\documentclass[12pt]{article}
\title{Digest of The New York Times}
\author{The New York Times}

% \definecolor{dark-green}{rgb}{0,0.25,0}
% \usepackage[utf8]{inputenc}

\usepackage{fontspec,xunicode,xltxtra}
\setmainfont[Ligatures=Required,Mapping=tex-text]{Georgia}
% \setmainfont[Ligatures=Required,Mapping=tex-text]{Palatino Linotype}
\setsansfont[Ligatures=Required,Mapping=tex-text]{Myriad Pro}
\setmonofont[Ligatures=Required,Mapping=tex-text]{Courier New}
\usepackage{xeCJK}
\setCJKmainfont{宋体}
\setCJKmonofont{黑体}

% \usepackage{indentfirst}


\usepackage[dvips,dvipsnames,svgnames]{xcolor}
\definecolor{light-gray}{gray}{0.95}

%生成PDF的链接
\usepackage[dvipdfmx]{hyperref}
\hypersetup{
    bookmarks=true,         % show bookmarks bar?
    bookmarksopen=true,
    pdfpagemode=UseNone,    % options: UseNode, UseThumbs, UseOutlines, FullScreen
    pdfstartview=FitH,
    pdfborder=1,
    pdfhighlight=/P,
    pdfauthor={wuxch},
    unicode=true,           % non-Latin characters in Acrobat’s bookmarks
    colorlinks,             % false: boxed links; true: colored links
    linkcolor=blue,         % color of internal links
    citecolor=blue,        % color of links to bibliography
    filecolor=magenta,      % color of file links
    urlcolor=cyan           % color of external links
}
% 使用目录
\usepackage{makeidx}
% 使用图形
\usepackage{graphicx}


\usepackage{titletoc}
\dottedcontents{section}[1.5em]{}{2.3em}{1pc}

% % 页眉页脚的设定
% \usepackage{fancyhdr}
% \pagestyle{fancy}
% \lhead{\includegraphics[height=0.2in]{zte_logo}}

% \usepackage{footmisc}
\usepackage{pfnote}
\usepackage{fnpos}

% % 嵌入的代码显示
% \usepackage{listings}
% \lstset{language=C++, breaklines, extendedchars=false}
% \lstset{keywordstyle=\color{blue},
%         commentstyle=\color{SeaGreen},
%         stringstyle=\ttfamily,
%         showstringspaces=false,
%         backgroundcolor=\color{light-gray}}
% \lstset{basicstyle=\small}

\usepackage{booktabs}

\usepackage{tipa}
\usepackage{phonetic}
\usepackage[gen]{eurosym}

% \usepackage{CJK}
\usepackage{makeidx}
% \usepackage{glossary}

\usepackage{fancyhdr}

\usepackage{lastpage}
\cfoot{\textsf{\thepage\ of \pageref*{LastPage}}}

\pagestyle{fancy}
\fancyhead{}
\renewcommand{\headrulewidth}{0pt}
\pagenumbering{arabic}

\usepackage{graphicx}
\usepackage{picinpar}

\setlength{\parindent}{0in}
\setlength{\parskip}{2.5ex plus4mm minus3mm}
% \setlength{\parskip}{.3\baselineskip}

\setcounter{section}{0}

% bibtex
\usepackage[super]{natbib}
% \usepackage[globalcitecopy,labelstoglobalaux]{bibunits}
% \makeatletter
% \renewcommand{\@bibunitname}{\jobname.\the\@bibunitauxcnt}
% \makeatother

% \renewcommand{\refname}{缩略语}
% \renewcommand{\refname}{}

\makeatletter                   % Reference list option change
\renewcommand\@biblabel[1]{#1:  } % from [1] to 1: 
\makeatother                    %

\newcommand{\citenumfont}[1]{\textsf{#1}}

% \usepackage{tipa}
% \usepackage{phonetic}

\newcommand{\phonetic}[1]{{\fontspec[Colour=CC00CC]{Microsoft Sans Serif} #1}}


%%% Local Variables:
%%% TeX-master: "medical_trivia"
%%% End:


\makeindex
\begin{document}
\date{}
\thispagestyle{empty}
\renewcommand\contentsname{\textsf{Digest of The New York Times}}
{\footnotesize\textsf{\tableofcontents}}
\clearpage
\setcounter{page}{1}

\section{The Teen Brain: The More Mature, the More Reckless}

\lettrine{T}{eenagers} are a famously reckless species. They floor the gas and experiment with drugs
and play with guns; according to Centers for Disease Control and Prevention figures, more than
16,000 young people die each year from unintentional injuries. The most common-sense explanation for
teens' carelessness is that their brains just aren't developed enough to know better. But new
research suggests that in the case of some teens, the culprit is just the opposite: the brain
matures not too slowly but, perhaps, too quickly.

In a paper just published in PLoS ONE -- a journal of the Public Library of Science -- a team led by
psychiatrist Gregory Berns of Emory University in Atlanta shows that adolescents who engage in more
dangerous activities have white-matter pathways that appear more mature than those of risk-averse
youths. White matter is essentially the brain's wiring -- the neural strands that connect the
various gray-matter regions, where the actual nerve cells reside, that are otherwise independent of
one another. Maturation of white matter is important because it increases the brain's processing
speed; nerve impulses travel faster in mature white matter. 

Berns and his colleagues recruited 91 kids ages 12 to 18 and asked them to fill out a questionnaire
about their tendency to engage in behaviors such as driving without a license, having unprotected
sex and using drugs. Then they had the kids undergo a relatively new kind of brain scan called
diffusion tensor imaging, a type of magnetic resonance imaging that is used to look at dense tissues
like white matter. After analyzing the scans, the authors found a strong correlation between how
risky the students described their behavior to be and how sophisticated their white matter was. The
more mature the look of the brain, the more risk-taking the teenager tended to report.

The results challenge the accepted notion that teens make dumb decisions because their brains are
immature. Although previous research has shown that most teens' gray-matter structures -- including
those involved in \\decision-making -- are less advanced than those of adults, as you would expect,
until now no one had studied teens' white matter, which works along with gray matter to produce
decisions. The key part of white matter is called myelin, a fatty substance that coats the
individual neural strands, or axons, that make up white matter. Myelination of axons begins during
childhood and is completed at the end of adolescence, around the mid-20s. Myelination in the frontal
lobe -- the brain region responsible for decision-making -- happens last, and it was in this region
that the brains of risk-seeking teens in the study showed greater development compared with the
frontal lobes of their more restrained peers. 

Why would kids who take more risks turn out to have more adultlike white matter than other kids? The
authors suggest that some risk-taking among adolescents is evidence that they are trying out more
adultlike roles. Having unsafe sex and driving too fast may be mistakes, but kids often have to
experiment with limits in order to learn how to live within them. Which, in turn, is a sign of
maturity. ``Adolescents who engage in \emph{risky} behaviors obtain more experience in a variety of
domains,'' the authors write. ``Their more conservative peers, in contrast, do not have as much
'life experience' and therefore might be expected to have more immature brains.''

Another possible explanation is that some teenagers whose brains develop more quickly than others
become uncomfortable with the gap between their biological capabilities and the social rules they
must follow as kids. ``Precocious development of these \emph{white-matter} tracts may predispose
some adolescents to engage in behaviors that society considers too adult in nature for their
chronological age,'' the authors write. In other words, having a more mature brain may actually spur
some kids to seek out new and potentially harmful experiences.

For now, these theories are just speculation, and the researchers concede that the interaction of
white and gray matter is so complex that hard conclusions remain elusive. ``We have a new piece to
the puzzle here,'' says Emory's Monica Capra, one of the study's authors. ``But we don't have it all
together.''

One further x-factor in the study: What if teens weren't being honest in their self-reports about
risky behavior? Berns' team addressed this question by drug-testing the 91 research subjects. Only
nine had actually done drugs -- in each case, marijuana -- but eight of those nine admitted their
drug use in the survey. No students who tested negative falsely claimed to have tried drugs. The
teen brain, it appears, can be often an honest thing -- even if it's not always a wise one.

\section{U.S. Is Seeing Policy Thorns\cite{thorn} in Japan Shift}

\lettrine{J}{apan}'s landmark election presents the Obama administration with an untested
government, creating a new set of imponderables for a White House already burdened by foreign policy
headaches in Afghanistan, Iran and North Korea.

Inside the administration, the historic change in Tokyo is raising concerns that Japan may back away
from supporting key American priorities like the war in Afghanistan or the redeployment of American
troops in Asia, according to senior officials.

Specifically, the newly elected Democratic Party says it may recall the Jap-anese naval forces from a
mission to refuel American warships near\\Afghanistan. And it wants to reopen an agreement to
relocate a Marine airfield on Okinawa, which requires Japan to pick up much of the cost for moving
thousands of Marines to Guam.

The victory of the Democrats on Sunday means the White House must deal, for the first time in
decades, with a Japanese government that is a complete stranger, and one that has expressed
blunt\cite{blunt} criticism of the United States. The party's leader and presumptive prime minister,
Yukio Hatoyama, recently spoke out against American-led globalization and called for a greater
Japanese focus on Asia.

Despite the party's campaign rhetoric, its leaders insist they will not threaten the alliance with
the United States, particularly when Japan faces a fast-rising China and a nuclear-armed North
Korea. Senior American officials said they expected Japan to remain a bulwark in Asia, even noting
that the new government, unburdened by history, could play a more central role in negotiating with
North Korea.

But for the most part, the United States is perplexed\cite{perplexed} by what one official described
as a ``seismic\cite{seismic} event,'' with unknown consequences for one of its most important
relationships.

``The election of a new party could produce new ways of doing things, which we will have to adjust
to,'' said a senior official, who spoke on condition of anonymity because of the delicacy of the
matter. ``You'll have this period of unpredictability.''

The big question many in Washington are asking is whether the vote was a harbinger\cite{harbinger}
of a deeper change in Japan, away from its historic dependence on the United States.

``There is a fear of dramatic change in the U.S.-Japan alliance,'' said Michael Auslin, an expert on
Japanese foreign policy at the American Enterprise Institute in Washington. ``No one knows what will
happen next, or even who to talk to for answers.''

The Democratic Party struck a chord with its talk of improving ties with China and other neighbors,
reflecting the fact that Japan's \$5 trillion economy has grown more dependent on commerce with its
neighbors.

Fears of Japanese drift seemed to be confirmed last week when an article by Mr.~Hatoyama, excerpted
and translated from a Japanese journal, appeared on the Web site of The New York Times. It stirred a
hornet's nest in Washington by casting Japan's embattled economy as the victim of American-inspired
free-market fundamentalism. Yet it also stressed the importance of the American alliance.

Mr.~Hatoyama's views sent many in Washington's diplomatic establishment scurrying\cite{scurry} to
learn more about him and the Democrats. That highlighted a problem: While American officials and
academics have spent decades cultivating close ties with the Liberal Democrats, who have governed
Japan for most of the last half century, they have built few links to the opposition.

Some Japan experts said it would be a mistake to read too much into Mr.~Hatoyama's remarks, and
Japanese officials privately conveyed that same caution to the Obama administration.

``It was an indication they still haven't figured out what they're going to do in power,'' said
Michael J.~Green, a professor of international relations at Georgetown University who served on the
National Security Council during the last Bush administration. ``This could get confused and
dysfunctional for a while.''

Stung by the reaction, Mr.~Hatoyama appears to be back-pedaling and engaging in damage control. On
Monday night, he said he had not intended for the article to appear abroad, and said it was being
misinterpreted. ``If you read the entire essay, you will understand that it is definitely not
expressing anti-American ideas,'' he said.

Professor Green noted that in many ways, relations between the United States and Japan were smoother
now than in years past because the trade disputes of the 1980s and 1990s were largely settled.

He said the new government would find that some of its proposals, like reopening talks on the
relocation of the Futenma Marine airfield on Okinawa, were unrealistic, given the years it took to
negotiate that deal. For the Obama administration, he said, the challenge will be to give Japan's
new leaders a face-saving way to back down.

Japan, experts said, could play a more muscular role in talks with North Korea if, as expected, the
Democrats turn down the heat on the issue of Japanese abducted by North Korea decades ago, a
perennial\cite{perennial} sticking point for the Liberal Democrats.

And Obama administration officials said they were eager to dispel\cite{dispel} perceptions in Japan
that a better relationship with China would somehow undermine its alliance with the United States.

``We have no desire to see our defense commitment tested by battle,'' a senior official said. ``We
see no contradiction between Japan reducing frictions with China and a strong Japan-U.S.~alliance.''

In recent years, many Japanese have thought the United States took the relationship for granted,
paying more attention to China.

Traditionally, the United States has sent high-powered diplomats or political figures to Tokyo. But
the Obama administration chose to send a big campaign donor, John Roos, as ambassador, passing over
a longtime Asia hand, Joseph S.~Nye Jr., who had been championed by Secretary of State Hillary
Rodham Clinton.

Administration officials counter that Mr.~Roos, a Silicon Valley lawyer, will be influential because
he has the ear of President Obama.

Political analysts and former diplomats say the Democrats are so sharply divided ideologically --
between pacifist former Socialists and flag-flying former Liberal Democrats -- that they will avoid
treading too heavily on divisive foreign policy issues for fear of splitting the party.

Policy analysts also say the Japanese public would turn against the Democrats if they undermined the
Washington alliance, pointing out that the opposition won because of anger with the incumbents'
failed economic policies, not because of a desire to change the nation's reliance on the United
States, which remains widely accepted here.

``They do not have a mandate for changing the alliance with the U.S.,'' said Yukio Okamoto, a former
adviser to several prime ministers on foreign affairs.

\section{Canadian Politician Charged in Death}

\lettrine{A}{s} the high-profile attorney general for Ontario, Michael Bryant had championed severe
and controversial traffic safety laws. On Tuesday, he was charged with criminal negligence causing
death and with dangerous driving causing death in an unusually violent episode of road rage
involving a bicyclist.

The arrest of Mr.~Bryant stemmed from a collision between a bicycle and an automobile in Toronto's
most prestigious shopping district late Monday evening. The episode started off as minor but swiftly
escalated, leading to the death of Darcy Allan Sheppard, 33, who was identified as a bicycle
courier.

After the collision, Mr.~Sheppard apparently grabbed the driver's side door and held on. Within
moments, the police received reports of a Saab convertible racing past the fashionable shops of
Bloor Street with a man clinging to its side. Two construction workers doing repairs along the road
told CTV, a Canadian television network, that the car accelerated, its tires squealing, before
veering into oncoming traffic on the left side of the street.

The workers said that the motorist repeatedly mounted the sidewalk and drove near lampposts in what
seemed to be an attempt to brush off the man hanging onto the side.

One of the workers said the driver was ``yelling pretty loud and he sounded very, very angry.'' The
other worker said, ``He meant to knock him off.''

Several witnesses said that the clinging man flew off of the car after striking a mailbox. Sgt.~Tim
Burrows of the Toronto police traffic division said Mr.~Sheppard died shortly after arriving at a
hospital.

The convertible\cite{convertible} raced into the covered driveway of a nearby luxury hotel, the
witnesses said. Not long after, Mr.~Bryant was filmed by television crews sitting in the back of a
police cruiser. A black Saab convertible near the hotel entrance had extensive damage on its
driver's side.

On Tuesday, Mr.~Bryant left a Toronto police station after being charged. He made a brief comment,
offering his ``deepest condolences'' to Mr.~Sheppard's family.

Until he quit politics this spring to start an economic development agency for the city of Toronto,
Mr.~Bryant, a Liberal, was among the highest-profile members of Ontario's government. In his last
role, as economic development minister, Mr.~Bryant negotiated the province's participation in the
bailout of Chrysler and General Motors.

But it was as the province's top law enforcement official that he produced the greatest controversy
and acclaim\cite{acclaim} of his 10-year political career.

Among his targets were street-racing motorists. In 2007 he gave the police the power to seize and
destroy cars modified for racing even if no charges were lodged against their owners.

After describing such cars as being as dangerous as explosives, he said, ``We will crush your car,
we will crush the parts.''

Later that year the province passed a bill to deem any vehicle traveling more than 50 kilometers an
hour, or 31 miles an hour, faster than the speed limit to be racing. The legislation, under which
more than 10,000 charges have been brought, allows the police to immediately seize vehicles and
suspend licenses.

\section{Changing Health Care by Steps}

\lettrine{A}{fter} Harry Truman repeatedly failed to persuade Congress to pass universal health
insurance, some Truman administration officials came up with a less ambitious idea. They suggested
covering only 8 percent of the population, and an especially sympathetic 8 percent at that:
everybody 65 and older.

Truman never really pushed the plan, however. John F.~Kennedy later did, yet was
stymied\cite{stymie} by Congress and the American Medical Association, which equated it with
Soviet-like socialism. So it fell to Lyndon Johnson. Even after he won a 23-point landslide in 1964,
he had to agree to some unseemly deal-making, as Jonathan Cohn of The New Republic has noted, that
handed a big payday to hospitals and doctors.

Only then, in the summer of 1965, was Medicare born.

Next week, Congress will return to session, and health care, of course, will be at the top of its
agenda. Passing a bill, it's clear, will be no easier than in previous decades. President Obama's
poll numbers have fallen, while untruths about death panels have made the rounds and members of
Congress have been subjected to town hall harangues\cite{harangue}.

But the job facing Mr.~Obama hasn't really changed: he will have to figure out how to end up more
like Johnson than like Truman or, more recently, Bill Clinton. He and Congress will have to figure
out how to make some progress toward fixing the country's troubled health care system.

Any bill they pass will inevitably be flawed. It will not do enough to reduce wasteful spending. It
probably will not result in universal coverage. Special interests -- like drug companies and, once
again, hospitals -- will get off too lightly.

But such is politics. Starting from scratch with a more purely liberal bill or a more conservative
one, as some have urged, won't change the messy reality of democracy. ``You're not going to fix
health care on one bill,'' Senator Mark Warner, a Virginia Democrat, said. ``Even if my most liberal
colleagues had their perfect bill, it's not going to fix health care. A bill is going to have to be
altered and amended later.''

Indeed, from an economic standpoint, the biggest risk seems to be that health care will be left
unchanged and that we'll simply pretend we are as healthy, wealthy and wise as can be.

We're not. The United States is the only rich country in the world without something resembling
universal health insurance. Medical costs already take a big whack out of workers' paychecks, and
Medicare's budget is growing at an unsustainable rate. Despite all this spending, we too often
receive inadequate care.

Emergency room care and certain cancer treatments are among the best in the world. But management of
chronic diseases isn't nearly as good. Senior citizens and children are less likely to receive
vaccinations here than in some other wealthy countries, a recent analysis found. Americans face
higher risks of being killed by a medical error or harmed by a medication error.

Health reform won't solve all these problems. But there is still reason to think that it can improve
things. A limited, imperfect, warts-and-all bill would, as in 1965, be a lot better than no bill at
all.

•

The main political problem with health reform is still the immediate cost of it. Covering the
uninsured will require \$1.2 trillion or so over 10 years, and Congress hasn't been able to agree on
\$1.2 trillion of spending cuts and tax increases to keep the deficit from growing.

This standard -- deficit neutrality -- may not be fair. George W.~Bush, after all, spent trillions
of dollars cutting taxes and adding a Medicare prescription drug benefit without offsetting the
costs. Yet given the current size of the deficit, Mr.~Obama and Congress have to do better.

So a scaled-down version of health reform will probably mean spending less money to subsidize health
insurance, thus leaving some of the uninsured without coverage. Such a compromise would not be the
end of the world. (Medicare, remember, initially applied to a measly 8 percent of the population.)
But it would create one tricky technical problem.

Any plan to expand insurance depends on creating a broad group of people known as a risk pool. A
company's work force is such a pool: it allows an insurer to balance the high cost of covering the
sick with the low cost of covering the healthy.

If Congress passed a universal health insurance plan, the 45 million or so newly insured would make
up their own risk pool, one that included both today's young, healthy uninsured as well as older,
sicker uninsured. But if Congress stops short of universal coverage, it then has a problem.

It can't allow just people to sign up for subsidies voluntarily. That would create an incentive for
the sick to buy insurance and the healthy to go without it. Average insurance costs would then soar
-- and the subsidies would have to soar too, undercutting the plan.

There are ways out of this problem. Congress could mandate that the young and, say, the near poor
buy insurance. Together, they would create a decent risk pool.

A second big question is whether health reform will include any serious attempt to reduce medical
costs and improve quality. Economically, this issue may matter more than any other. Politically, it
would make the White House seem serious about reducing the budget deficit.

``Whatever new bill is going to come about has to have a far greater emphasis on cutting costs and
making the delivery system far more efficient,'' said Senator Charles E.~Schumer, Democrat of New
York. ``That's true substantively and politically.''

Right now, we pay hospitals, doctors and drug companies to give us more care rather than better
care, and we get what we pay for. I recently described how this system undermines the treatment of
prostate cancer, but there are dozens of examples.

Unfortunately, none of the major proposals in the House of Representatives address the problem. The
House's liberal leaders have been mostly uninterested in changing the incentives in health care. So
have conservative Blue Dog Democrats, despite their fiscal discipline rhetoric\cite{rhetoric}.

Republicans, for their part, have signaled that their top priority is defeating a bill; they even
rejected Mr.~Obama's offer to include tort reform, a pet Republican issue. It looks as if any attack
on fee-for-service medicine will fall to Senate Democrats.

And then there is the president.

Persuading Americans that health care needs reforming was always going to be hard. It depends on
making people see the logic in a series of hypotheticals. If we don't reduce the growth in costs, we
will leave our children with a crushing tax burden. If we do reduce costs, our paychecks will be
fatter. If we reward quality of care instead of quantity of care, fewer people will die an early
death. This isn't easy stuff.

But Mr.~Obama didn't get elected because he promised to do the easy jobs better than Hillary Clinton
or John McCain. He won the election by vowing to be a Reagan-like communicator who could rally
Americans to deal with daunting\cite{daunt}, intractable\cite{intractable} problems -- the kind of
problem that bedeviled Bill Clinton, John F.~Kennedy and Harry Truman.

Mr.~Obama and his advisers have spent much of this week trying to come up with a post-Labor Day
strategy for selling health reform. They understand very well that incremental success is a lot
better than nothing, and the Democratic-controlled Congress means that their odds are better than
even. But they're not there yet.

\section{In a Sale, Skype Wins a Chance to Prosper}

\lettrine{W}{ith} its sale to private investors, the online calling service Skype has thrown off the
last of the shackles\cite{shackle} that limited its growth and potential as a unit of eBay. Now its
challenge is to turn its global popularity into bigger profits.

EBay announced on Tuesday that it was selling Skype to a group led by Silver Lake Partners, a
private equity firm in Silicon Valley. As part of the deal, which values Skype at \$2.75 billion,
the buyers agreed to pay \$1.9 billion in cash, which includes a loan from eBay of \$125 million,
for 65 percent of the company. EBay, which is based in San Jose, Calif., will retain a 35 percent
stake.

Skype offers free software for computers and smartphones that lets people make free voice and video
calls to other Skype users over the Internet. It makes money primarily by charging for calls to
landlines and cellphones around the world, although its rates are generally far below those of
traditional phone companies.

When it acquired Skype in 2005, eBay said it hoped the service would support its auctions and its
PayPal payment service by letting buyers and sellers discuss transactions. But eBay users were not
so chatty.

The deal announced Tuesday would essentially allow Skype to go back into start-up mode. For example,
it will be able to conceal investments in projects and new technologies from the public -- and from
rivals -- instead of disclosing them in eBay's public regulatory filings.

The company will also be able to lure new employees with its own stock, which could become valuable
if Skype's buyers decide to sell shares in an initial public offering.

``This gives us a great set of investors who are going to add a lot of value to the business,'' said
Josh Silverman, Skype's chief executive. ``And as a stand-alone company we are focused solely on
communications, and there is always benefit to focus.''

The group buying Skype includes the London venture capital firm Index Ventures and Andreessen
Horowitz, a new venture capital firm co-founded by Marc Andreessen, the Netscape co-founder.

The deal completes a journey that began early last year, when Meg Whitman, eBay's longtime chief
executive, left the company and her deputy, John Donahoe, took over. Mr.~Donahoe moved Mr.~Silverman
over from Shopping.com, another eBay division, and gave Skype, which is based in Luxembourg, wide
latitude to operate independently.

``All the previous presidents had a really short leash back to San Jose, and the company was
basically a bureaucratic mess,'' said Phil Wolff, editor of Skype Journal, a blog covering the
service. ``With the changes last year, the company got a strong leadership team and a strong sense
of direction.''

The result has been steady, solid growth, even as competing Internet calling services like Google
Voice came on the scene. Skype, which had 276 million registered users in the first quarter of last
year, ended July with 480 million.

Revenue rose to \$170 million in the second quarter of this year from \$136 million in the same
quarter a year earlier. EBay does not break out Skype's profits but says it has been profitable for
10 consecutive quarters.

In a sign that Skype may have discovered new opportunities outside the personal computer, the Skype
application for the iPhone has been one of the most popular programs for the device since it was
released in March. Skype has also struck deals to place its service on Nokia phones.

Egon Durban, managing director at Silver Lake Partners, said the iPhone application was ``a great
example of what we feel are the attractive opportunities for the company to develop.''

``This is one of the leading Internet franchises with terrific growth prospects,'' Mr.~Durban said.
He gave no specifics on features Skype might offer, but said it was easy to imagine possibilities.
He also said the buyers had no plans for a public offering.

Some of Skype's newest features may suggest directions for the company. The latest version of its
software, released this year, emphasizes face-to-face video chats. Skype now says that 34 percent of
calls between Skype users include video, and such chats have become a popular way for people around
the world to connect with one another.

In addition to working on allowing outside programmers to weave Skype's features and infrastructure
into their own programs or Web sites, the company is also working on replacing the service's
underlying peer-to-peer technology. That is partly out of legal necessity. The rights to that
technology remain with Skype's founders, Niklas Zennstrom and Janus Friis, who have sued eBay in a
British court over some changes eBay made to it. That case is to go to trial next year.

For eBay, selling Skype -- particularly at a valuation higher than many analysts had thought
possible -- offers partial redemption for a deal that many Internet analysts said was an awkward
fit. EBay paid \$2.6 billion for Skype, and performance incentives lifted the final price to \$3.1
billion.

``The purchase was a serious mistake by Meg Whitman. It was an attempt to buy growth, which
investors saw through instantly,'' said Jeffrey Lindsay, an analyst at Sanford C.~Bernstein. The new
eBay management found ways for Skype to generate revenue, he said, ``and now they sold it at a great
price.''

Mr.~Donahoe said eBay did not regret having bought Skype when company executives believed eBay was
in a mortal struggle with Google, which was also pursuing the service. He said the spinoff would
allow eBay to focus on its core e-commerce and online payment businesses and avoid extra
distractions.

``We don't regret having done this at all. We compete in a dynamic market, and you have to move
quickly and take risks,'' Mr.~Donahoe said. ``When we bought Skype we thought it had synergies with
our other two businesses, and it turns out it did not. But it also turned out that it's a great
stand-alone business.''

\section{The Race to Be an Early Adopter of Technologies Goes Mainstream, a Survey Finds}

\lettrine{F}{or} decades, the adoption and use of the latest technologies was limited to a
subculture: Whether called ``tech enthusiasts'' or ``gadget geeks,'' the implication was that most
of the world got along fine with older, established products and services, while a smaller group
pursued the most leading-edge technology.

But according to a study released Wednesday by Forrester Research, a marketing firm based in
Cambridge, Mass., a shift has taken place. What used to be the pursuit of a few has become decidedly
mainstream. We're all gadget geeks now.

According to the study, which surveyed 53,668 households in the United States and Canada by mail,
half of all American adults are gamers. Sixty-three percent of American households have a broadband
Internet connection. Three-quarters of American households have cellphones and PCs. And nearly 10
million American households, out of nearly 118 million, added an HDTV in the last year, a jump of 27
percent over 2007.

``There's really no group out of the tech loop,'' said Jacqueline Anderson, an analyst with
Forrester who was one of the study's authors. ``America is becoming a digital nation. Technology
adoption continues to roll along, picking up more and more mainstream consumers every year.''

High-definition television sets were one of the fastest-growing consumer technologies in 2008. Over
the next five years, the company forecasts, nearly 39 million households in the United States will
get their first high-definition set, bringing total market penetration for HDTV to nearly 70
percent.

The study also found that despite the recession, online spending remained strong, with older
consumers leading the charge. On average, those consumers spent \$560 in the last three months,
although ``20 percent of that group spent more than \$1,000 online in the last three months,''
Ms.~Anderson said. Given the tumultuous economic climate, ``that's a lot of money,'' she said.

Ms.~Anderson also pointed out that families were a big driver behind the widespread adoption of
technologies. The popularity of video game consoles like the Nintendo Wii, which took a decidedly
different approach from other game-console makers by appealing to nongamers and families, created an
opening for more digital entertainment to enter the home.

Families are also more likely to have gadgetry like MP3 players, digital cameras and digital
camcorders. ``They have little kids so they want to catalog those memories,'' Ms.~Anderson said. In
addition, 86 percent of families with children had mobile phones but were also more likely to use
mobile phones with more features like music and video playback.

The study also suggests a growing reliance on the Internet for commerce, communication,
entertainment and social lives, said Charles S.~Golvin, an analyst with Forrester Research, and a
co-author on the study.

``The digitization of our daily lives has been steadily ramping up over the past decade,''
Mr.~Golvin said.

One area that appears to be slower to catch on is home networks. The survey found that 33 percent of
households in the United States with an Internet connection reported having a home network, up from
28 percent a year ago. Although Ms.~Anderson says that figure is relatively high, the adoption is
still lower compared with the adoption of other home technologies.

``The barrier to entry for a home network is a lot higher than for an HDTV, where all you have to do
is buy one,'' she said. ``There are more components and you have to understand how to connect them.
Many people had the components for a home network before but didn't necessarily understand what it
meant to put them together or why they'd want to,'' she said. In the next five years, the company
forecasts, more than 30 million households will install a home network, bringing market penetration
to just over 50 percent.

Already, Mr.~Golvin says, more people are migrating away from the home and office to use the Web and
turning toward their smartphones. About 15 percent of cellphone owners were using the Internet on
their phones in 2008, the study found, showing that, for a growing number of Americans, there is an
increasing ``expectation that all the same services and resources are available to us no matter
where we are,'' he said.

\section{Filmmakers Barred From Chinese Festival}

\lettrine{W}{hen} the American filmmakers Jon Alpert and Matthew O'Neill traveled around Sichuan
Province last year to document the anger of parents whose children had died in school collapses
during the earthquake in May, they ran into a chilly reception from officials.

Police officers harassed the two men and their co-workers, detained them and interrogated them for
eight hours, they said.

Now, the Chinese government has denied both of them visas, blocking them from presenting their
documentary, ``China's Unnatural Disaster: The Tears of Sichuan Province,'' at the Beijing
Independent Film Festival this week. The two men, who made the film for HBO with the co-producer
Peter Kwong, said their visa applications were rejected late last week.

``We are extremely disappointed that the Chinese government denied our request for visas and that we
will not be permitted to discuss this film with a Chinese audience in Beijing,'' Mr.~Alpert and
Mr.~O'Neill said in a joint e-mail message. ``The denial of our visas fits in with a pattern of what
seems to be a complete commitment on the part of this Chinese government to crush any inquiry into
the possibility of wrongful deaths during the earthquake in Sichuan.''

Chen Cong, a vice consul in the press office of the Chinese Consulate in New York, declined to
explain the rejection, saying that diplomatic organizations had ``the right not to give a reason for
why the visa was denied.''

Mr.~Alpert and Mr.~O'Neill have both won Emmy Awards and have worked together on highly praised
documentaries, including ``Baghdad ER.'' The Sichuan documentary was shown on HBO in May, one year
after the earthquake, and got positive reviews. The official Web site of the film is blocked in
China.

The Chinese government has gone to great lengths to silence any mention of the collapsed schools
and, according to an official count, the 5,335 children who died or remain missing. In the weeks
after the earthquake, which left nearly 87,000 people dead or missing, parents took to the streets
to demand official investigations into why so many school buildings had collapsed even though other
buildings around them remained standing. The parents said shoddy construction and corruption were
the obvious causes.

Local officials ordered security forces to detain the parents or tried to buy the silence of the
parents with compensation money. Meanwhile, journalists who tried approaching the schools were
stopped, and two rights advocates who pressed for official inquiries were detained. The two
advocates, Huang Qi and Tan Zuoren, were put on trial last month.

Artists trying to raise consciousness over the collapsed schools have been similarly harassed. The
Chinese filmmaker Pan Jianlin was tracked by security officials after his documentary on the deaths,
``Who Killed Our Children?'' was shown last year at the Pusan International Film Festival in South
Korea. Ai Weiwei, a prominent artist who often criticizes the Communist Party, had his Web site
blocked after he tried to compile online a comprehensive tally of dead schoolchildren. He was
temporarily detained in Sichuan last month when he tried to attend the trial of Mr.~Tan.

A person helping to organize the film festival in Beijing said the HBO documentary would be shown on
Thursday even though the filmmakers would not be able to attend. The festival is showcasing more
than 80 films, and each one is generally shown once.

\section{Beijing Limits Information on Burmese Refugees Remaining in China}

\lettrine{C}{hinese} officials imposed an information blackout on Tuesday on the situation along its
border with Myanmar and began taking down tents that had sheltered an estimated 30,000 refugees who
fled into China to escape recent fighting between Myanmar's military and ethnic rebels.

But news reports stated that many thousands of refugees remained in China, unwilling or unable to
return to Myanmar, formerly called Burma, and it was not clear how the Chinese government intended
to address their plight.

The Chinese authorities withheld comment on the border situation on Tuesday, aside from saying, in a
Foreign Ministry briefing, that ``necessary humanitarian assistance'' was being provided. And they
began ordering foreign journalists to leave the area around Nansan and Genma, Chinese towns on the
mountainous border where the refugees have been housed in seven separate camps.

While about 4,000 refugees had returned to Myanmar on Monday, the day after the fighting ended, the
pace has since slowed significantly. Only about 30 people crossed the border into Myanmar in a
half-hour period on Tuesday morning, The Associated Press reported.

``It seems to be slowing down,'' one foreigner near Nansan said in a telephone interview on Tuesday.
``There's still a large number of refugees in and around Nansan, both in the camps and hanging
around.'' The foreigner, who asked not to be identified, said Chinese Army troops had stepped up
patrols in the area.

An unknown number of those who fled to China during the fighting are Chinese citizens who have been
conducting business in Myanmar, where China is building dams and other projects and has extensive
mining ventures. They are unlikely to return soon.

China has insisted that the northern Myanmar region of Kokang is safe and stable after the fighting
last week, in which hundreds of government troops overwhelmed an armed ethnic group, breaking a
cease-fire that had prevailed for two decades. Human rights groups and others have warned that the
junta's actions could ignite a wider conflict in the area, where other, better armed, ethnic groups
also are resisting government control.

Thai newspapers and The Irrawaddy, an independent magazine that focuses on Myanmar, have reported
that the government is sending fresh troops into the northern state of Shan in an attempt to
consolidate its control there. The army wants the rebels to disarm and join a government border
patrol force, as required under a new Constitution. Most of the rebels have resisted the order,
which would effectively place them under government control.

Myanmar's military junta\cite{junta} apparently seeks to take control of the region before
elections, the first in almost 20 years, that are scheduled for next year. Outside monitors accuse
the military junta of brutal human rights violations as part of its effort to stay in power. The
Myanmar government has said that 26 of its soldiers and at least 8 rebels died in three days of
battles.

The Myanmar conflict has thrust the Chinese government, one of Myanmar's only staunch\cite{staunch}
backers, into an awkward situation. China has provided diplomatic support to the junta in exchange
for access to its considerable mineral wealth and cooperation in efforts to suppress a growing
cross-border trade in heroin and other illicit drugs. The flood of refugees prompted the Chinese to
issue muted criticism of the junta, on Friday calling for it to secure Myanmar's borders.

\section{Heavy Data Use Puts a Strain on AT\&T Service}

\lettrine{S}{lim} and sleek as it is, the iPhone is really the Hummer of cellphones.

It's a data guzzler. Owners use them like minicomputers, which they are, and use them a lot. Not
only do iPhone owners download applications, stream music and videos and browse the Web at higher
rates than the average smartphone user, the average iPhone owner can use 10 times the network
capacity used by the average smartphone user.

``They don't even realize how much data they're using,'' said Gene Munster, a senior securities
analyst with Piper Jaffray.

The result is dropped calls, spotty service, delayed text and voice messages and glacial download
speeds as AT\&T's cellular network strains to meet the demand. Another result is outraged customers.

Cellphone owners using other carriers may gloat now, but the problems of AT\&T and the iPhone
portend their future. Other networks could be stressed as well as more sophisticated phones
encouraging such intense use become popular, analysts say.

Taylor Sbicca, a 27-year-old systems administrator in San Francisco, checks his iPhone 10 to 15
times a day. But he is not making calls. He checks the scores of last night's baseball game and
updates his Twitter stream. He checks the local weather report to see if he needs a coat before
heading out to dinner -- then he picks a restaurant on Yelp and maps the quickest way to get there.

Or at least, he tries to.

``It's so slow, it feels like I'm on a dial-up modem,'' he said. Shazam, an application that
identifies songs being played on the radio or TV, takes so long to load that the tune may be over by
the time the app is ready to hear it. On numerous occasions, Mr.~Sbicca says, he missed invitations
to meet friends because his text messages had been delayed.

And picking up a cell signal in his apartment? ``You hit the dial button and the phone just sits
there, saying it's connecting for 30 seconds,'' he said.

More than 20 million other smartphone users are on the AT\&T network, but other phones do not drain
the network the way the nine million iPhones users do. Indeed, that is why the howls of protest are
more numerous in the dense urban areas with higher concentrations of iPhone owners.

``It's almost worthless to try and get on 3G during peak times in those cities,'' Mr.~Munster said,
referring to the 3G network. ``When too many users get in the area, the call drops.'' The problems
seem particularly pronounced in New York and San Francisco, where Mr.~Munster estimates AT\&T's
network shoulders as much as 20 percent of all the iPhone users in the United States.

Owners of the iPhone 3GS, the newest model, ``have probably increased their usage by about 100
percent,'' said Chetan Sharma, an independent wireless analyst. ``It's faster so they are using it
more on a daily basis.''

Mr.~Sharma compares the problem to water flowing through a pipe. ``It can only funnel so much at a
given time,'' he said. ``It comes down to peak capacity loads, or spikes in data usage. That's why
you see these problems at conferences or in large cities with high concentration of iPhone users.''

When thousands of iPhone owners descended on Austin, Tex., in March during South by Southwest, an
annual technology and music conference, attendees were unable to send text messages, check their
e-mail or make calls until AT\&T installed temporary cell sites to amplify the service.

AT\&T's right to be the exclusive carrier for iPhone in the United States has been a golden ticket
for the wireless company. The average iPhone owner pays AT\&T \$2,000 during his two-year contract
-- roughly twice the amount of the average mobile phone customer.

But at the same time the iPhone has become an Achilles' heel for the company.

``It's been a challenging year for us,'' said John Donovan, the chief technology office of AT\&T.
``Overnight we're seeing a radical shift in how people are using their phones,'' he said. ``There's
just no parallel for the demand.''

AT\&T says that the majority of the nearly \$18 billion it will spend this year on its networks will
be diverted into upgrades and expansions to meet the surging demands on the 3G network. The company
intends to erect an additional 2,100 cell towers to fill out patchy coverage, upgrade existing cell
sites by adding fiber optic connectivity to deliver data faster and add other technology to provide
stronger cell signals.

As fast as AT\&T wants to go, many cities require lengthy filing processes to erect new cell towers.
Even after towers are installed, it can take several months for software upgrades to begin operating
at faster speeds.

The company has also delayed bandwidth-heavy features like multimedia messaging, or text messages
containing pictures, audio or video. It is also postponing ``tethering,'' which allows the iPhone to
share its Internet connection with a computer, a standard feature on many rival smartphones. AT\&T
says it has no intention of capping how much data iPhone owners use.

The upgrades are expected to be completed by next year and the company has said it is already seeing
improvements.

But AT\&T faces another cost -- to its reputation. AT\&T's deal with Apple is said to expire as
early as next year, at which point other carriers in the United States would be able to sell the
popular Apple phones. Indeed, a recent survey by Pricegrabber.com found that 34 percent of
respondents pinpointed AT\&T as the primary reason for not buying an iPhone.

``It's a P.R. nightmare,'' said Craig Moffett, a senior analyst with Sanford C.~Bernstein \&
Company.

AT\&T might be in the spotlight now, analysts say, but other carriers will face similar problems as
they sell more smartphones, laptop cards and eventually tablets that encourage high data usage.

Globally, mobile data traffic is expected to double every year through 2013, according to Cisco
Systems, which makes network gear. ``Whether an iPhone, a Storm or a Gphone, the world is
changing.'' Mr.~Munster said. ``We're just starting to scratch the surface of these issues that
AT\&T is facing.''

In preparation for the next wave of smartphones and data demands, all the carriers are rushing to
introduce the next-generation of wireless networks, called 4G.

Analysts expect that in a year or so, AT\&T's network will have improved significantly -- but it may
not be soon enough for some iPhone owners paying for the higher-priced data plans, like Mr.~Sbicca,
who says he plans to switch carriers as soon as the iPhone becomes available on other networks.

``What good is having all those applications if you don't have the speed to run them?'' he said.
``It's not exactly rocket science here. It's pretty standard stuff to be able to make a phone
call.''

\section{YouTube Said to Consider Pay Movies}

\lettrine{Y}{ouTube}, the largest video site, is in negotiations with major Hollywood studios for a
deal that would let its visitors pay to watch full-length movies, according to two people briefed on
the negotiations.

If an agreement is reached, it would be a major change for YouTube, which has built a huge audience
by offering an eclectic collection of free video clips and earns most of its revenue from
advertising. It would also put YouTube, which is owned by Google, in direct competition with
services from Netflix, Amazon and Apple, which allow users to buy or rent movies online.

YouTube, which already offers some older movies on its site free, is talking with Lionsgate
Entertainment, Sony and Warner Brothers about making newer ones available, said one person briefed
on the talks, which were intended to be confidential.

YouTube, which has long sought to add more professionally produced video, said in a statement:
``While we don't comment on rumor and speculation, we hope to expand both on our great relationships
with movie studios and on the selection and types of videos we offer our community.''

Scott Rowe, a spokesman for Warner Brothers, and Jim Kennedy, a spokesman for Sony Pictures,
declined to comment. Peter Wilkes, a spokesman for Lionsgate, also would not comment specifically
but said that the company was ``always exploring alternatives'' that would help it make more money
from its films.

Lionsgate, Mr.~Wilkes said, has enjoyed considerable success with offering movies and television
shows through Apple. According to Mr.~Wilkes, the studio's ``Mad Men'' series has been downloaded
about two million times through Apple's iTunes store.

The negotiations between YouTube and the studios were first reported on the Web site of The Wall
Street Journal.

Movie studios have been pushing YouTube to consider charging for certain types of content, said a
person briefed on the discussions. YouTube appears willing to do so if the studios agree to give it
access to enough of their newer films close to when they come out on DVD, the person said.

One studio executive -- who was briefed on the talks but spoke on the condition of anonymity to
minimize disruption to the discussions -- said the issues still to be resolved were pricing and the
timing of YouTube releases. Though DVD sales have dropped, studios remain protective of the period
during which films and shows are available on DVD but not in other formats.

Analysts said that without knowing the terms of the deal, it would be impossible to gauge its
financial impact on YouTube, but said that, in general, agreements with major content providers
would help YouTube over time.

``Wall Street is under the assumption that YouTube will continue to lose money,'' said Ross Sandler,
an analyst with RBC Capital Markets. Mr.~Sandler said that with deals like these, ``YouTube's
potential long-term upside goes up.''

YouTube overwhelmingly dominates the world of online video. In July, users in the United States
watched nearly nine billion clips on the site, more than 10 times as many clips as they watched on
the sites owned by its nearest competitor, Viacom, according to comScore.

But much of YouTube's audience visits the site to watch a random mix of clips generated by amateurs,
which advertisers view with trepidation\cite{trepidation}. As a result, YouTube has been on a long
quest to obtain more professionally produced video that it can use to generate revenue and offset
the enormous cost of streaming billions of free clips.

As it has sought to accommodate professional content providers, YouTube has shown itself to be more
willing to change. In April, for instance, it announced an agreement with Universal Music to create
a separate site, called Vevo, to showcase music videos and related content. Sony Music later joined
the venture, which has yet to open its site.

Vevo was widely seen as an attempt to emulate the success of Hulu, an online joint venture of NBC,
Fox and Disney that offers users free television shows and movies. While Hulu's audience is much
smaller than YouTube's, the site has been able to attract major advertisers.

Some of YouTube's efforts may be starting to pay off. In recent months, Google executives have said
that thanks to a host of new advertising efforts tied to professional clips, YouTube was inching
closer to profitability. But they declined to say when the site might cross that line.

Google has charged for video rentals and downloads in the past through Google Video. But less than a
year after acquiring YouTube in October 2006, it stopped offering the paid video service.

\section{Facebook Exodus\cite{Exodus}}

\lettrine{T}{hings} fall apart; the center cannot hold. Facebook, the online social grid, could not
command loyalty forever. If you ask around, as I did, you'll find quitters. One person shut down her
account because she disliked how nosy it made her. Another thought the scene had turned desperate. A
third feared stalkers. A fourth believed his privacy was compromised. A fifth disappeared without a
word.

The exodus is not evident from the site's overall numbers. According to comScore, Facebook attracted
87.7 million unique visitors in the United States in July. But while people are still joining
Facebook and compulsively visiting the site, a small but noticeable group are fleeing -- some of
them ostentatiously.

Leif Harmsen, once a Facebook user, now crusades against it. Having dismissed his mother's snap
judgment of the site (``Facebook is the devil''), Harmsen now passionately agrees. He says, not
entirely in jest, that he considers it a repressive regime akin to North Korea, and sells T-shirts
with the words ``Shut Your Facebook.'' What especially galls him is the commercialization and
corporate regulation of personal and social life. As Facebook endeavors to be the Web's headquarters
-- to compete with Google, in other words, and to make money from the information it gathers -- it's
inevitable that some people would come to view it as Big Brother.

``The more dependent we allow ourselves to become to something like Facebook -- and Facebook does
everything in its power to make you more dependent -- the more Facebook can and does abuse us,''
Harmsen explained by indignant e-mail. ``It is not 'your' Facebook profile. It is Facebook's profile
about you.''

The disillusionment with Facebook has come in waves. An early faction lost faith in 2008, when
Facebook's beloved Scrabble application, Scrabulous, was pulled amid copyright issues. It was
suddenly clear that Facebook was not just a social club but also an expanding force on the Web,
beholden to corporate interests. A later group, Harmsen's crowd, grew frustrated last winter when
Facebook seemed to claim perpetual ownership of users' contributions to the site. (Facebook later
adjusted its membership contract, but it continues to integrate advertising, intellectual property
and social life.) A third wave of dissenters appears to be bored with it, obscurely sore or just
somehow creeped out.

My friend Alex joined four years ago at the suggestion of ``the coolest guy on the planet,'' she
told me in an e-mail message. For a while, they cultivated a cool-planet online gang. But then
Scrabulous was shut down, someone told her she was too old for Facebook, her teenage stepson seemed
to be losing his life to it and she found the whole site crawling with mercenaries trying to sell
books and movies. ``If I am going to waste my time on the Internet,'' she concluded, ``it will be
playing in online backgammon tournaments.''

Another friend, who didn't want his name used, found that Facebook undermined his whole notion of
online friendship. ``It's easy to think of your circle of 'Friends' as a coherent circle, clear and
moated, when in fact the splay of overlap/network makes drip/action painting a better (visual)
analogy.'' Something happened to this drip painting that he won't discuss. He said, ``Postings that
seem private can scatter and slip unpredictably into a sort of semipublic status.''

That friend was not the only Facebook dissenter who was reticent about specifics. Many seem to have
just lost their appetite for it: they just stopped wanting to look at other people's photos and
r\'esum\'es and updates, or have their own subject to scrutiny. Some ex-users seemed shaken, even
heartbroken, by their breakups with Facebook. ``I primarily left Facebook because I was wasting so
much time on it,'' my friend Caroline Harting told me by e-mail. ``I felt fairly detached from my
Facebook buddies because I rarely directly contacted them.'' Instead, she felt as if she stalked
them, spending hours a day looking at their pages without actually saying hello.

But then came the truly weird part: ``Facebook was stalking me,'' Harting wrote. One day, on another
Web site, she responded to an invitation to rate a movie she saw. The next time she logged on to
Facebook, there was a message acknowledging that she had made the rating. ``I didn't appreciate
being monitored so closely,'' she wrote. She quit.

Julie Klam, a writer and prolific and eloquent Facebook updater, said in her own e-mail message, ``I
have noticed the exodus, and I kind of feel like it's kids getting tired of a new toy.'' Klam, who
still posts updates to Facebook but now prefers Twitter for professional networking, added,
``Facebook is good for finding people, but by now the novelty of that has worn off, and everyone's
been found.'' As of a few months ago, she told me, Facebook ``felt dead.''

Is Facebook doomed to someday become an online ghost town, run by zombie users who never update
their pages and packs of marketers picking at the corpses of social circles they once hoped to
exploit? Sad, if so. Though maybe fated, like the demise of a college clique.

Points of Entry: This Week's Recommendations

THE QUIT Put ``Why I Quit'' into Google, and the search engine proposes you look into both ``Why I
Quit Facebook'' and ``Why I Quit Church.'' If you aim to be a lapsed social networker, wikiHow, the
collaborative how-to guide, provides a useful step-by-step way to disengage, emotionally and
practically: wikihow.com/quit-facebook.

AN INQUIRY You're not the first to think it's creepy to have your personal life commercialized.
Jürgen Habermas has been especially eloquent about this. Start with ``The Theory of Communicative
Action.'' Copies are available on AbeBooks.com. Also interesting on this score: ``The Purchase of
Intimacy,'' by Viviana Zelizer.

GET BOARD ONLINE Scrabble is alive and well in cyberspace. If you like Scrabble, try lexulous.com.
For backgammon: ItsYourTurn.com.

\section{For Profit, Industry Seeks Cancer Drugs}

\lettrine{P}{fizer}'s fortunes in the past were built on cardiovascular drugs, like the cholesterol
buster Lipitor and the blood pressure pill Norvasc.

But the future of Pfizer, the world's largest pharmaceutical company, may rest in a cluster of
buildings on a bluff not far from the Pacific Ocean. It is here that Pfizer has amassed about 1,000
researchers for an all-out effort to develop drugs for cancer, a disease the company once largely
ignored.

Virtually every large pharmaceutical company seems to have discovered cancer, and a substantial
portion of the smaller biotechnology companies are focused on it as well. Together, the companies
are pouring billions of dollars into developing cancer drugs.

Two industry trends are driving the push. Recent scientific discoveries have suggested new targets
for cancer drug researchers to attack. And as drug companies see profits beginning to wane from
mainstays like Lipitor, the high prices that cancer drugs can command have become an irresistible
lure.

About 860 cancer drugs are being tested in clinical trials, according to the pharmaceutical
industry's main trade group. That is more than twice the number of experimental drugs for heart
disease and stroke combined, nearly twice as many as for AIDS and all other infectious diseases
combined, and nearly twice as many as for Alzheimer's and all other neurological diseases combined.

But for all the industry's spending and effort, only a trickle\cite{trickle} of new cancer drugs
make it to market. Last year there were two, and this year there has been only one.

And even some of those drugs offer only a few months at most of extra life or tumor stabilization
despite prices that often reach thousands of dollars a month. The drug Tarceva, which costs about
\$3,500 a month, was approved as a treatment for pancreatic cancer because it improved survival by
12 days.

The battle to treat cancer has become, as a commentary in a leading journal put it, a ``grinding war
of the trenches.''

Why? Experts say the same factors that attract drug companies to the cancer business help explain
the slow progress.

One reason is scientific. Studies are rapidly revealing the genetic changes in cells that cause
cancer and spur its growth. That is providing drug companies with dozens of molecules, or
``targets,'' that drugs could block.

But those same studies have shown that cancer is devilishly complicated. There are so many aberrant
molecules in a tumor that blocking just one or two is like trying to stop all traffic in Manhattan
with a roadblock at a single intersection.

Tumor cells, like bacteria, can develop resistance to drugs. Some experts believe that drugs that
kill most tumor cells do not affect cancer stem cells, which can regenerate the tumor.

And even two people with breast cancer, or two people with lung cancer, might have two very
different diseases on the molecular level, so a drug that works for one might not work for the
other.

``Cancer is not a single disease,'' said Robert A.~Weinberg, a cancer biologist at the Whitehead
Institute and the Massachusetts Institute of Technology. ``It's really dozens, arguably hundreds of
diseases.''

The other reason for the drug makers' interest is financial. Patients are often desperate, and
insurers risk outrage by denying payments for a cancer drug, even if the odds say it will have
little benefit. That has allowed pharmaceutical companies to charge thousands of dollars a month for
cancer medicines. Such prices can make drugs for even rare cancers, or drugs that do not work very
well, into big moneymakers.

Take Erbitux, developed by ImClone Systems, which costs \$10,000 a month. A study in Canada showed
that as a last-ditch treatment for colorectal cancer, Erbitux lengthened lives by an average of
about one and a half months compared with not treating the cancer at all. Using the price of the
drug in the United States and the average length of treatment, the extra cost per patient was about
\$50,000.

Erbitux, which is also approved to treat head and neck cancers, recorded global sales of \$1.6
billion last year, higher than all but about 70 other drugs. Last year, as part of the industry
scramble into cancer drugs, Eli Lilly \& Company outbid Bristol-Myers Squibb to acquire ImClone for
\$6.5 billion.

In 1998, there were only 12 cancer drugs on the list of the world's 200 medicines with the highest
sales, compiled by the trade magazine Med Ad News. Taxol, No.~21, was the only cancer drug among the
30 drugs with sales of at least \$1 billion.

The same list for last year contained 23 cancer drugs among the top 200 -- and three in the top 10.
Of the 126 drugs with \$1 billion in sales, 20 were for cancer.

Cancer drugs have been the biggest category of drugs in terms of sales worldwide since 2006 and in
the United States since 2008, according to the market researcher IMS Health.

Such money attracts companies. ``Cancer is such an emotional issue that the free market doesn't work
like it does for bicycle wheels and umbrellas,'' said Robert L.~Erwin, a biotechnology industry
executive who heads the Marti Nelson Cancer Foundation, a patient advocacy group. ``As long as the
health care system will pay the price, the money will flow in that direction.''

But Mr.~Erwin and some other experts say that is not always a good thing for patients because it can
set the bar too low for drug companies.

``As long as the marketplace does not distinguish between modestly effective drugs and dramatically
effective drugs, there won't be an incentive to shift resources to a greater emphasis on a larger
benefit,'' said Dr.~Neal J.~Meropol, an oncologist at the Fox Chase Cancer Center in Philadelphia
who has been studying drug prices.

Many executives dispute this, saying they would produce drugs offering bigger gains if they knew
how. But they must balance their portfolio of experimental drugs between long shots and some drugs
that have a better chance of making it to market and sustaining the enterprise.

``If you always swing for home runs, you strike out a lot,'' said George A.~Scangos, chief executive
of Exelixis, a biotechnology company with 11 cancer drugs in clinical trials. ``It's not the
companies' profit motives,'' he said. ``It's largely the difficulty of hitting home runs.''

With health care costs rising, there is new pressure on companies to be more selective in drugs they
develop. Some experts now talk about ``financial toxicity'' as a side effect of cancer drug
treatment, along with nausea and hair loss.

``A question is how the system can tolerate 400 new drugs on the market, all at the same price'' of
\$50,000 a year, said Dr.~Lee Newcomer, senior vice president for oncology at United Healthcare, a
big insurer.

Such cost pressures, and the fact that only a handful of cancer drugs get to market each year, mean
the big investments now being made into cancer drugs are likely to turn sour for many companies.

``It's the biggest bubble you've ever seen,'' said Dr.~Mark Ratain, an oncologist at the University
of Chicago.

But Pfizer is counting on cancer to help save the company. It hopes to reach \$11 billion in sales
of cancer drugs by 2018. That would be more than four times the category's sales last year of \$2.5
billion, which represented only 5 percent of Pfizer's revenue.

Cancer was once unattractive for big pharmaceutical companies like Pfizer. There were relatively few
patients with any one type of cancer, and they died fairly quickly. By contrast, there were millions
of patients with chronic diseases like hypertension who would take drugs for life.

Indeed, the three main cancer drugs Pfizer now sells came to it with its 2003 acquisition of a
rival, Pharmacia, a deal done mainly to acquire the arthritis drug Celebrex.

But there are now many good cardiovascular drugs. Lipitor, the world's best-selling drug, will lose
patent protection in 2011, and Pfizer failed to develop a successor.

So Pfizer is scaling back cardiovascular research and has made cancer drugs one of its six focus
areas. About 20 percent of Pfizer's more than \$7 billion in research and development spending is on
cancer, and 22 of the roughly 100 drugs in clinical trials are cancer drugs.

``I've taken a lot of personal interest in this business unit,'' said Jeffrey B.~Kindler, Pfizer's
chief executive. ``We think we are positioned to be a top leader in oncology.''

Cancer research is concentrated here in San Diego, in a cluster of buildings once owned by one of
the many biotechnology companies in the region. Pfizer has tried to retain some of the looser
culture of entrepreneurial start-ups, like Friday afternoon beer parties. The head of the site,
Catherine Mackey, a transplant from Pfizer's laboratories in Connecticut, has become an avid
early-morning surfer.

The clinical trials for cancer are being overseen by Dr.~Mace L.~Rothenberg, an oncologist recruited
this year from Vanderbilt University. He hopes Pfizer can develop those home-run drugs. ``Having
treated patients for 20 years,'' Dr.~Rothenberg said, ``I know their needs are not for singles.''

The big thrust in cancer drug development for the last few years has been so-called targeted
therapies. These drugs aim, so far with modest success, to block aberrant molecules in tumor cells
while leaving normal cells unscathed.

But even most targeted therapies have limited impact. One reason is that most tumors are fueled by
numerous, often redundant, genetic anomalies. That means that drugs with different targets need to
be used in combination. But combinations increase both the costs and side effects of therapy. And it
is difficult to test two experimental drugs in combination because the regulatory system is geared
to assessing a single drug at a time.

Another reason is that tumors differ among people. Dr.~Bert Vogelstein, a cancer geneticist at Johns
Hopkins, said a typical tumor might have 50 to 100 genetic mutations. But two patients with the same
type of cancer might have only five mutations in common.

So even though a drug might work well for patients whose tumors have a particular mutation, when the
drug is used for a broader population, it shows only a small effect.

One solution is to try to determine which patients should get which drug based on the genetic
profile of the tumor.

Pfizer is moving in that direction. It plans soon to start a late-stage clinical trial of a drug for
lung cancer. But the only patients in the trial will be from the 5 percent or so of lung cancer
patients with a mutation in a gene called A.L.K.

``What we're looking for,'' said Dr.~Rothenberg of Pfizer, ``is not a small benefit in a large
group, but a larger benefit in a smaller group.''

For now at least, making the ballpark smaller may be the industry's best chance to hit home runs.

\section{At ABC, an Anchor Shift; for TV, an Image Shift}

\lettrine{O}{n} Tuesday, Diane Sawyer was on stage with Whitney Houston before a crowd in Central
Park. On Wednesday, she was crowned the leader of ABC's network news division.

Ms.~Sawyer, the longtime -- some would say long-suffering -- co-host of ABC's ``Good Morning
America,'' was named the successor to Charles Gibson, who is stepping down as the anchor of ``World
News'' on ABC. The network announced the moves on Wednesday; they will be effective in January.

The arrival of Ms.~Sawyer will comprehensively alter the long-established image of an avuncular male
nightly news anchor. With Katie Couric, who took the CBS anchor position in 2006, two of the three
main network news voices will be female, a role that in the past has punished others, like Barbara
Walters and Connie Chung.

``You're going to have, for the first time ever, two women competing as solo anchors in a television
framework that just -- within living memory -- sort of destroyed every woman who tried to do it,''
said Richard Wald, a former news executive at ABC and NBC.

For Ms.~Sawyer, the promotion is the culmination of years of waiting. She took the morning job on a
temporary basis, then stayed for 10 years, reviving ABC's morning franchise. But she was passed over
twice for the evening anchor job.

Her wait came to an abrupt and unexpected end when Mr.~Gibson announced his intention to step down.
David Westin, the president of ABC News, characterized the decision as entirely Mr.~Gibson's, after
the two men held conversations that extended through the summer.

Mr.~Westin said in an interview that if he had had his way, none of this would be happening. He said
he had tried to talk Mr.~Gibson out of the decision and to give him some time to think it over.

``This was not a result I wanted,'' Mr.~Westin said.

Mr.~Westin said he finally assented last week and asked Ms.~Sawyer if she would accept the position
last Thursday. She agreed Tuesday night, he said.

Other senior ABC News executives, including Jon Banner, the executive producer of ``World News,''
said that Mr.~Gibson, who is 66, had indeed initiated the move out of a desire to pull back from
daily television work.

The decision was sudden enough that ABC does not have a plan in place to replace Ms.~Sawyer, who is
63, on ``Good Morning America,'' which is by far the most profitable program in the news division
and where she is the biggest attraction for viewers. One executive estimated that it brings in about
\$50 million a year for ABC.

But ``Good Morning America'' has been second in its time slot, behind ``Today'' on NBC, for more
than a decade, and it has fallen further behind in recent months. All along, Ms.~Sawyer has remained
its strongest on-camera presence. Her departure leaves a serious gap, as Mr.~Westin acknowledged.

In the face of 24-hour news on cable and on the Web, the ratings for all three evening newscasts
continue to erode. The combined newscasts now struggle to reach 20 million viewers some nights.
During the last week of August, the ``NBC Nightly News With Brian Williams'' averaged 7.9 million
viewers, about 800,000 more than Mr.~Gibson's ``World News.'' Ms.~Couric's ``CBS Evening News''
remained a distant third with 5.4 million viewers.

Neither Mr.~Gibson nor Ms.~Sawyer would give interviews on Wednesday about the moves, an unusual
decision. One ABC News executive said it reflected how surprised many were in the news division to
hear of the development. Both issued statements instead.

But ABC News executives said the choice was inevitable, both because Ms.~Sawyer had expressed
interest in succeeding Mr.~Gibson whenever he stepped down and because she was the most prominent
name in the news division.

``Diane was the obvious choice,'' Mr.~Westin said. ``She always said this was Charlie's time, but
yes, she always expressed interest in the job.'' Ms.~Sawyer only recently signed a new contract with
ABC.

Both Mr.~Westin and Mr.~Banner used the same expression, ``the DNA of the newscast will not
change,'' in saying Ms.~Sawyer would make the newscast her own without changing it in any
fundamental ways. But they also cited her eagerness to chase stories all over the globe: she has
traveled to Afghanistan and North Korea on assignment.

Ms.~Sawyer started her career in journalism in her home state, Kentucky, but her big break came
through politics when in 1970 she was offered a job at the White House by Ron Ziegler, press
secretary to President Richard Nixon. She joined CBS in 1978 as a correspondent, became co-anchor of
the morning news and later a correspondent on ``60 Minutes.'' She moved to ABC in 1989, co-anchoring
the news programs ``Primetime Live'' and ``20/20.''

Despite her experience in hard news and interviews with world leaders, Ms.~Sawyer, as Ms.~Couric
did, comes to the evening anchor job better known for soft news and celebrity interviews, like the
one she conducted with Michael Jackson and Lisa Marie Presley, then his wife, in 1995. She inquired
about their sex life with the preface, ``I didn't spend my life as a serious journalist to ask these
kinds of questions.''

Ms.~Sawyer's presence in the evening news competition will surely revive stories of personal
competition between her and Ms.~Couric. They previously faced off as the leading figures on the
morning programs.

``We'll find out whether people don't want to watch a woman anchor the news or whether they just
didn't want to watch Katie,'' said one veteran television news industry executive, asking for
anonymity to speak candidly about the competition. Ms.~Couric's CBS broadcast is consistently in
third place.

Others hailed the move as positive for the business and for American society. ``I think it says a
lot about where we've come as a nation,'' Mr.~Westin said.

Ms.~Walters, the ABC host who was the first woman to co-anchor a network newscast, in 1976, said in
an interview on Wednesday, ``I think it's just a great day.''

Mr.~Gibson, who had his own long run as a co-host, with Ms.~Sawyer, of ``Good Morning America,''
took over the anchor job after a series of unexpected events -- the sudden death of the longtime
anchor Peter Jennings and the serious injuries suffered in Iraq by Bob Woodruff, one of the
co-anchors Mr.~Westin had chosen to replace Mr.~Jennings.

Mr.~Gibson was almost immediately successful. He even ascended briefly to first place in the evening
news ratings in spring 2007 and has since been somewhat competitive with the leader, Brian Williams
at NBC.

But he lost the most recent contest: on his Wednesday night newscast, Mr.~Williams reported the
switch at ABC, assuring viewers that he would keep them posted on the competition. Mr.~Gibson,
however, did not mention it.

\section{Kennedy's Rough Waters and Still Harbors}

\lettrine{A}{t} the end of his deeply affecting memoir, the late Senator Edward M.~Kennedy writes
about his grandson ``Little Teddy'' -- the son of his son ``Medium Teddy'' who delivered such a
heartbreaking eulogy\cite{eulogy} at the senator's funeral on Saturday -- and his difficulties
mastering the family tradition of sailing. The senator told the 10-year-old ``we might not be the
best,'' but ``we can work harder than anyone,'' and Little Teddy stayed with it, grew eager to learn
and started winning races. That, the senator writes, ``is the greatest lesson anyone can learn'':
that if you ``stick with it,'' that if, as the title of his book suggests, you keep a ``true
compass'' and do your best, you will eventually ``get there.''

And that, in a sense, is the theme of this heartfelt autobiography: that persistence, perseverance
and patience in pursuit of a cause or atonement\cite{atonement} for one's failures can lead to
achievement and the possibility of redemption. It's the story of how this youngest and most
underestimated of siblings slowly, painfully, incrementally found genuine purpose of his own in
shouldering the weighty burden of familial expectations and the duty of carrying on his slain
brothers' work. He found a purpose, not as they did in the high-altitude pursuit of the presidency
but in the dogged, daily grind of being a senator -- of laboring over bills, of sitting through
endless committee meetings, of wading through briefing books and making deals with members across
the aisle. The resulting legislation -- including the Civil Rights Act of 1964 and the State
Children's Health Insurance Program of 1997 -- would help the poor and the disenfranchised and those
with disabilities, and win him recognition as one of the foremost legislators in American history.

Mr.~Kennedy is not a particularly introspective writer -- he acknowledges in these pages that he
coped with the assassinations of his brothers Jack and Bobby by pushing his grief down, by trying to
keep moving forward so as to stay ahead of the darkness and not to be engulfed by despair. But he
writes in these pages with searching candor\cite{candor} about the losses, joys and lapses of his
life; the love and closeness of his family; the solace he found in sailing and the sea; his complex
relationships with political allies and rivals. Mr.~Kennedy's conversational gifts as a storyteller
and his sense of humor -- so often remarked on by colleagues and friends -- shine through here, as
does his old-school sense of public service and his hard-won knowledge, in his son Teddy Jr.'s
words, that ``even our most profound losses are survivable.''

In these pages (Ron Powers is credited as a collaborator) Mr.~Kennedy draws some telling portraits
of other politicians. Of Jimmy Carter, he writes, ``He baffled\cite{baffle} many potential allies in
his own party,'' but ``I believed then and now that he reserved a special place in his
animus\cite{animus} toward me.'' He writes that his objections to Ronald Reagan's policies are ``far
too vast to enumerate'' but that he admired the optimism Reagan brought to the country after the
Carter era. More revealingly, Mr.~Kennedy says that he is convinced that had his brother Jack lived,
he would have sought a way out of Vietnam (``He had spoken with McNamara,'' referring to Defense
Secretary Robert S.~McNamara, ``about a plan for withdrawal within two or three years,'' he writes)
and that he, Teddy, is ``satisfied that the Warren Commission got it right.''

But it is Mr.~Kennedy's personal rather than political reminiscences that are most memorable in
``True Compass.'' He was a talented amateur painter, and there is a vivid, almost pictorial
evocation of his privileged but pressured and sometimes lonely childhood and youth. He sometimes
felt, he says, that his life was ``a constant state of catching up'' to his glamorous,
larger-than-life brothers, whom he hero-worshipped as a boy: Joe Jr., who died young in war; Jack,
who Teddy believed would always win even when the odds were against him, who ``could do anything he
wanted''; and Bobby, who was not ``cold, calculating'' as some of his critics charged, but who
``lived and made decisions in the moment,'' completely absorbed in whatever he was doing.

Teddy is always the one who through sheer\cite{sheer} will and fortitude -- whether passing a piece
of legislation, finishing a perilous mountain climb or gritting his teeth through the pain of kidney
stones to deliver a speech -- keeps on keeping on, telling himself ``I can handle this,'' ``I can
handle this.''

In trying to decide whether to run for the Senate in 1962, Mr.~Kennedy writes, he remembered his
father's words to him as a boy: ``You can have a serious life or a nonserious life, Teddy. I'll
still love you whichever choice you make. But if you decide to have a nonserious life, I won't have
much time for you. You make up your own mind. There are too many children here who are doing things
that are interesting for me to do much with you.''

Later in this volume, Mr.~Kennedy addresses his own failings and regrets. He writes about how his
actions in 1969 at Chappaquiddick were ``inexcusable,'' how Mary Jo Kopechne's death ``haunts me
every day of my life'' and how ``atonement is a process that never ends.'' When his father died four
months later, he says, he ``wondered whether I had shortened'' his life ``from the shock I had
visited on him with my news of the tragic accident on Chappaquiddick Island. The pain of that burden
was almost unbearable.''

In another chapter, Mr.~Kennedy observes that ``with all of the background noise about Palm Beach''
-- surrounding the rape trial of his nephew William K.~Smith (who was found not guilty) -- ``and my
bachelor lifestyle, I would have been the wrong person to lead the questioning'' of Clarence Thomas
during the Anita Hill portion of his confirmation hearings and that ``many people were disappointed
that I was unable to succeed in making a persuasive case against Thomas's confirmation.''

After his brother Robert's assassination, Mr.~Kennedy recalls how his anguish led him to drive his
car at high speeds, to sometimes drive his ``capacity for liquor to the limit'' and how years
earlier, Bobby's grief over their brother Jack's assassination ``veered close to being a tragedy
within the tragedy,'' with their mother, Rose, and Bobby's wife, Ethel, fearing for ``his psychic
survival.'' His brother Bobby's ``blossoming idealism'' -- about Vietnam, about taking on the fight
against poverty and urban violence -- was in fact, he suggests, ``provoked by Jack's death.''

The murders of Jack and Bobby not only devastated Teddy and left him with an abiding sorrow and
loneliness, but they also, on a subconscious level at least, made him fearful for his own life. He
writes that he flinched at 21-gun salutes at Arlington to honor the fallen in Iraq, once dived for
the pavement when a car backfired in the street. In 1982, he says, his children's hope that he would
not run for the presidency and their unspoken fears for his safety were crucial in his decision not
to enter the race.

The tribulations of the Kennedy family have frequently been likened to something out of Shakespeare
or a Greek tragedy, but Teddy Kennedy manages the difficult task in these pages of conveying the
profoundly ordinary, human dimensions of his and his family's losses, the day-to-day reality of
losing two siblings (Joe Jr.~and Kathleen) when he was still a youth; of becoming a paterfamilias at
the age of 37 for his nephews and nieces after Jack's and Bobby's deaths; of shepherding two of his
children through the ordeal of cancer; of coming to a point in his life after all this sorrow and
struggle of having ``stopped looking forward to things,'' of retreating from the risk of ``new
personal commitments.''

In the end, however, Mr.~Kennedy was able to write a happy ending to his own life. He fell in love
with and in 1992 married Victoria Reggie, whose ``acute understanding and love of me'' gave him a
new sense of stability and tranquillity. He found renewed meaning in his work in the Senate. And he
found, in Barack Obama, whom he helped elect, a new incarnation of the idealism and sense of public
service he and his brothers had embraced as young men so many years ago.

In his last months, Mr.~Kennedy says, he found that ``simple pleasures fill me with happiness,''
that gazing out at the sea and his beloved boat Mya left him with a sense of peace. Sailing, he
writes, always helped ``displace the emptiness inside me with the awareness of direction'': ``an
awareness that there is a beginning to the voyage\cite{voyage} and an end to the voyage, and that
this beginning and ending is part of the natural order of things.''

\section{Politics Permeates\cite{permeate} Anti-Corruption Drive in China}

\lettrine{S}{hortly} after Huang Guangyu was named the second richest man in China by Forbes last
October, the 39-year-old entrepreneur disappeared.

Nearly a week later, Beijing's Public Security Bureau issued a brief notice saying he had been
detained on corruption charges, the first indication of his whereabouts.

Soon, other prominent individuals were arrested or charged with corruption or bribery: Rixin Kang,
the former head of China's nuclear power agency; Chen Tonghai, the former chairman of Sinopec, the
state-owned oil company; the head of Beijing's Capital Airport (who was executed last month); and
the former mayor of Shenzhen, one of the country's biggest manufacturing centers.

Chinese authorities say the arrests are part of the Communist Party's latest anticorruption campaign
-- and they include the arrest last month of four employees of the British-Australian mining giant,
Rio Tinto, on bribery charges.

But analysts say that prominent corruption cases in China are often the outgrowth of power struggles
within the Communist Party, with competing factions using the ``war on corruption'' as a tool to
eliminate or weaken rivals and their corporate supporters.

In the case of Mr.~Huang, the electronics billionaire, for example, state-run media say a number of
other high-ranking officials with longstanding ties to him have also been dismissed and arrested in
what looks to be a Communist Party power shuffle.

China does not have an independent police or judicial system; party leaders order investigations.
``It's a very politicized process,'' says Minxin Pei, a professor at Claremont McKenna College in
California and author of numerous studies on corruption in China. ``If your patrons do not protect
you, you're toast.'' This may help explain one of the enduring contradictions of China's political
and economic system: the government regularly publicizes an astounding number of corruption cases,
yet little progress seems to be made in uprooting corruption.

But as China prepares to celebrate the 60th anniversary of the Communist victory here, Beijing is
growing increasingly worried that rampant\cite{rampant} corruption is eroding\cite{erode} public
trust in the party, hurting the economy and threatening social stability.

In a recent government poll, 75 percent of those surveyed listed corruption as their No.~1 concern.
And this year, public outrage over corruption has caused mass demonstrations -- like one in June,
when thousands of people in the eastern city of Shishou rioted after a cook died mysteriously in a
hotel believed to be controlled by corrupt officials.

``The public is fed up with corruption,'' says Gao Quanxin, a senior fellow at the Chinese Academy
of Social Sciences in Beijing.

China's top leaders, including President Hu Jintao, have called on the Communist Party to strengthen
its fight against graft\cite{graft}.

Experts say one of Beijing's objectives in the current crackdown is to send a message to officials
who overstep some unknown acceptable level of graft or who too ostentatiously flaunt its rewards
and, second, to reassure an increasingly angry public.

The problem is apparently so large that even selective enforcement results in about 150,000
officials being punished every year for bribery, corruption and other offenses.

China is paying a hefty price for corruption, analysts say, through the misallocation of resources,
and the health and safety threats that can arise when regulators, for instance, are paid to look the
other way.

A 2007 study by the Carnegie Endowment of International Peace estimated that in 2003, corruption
cost China about \$86 billion, or about 3 percent of gross domestic product at the time.

``Corruption has not derailed China's economic rise,'' says Professor Gao at the Chinese Academy.
``But it's rotting the establishment of a rule of law. The Chinese government has more than 1,200
laws, rules and directives against corruption, but implementation is ineffective.''

Experts say corruption is thriving here because relatively low-paid government officials wield
enormous power over business and resources. ``The key variable is the extent to which the government
gets involved in business in China,'' says Professor Pei at Claremont McKenna.

Many officials come up with ever more imaginative ways to gain wealth, like setting up private
foundations to accept bribes or encouraging entrepreneurs and companies to bankroll the education
overseas of family members.

``It's no longer bags of cash,'' says Violet Ho, an investigator working for Kroll, the risk
consulting firm. ``There are sophisticated games going on.''

Chinese lawmakers are considering a law that would give investigators larger powers, including
criminal penalties on the ``bribe-taking relatives and spouses'' of corrupt officials.

No one knows how many Western companies participate in corruption here or how often, but many
multinational corporations are worried, particularly American companies that must comply with the
Foreign Corrupt Practices Act, which forbids bribing foreign officials.

This year, Morgan Stanley dismissed a real estate executive in Shanghai after discovering he might
have violated the bribery law. And in 2007, Lucent Technologies agreed to pay a \$2.5 million fine
after the United States government accused the company of violating the act.

In a filing, the Justice Department said Lucent spent more than \$10 million from 2000 to 2003 to
bring 1,000 officials to the United States. The officials were supposed to visit Lucent operations,
but instead Lucent often paid for them to travel to Las Vegas, Disneyland and the Grand Canyon.

Still, when major corruption cases occur, the origins and nature of the cases are often shrouded in
mystery.

In 2006, when Chen Liangyu, Shanghai's party secretary and a Politburo member, was ousted, many
experts said there was evidence that he -- a major regional figure -- had lost a power struggle with
leaders in Beijing.

Many China experts believe the Rio Tinto case was politically motivated because Rio Tinto battled
with the government steel industry over iron ore prices.

Those caught in the party's campaigns are usually humiliated and denounced for taking bribes,
leading ``decadent lifestyles'' and, sometimes, for taking multiple mistresses.

But so many officials are engaging in the same practices that in Hong Kong there is a cottage
industry in books that detail the private lives of corrupt officials, with titles like, ``The Sexy
Record of High-Ranking Communist Party Officials,'' ``Power and Money: How They Steal,'' and
``Shared Mistresses: Lust and Caution Among Chinese Officials.''

Whatever the public outrage, many experts say they doubt corruption will go away anytime soon,
largely because China does not have an independent police force or judicial system.

There are limits to where the party is willing to go. When President Hu's son became embroiled in a
corruption investigation in the southern African nation Namibia involving a Chinese company he once
controlled, Beijing asked Chinese Internet portals and state-run outlets not to make any mention of
the case.

Analysts say there is another reason change will come slowly: because officials seek power precisely
so that they can form alliances and make the decisions that will eventually make them rich.

``It's now a recruiting tool,'' Professor Pei says of the promise of future graft. ``Corruption is
the glue that keeps the party stuck together. Getting rid of it is not possible as long as they keep
this system.''

\section{China Fails to Prevent Myanmar's Ethnic Clashes}

\lettrine{C}{hina} is Myanmar's closest ally -- almost its only one. It is Myanmar's chief defender
in international forums, its major weapons supplier, its largest foreign investor and a crucial
backer of its ruling military junta.

But in the wake of a recent clash between Myanmar's army and ethnic rebels, a rout\cite{rout} that
sent thousands of people streaming over the mountainous border into China, analysts have begun to
question how much influence China has.

The answer may determine whether that brief battle grows into a much bigger and deadlier war.

Although it tried, analysts, journalists and other experts say, China was unable to dissuade
Myanmar's junta last month from sending thousands of troops into the nation's northern Kokang
region, where they easily routed about 1,500 armed rebels. The rebels had observed a cease-fire with
Myanmar's government for nearly 20 years.

Now news reports say that the junta has sent 7,000 troops and 20 tanks into a neighboring region
known as Wa State, where a much larger rebel force, the United Wa State Army, has been observing the
same cease-fire. The Wa forces, at least 20,000 strong and heavily armed, promise a fight if
attacked.

``We want peace, but we are not going to lay down our weapons and surrender,'' a rebel spokesman who
called himself Su said by telephone on Thursday. ``We will not become the second Kokang.''

The recent fighting is the result of demands by the junta that the rebels disarm and join a
government-run border patrol. The ultimatum is widely seen as an effort by the military rulers to
defang their opponents in advance of an election next year that they are billing as the first
democratic vote in more than 20 years.

China's Foreign Ministry spokesman said Tuesday that Myanmar, long called Burma, had ``promised to
restore peace and stability along the border,'' and some local news reports suggest that the
confrontation in Wa State may yet be defused. But there are also signs that China and Myanmar, so
close for two decades, are having differences.

``I've spoken to Chinese Foreign Ministry people, and they're very concerned about this hostile
attitude Burma has,'' said Aung Zaw, a Burmese exile who is editor of The Irrawaddy, a magazine
based in Thailand. ``China has given them political and diplomatic support. But when Burma wants to
put a stop to its own internal matters, they don't care about anybody else.''

That view is echoed by a number of Beijing-based political analysts and scholars, some of whom have
worried publicly that China may not have enough clout to ward off a larger war that could send many
more refugees pouring into China. ``They don't always heed China's advice,'' said Shi Yinhong, a
professor of international relations at Renmin University in Beijing. ``China has so little
leverage\cite{leverage} against them because China, in some sense, depends on them.''

For decades, Myanmar has traded access to its ample\cite{ample} natural resources and to the Indian
Ocean for political support from China. This month, Chinese companies are set to start construction
of a \$2.5 billion oil-and-natural-gas pipeline project that would run from the ocean to Kunming,
the capital of China's neighboring province, Yunnan.

But China's relations with Myanmar are not so straightforward. In an earlier era, the Chinese gave
money and arms to ethnic groups, including the Wa and Kokang rebels, on Myanmar's side of the border
that were allied with the Burmese Communist Party. Factionalism sank the party in the 1980s, and in
1989 Myanmar's government struck a cease-fire with the ethnic groups that has lasted until now.

While staying close to the military junta, the Chinese have also kept in close touch with the ethnic
groups, leaving the border open for trade, family visits and no small amount of smuggling of arms
and other contraband.

As early as June, said one Beijing analyst, the Chinese government told a ranking official in
Myanmar's government during a visit to Beijing that it wanted to avoid conflict on the border. The
warning was repeated in July, when a team of officials from Myanmar traveled to Kunming to meet with
Yunnan provincial leaders, according to the analyst, who refused to be identified for fear of
retaliation from the Chinese government.

But some experts say they believe that the junta was nettled\cite{nettle} by China's
tacit\cite{tacit} support for the ethnic groups, many of whom are ethnic Chinese, and its refusal to
close the border and cut off the groups' economic lifeline. Other experts say the junta's internal
political calculations, geared toward a sweeping victory in the 2010 elections, trumped\cite{trump}
any diplomatic concerns.

When battles broke out in Kokang, the Chinese government reacted with unusual force, issuing a
statement asking the junta to restore ``regional stability.'' Unconfirmed new reports suggest that a
senior Chinese military official traveled to Myanmar early this week to assess the situation, and
that a delegation of five officials from Myanmar traveled Monday to Kunming to meet with
unidentified Chinese officials.

It is unclear whether China's forcefulness will cause the junta to stand down. But some believe
Myanmar's government sent its own signal in this week's edition of The Myanmar Times, a newspaper
that, like all the press there, reflects the leaders' thinking. The paper carried an article
recording a visit to Taiwan by the Dalai Lama, the Tibetan spiritual leader whom Beijing accuses of
trying to foment rebellion against the Chinese government.

The Asia Times, which first reported on the article, said it was the first time in at least 20 years
that the Dalai Lama's name had appeared in Myanmar's press.

\section{Massachusetts Waits for a Kennedy to Decide on Race}

\lettrine{N}{early} everyone in Massachusetts is waiting for Joseph P.~Kennedy II to make up his
mind.

It is hard to imagine this state without a Kennedy in the United States Senate. But it seems that
Mr.~Kennedy, 56, an entrepreneur and a former congressman who has avoided politics for more than a
decade, is the only family member seriously mulling a run for the seat his uncle, Edward M.~Kennedy,
held for 47 years.

For Massachusetts, with its top hospitals, universities and research centers that counted on the
federal dollars that flowed from Senator Kennedy's influence on Capitol Hill, the stakes extend far
beyond its deep emotional connection with him. People here are starting to grapple with the big
political question: how to replace their irreplaceable senator, whose decades of relationships in
Washington made Massachusetts a prime beneficiary of the things he believed in.

``With Joe, a lot of those ties would come back quickly,'' said David Gergen, a professor at the
Kennedy School of Government at Harvard who has advised both Democratic and Republican presidents.
``There would be an army of people who worked for his uncle who would want to work for him.''

Several people close to the family said pressure on Joe Kennedy was not coming from the family,
still grieving over the loss of its patriarch, but from within Joe himself, and, they speculated,
from various stakeholders in Massachusetts -- mayors, labor leaders and others -- who view him as
the candidate most likely to inherit enough of his uncle's monumental legacy to take care of their
needs.

Friends said that Joe Kennedy, who is the son of Robert F.~Kennedy and who has the trademark Kennedy
toothy smile, call to public service and complicated past, is agonizing\cite{agonizing} over whether
he truly wants the job, and depending almost exclusively on his wife, Beth, for counsel.

``It's not about ensuring that someone named Kennedy stays in the seat,'' said one friend of the
Kennedy family, who requested anonymity, citing the private nature of the discussions. ``It's going
to be a very personal decision based on whether or not this fits into Joe's life right now.''

And until he makes a move, political observers say, no other possible contender -- with the
exception of Martha Coakley, the state attorney general who announced Thursday -- is going to
declare for the primary, scheduled Dec.~8.

``This is almost a unique situation in which Joe Kennedy has the right of first refusal to the
seat,'' Mr.~Gergen said. ``He'll be a powerful favorite not only for the nomination but for the
election because of who he is -- the family name.''

Mr.~Gergen and others said Joe Kennedy had shown no sign of national political ambition until last
week, when he delivered a stirring tribute to his uncle at a memorial service here.

``Nobody can tell what's in Joe's head right now,'' said Robert Healy, a former executive editor of
The Boston Globe who was the newspaper's Washington bureau chief during John F.~Kennedy's
presidency. ``But I have to think that he has to have been affected -- as much as the people who are
for him and against him -- about the turnout for Teddy last week, the power of it.''

The bond between the Kennedys and Massachusetts is stronger than any other between a family and a
state in recent history, and Joe Kennedy has already played a role in perpetuating\cite{perpetuate}
it. He was a member of the United States House of Representatives from 1986 to 1998, occupying the
same seat that his Uncle Jack won in 1946.

He was a popular congressman, focusing on low-cost housing and veterans' health care and easily
winning re-election. But after his first wife publicly fought his attempt to annul their marriage
and his brother Michael was killed in a skiing accident in 1997, Mr.~Kennedy decided to leave
office.

For the previous two years, he had suffered from the fallout of charges of improper sexual relations
between Michael and a teenage baby sitter. Michael had been his campaign manager and had run his
company.

``There are times when you have to get your priorities right,'' Mr.~Kennedy said in 1998 when he
announced he would not seek re-election.

Political observers here say that while Mr.~Kennedy has plenty of charisma and political
savvy\cite{savvy}, he has not yet demonstrated his uncle's patience for the tedious, incremental
process of legislating.

Edward Kennedy apparently agreed: according to The Boston Herald, the senator told editors and
reporters there in 2005 that Joe was ``made intuitively and instinctively to be an executive, much
more than the legislative.''

Since leaving Congress, Joe Kennedy has run the Citizens Energy Corporation, a nonprofit company he
founded in 1979 to provide low-cost heating oil to the poor. The corporation has grown to include
seven subsidiaries\cite{subsidiary}, some of them for-profit; Mr.~Kennedy earned a salary of
\$544,792 in 2007, federal tax filings show.

He has won legions\cite{legion} of working-class admirers for his work distributing free oil, a
service he promotes in television advertisements urging viewers to call 1-877-JOE-4-OIL. The
Citizens Energy Web site features wintertime photos of Mr.~Kennedy embracing bundled-up elderly and
poor people, and the company's motto -- ``No one should be left out in the cold'' -- suggests the
kind of populist touch that made Edward Kennedy so popular here.

Joe Kennedy has been ambivalent\cite{ambivalent} about whether to enter races before, like in 2002,
when he signaled for months that he might run for governor but ultimately opted not to do so. Now,
as then, friends say, he is weighing whether to expose his wife and twin sons, now 28, to the
rigors\cite{rigor} of a campaign and the questions and attacks that surely would come with it.

An early dose came Wednesday, when Scot Lehigh, a columnist for The Boston Globe, called Mr.~Kennedy
a ``knucklehead'' prone to tantrums\cite{tantrum} and impetuousness.

One of Mr.~Kennedy's friends said: ``He's been very content. So there's some feeling of, why mess
around with that?''

\section{Obama Aides Aim to Simplify and Scale Back Health Bills}

\lettrine{P}{resident} Obama plans to address a joint session of Congress next week in an effort to
rally support for health care legislation as White House officials look for ways to simplify and
scale back the major Democratic bills, lower the cost and drop contentious but nonessential
elements.

Administration officials said Wednesday that Mr.~Obama would be more specific than he has been to
date about what he wants included in the plan. Doing so amounts to an acknowledgment that the
president's prior tactic of laying out broad principles and leaving Congress to fill in the details
was no longer working and that Mr.~Obama needed to become more personally involved in shaping the
outcome.

But the officials said Mr.~Obama was unlikely to unveil a detailed legislative plan of his own. And
they insisted that Mr.~Obama had not given up on the provision that has attracted the most fire from
the right, a proposal for a government-run competitor to private insurers, although many Democrats
say the proposal may eventually be jettisoned.

Rahm Emanuel, the White House chief of staff, said Mr.~Obama would be ``more prescriptive than he
has been to date.'' And he added, ``We have a tremendous amount of consensus in Congress to build
off of.''

In his address, Mr.~Obama is expected to emphasize areas of potential agreement. One is the need for
federal regulation of health insurance companies to prohibit them from denying coverage, or charging
higher premiums, because of a person's medical history or current condition. Another is the need for
federal subsidies to make insurance affordable to millions of lower-income people.

By signaling that they would seek to revise existing versions of legislation moving through the
House and Senate, administration officials and Democratic leaders in Congress -- many of whom had
said earlier in the summer that they saw no need to scale back their ambitions -- made clear that
their political calculations had changed. With Congressional Republicans standing almost unanimously
in opposition to the Democratic approach, the target now for Mr.~Obama is primarily a handful of
moderate Democrats and the one Republican who seems open to a deal, Senator Olympia J.~Snowe of
Maine.

``It's so important to get a deal,'' a White House official said, speaking on the condition of
anonymity in order to be candid about strategy. ``He will do almost anything it takes to get one.''

In scheduling a prime-time speech for next Wednesday before members of the House and Senate and a
national television audience, Mr.~Obama chose to put his political standing on the line more
directly and dramatically than he has so far on health care, his signature domestic initiative.

He will deliver the address 16 years after President Bill Clinton outlined his plan for universal
insurance coverage in a speech to Congress on Sept.~22, 1993. A year later, in September 1994, the
legislation was declared dead, after withering attacks by Republicans and insurance companies. Some
Obama advisers, wary of parallels between that effort and Mr.~Obama's push for an overhaul of the
health system, had argued that the president should give a televised speech from the Oval Office
instead of the House chamber.

Timed to coincide with Congress's return from its summer recess, the president's address would
follow a tumultuous month in which opponents of sweeping health legislation disrupted lawmakers'
town-hall-style meetings and the White House struggled to regain control of the debate. Mr.~Obama is
also scheduled to travel to Cincinnati on Monday to speak at a large Labor Day picnic organized by
the A.F.L.-C.I.O. He will have a receptive audience, as labor unions have been among the strongest
supporters of his effort to expand coverage and rein in health costs.

So far, the administration's ideas of concessions are likely to fall far short of the fundamental
changes that Congressional Republicans seek.

For now, White House officials said, Mr.~Obama remains committed to the goal of insuring all
Americans and still prefers to foster competition for insurance companies by creating a new
government insurance program, or public option.

White House officials are combing the versions of health care legislation approved by four of the
five Congressional committees with jurisdiction on the issue, both to find common ground and to
jettison provisions -- some relatively minor -- that have drawn fire from critics on the political
right.

To avoid some of the most heated criticism voiced in recent weeks, White House officials said they
would have no objection if Congress scrapped proposals to have Medicare pay for counseling on
end-of-life care.

Critics said such counseling could lead to pressure on patients to forgo expensive treatments for
terminal illnesses. Mr.~Obama has said it is ludicrous to suggest that ``we want to set up death
panels to pull the plug on Grandma.''

White House officials said Congress could also drop proposals requiring the government to create
school-based health clinics and collect nationwide data on health and health care by race, sex,
sexual orientation and ``gender identity.''

Supporters of the House bill said such data would help reduce ``health disparities,'' but critics
said they feared the government could assemble a database that posed a threat to personal privacy.

If Mr.~Obama does not gain traction by making these concessions, his allies on Capitol Hill said,
they may have to consider bigger changes. For example, they said, rather than requiring all
Americans to carry health insurance, Congress might start by requiring coverage of children, or
families with children.

While such a change would deeply disappoint many of Mr.~Obama's supporters, it could have two
potential political benefits, reducing the initial cost of any bill and reducing the size of cuts
needed in the future growth of Medicare.

Senator Charles E.~Schumer of New York, the third-ranking Democrat in the Senate, welcomed
Mr.~Obama's plan for an address to Congress.

``This level of involvement from the president could well be a game-changer,'' Mr.~Schumer said.
``There is no better way to turn public opinion around than to have someone as popular as President
Obama addressing the American people directly, without intermediaries interpreting, or
misinterpreting, his ideas.''

Republicans' reactions showed that they had become emboldened in their opposition since Congress
last met.

The House Republican leader, Representative John A.~Boehner of Ohio, said: ``House Republicans want
to hear what the president has to say, but after the public outcry this August, it's clear the
American people don't want a new speech. They want a new plan.''

Mr.~Boehner said the Democrats should scrap their current proposals and start over.

\section{New Snag for Oracle in Sun Deal}

\lettrine{E}{uropean} regulators delayed the proposed takeover of Sun Microsystems by the software
company Oracle on Thursday, indicating that the combination could squelch\cite{squelch} the growth
of a popular, free corporate database program owned by Sun.

The decision by the European Commission to extend its investigation into the deal, worth \$7.4
billion, is especially sensitive because the Justice Department has already approved the merger.
Regulators in the United States questioned Oracle's market power in some areas of its business but
raised fewer concerns than the Europeans about open-source software.

The European Commission's assertiveness has conflicted in the past with the Justice Department's
judgment. It has objected to mergers of American companies on several occasions, but in 2003, it
outright rejected the merger of General Electric and Honeywell after the American authorities
approved it. Mario Monti, competition commissioner at the time, said that G.E. would become too
dominant in markets for aircraft engines.

In recent years, Mr.~Monti and Neelie Kroes, the current European Union competition commissioner,
have found themselves at odds with some of their American counterparts over whether to force
Microsoft to change its Windows operating system.

American and European antitrust officials, legal experts say, agree far more often than they differ,
other than in a handful of cases. ``But there are somewhat different sensitivities,'' noted Andrew
I.~Gavil, a law professor at Howard University. ``Even after the change in administrations in
Washington, the American level of concern about postmerger price increases tends to be less than in
Europe. And European antitrust officials are more protective of consumers and more confident of the
beneficial consequences of intervention.''

Samuel R.~Miller, a partner at Sidley Austin in San Francisco who acted as special trial counsel for
the Justice Department's first antitrust case against Microsoft, said, ``This action reflects a
continued pattern of aggressive antitrust enforcement regardless of whether the companies are based
in Europe, the U.S.~or Asia.''

Oracle has tussled with European Union regulators in the past over its ambitions but has overcome
initial opposition. In 2004, the commission approved Oracle's acquisition of PeopleSoft without
conditions after subjecting the deal to the kind of in-depth inquiry now under way over its purchase
of Sun. The Justice Department had opposed that merger, but lost in a legal battle with Oracle. The
following year the European commission also approved Oracle's acquisition of Siebel Systems, again
without conditions.

In announcing the decision Thursday, Ms.~Kroes warned that the acquisition could hamper development
of an important software product owned by Sun, which specializes in computer hardware. The product,
MySQL, is the most widely used corporate database software in the world, and it competes with
software produced by Oracle.

Ms.~Kroes said preserving access to open-source software was vital when much of the world, including
Europe, might just be emerging from a deep slump.

``In the current economic context, all companies are looking for cost-effective IT solutions and
systems based on open-source software are increasingly emerging as viable alternatives to
proprietary solutions,'' Ms.~Kroes said. She said a longer investigation was needed ``to ensure that
such alternatives would continue to be available.''

Oracle had no comment on the action.

The commission has at least three months, or until Jan.~19, to decide whether to clear the deal or
issue an order blocking it.

European and American antitrust authorities have sought to narrow their differences in recent years
in an effort to avoid disputes that have marred trans-Atlantic relations this decade.

Antitrust experts said the decision to investigate the effects on open-source software of Oracle's
acquisition of Sun showed that differences persisted between American and European regulators.

``Europeans still have a lot more concerns than Americans about companies using strong or dominant
positions to create bottlenecks for competitors in the information and technology sectors,'' said
Peter Alexiadis, a partner at the law firm Gibson, Dunn \& Crutcher who is based in Brussels.

``Any whiff\cite{whiff} of dominance over different platforms used to deliver information raises
particular concerns,'' he said. ``This may in part explain why Europeans, who are used to multiple
business traditions, might be less inclined to view Oracle's traditional strengths in databases as
not posing competitive concerns.''

Mr.~Alexiadis also noted that a German company, SAP, was an example of a European software company
that fiercely competes with Oracle on a number of relevant markets affected by the deal.

In the beginning, writers of open-source software did not see it as a source of profit. More
recently, companies with open-source operations have begun making large amounts of money on those
products by providing support services to go along with them. That bundle of software and service is
still often priced at a discount to software sold by companies like Oracle.

A major concern cited by European investigators was what would happen to MySQL once Oracle took it
over.

They expressed concern that Oracle would have an incentive to stymie the development of MySQL as a
way of improving the sales of its competing database products.

Some experts said the concerns of Ms.~Kroes were probably unfounded.

Bo Lykkegaard, an analyst with IDC in Copenhagen, said Oracle had bought Sun for a variety of
reasons and that MySQL was not among Oracle's priorities.

Even so, Mr.~Lykkegaard said that Oracle, by keeping MySQL's open-source status, would be able to
develop parts of its business by reaching out to small companies or departments within large
companies that were seeking value-priced software for running operations that were not necessarily
``mission critical.''

Another issue that may have led the Europeans to take more time with the case is the way that Oracle
has handled regulators on both sides of the Atlantic.

Oracle notified European Union regulators of its deal in late July, more than two months after it
informed American officials.

European merger watchdogs can take a dim view if companies spread out their notifications between
jurisdictions over long periods of time, and they have said in the past that such tactics might be
aimed at pressuring the Europeans to give the green light to takeovers already approved in the
United States.

\section{New Protests Reported in Restive Chinese Region}

\lettrine{T}{housands} of Han Chinese protesters swarmed around government buildings in the capital
of the restive Xinjiang region on Thursday to demand a crackdown on Uighurs after rumors spread that
they were sticking hundreds of Hans with H.I.V.-tainted hypodermic needles.

The fresh conflict in the capital, Urumqi, showed the resilient hostility between the Han, China's
dominant ethnic group, and the Uighurs, a Turkic-speaking people who mostly follow Sunni Islam. The
Uighurs are the largest ethnic group in Xinjiang, and long-simmering tensions between the two groups
broke out in deadly rioting in Xinjiang two months ago.

On Thursday, witnesses described chanting crowds marching with Chinese flags, as well as moments of
violence. At least one Uighur was beaten by a crowd as paramilitary police officers watched, said
one witness. The police erected roadblocks, and schools were shut down.

Many protesters yelled, ``Wang Lequan, step down!'' referring to the powerful regional secretary of
the Communist Party who has run Xinjiang for 15 years. Mr.~Wang, a member of the ruling Politburo,
has been widely criticized by Uighurs and foreign scholars of Xinjiang as a hard-liner whose
policies have widened the divides among Hans, Uighurs and other ethnicities. But the Han frustration
with him is based on criticism that he has not been harsh enough with the Uighurs.

Officials declined to discuss the protests. An employee at a government news media center in a hotel
in Urumqi said: ``You guys should not be so eager competing for the news. Just wait one day until
the official news comes out.''

The latest protests took place in the middle of a five-day trade fair in Urumqi that is aimed at
attracting overseas investment to Xinjiang. One government employee said regional officials feared
that the Xinjiang economy would continue to suffer in the aftermath of the July violence, which left
at least 197 people dead and 1,721 wounded, and brought ethnic tensions into the spotlight. The
government said most of the casualties were Han civilians.

On Wednesday, China Daily, the official English language newspaper, ran a front-page article on the
trade fair with the headline ``Xinjiang Emerging From Riot Shadow.''

Rumors of people trying to spread H.I.V., the virus that causes AIDS, with needle attacks have
circulated in China for years. Talk that Uighurs were on a needle-sticking spree built rapidly in
Urumqi. A small protest on Wednesday led to the much larger one on Thursday.

Media accounts might have helped fan the anxiety. On Thursday, state-run Xinjiang Television said
that 476 people, most of them Hans, had sought treatment since Aug.~20 for needle stab wounds,
though only 89 had obvious signs they had been pricked.

Also, the Web site of China Daily published an article from Xinhua, the state news agency, that said
police officers had detained 15 people in Xinjiang for needle attacks, though it did not say when.
Four had been prosecuted, the account said.

The account quoted Yu Yunlin, a health official, as saying no one had been infected or poisoned so
far.

After the large protests, government censors rushed to delete the news accounts on the Internet. By
late afternoon, the Xinhua article published by China Daily was still on the newspaper's Web site,
but a news release on the same subject, posted on the Web site of the information office of the
State Council, China's cabinet, had been deleted.

By then, the rumors of stabbings had already struck fear into Han civilians across Urumqi.

``It's the Uighurs!'' said a retired woman living in an apartment near People's Square in Urumqi,
where protesters had gathered, when reached by telephone. ``They hate the Han! Hundreds have been
stabbed. Old people, young children, pregnant women.''

The woman said she had bought a week's worth of food because she feared she might be stuck with a
needle if she went outside. People were taking taxis because they feared being stabbed on buses, she
said.

``The government hasn't done anything,'' she said. ``They haven't told us anything. They haven't
kept order. We're all so angry. The July 5 incident was so brutal that we still haven't recovered
from that yet. And now this?''

\section{Out of Work, Too Down to Search On, and Uncounted}

\lettrine{T}{hey} were left out of the latest unemployment rate, as they are every month: millions
of hidden casualties\cite{casualty} of the Great Recession who are not counted in the rate because
they have stopped looking for work.

But that does not mean these discouraged Americans do not want to be employed. As interviews with
several of them demonstrate, many desperately long for a job, but their inability to find one has
made them perhaps the ultimate embodiment of pessimism as this recession wears on.

Some have halted their job searches out of sheer frustration. Others have decided it makes more
sense to become stay-at-home fathers or mothers, or to go back to school, until the job market
improves. Still others have chosen to retire for now and have begun collecting Social Security or
disability benefits, for which claims have surged.

Rick Alexander, a master carpenter in Florida who has given up searching after months of effort,
said the disappointment eventually became unbearable.

``When you were in high school and kept asking the head cheerleader out for a date and she kept
saying no, at some point you stopped asking her,'' he said. ``It becomes a 'why bother?' scenario.''

The official jobless rate, which garners the bulk of attention from politicians and the public, was
reported on Friday to have risen to 9.7 percent in August. But to be included in that measure, which
is calculated by the Bureau of Labor Statistics from a monthly nationwide survey, a worker must have
actively looked for a job at some point in the preceding four weeks.

For an increasing number of people in this country who would prefer to be working, that is not the
case.

It is difficult to assign an exact figure, because of limitations in the data collected by the
bureau, but various measures that capture discouragement have swelled in this recession.

In the most direct measure of job market hopelessness, the bureau has a narrow definition of a group
it classifies as ``discouraged workers.'' These are people who have looked for work at some point in
the past year but have not looked in the last four weeks because they believe that no jobs are
available or that they would not qualify, among other reasons. In August, there were roughly 758,000
discouraged workers nationally, compared with 349,000 in November 2007, the month before the
recession officially began.

The bureau also has a broader category of jobless it calls ``marginally attached to the labor
force,'' which includes discouraged workers as well as those who have stopped looking because of
other reasons, like school, family responsibilities or health issues. But economists agree that many
of these workers probably would have found a way to work in a good economy.

There were roughly 2.3 million people in this group in August, up from 1.4 million in November 2007.
If the unemployment rate were expanded to include all marginally attached workers, it would have
been 11 percent in August.

But even this figure is probably an undercount of the extent of the jobless problem in this country.
There are about 1.4 million more people who are not in the labor force than when the recession
began. Some of these are retirees, stay-at-home parents, people on disability and students. But it
is also rather likely that many of these people have given up looking for work at least partly
because of economic reasons as well.

Here are four people's stories:

Rick Alexander: A Builder by Trade, With Too Much Time

In the worst case, Rick Alexander figured, he could scrounge up a job at Home Depot.

He was a master carpenter, after all. He had skills. He had run his own successful home-restoration
business for 28 years.

In early 2008, however, he moved to Florida to take care of his ailing parents, leaving his business
in Connecticut to his daughter.

After helping his parents into an assisted-living facility, he began applying for jobs. He devoted
eight hours a day to the task, sometimes sending out three or four applications a day.

``It was a full-time job,'' he said.

At first, he focused on jobs in construction, applying to be a site supervisor. He looked for
anything within an hour's commute of where he was living in Jensen Beach.

But the real estate industry had fallen off precipitously, bringing building to a near standstill.
Mr.~Alexander, 58, began branching out to suppliers, applying at lumberyards and other wholesalers.
Eventually, he expanded his search to Home Depot, Lowe's and mom-and-pop hardware stores. Finally,
he began applying for ``everything under the sun,'' even the overnight shift at convenience stores.

By that summer, he had still received no callbacks for interviews. He went back to Connecticut for
several weeks to do a renovation for an old client to earn some cash. When he returned to Florida in
August 2008, he tried to start his own business, selling advertising on video displays mounted in
coffee shops and other places.

He networked furiously with local businesses, but by then the economy had nose-dived. Mr.~Alexander
said he grossed a total of \$150. He sank into a funk and stopped looking.

``There are thousands of people applying for every job I'm looking at, and potential employers won't
even give me the courtesy of acknowledging I applied,'' he said. ``The entirety of that causes me
not to bother. It's a waste of my time and theirs.''

He has applied to just two jobs this year, both several months ago. The unemployment rate in his
area, Martin County, now exceeds 11 percent. After prodding from his companion, Dona Olinger, he
went down to Home Depot a little over a month ago to re-activate his application there.

His savings are gone. He lives with Ms.~Olinger, who makes \$10 an hour as a volunteer coordinator
at a food pantry, Harvest Food and Outreach Center, where they also get groceries every week. It is
her salary that pays their rent.

Mr.~Alexander's parents have since moved out of the assisted-living facility and back into their
home, so he tends to them most days. He reads Robert Ludlum novels. He sleeps. To fill his time, he
is looking into volunteer work. The other day, he cut the grass on his small lawn using just a pair
of clippers.

Ray Rucker: Feeling Counted Out With Years Still Left

Ray Rucker came home from a job interview several months ago, sat down in his living room with his
suit still on and wept.

The meeting with the interviewer had lasted 10 minutes. The man did not even open a folder in front
of him to study Mr.~Rucker's r\'esum\'e. It was just ``jibber jabber,'' Mr.~Rucker said later.

Mr.~Rucker, who lives in Overland Park, Kan., had little doubt about what had happened. He is 62
years old and, as he puts it, ``I look 62.''

He lost his job as a facilities manager for Starbucks in Kansas City and Wichita, Kan., last
November, when the company closed hundreds of stores across the country. He had done similar work
for years for other national restaurant chains and retail outlets.

He landed his first interview within a month, with a retail chain. He was invited back to talk to
the vice president of operations and to the director of operations. He was also invited to meet with
the company's chief executive.

But as Mr.~Rucker was finishing with the director of operations, she asked him straight out whether
he was retiring soon. Shocked, Mr.~Rucker answered, truthfully, that he planned to work at least 10
more years.

The meeting with the chief executive never came. Mr.~Rucker said he thinks his interviewer simply
did not believe he planned to continue working.

A month ago, he found a job posting that seemed tailored for him, a facilities manager for a
national restaurant chain. He sent in his r\'esum\'e and three days later got called for an
interview. The company official said he was in a hurry to fill the position. But Mr.~Rucker soon
learned that this one, too, had slipped from his grasp.

``That's the one when I kind of threw in the towel,'' he said.

Mr.~Rucker said he was done looking. His wife, who works at a small nonprofit organization,
protested, saying there was more he could do to look.

``You don't know what I'm going through,'' Mr.~Rucker said he told her.

``You send out so much, and you don't get responses,'' he said. ``Then when you get called in,
you're treated like you're too old. Why am I doing this?''

So he made an appointment with the local Social Security office to begin claiming benefits. He might
try to get some kind of hourly job to help make ends meet. He has mapped out some home renovation
projects he wants to do.

The Social Security checks will not equal even a third of what he used to make. But he is now
preparing for semiretirement.

Jenny Salinas: From a Nonstop Career to a Focus on the Home

Jenny Salinas never envisioned being a stay-at-home mother, taking care of the children and keeping
house. She was the one with the high-powered career, the six-figure salary, always jetting off to
Russia or China.

She put her 5-year-old daughter, Mia, in day care when she was three months old. Mia got so used to
her mother going away she would simply say, ``Mommy's on a trip,'' and blow her kisses when she
left.

But after searching unsuccessfully since January for a job, Mrs.~Salinas, 37, said her priorities
had shifted. She is now content to stay home and focus on her family. She and her husband are even
talking about having more children.

``It's just amazing how it changes your perspective on what's important,'' she said.

Mrs.~Salinas had been a manager of corporate marketing and media relations at an oil and gas company
in Houston, where she lives. She was so focused on her career, she said, that she never noticed her
daughter had a lazy eye. Mrs.~Salinas's mother mentioned something to her, but only after
Mrs.~Salinas was laid off did she realize that her daughter needed to see an ophthalmologist.

``That's how much I was on my BlackBerry,'' Mrs.~Salinas said.

Mrs.~Salinas was initially confident that she would land somewhere quickly. She seemed to be doing
well, too, scoring interview after interview for senior-level corporate marketing positions. But
each of those prospects dried up, usually because of a hiring freeze or some other obstacle.

So, for the last two months, she has not looked at all. Partly, she has been busy, selling their old
house, moving into a new one they are renting at half the monthly expense, seeing her daughter off
to kindergarten.

She is helped by the fact that her husband, a vice president at an advertising agency, still has his
job. After the couple realized that her job search might take time, they decided to cut back on
their spending.

She has in mind a specific set of companies, but they are all still not hiring. Unwilling to settle
for just any job, she said, she would rather bide her time.

But the process of searching for work and coming up empty has also left her feeling spent.

``I was just discouraged, fed up and angry, feeling like my career had betrayed me,'' she said.

Her daughter used to be in day care or preschool from 7 a.m. to 6 p.m., but Mrs.~Salinas began
dropping her off later and picking her up earlier. Some days, they skip day care completely and
while away the day together.

Tatjana Jovanovic-Grove: Moving From Serbia, Scraping By Online

Tatjana Jovanovic-Grove now occupies her days with arts and crafts projects. She makes a little
money selling them online -- \$10 here, \$50 there -- but mostly it beats the sense of futility that
used to envelop her each day during her quest to find a job.

``I stopped looking because that feeling of being rejected again and again is hard,'' she said.
``It's just like somebody punching you in the face.''

Ms.~Jovanovic-Grove, 41, has struggled to find work since she immigrated in late 2005 to the United
States from her native Serbia, where she was a biology researcher at a prestigious research
institute in Belgrade.

She had married an American, Doug Grove, 42, a Wal-Mart mechanic she met over the Internet. The
couple initially lived in Glendale, Ariz., with their three children from previous marriages, but
they moved to Winston-Salem, N.C., in late 2007.

They were attracted by the weather and the low crime rate. They also thought Ms.~Jovanovic-Grove,
who earned a master's degree in Serbia in environmental protection and zoology, would have an easier
time finding a job in an area rich with universities.

``I was really thinking I would have no problem,'' she said.

The need for her to find work became more urgent after the couple took on thousands of dollars in
additional debt after they turned their Arizona home over to a bank in lieu of a foreclosure
settlement. They had been unable to sell it amid the state's collapsing real estate market.

But aside from a few temporary jobs, Ms.~Jovanovic-Grove has come up empty on everything from
research assistant positions to retail jobs. Meanwhile, her husband's hours at Wal-Mart, where he is
paid a little more than \$14 an hour, have been cut back.

In May, she stopped looking completely, concluding that the job market was saturated.
Winston-Salem's unemployment rate exceeded 10 percent.

``You figure out it's just like when you toss a piece of meat at a pack of hungry cats,'' she said.
``I just gave up because I could not compete.''

Instead, she has turned to making wood handicrafts and selling them on Etsy.com, an online
marketplace. The small payments she gets often mean she earns less than fifty cents an hour for her
effort. But she reasoned it is better than wasting gas driving around applying for jobs she believes
she cannot get.

\section{Obama Faces a Critical Moment for His Presidency}

\lettrine{P}{resident} Obama returned to the White House from his summer break on Sunday determined
to jump-start his struggling presidency by reasserting command of the health care debate and
recalibrating expectations that some advisers believe got away from him.

With his honeymoon seemingly over and his White House on the defensive, Mr.~Obama faces what friends
and foes alike call a make-or-break moment in his young administration. Because he has elevated
health care to such a singular priority, advisers said he must force through a credible plan or risk
crippling his presidency.

``It goes without saying that a lot is riding now on his ability to re-energize the health care
debate and bring it home to a successful conclusion,'' said John D.~Podesta, who ran Mr.~Obama's
transition and still advises him on health care, energy and other issues. ``Nothing will influence
the perception of the presidency more than whether he can be successful in getting a health care
bill through the Congress.''

Recognizing the stakes, Mr.~Obama has worked on a strategy for autumn to regain the initiative. He
talked on Thursday from Camp David with Nancy Pelosi, the House speaker, and Harry M.~Reid, the
Senate majority leader. He spent part of Sunday working on this week's speech to the nation and
dispatched top surrogates to the talk shows to try to reframe the health care debate. And he has two
meetings scheduled for Monday with his health policy and political advisers planned around a trip to
Cincinnati to observe Labor Day.

As much as health care has consumed the president, other vexing issues await him in the fall. In the
coming weeks, he will decide whether to order thousands more troops to Afghanistan and pursue new
sanctions against Iran. He will host a meeting of the Group of 20 nations to spur the world economy
and push forward with arms control negotiations with Russia.

Now, as he prepares for Wednesday's address before a joint session of Congress, Mr.~Obama and his
team are simultaneously trying to figure out how they got into this dilemma and how to get out of
it. An administration that swept into office just seven months ago on a wave of hope and optimism
has burned through good will and public patience in swift fashion and now finds itself under fire
from both the left and the right.

He faces a crisis of expectations tough to manage. Can he form a health care compromise that
satisfies both his liberal base and fiscal conservatives in his own party, much less the other one?
Can he stanch the slide in support for the war in Afghanistan even as he considers sending more
troops? Can he soothe discontent with an economy that appears to have bottomed out but remains
moribund? Can he change the tenor of debate in a capital that seems as polarized as ever?

``To govern is to make choices, and to make choices is to make some unhappy,'' David Axelrod, the
president's senior adviser, said in an interview. ``He made some very tough decisions that pulled us
away'' from a new Great Depression. ``But he had to expend some political capital to do that. He's
expending some capital to do something that's very important, which is to bring security and health
care to people who don't have it.''

Some Republicans said Mr.~Obama's fundamental mistake was believing his election presaged a larger
ideological shift in the country. ``If they thought that his popularity and the good will he had
would support liberal policies, they were wrong,'' said Charles R.~Black, a Republican strategist
who worked last year for Senator John McCain of Arizona, Mr.~Obama's Republican presidential
opponent.

White House officials have signaled that they are prepared to scale back their aspirations for the
health care legislation. In private conversations, some said they would be happy even if they end up
with a pared-back program that can serve as a basis for future efforts.

One element clearly on the table is a proposed government-backed health insurance plan to compete
with private insurers. Just as they have in recent weeks, White House officials indicated Sunday
that Mr.~Obama would continue to push for the so-called public option but they did not make it a
condition of signing whatever bill lands on his desk.

Mr.~Axelrod, appearing on NBC's ``Meet the Press,'' said the public option ``is a valuable tool''
but added that ``it shouldn't define the whole health care debate.'' Robert Gibbs, the White House
press secretary, said on ABC's ``This Week With George Stephanopoulos'' that Mr.~Obama would ``draw
some lines in the sand'' on Wednesday but ``I doubt we're going to get into heavy veto threats.''

The conundrum for the president, though, was on display during a roundtable discussion later on the
same program. Robert Dole, the former Republican senator from Kansas and onetime presidential
nominee, said a public option would never pass the Senate. Representative Maxine Waters, Democrat of
California and a leader of Congressional liberals, said no plan could pass with House without a
public option.

Mr.~Obama is hardly the first president to run into trouble after the bunting and balloons have
vanished, but his slipping support has fueled a narrative about a young and relatively inexperienced
president who overinterpreted his mandate and overreached in his policies. His job approval rating
has fallen to 56 percent from 62 percent since February in polls taken by The New York Times and CBS
News. Other surveys register an even sharper drop.

But his overall standing with the public is still healthy, and his first seven months have not been
as rocky as those of Bill Clinton or Gerald R.~Ford. Mr.~Clinton, at least, later recovered enough
to win re-election. And Mr.~Obama showed during last year's campaign that he has the capacity to
ride out rough moments. If he ultimately gets some form of health care program passed that he can
call a victory, this turbulence may ultimately be forgotten.

Vin Weber, a former Republican congressman from Minnesota, said the backlash to Mr.~Obama's spending
and health care proposals had eroded his support but not fundamentally shifted the nation's
politics. ``The American people are sort of returning to where they were,'' he said. ``I don't think
they've made a big swing to the right. They're returning to their centrist moorings.''

Of all the challenges Mr.~Obama faces this fall, health care has come to dominate so much that the
fate of the rest of his domestic program, particularly climate change legislation and new
regulations on the financial industry, may depend in part on whether he wins this fight.

``He's gone all in,'' said Matt Bennett, vice president of Third Way, a Democratic-oriented advocacy
organization, using a poker term. ``Everyone's watching. The bets are all on the table. And we're
just waiting to see what the cards say.''

\section{Sudan Court to Define Indecent Dress for Women}

\lettrine{T}{his} is not about pants, Lubna Hussein insists. It is about principles.

A woman should be able to wear what she wants and not be publicly whipped for it, says Mrs.~Hussein,
a defiant Sudanese journalist, and on Monday her belief will be put to the test.

Mrs.~Hussein has been charged in Khartoum, Sudan's capital, with indecent dress, a crime that
carries a \$100 fine and 40 lashings. She was arrested in July, along with 12 other women, who were
caught at a cafe wearing trousers.

Sudan is partially ruled by Islamic law, which emphasizes modest dress for women. Mrs.~Hussein, 34,
has pleaded not guilty and is daring the Sudanese authorities to punish her.

``I am Muslim; I understand Muslim law,'' Mrs.~Hussein said in an interview. ``But I ask: what
passage in the Koran says women can't wear pants? This is not nice.''

Mrs.~Hussein even printed up invitation cards for her initial court date in July and sent out e-mail
messages asking people to witness her whipping, if it came to that. She said she wanted the world to
see how Sudan treated women.

Hundreds of Sudanese women -- many wearing pants -- swarmed in front of the court where the trial
was supposed to take place, protesting that the law was unfair. Twice now, the trial has been
postponed. Some of the other women arrested with Mrs.~Hussein have pleaded guilty and were lashed as
a result. Past floggings have been carried out with plastic whips that leave permanent scars.

``The flogging, yes, it causes pain,'' Mrs.~Hussein said. ``But more important, it is an insult.
This is why I want to change the law.''

The law in contention here is Article 152 of Sudan's penal code. Concisely stated, the law says that
up to 40 lashes and a fine should be assessed anyone ``who commits an indecent act which violates
public morality or wears indecent clothing.''

The question is: what exactly is indecent clothing?

In Sudan, some women wear veils and loose fitting dresses; others do not. Northern Sudanese, who are
mostly Muslim, are supposed to obey Islamic law, while southern Sudanese, who are mostly Christian,
are not. Mrs.~Hussein argues that Article 152 is intentionally vague, in part to punish women.

Rabie A.~Atti, a Sudanese government spokesman, said the law was meant for the opposite reason, to
``protect the people.''

``We have an act controlling the behavior of women and men so the behavior doesn't harm others,
whether it's speech or dress or et cetera,'' he said.

But, he insisted, Mrs.~Hussein must have done something else to run afoul of the authorities,
besides wearing pants.

``You come to Khartoum and you will see for yourself,'' he said. ``Many women, in offices and
wedding ceremonies, wear trousers.''

``Thousands of girls wear the trousers,'' he added.

Asked what other offenses Mrs.~Hussein may have committed, Mr.~Atti said that the case file was
secret and that he did not know.

Mrs.~Hussein countered that she did not do anything else that might have violated the law, and that
countless people from inside and outside Sudan are supporting her.

``It's well known that Sudanese women are pioneers in the history of women's rights in this region,
and that we won our rights a long time ago because of our awareness, open mind, good culture and
struggle,'' she said.

The last time Sudan's courts handled a case that attracted such international attention, they found
a compromise solution. A British schoolteacher faced up to 40 lashes and six months in prison for
allowing her students to name a class teddy bear Muhammad, which was perceived as an insult to
Islam. But after being sentenced to 15 days in jail, she was soon pardoned by the Sudanese
president.

A widow with no children, Mrs.~Hussein is a career journalist who recently worked as a public
information assistant for the United Nations in Sudan. She quit, she said, because she did not want
to get the United Nations embroiled in her case. But Sudan, given its renewed interest in
normalizing relations with the United States, might be reluctant to draw much international ire by
harshly punishing her.

Protesters are expected to come out on her behalf again when Mrs.~Hussein returns to court Monday
morning. She says her family is also behind her.

``My mother supports me,'' she said, ``but she is worried for me and prays for me.''

\section{Secretary Endorses Obama's Talk for Students}

\lettrine{C}{riticism} of President Obama's plans to address the nation's students, which has
bubbled up on conservative radio talk shows and in calls from parents to hundreds of school
districts in recent days, has been ``just silly,'' Education Secretary Arne Duncan said Sunday.

Mr.~Obama's back-to-school address, scheduled for Tuesday morning, will encourage students to study
hard and will steer clear of political themes, Mr.~Duncan said in an appearance on ``Face the
Nation'' On CBS. In any event, he said, whether students watch the speech is entirely up to parents
and local school officials.

``The whole message is about personal responsibility and challenging students to take their
education very, very seriously,'' Mr.~Duncan said. The speech text was to be posted on the White
House Web site on Monday.

Last week, conservative talk show hosts accused the president of sinister motives, and some parents
vowed to keep their children out of school on Tuesday.

``That's just silly,'' Mr.~Duncan said. He acknowledged, however, that a poorly-worded message sent
to public schools across the nation may have contributed to public misunderstanding.

The Department of Education sent out a list of suggested classroom activities that teachers could
use to accompany the speech. One among dozens suggested that teachers could have students ``write
letters to themselves about what they can do to help the president.''

Some conservatives focused on the letter-writing activity as an effort by the administration to use
classrooms to build political support for the president. Department officials on Wednesday replaced
that initial suggested activity with one in which students would ``write letters to themselves about
how they can achieve their short-term and long-term educational goals.''

The original activity, Mr.~Duncan said Sunday, ``wasn't worded quite correctly.''

Among those criticizing plans for the speech was the chairman of the Florida Republican Party, who
accused Mr.~Obama of seeking to spread what he called ``socialist ideology.''

But two prominent Republicans distanced themselves from such talk, defending the administration's
plans.

Newt Gingrich, the former House speaker, said on ``Fox News Sunday,'' ``It is good to have the
president of the United States say to young people across America, 'Stay in school, study and do
your homework.' ''

And Senator Lamar Alexander, Republican of Tennessee, said on the same show, ``Of course the
president of the United States should be able to address students.''

``If I were a teacher,'' Mr.~Alexander added, ``I'd take advantage of it.''

\section{Google's Internet Techniques Inspire Studies of Food Webs}

\lettrine{A}{ major} reason Google's search engine is so successful is its PageRank algorithm, which
assigns a pecking order to Web pages based on the pages that point to them. A page is important,
according to Google, if other important pages link to it.

But the Internet is not the only web around. In ecology, for instance, there are food webs -- the
often complex networks of who eats whom.

Inspired by PageRank, Stefano Allesina of the University of Chicago and Mercedes Pascual of the
University of Michigan have devised an algorithm of their own for the relationships in a food web.
As described in the online open-access journal PLoS Computational Biology, the algorithm uses the
links between species in a food web to determine the relative importance of species in a food web,
which will have the most impact if they become extinct.

Dr.~Allesina, who studies network theory and biology, was reading a paper about Google's algorithm
one day while at the University of California, Santa Barbara. ``I said, 'This reminds me of
something,' '' he recalled.

One key to PageRank's success is that its developers introduced a small probability that a Web user
would jump from one page to any other. This in effect makes the Web circular, and makes the
algorithm solvable. But in food webs, Dr.~Allesina said, ``you can't go from the grass to the lion
-- the grass has to go through the gazelle first.

``We could not use the same trick to make food webs circular,'' he went on.

So they used another trick, he said. Since all organisms die and decompose, they created a
``detritus pool'' that all species link to. The pool also links to primary producers in a food web,
which make use of the decomposed matter.

Their algorithm differs also in that it determines the relative importance of species through
reverse engineering -- by seeing which species make the food web collapse fastest if they are
removed. The researchers found that the algorithm produces results that were as accurate as much
more complex (and computationally costly) software that builds webs from the ground up, simulating
evolution.

The next step, Dr.~Allesina said, is to refine the algorithm so that it will work with more complex
webs. There are many other factors that affect extinctions, including pollution and habitat loss.
The goal is to create an algorithm that can take these and other elements into account as well.

\section{Panels of Light, Not Just Points, Fascinate Designers and Industry}

\lettrine{L}{ED} light bulbs, with their minuscule energy consumption and 20-year life expectancy,
have grabbed the consumer's imagination.

But an even newer technology is intriguing the world's lighting designers: OLEDs, or organic
light-emitting diodes, create long-lasting, highly efficient illumination in a wide range of colors,
just like their inorganic LED cousins. But unlike LEDs, which provide points of light like standard
incandescent bulbs, OLEDs create uniform, diffuse light across ultrathin sheets of material that
eventually can even be made to be flexible.

Ingo Maurer, who has designed chandeliers of shattered plates and light bulbs with bird wings, is
using 10 OLED panels in a table lamp in the shape of a tree. The first of its kind, it sells for
about \$10,000.

He is thinking of other uses. ``If you make a wall divider with OLED panels, it can be extremely
decorative. I would combine it with point light sources,'' he said.

Other designers have thought about putting them in ceiling tiles or in Venetian blinds, so that
after dusk a room looks as if sunshine is still streaming in.

Today, OLEDs are used in a few cellphones, like the Impression from Samsung, and for small,
expensive, ultrathin TVs from Sony and soon from LG. (Sony's only OLED television, with an 11-inch
screen, costs \$2,500.) OLED displays produce a high-resolution picture with wider viewing angles
than LCD screens.

In 2008, seven million of the one billion cellphones sold worldwide used OLED screens, according to
Jennifer Colegrove, a DisplaySearch analyst. She predicts that next year, that number will jump more
than sevenfold, to 50 million phones.

But OLED lighting may be the most promising market. Within a year, manufacturers expect to sell the
first OLED sheets that one day will illuminate large residential and commercial spaces. Eventually
they will be as energy efficient and long-lasting as LED bulbs, they say.

Because of the diffuse, even light that OLEDs emit, they will supplement, rather than replace, other
energy-efficient technologies, like LED, compact fluorescent and advanced incandescent bulbs that
create light from a single small point.

Its use may be limited at first, designers say, and not just because of its high price. ``OLED
lighting is even and monotonous,'' said Mr.~Maurer, a lighting designer with studios in Munich and
New York. ``It has no drama; it misses the spiritual side.''

``OLED lighting is almost unreal,'' said Hannes Koch, a founder of rAndom International in London, a
product design firm. ``It will change the quality of light in public and private spaces.''

Mr.~Koch's firm was recently commissioned by Philips to create a prototype wall of OLED light, whose
sections light up in response to movement.

Because OLED panels could be flexible, lighting companies are imagining sheets of lighting material
wrapped around columns. (General Electric created an OLED-wrapped Christmas tree as an experiment.)
OLED can also be incorporated into glass windows; nearly transparent when the light is off, the
glass would become opaque when illuminated.

Because OLED panels are just 0.07 of an inch thick and give off virtually no heat when lighted, one
day architects will no longer need to leave space in ceilings for deep lighting fixtures, just as
homeowners do not need a deep armoire for their television now that flat-panel TVs are common.

The new technology is being developed by major lighting companies like G.E., Konica Minolta, Osram
Sylvania, Philips and Universal Display.

``We're putting significant financial resources into OLED development,'' said Dieter Bertram,
general manager for Philips's OLED lighting group. Philips recently stepped up its investment in
this area with the world's first production line for OLED lighting, in Aachen, Germany.

Universal Display, a company started 15 years ago that develops and licenses OLED technologies, has
received about \$10 million in government grants over the last five years for OLED development, said
Joel Chaddock, a technical project manager for solid state lighting in the Energy Department.

Armstrong World Industries and the Energy Department collaborated with Universal Display to develop
thin ceiling tiles that are cool to the touch while producing pleasing white light that can be
dimmed like standard incandescent bulbs. With a recently awarded \$1.65 million government contract,
Universal is now creating sheetlike undercabinet lights.

``The government's role is to keep the focus on energy efficiency,'' Mr.~Chaddock said. ``Without
government input, people would settle for the neater aspects of the technology.''

G.E. is developing a roll-to-roll manufacturing process, similar to the way photo film and food
packaging are created; it expects to offer OLED lighting sheets as early as the end of next year.

``We think that a flexible product is the way to go,'' said Anil Duggal, head of G.E.'s 30-person
OLED development team. OLED is one of G.E.'s top research priorities; the company is spending more
than half its research and development budget for lighting on OLED.

Exploiting the flexible nature of OLED technology, Universal Display has developed prototype
displays for the United States military, including a pen with a built-in screen that can roll in and
out of the barrel.

The company has also supplied the Air Force with a flexible, wearable tablet that includes GPS
technology and video conferencing capabilities.

As production increases and the price inevitably drops, OLED will eventually find wider use, its
proponents believe, in cars, homes and businesses.

``I want to get the price down to \$6 for an OLED device that gives off the same amount of light as
a standard 60-watt bulb,'' said Mr.~Duggal of G.E. ``Then, we'll be competitive.''

\section{Like Apple, TV Explores Must-Have Applications}

\lettrine{T}{he} cable and satellite TV business has a big case of iPhone envy.

Apple has been able to popularize its cellphone in a crowded field by giving away or selling
specialized applications that make the phone more useful. So far, independent developers have
written more than 65,000 apps.

DirecTV and the FiOS service from Verizon Communications have recently announced app stores modeled
directly on Apple's App Store. Just a few applications have shown up so far, but already these few
-- Bible verses, Facebook updates and fantasy sports team updates -- suggest that people may not be
content to sit back while watching TV but rather want to lean forward and interact and customize
their TVs.

Most of the other cable, satellite and phone companies are also developing technology that will let
their set-top boxes run more complex applications, including those written by outside developers.

But the companies are still wrestling with how open they want their systems to be to outside
developers, what business arrangement to make with developers and what sorts of things people want
to do while watching their TV from their couches.

TV systems, after all, have long been tightly controlled by their operators, who send squadrons of
lawyers to negotiate deals with even the most obscure channel. To them, the prospect of emulating
Apple's sprawling marketplace is frightening, yet still increasingly appealing.

``The beauty of the iPhone is that there are a lot of applications that Apple would not have
imagined people want,'' said Sree Kotay, chief software architect for Comcast. ``We want people
engaged with television in ways we haven't thought of yet.''

The apps idea goes beyond interactive television.

For years, many cable and satellite customers have been able to get news articles, weather forecasts
and sports scores through their set-top boxes. Some boxes also insert on-screen messages during some
commercials, urging viewers to use their remote to get more information on screen.

But very little of this interactivity actually engaged customers, leading many in the industry to
conclude that all most people want to do with their TVs is watch.

With people getting used to sophisticated electronic interactions, not just on computer screens but
on devices used all day -- car dashboards, cameras and especially cellphones -- the industry is
rethinking that theory.

While viewers remain passive most of the time, they increasingly want the capabilities and
information they come to depend on from one gadget with a screen to be available to them on all the
other screens they use every day.

On FiOS, the Facebook application lets people see photos on their television screens that have been
shared by their friends. The Twitter application shows a running stream of tweets about the show
being viewed on the left side of the screen.

``The shows become a lot more fun to watch because there is a whole conversation going on,'' said
Shadman Zafar, Verizon's senior vice president for product development. ``For sports, it's like
bringing the rest of the sports bar into your home.''

The Twitter service has been a hit, with more than a million FiOS customers using it within the
first three days. But Mr.~Zafar mistakenly believed that people wanted to read tweets on the TV and
use their computers to send them.

``We didn't allow people to send tweets from their remotes,'' he said. ``We thought you mostly lean
back and want to see what other people say. In four hours after the product went live, we had
hundreds of requests, 'Can I tweet on the TV now?' ''

Within two days, Verizon added that feature; users tap out messages on the number pad on the remote
as if they were sending a text message on an older cellphone.

Another lesson learned: one of the biggest impediments to adding more features to televisions is the
remote control. It is difficult to enter text or even move quickly through complex menus.

So the cable industry would have to redesign and replace old remotes as well as old set-top boxes
that have limited memory and processing capabilities in tens of millions of homes.

``The lowest end set-top box is the equivalent of a Mac II from 1991,'' said Mr.~Kotay of Comcast.
Even recent models are hardly able to handle all the interactive tasks that developers have started
to imagine, especially combining text, graphics and video from cable channels and the Internet. And
the next generation will go further.

For example, a box Echostar plans to introduce this year will have a full Internet browser that can
play video from nearly any site on the Web.

AT\&T and Verizon, which only recently began competing directly with Comcast and Time Warner Cable,
have the benefit of a technical fresh start. Their boxes have more computerlike features and faster
Internet connections. AT\&T's U-verse service, for example, recently offered an interactive version
of the Professional Golfers Association golf tournament that lets users switch between cameras aimed
at different parts of the course.

Comcast, Time Warner Cable and most of the other major cable companies are countering the phone
companies by backing a standard called the Enhanced TV Binary Interchange Format. When it is put
into place next year, it will allow all but the oldest set-top boxes to display simple interactive
applications that are set off by certain programs and advertisements. One of the first uses of this
software is to put the caller ID information on the TV screen when the phone rings.

In theory, this standard will also let outside developers write applications that could run in tens
of millions of homes. But the cable companies have not figured out how developers will submit
applications and what sort of business models it will support. (Even Verizon has not published the
terms of its new app store yet, but it hopes to split ad revenue or app purchase fees with
developers.)

Cable companies are betting that viewers may well want to see information and photographs related to
the programs they watch and the actors in them. Other areas discussed include casual games, sports
information and online shopping.

Steve Necessary, the vice president for video product development at Cox Communications, said the
company is wary of opening up too far because that would put the burden on the cable companies to
make sure that applications do not display things that children should not see.

``We want to provide tools to make sure parents can provide he proper level of monitoring,''
Mr.~Necessary said. ``By no means do we want to be the troll under the bridge keeping applications
from coming in.''

Some people, however, wonder if one difference in the use of cellphones, which usually have a single
user, and TVs, which are more often shared among family members, might create problems.

Jeffrey G.~Weber, AT\&T's vice president for video products, worries that opening the TV to
applications like e-mail might encourage domestic disputes.

``E-mail is a private thing,'' he said. ``What happens when the wife wants to read e-mail from her
mother on the set when you want to watch the game?''

\section{China and India Dispute Enclave\cite{enclave} on Edge of Tibet}

\lettrine{T}{his} is perhaps the most militarized Buddhist\cite{Buddhist} enclave in the world.

Perched above 10,000 feet in the icy reaches of the eastern Himalayas, the town of Tawang is not
only home to one of Tibetan Buddhism's most sacred monasteries, but is also the site of a huge
Indian military buildup. Convoys of army trucks haul howitzers along rutted mountain roads. Soldiers
drill in muddy fields. Military bases appear every half-mile in the countryside, with watchtowers
rising behind concertina wire. A road sign on the northern edge of town helps explain the reason for
all the fear and the fury: the border with China is just 23 miles away; Lhasa, the Tibetan capital,
316 miles; and Beijing, 2,676 miles.

``The Chinese Army has a big deployment at the border, at Bumla,'' said Madan Singh, a junior
commissioned officer who sat with a half-dozen soldiers one afternoon sipping tea beside a
fog-cloaked road. ``That's why we're here.''

Though little known to the outside world, Tawang is the biggest tinderbox in relations between the
world's two most populous nations. It is the focus of China's most delicate land-border dispute, a
conflict rooted in Chinese claims of sovereignty over all of historical Tibet.

In recent months, both countries have stepped up efforts to secure their rights over this rugged
patch of land. China tried to block a \$2.9 billion loan to India from the Asian Development Bank on
the grounds that part of the loan was destined for water projects in Arunachal Pradesh, the state
that includes Tawang. It was the first time China had sought to influence the territorial dispute
through a multilateral institution. Then the governor of Arunachal Pradesh announced that the Indian
military was deploying extra troops and fighter jets in the area.

The growing belligerence has soured relations between the two Asian giants and has prompted one
Indian military leader to declare that China has replaced Pakistan as India's biggest threat.

Economic progress might be expected to bring the countries closer. China and India did \$52 billion
worth of trade last year, a 34 percent increase over 2007. But businesspeople say border tensions
have infused business deals with official interference, damping the willingness of Chinese and
Indian companies to invest in each other's countries.

``Officials start taking more time, scrutinizing things more carefully, and all that means more
delays and ultimately more denials, `` said Ravi Bhoothalingam, a former president of the Oberoi
Group, the luxury hotel chain, and a member of the Institute of Chinese Studies in New Delhi.
``That's not good for business.''

The roots of the conflict go back to China's territorial claims to Tibet, an enduring source of
friction between China and many foreign nations. China insists that this section of northeast India
has historically been part of Tibet, and should be part of China.

Tawang is a thickly forested area of white stupas and steep, terraced hillsides that is home to the
Monpa people, who practice Tibetan Buddhism, speak a language similar to Tibetan and once paid
tribute to rulers in Lhasa. The Sixth Dalai Lama was born here in the 17th century. The Chinese Army
occupied Tawang briefly in 1962, during a war with India fought over this and other territories
along the 2,521-mile border.

More than 3,100 Indian soldiers and 700 Chinese soldiers were killed and thousands wounded in the
border war. Memorials here highlighting Chinese aggression in Tawang are big draws for Indian
tourists.

``The entire border is disputed,'' said Ma Jiali, an India scholar at the China Institutes of
Contemporary International Relations, a government-supported research group in Beijing. ``This
problem hasn't been solved, and it's a huge barrier to China-India relations.''

In some ways, Tawang has become a proxy battleground, too, between China and the Dalai Lama, the
exiled spiritual leader of the Tibetans, who passed through this valley when he fled into exile in
1959. From his home in the distant Indian hill town of Dharamsala, he wields enormous influence over
Tawang. He appoints the abbot of the powerful monastery and gives financial support to institutions
throughout the area. Last year, the Dalai Lama announced for the first time that Tawang is a part of
India, bolstering the India's territorial claims and infuriating China.

Traditional Tibetan culture runs strong in Tawang. One morning in June, the monastery held a
religious festival that drew hundreds from the nearby villages. As red-robed monks chanted sutras,
blew horns and swung incense braziers in the monastery courtyard, the villagers jostled each other
to be blessed by the senior lamas.

At the monastery, an important center of Tibetan learning, monks express rage over Chinese rule in
Tibet, which the Chinese Army seized in 1951.

``I hate the Chinese government,'' said Gombu Tsering, 70, a senior monk who watches over the
monastery's museum. ``Tibet wasn't even a part of China. Lhasa wasn't a part of China.''

Few expect China to try to annex Tawang by force, but military skirmishes are a real danger,
analysts say. The Indian military recorded 270 border violations and nearly 2,300 instances of
``aggressive border patrolling'' by Chinese soldiers last year, said Brahma Chellaney, a professor
of strategic studies at the Center for Policy Research, a research organization in New Delhi.
Mr.~Chellaney has advised the Indian government's National Security Council.

``The India-China frontier has become more `hot' than the India-Pakistan border,'' he said in an
e-mail message.

Two years ago, Chinese soldiers demolished a Buddhist statue that Indians had erected at Bumla, the
main border pass above Tawang, a member of the Indian Parliament, Nabam Rebia, said in a session of
Parliament.

Tawang became part of modern India when Tibetan leaders signed a treaty with British officials in
1914 that established a border called the McMahon Line between Tibet and British-run India. Tawang
fell south of the line. The treaty, the Simla Convention, is not recognized by China.

``We recognize it because we agreed to it,'' said Samdhong Rinpoche, prime minister of the Tibetan
government-in-exile. ``If China agreed to it now, it would be a recognition of the power of the
Tibet government at that time.''

China has grown increasingly hostile to the Dalai Lama after severe ethnic unrest in Tibet in 2008.
This year, it turned its diplomatic guns on India over the Tawang issue. China moved in March to
block a \$2.9 billion loan to India from the Asian Development Bank, a multination group based in
Manila that has China on its board, because \$60 million of the loan had been earmarked for
flood-control projects in Arunachal Pradesh. The loan was approved in mid-June over China's heated
objections.

``China expresses strong dissatisfaction to the move, which can neither change the existence of
immense territorial disputes between China and India, nor China's fundamental position on its border
issues with India,'' Qin Gang, the Foreign Ministry spokesman, said in a written statement.

In May, weeks after China first tried to block the loan, the chief of the Indian Air Force, Air
Chief Marshal Fali Homi, now retired, told a prominent Indian newspaper that China posed a greater
threat than Pakistan.

Another official, J.~J.~Singh, the governor of Arunachal Pradesh and a retired chief of the Indian
Army, said the next month that the Indian military was adding two divisions of troops, totaling
50,000 to 60,000 soldiers, to the border region over the next several years. Four Sukhoi fighter
jets were immediately deployed to a nearby air base.

Since 2005, when Prime Minister Wen Jiabao of China visited India, the two countries have gone
through 13 rounds of bilateral negotiations over the issue. A round was held just last month, with
no results.

``The China-India border has got to be one of the most continuously negotiated borders in modern
history,'' said M.~Taylor Fravel, an associate professor of political science at the Massachusetts
Institute of Technology who is a leading expert on China's borders. ``That shows how intractable
this dispute is.''

\section{China Web Sites Seeking Users' Names}

\lettrine{N}{ews} Web sites in China, complying with secret government orders, are requiring that
new users log on under their true identities to post comments, a shift in policy that the country's
Internet users and media have fiercely opposed in the past.

Until recently, users could weigh in on news items on many of the affected sites more anonymously,
often without registering at all, though the sites were obligated to screen all posts, and the posts
could still be traced via Internet protocol addresses.

But in early August, without notification of a change, news portals like Sina, Netease, Sohu and
scores of other sites began asking unregistered users to sign in under their real names and
identification numbers, said top editors at two of the major portals affected. A Sina staff member
also confirmed the change.

The editors said the sites were putting into effect a confidential directive issued in late July by
the State Council Information Office, one of the main government bodies responsible for supervising
the Internet in China.

The new step is not foolproof, the editors acknowledged. It was possible for a reporter to register
successfully on several major sites under falsified names and ID and cellphone numbers.

But the requirement adds a critical new layer of surveillance to mainstream sites in China, which
were already heavily policed. Further regulations of the same nature also appeared to be in the
pipeline.

And while the authorities called the measure part of a drive to forge greater ``social
responsibility'' and ``civility'' among users, they moved forward surreptitiously and suppressed
reports about it, said the editors and others in the media industry familiar with the measure, who
spoke on condition of anonymity to avoid putting their jobs at risk.

Asked why the policy was pushed through unannounced, the chief editor of one site said, ``The
influence of public opinion on the Net is still too big.''

Government Internet regulators have been trying to usher in real-name registration controls since
2003, when they ordered Internet cafes around China to demand that customers show identification,
nominally to keep out minors. Last year, lawmakers and regulators began discussing legislation on a
more extensive ``real name system,'' as it is known.

But such proposals have aroused heated debate over the purview of the state to restrict China's
online community, which is the largest in the world at about 340 million people and growing.

Proponents, led by officials and state-connected academics in the information security field, argue
that mandatory controls are necessary to help subdue inflammatory attacks, misinformation and other
illegal activity deemed to endanger social order. They often note registration requirements on large
sites in South Korea to support their point.

Critics counter that government regulation represents an incursion on free speech, individual
privacy and the watchdog role of the Web in China.

The critics say sites and users should retain the right to discipline themselves. Given the
country's huge population of Internet users and its failure to guarantee freedom of expression, they
argue, the case of China is hardly analogous to that of South Korea.

In 2006, Internet users and the news media rebuffed one official proposal to require real-name
registration on blog hosting sites. Star bloggers denounced the notion, while ordinary users
overwhelmingly rejected it in surveys conducted on sites like Sina.

In another key test of the policy earlier this year, the legislature in Hangzhou, near Shanghai,
passed a regulation that would have placed the requirement on users who comment, blog or play games
on sites based there. Amid a popular outcry, however, the city shied away from enforcing the
regulation.

Central authorities have gone to new lengths to tame online activity in 2009, a year peppered with
politically delicate anniversaries.

Government censors have closed thousands of sites in a continuing war on ``vulgarity,'' closed
liberal forums and blogs for spreading ``harmful information,'' blocked access to YouTube, Facebook
and Twitter, and cut off Internet service where serious unrest has erupted, notably in the Xinjiang
region of the west after deadly clashes between ethnic Uighurs and Han in July. Increasingly,
officials have defended the Web shutdowns on the grounds of national security.

The government recently set off an international furor when it ordered that all computers sold in
China come prepackaged with pornography filtering software that authorities could remotely control.
Officials were forced to retreat from the order after international companies and trade bodies
protested and Chinese hackers showed that the software was designed to block politically offensive
content as well.

The authorities had aimed to avoid a similar showdown over the new real-name requirement. ``We had
no recourse to challenge it,'' said the news editor of another portal.

Ta Kung Pao, a Hong Kong-based newspaper loyal to Beijing, first leaked news of the State Council
edict in late July. But the report was scrubbed from the paper's Web site within a few days.

Another state newspaper tried to follow up on the Ta Kung Pao report soon thereafter, the paper's
editors said, but they were forced to abort their article because they were warned that the order
was a state secret.

The State Council Information Office had yet to respond to a list of submitted questions about the
move.

The new mandate did not appear to affect formerly registered users of the portals. Nor did it affect
blog hosts, forums or government news sites like People's Daily or Xinhua.

Whether because it had an impact mainly on rookie users or because of the void of news about it,
bloggers in China were unusually slow to recognize the measure. But those who did were critical.

One commentator on the popular forum Tianya wrote, ``Not daring to write one's real name, in truth,
is a form of self-protection for the weak.''

There were signals in the state media in recent weeks that more name registration measures would
follow.

An influential advocate of the policy, Fang Bingxing, the president of Beijing University of Posts
and Telecommunications, told a forum in August that the ``time was ripe'' to roll it out widely to
bolster information security, newspapers reported.

A trail of comments on Sina thrashed the report.

Late last month, the Communist Party-run Guangming Daily ran a positive story about a city
government portal in western China that imposed the requirement on new bloggers, calling it a
``forerunner.''

Hu Yong, a new media specialist at Peking University, said government-enforced registration
requirements carried long-term side effects.

``Netizens will have less trust in the government, and to a certain extent, the development of the
industry will be impeded,'' he said.

From a comparison of the most commented-on articles in July and August on a number of portals it was
hard to determine whether the volume of posts had been affected so far.

But both editors at two of the major portals affected said their sites had shown marked drop-offs.

\section{China Ousts\cite{oust} Top Official After Protests}

\lettrine{T}{he} top Communist official in Urumqi in western China was dismissed on Saturday as a
large deployment of the military police appeared to have brought a measure of peace to the city
after two days of large street protests.

Li Zhi, the party secretary of Urumqi, lost his post, the official Xinhua news agency reported on
Saturday evening. He became the most senior person to be removed since ethnic tensions erupted there
in rioting in July.

Beijing officials also sent to Urumqi a special medical inspection unit from the People's Liberation
Army to investigate reports that people had been stabbed with needles.

It is somewhat unusual for China's leaders to replace a senior local official so quickly after
protests -- in this case, while large deployments of armed police officers are still blocking
intersections in Urumqi and most shops are still closed. The Beijing leadership has often sought to
avoid giving the impression of giving in to public pressure.

The removal of Mr.~Li ``shows that Xinjiang is viewed as a strategic region where there cannot be
the kind of social protests we have seen in recent days,'' said Nicholas Bequelin, a senior
researcher at Human Rights Watch.

The latest protests were notable for including large crowds of people who specifically called on
Friday for the removal of Mr.~Li's boss, Wang Lequan, the powerful party secretary of restive
Xinjiang region, of which Urumqi is the capital. Mr.~Wang, a member of the Politburo believed to be
a close ally of President Hu Jintao, has run the nominally autonomous region for 15 years and is
famous within China for taking a hard line toward minorities.

``They want to protect Wang Lequan, because firing a Politburo member would send a message they do
not want to send,'' namely that hard-line policies toward ethnic minorities can be questioned, said
Li Cheng, the research director of the Brookings Institution's China Center.

The Urumqi party chief seems to have been sacrificed partly to protect Mr.~Wang, but also because he
proved incapable of keeping order.

``Usually if a local leader could not deal with two crises that occur in a short period of time,
that leader will be fired,'' Li Cheng said.

Also forced out on Saturday was the director of Xinjiang's public security department, Liu Yaohua.
The party chief of the region's Aksu Prefecture, Zhu Changjie, replaced him.

The crowds in Urumqi last week have been accusing Mr.~Wang and his aides of not being tough enough.
The violence last week involves renewed tensions between Uighurs, the dominant ethnic group in
Xinjiang, and Han Chinese, the dominant ethnic group for China over all.

Han have been moving in large numbers to Xinjiang since the 1960s, occupying many of the best jobs;
according to the authorities, most of the victims in the rioting in July, in which at least 197
people died, were Han.

Tens of thousands of Han protesters took to the streets on Thursday as word spread that young
Uighurs had reportedly been stabbing people with needles. Zhang Hong, the deputy mayor of Urumqi,
announced Friday that the unrest on Thursday had resulted in five deaths with 14 people injured,
Xinhua said.

Friday's protests were much smaller and more peaceful than Thursday's.

Through Thursday, 531 people had gone to hospitals claiming to have been stabbed, and 106 of them
had physical wounds consistent with having suffered such a stabbing, Xinhua said. Some Han have
voiced worry that the needles could carry H.I.V.

Before his fall from power, Mr.~Li had been the most visible Chinese official seeking to respond to
the unrest in his city last week. He had raced to a succession of gathering places and tried to
persuade residents to return home.

Zhu Hailun, the secretary of the Communist Party's Xinjiang legislative and political affairs
committee, will replace Mr.~Li.

A troubled quiet returned on Saturday to Urumqi, residents said in telephone interviews. ``The
military police are everywhere, and they have a very firm hold on the city,'' said a taxi driver who
gave his family name, Li.

Most of the downtown area, the scene of some of the largest protests, was blocked off to traffic.

``There's no one out there, every store is still closed,'' said a downtown resident who gave only
her family name, Liu.

\section{China Says Five Dead in Latest Xinjiang Unrest}

\lettrine{A}{t} least 5 people have been killed and 14 wounded during days of protests by ethnic Han
demanding a government clampdown on Uighurs in the capital of the restive western region of
Xinjiang, according to the state news agency, which quoted the city's deputy mayor.

The deputy mayor, Zhang Hong, gave the casualty toll at a news conference late Friday night in
Urumqi, the regional capital. The news agency, Xinhua, did not report the ethnicities of the victims
and had no other details.

Urumqi remained paralyzed Friday as Han protesters continued to rally in parts of the city to call
for greater security measures, forcing the government to close schools, offices and shops to keep
people off the streets and to block traffic from the city center. The police used tear gas to
disperse protesters, and loud booms from the firing of gas canisters could be heard from the city
center, witnesses said.

Helicopters were seen swooping overhead in the morning, a day after tens of thousands of Han, the
dominant ethnic group in China, took to the streets to criticize the government for failing to
prevent what they believe to be needle-stabbing attacks by Uighurs, according to an earlier report
by Xinhua. Some of the Han said the needles were infected with H.I.V.

Tensions between Uighurs, a Turkic-speaking people who are largely Muslim, and Han have remained
high since rioting broke out in early July, in China's worst ethnic violence in decades. At least
197 people were killed and 1,721 wounded, according to the government, and most of the dead were
Han.

The violence brought to the surface the resentments that have long coursed through society in
Xinjiang. Uighurs complain that the ruling Han discriminate against them, and Han say the Uighurs
are ungrateful for the benefits that Han governance has brought.

On Friday, groups of protesters continued to gather in front of the offices of the regional
government to demand the resignation of Wang Lequan, the top Communist Party official in Xinjiang,
an oil-rich and ethnically fractious region that makes up one-sixth of China. The Han protesters
said that Mr.~Wang had failed to make the city safe for Han residents after the July riots.
Mr.~Wang, a proponent of hard-line ethnic minority policies who has ruled Xinjiang for 15 years, is
a member of the Communist Party's Politburo, and the outcry is a rare attack against a member of the
country's governing elite.

Talk of the hypodermic needle attacks spread quickly in recent days, fueled by official news reports
on the assaults. The local authorities told Chinese reporters that the police had arrested 21
suspects. Then, as outrage poured into the streets, starting with a small protest Wednesday, the
government deployed more security forces around the city and tried to delete mentions of the needle
attacks from the Internet.

Li Zhi, the party secretary of Urumqi, was at various sites on Thursday trying to calm protesters.
Mr.~Li said that someone had tried to attack a woman with a needle on Wednesday and was detained by
the police, according to Xinhua.

``The current situation is that a very few hostile elements are maliciously sabotaging the capital's
stability,'' he said. ``Their goal is to create ethnic estrangement and incite ethnic opposition.''

On Friday, groups of angry Han young men shoved against lines of police officers to try to get to
the main Uighur quarter. ``The situation is even more tense than it was yesterday,'' a Han driver in
Urumqi said Friday. ``The entire downtown has been cut off from incoming traffic. There are no buses
running, either.''

The driver, who spoke on the condition that his name not be used, said he had seen a Han mob beating
a Uighur woman at 6 p.m. Thursday. ``These people are just so angry,'' he said.

\section{Joseph Kennedy Will Not Seek Uncle's Senate Seat}

Joseph P.~Kennedy II announced Monday that he would not seek the Senate seat that his uncle Edward
M.~Kennedy held for almost 50 years, probably ending for now the family's long and prominent role in
Massachusetts politics.

Mr.~Kennedy, 56, said in a statement that ``after much consideration'' he had decided instead to
keep running the Citizens Energy Corporation, a nonprofit company he founded in 1979 to provide
heating oil to the poor. His decision clears the way for other Democrats to enter the special
election for the Senate seat in the coming days.

``Given all that my uncle accomplished, it was only natural to consider getting back involved in
public office,'' Mr.~Kennedy, the eldest son of Robert F.~Kennedy, said in the statement, which was
posted Monday afternoon on the Citizens Energy Web site.

``My father called politics an honorable profession, and I have profound respect for those who
choose to advance the causes of social and economic justice in elective office,'' Mr.~Kennedy said.
``After much consideration, I have decided that the best way for me to contribute to those causes is
by continuing my work at Citizens Energy Corporation.''

Mr.~Kennedy would not have been a newcomer to politics; he was a member of the House of
Representatives from 1987 to 1999, occupying the same seat another famous uncle, John F.~Kennedy,
won in 1946.

Joseph Kennedy was a popular congressman, focusing on low-cost housing and veterans' health care and
easily winning re-election. But after his first wife publicly fought his effort to annul their
marriage and his brother Michael was killed in a skiing accident in 1997, Mr.~Kennedy decided to
leave office.

For the previous two years, he had struggled with the fallout of charges of improper sexual
relations between Michael and a teenage baby sitter. Michael had been his campaign manager and had
run his company.

Friends said that since his uncle's death on Aug.~25, Mr.~Kennedy had been agonizing over whether to
expose his second wife and twin sons, now 28, to the rigors of a campaign and the questions and
attacks that would surely come with it.

Only one Democrat -- Martha Coakley, the state's attorney general -- has entered the Senate race,
but several others will probably jump in now that Mr.~Kennedy is not running. Representative Stephen
F.~Lynch took the first step last week, picking up nominating papers from the secretary of state's
office for the election, scheduled for Jan.~19.

Other possible Democratic candidates include Representatives Michael E.~Capuano and Edward
J.~Markey, and Martin T.~Meehan, a former congressman who is chancellor of the University of
Massachusetts, Lowell. Mr.~Kennedy's widow, Victoria Reggie Kennedy, was considered a strong
candidate, but family friends say she has ruled out running. A primary is set for Dec.~8.

On the Republican side, several state legislators are considering a run. Also, Curt Schilling, a
former Red Sox pitcher and an independent who has campaigned for Republicans, has expressed mild
interest.

It remains unclear whether the state legislature will honor a request that Edward Kennedy made just
before his death: to change the law and allow Gov. Deval Patrick, a Democrat, to appoint an interim
successor so the seat would not remain empty until the special election. A public hearing on the
proposed change is scheduled for Wednesday.

\section{In Michigan, the Disappearing State Fair}

Since 1849, Michigan residents have attended the annual state fair, where urban, suburban and farm
dwellers can pet Clydesdales, attend rock concerts and watch heifers being born. Barring last-minute
support, however, Michigan's State Fair, the nation's oldest, held its last event when it ended its
11-day run Monday.

Facing an approximate \$1.8 billion state budget deficit, Gov. Jennifer M.~Granholm proposed in
February ending fair financing after this year, saying the state could no longer absorb its
operating losses. Attendance has dropped in recent years, and corporate support has dwindled.

Pummeled by the auto industry's economic woes, Michigan had a 15.2 percent unemployment rate, the
nation's highest, as of July 17.

``We're facing very, very difficult decisions here in Michigan, and at a time of declining revenues,
we cannot be all things to all people,'' Ms.~Granholm's spokeswoman, Liz Boyd, said on Monday.

Under a bill before the House, the state would continue to cover any fair deficit up to an
unspecified amount. Ms.~Boyd would not say what the governor would do if the bill passed the
Legislature.

Last year's fair losses came to \$362,000, the former fair manager, Steve Jenkins, told The Detroit
Free Press. Mr.~Jenkins, who resigned his post after Ms.~Granholm's announcement, led a weekend
rally to keep the fair open.

Other efforts are afoot to keep the fair running. Members of the Michigan State Fair Authority, the
fair's governing body, have asked state officials to consider transferring operations to a nonprofit
organization. The authority's chairman, Mark Gaffney, said he was optimistic that some private
entity would create a redevelopment plan for the fairgrounds that would include keeping the fair
open. ``We just need more time,'' he said.

Ms.~Boyd said that while Ms.~Granholm was open to private efforts to continue the fair, she also
recognized the value of the property on which it is located -- a 164-acre site along Woodward Avenue
near Detroit's northern border, where the fair has been held since 1905. ``It's property that could
be developed to create jobs in southeast Michigan,'' Ms.~Boyd said.

Jaynemarie Hubanks, the finance director for suburban Ferndale, said she had gone to 50 consecutive
fairs and particularly enjoyed the crafts and agricultural products. This year's visit was
especially poignant. ``I think it's tragic,'' said Ms.~Hubanks, 51. ``It's just a huge loss.''

Mr.~Gaffney put it another way. ``Kids here grow up in the city and never see a cow,'' he said.
``Here, they not only can see a cow, they can milk it.''

\section{Where Did All the Flowers Come From?}

Throughout his life, Charles Darwin surrounded himself with flowers. When he was 10, he wrote down
each time a peony bloomed in his father's garden. When he bought a house to raise his own family, he
turned the grounds into a botanical field station where he experimented on flowers until his death.
But despite his intimate familiarity with flowers, Darwin once wrote that their evolution was ``an
abominable\cite{abominable} mystery.''

Darwin could see for himself how successful flowering plants had become. They make up the majority
of living plant species, and they dominate many of the world's ecosystems, from rain forests to
grasslands. They also dominate our farms. Out of flowers come most of the calories humans consume,
in the form of foods like corn, rice and wheat. Flowers are also impressive in their sheer diversity
of forms and colors, from lush\cite{lush}, full-bodied roses to spiderlike orchids\cite{orchid} to
calla lilies shaped like urns.

The fossil record, however, offered Darwin little enlightenment about the early evolution of
flowers. At the time, the oldest fossils of flowering plants came from rocks that had formed from
100 million to 66 million years ago during the Cretaceous period.
Paleontologists\cite{paleontologist} found a diversity of forms, not a few primitive forerunners.

Long after Darwin's death in 1882, the history of flowers continued to vex\cite{vex} scientists. But
talk to experts today, and there is a note of guarded optimism. ``There's an energy that I haven't
seen in my lifetime,'' said William Friedman, an evolutionary biologist at the University of
Colorado, Boulder.

The discovery of new fossils is one source of that new excitement. But scientists are also finding a
wealth of clues in living flowers and their genes. They are teasing apart the recipes encoded in
plant DNA for building different kinds of flowers. Their research indicates that flowers evolved
into their marvelous diversity in much the same way as eyes and limbs have: through the recycling of
old genes for new jobs.

Until recently, scientists were divided over how flowers were related to other plants. Thanks to
studies on plant DNA, their kinship\cite{kinship} is clearer. ``There was every kind of idea out
there, and a lot of them have been refuted\cite{refute},'' said James A.~Doyle, a paleobotanist at
the University of California, Davis.

It is now clear, for example, that the closest living relatives to flowers are flowerless species
that produce seeds, a group that includes pine trees and gingkos. Unfortunately, the plants are all
closely related to one another, and none is more closely related to flowers than any of the others.

The plants that might document the early stages in the emergence of the flower apparently became
extinct millions of years ago. ``The only way to find them is through the fossils,'' Dr.~Doyle said.

In the past few years scientists have pushed back the fossil record of flowers to about 136 million
years ago. They have also found a number of fossils of mysterious extinct seed plants, some of which
produce seeds in structures that look faintly like flowers. But the most intriguing\cite{intrigue}
fossils are also the most fragmentary, leaving paleobotanists deeply divided over which of them
might be closely related to flowers. ``There's no consensus\cite{consensus},'' Dr.~Doyle said.

But there is a consensus when it comes to the early evolution of flowers themselves. By studying the
DNA of many flowering plants, scientists have found that a handful of species represent the oldest
lineages\cite{lineage} alive today. The oldest branch of all is represented by just one species: a
shrub called Amborella that is found only on the island of New Caledonia in the South Pacific. Water
lilies and star anise represent the two next-oldest lineages alive today.

If you could travel back to 130 million years ago, you might not be impressed with the earliest
flowers. ``They didn't look like they were going anywhere,'' Dr.~Doyle said.

Those early flowers were small and rare, living in the shadows of far more successful nonflowering
plants. It took many millions of years for flowers to hit their stride. Around 120 million years
ago, a new branch of flowers evolved that came to dominate many forests and explode in diversity.
That lineage includes 99 percent of all species of flowering plants on Earth today, ranging from
magnolias to dandelions to pumpkins. That explosion in diversity also produced the burst of flower
fossils that so puzzled Darwin.

All flowers, from Amborella on, have the same basic anatomy. Just about all of them have petals or
petal-like structures that surround male and female organs. The first flowers were probably small
and simple, like modern Amborella flowers.

Later, in six lineages, flowers became more complicated. They evolved an inner ring of petals that
became big and showy, and an outer ring of usually green, leaflike growths called sepals, which
protect young flowers as they bud.

It would seem, based on this recent discovery, that a petal is not a petal is not a petal. The
flowers of, say, the paw-paw tree grow petals that evolved independently from the petals on a rose.
But the genes that build flowers hint that there is more to the story.

In the late 1980s, scientists discovered the first genes that guide the development of flowers. They
were studying a small plant called Arabidopsis, a botanical lab rat, when they observed that
mutations\cite{mutation} could set off grotesque changes. Some mutations caused petals to grow where
there should have been stamens, the flower's male organs. Other mutations transformed the inner
circle of petals into sepals. And still other mutations turned sepals into leaves.

The discovery was a remarkable echo of ideas first put forward by the German poet Goethe, who not
only wrote ``Faust'' but was also a careful observer of plants.

In 1790, Goethe wrote a visionary essay called ``The Morphology\cite{morphology} of Plants,'' in
which he argued that all plant organs, including flowers, started out as leaves. ``From first to
last,'' he wrote, ``the plant is nothing but a leaf.''

Two centuries later, scientists discovered that mutations to genes could cause radical
transformations like those Goethe envisioned. In the past two decades, scientists have investigated
how the genes revealed through such mutations work in normal flowers. The genes encode proteins that
can switch on other genes, which in turn can turn other genes on or off. Together, the genes can set
off the development of a petal or any other part of an Arabidopsis flower.

Scientists are studying those genes to figure out how new flowers evolved. They have found versions
of the genes that build Arabidopsis flowers in other species, including Amborella. In many cases,
the genes have been accidentally duplicated in different lineages.

Finding those flower-building genes, however, does not automatically tell scientists what their
function is in a growing flower. To answer that question, scientists need to tinker\cite{tinker}
with plant genes. Unfortunately, no plant is as easy to tinker with as Arabidopsis, so answers are
only beginning to emerge.

Vivian Irish, an evolutionary biologist at Yale, and her colleagues are learning how to manipulate
poppies\cite{poppy} because, Dr.~Irish points out, ``poppies evolved petals independently.''

She and her colleagues have identified flower-building genes by shutting some of them down and
producing monstrous flowers as a result.

The genes, it turns out, are related to the genes that build Arabidopsis flowers. In Arabidopsis,
for example, a gene called AP3 is required to build petals and stamens. Poppies have two copies of a
related version of the gene, called paleoAP3. But Dr.~Irish and her colleagues found that the two
genes produced different effects. Shutting down one gene transforms petals. The other transforms
stamens.

The results, Dr.~Irish said, show that early flowers evolved a basic tool kit of genes that marked
off different regions of a stem. Those geography genes made proteins that could then switch on other
genes involved in making different structures. Over time, the genes could switch control from one
set of genes to another, giving rise to new flowers.

Thus, the petals on a poppy evolved independently from the petals on Arabidopsis, but both flowers
use the same kinds of genes to control their growth.

If Dr.~Irish is right, flowers have evolved in much the same way our own anatomy evolved. Our legs,
for example, evolved independently from the legs of flies, but many of the same ancient
appendage-building genes were enlisted to build those different limbs.

``I think it is pretty cool that animals and plants have used similar strategies,'' Dr.~Irish said,
``albeit with different genes.''

Dr.~Irish said, however, that her studies of petals were only part of the story. ``Lots of things
happened when the flower arose,'' she said. Flowers evolved a new arrangement of sex organs, for
example. ``A pine tree has male cones and female cones,'' she said, ``but flowers have male and
female organs on the same axis.''

Once the sex organs were gathered together, they underwent a change invisible to the naked eye that
might have driven flowers to their dominant place in the plant world.

When a pollen grain fertilizes an egg, it provides two sets of DNA. While one set fertilizes the
egg, the other is destined for the sac that surrounds the egg. The sac fills with endosperm, a
starchy material that fuels the growth of an egg into a seed. It also fuels our own growth when we
eat corn, rice or other grains.

In the first flowers, the endosperm ended up with one set of genes from the male parent and another
set from the female parent. But after early lineages like Amborella and water lilies branched off,
flowers bulked up their endosperm with two sets of genes from the mother and one from the father.

Dr.~Friedman, of the University of Colorado, Boulder, has documented the transition and does not
think it was a coincidence that flowering plants underwent an evolutionary explosion after gaining
an extra set of genes in their endosperm. It is possible, for example, that with extra genes, the
endosperm could make more proteins.

``It's like having a bigger engine,'' Dr.~Friedman said.

Other experts agree that the transition took place, but they are not sure it is the secret to
flowers' success. ``I don't know why it should be so great,'' Dr.~Doyle said.

As Dr.~Friedman has studied how the extra set of genes evolved in flowers, he has once again been
drawn to Goethe's vision of simple sources and complex results.

Flowers with a single set of female DNA in their endosperm, like water lilies, start out with a
single nucleus at one end of the embryo sac. It divides, and one nucleus moves to the middle of the
sac to become part of the endosperm.

Later, a variation evolved. In a rose or a poppy, a single nucleus starts out at one end of the sac.
But when the nucleus divides, one nucleus makes its way to the other end of the sac. The two nuclei
each divide, and then one of the nuclei from each end of the sac moves to the middle.

Duplication, a simple process, led to greater complexity and a major change in flowers.

``Nature just doesn't invent things out of whole cloth,'' Dr.~Friedman said. ``It creates novelty in
very simple ways. They're not radical or mysterious. Goethe already had this figured out.''

\section{Birth Order: Fun to Debate, but How Important?}

The older girl was smart, neat and perfectly behaved in school; in her spare time, she won dance
trophies\cite{trophy}. At every checkup, her mother would tell me what a good girl she was.

She is the oldest, her mother would say, so she gets lots of attention, and she works very hard.
When her younger sister turned out to be an equally good student, the proud mother explained that
naturally she wanted to be just like her older sister.

Then a long-looked-for baby boy was born. When he was a toddler, I began to worry that his speech
seemed a little slow in coming. His mother was perfectly calm about it. He is the only boy, she
said, so he gets lots of attention, and he doesn't have to work very hard.

Everyone takes it personally when it comes to birth order. After all, everyone is an oldest or a
middle or a youngest or an only child, and even as adults we revert almost inevitably to a joke or
resentment or rivalry that we've never quite outgrown.

Children and parents alike are profoundly affected by the constellations of siblings; it is said
that no two children grow up in the same family, because each sibling's experience is so different.

But that doesn't mean the effects of birth order are as clear or straightforward as we sometimes
make them sound. Indeed, birth order can be used to explain every trait and its precise opposite.
I'm competitive, driven -- typical oldest child! My brother, two years younger, is even more
competitive, more driven -- typical second child, always trying to catch up!

I surveyed some experienced pediatricians about when parents are likely to bring up birth order.
Many cited the issue of speech, especially when a second child doesn't talk as well or as early as
the first.

And parents are likely to talk about mistakes they think they made the first time around. This time,
we're going to solve the sleep thing good and early. This time, we're going to get it right with
potty training. This time, we're going to sign the child up for soccer.

``Too many parents are haunted by experiences both good and bad that they identify with their birth
order,'' said Dr.~Peter A.~Gorski, a professor of pediatrics, public health and psychiatry at the
University of South Florida. And that might lead them to classify their own children according to
birth order, he went on, which in turn can lead to a sense of identification or even rejection and
to ``self-fulfilling prophecies.''

Frank J.~Sulloway, a visiting scholar at the University of California, Berkeley, and the author of
``Born to Rebel: Birth Order, Family Dynamics and Creative Lives'' (Pantheon, 1996), points out that
second-born children tend to be exposed to less language than eldest children. ``The best
environment to grow up in is basically two parents who are chattering away at you with fancy
words,'' Dr.~Sulloway said.

He cited a huge and well-publicized Norwegian study, published in 2007, which found that eldest
siblings' I.Q.'s averaged about three points higher than their younger brothers'. (The study made
use of Norwegian military records, so all the subjects were male.)

Those differences in verbal stimulation, like the differences in I.Q., are ``relatively modest,''
Dr.~Sulloway continued, and unlikely to result in clinical speech delays. But in a child who is
already vulnerable, a child who may be temperamentally less likely to evoke adults' attention, or a
child growing up in a less stimulating home -- well, then, being the second child might be the added
risk that makes the difference, he said.

``Birth order doesn't cause anything,'' Dr.~Sulloway said. ``It's simply a proxy for the actual
mechanisms that go on in family dynamics that shape character and personality.''

We can all cite examples and counterexamples, from our own families, our friends, history and
literature. There are plenty of families where the younger child is brighter or more academic, and
plenty of literary and historical examples (Jane and Elizabeth Bennet, Meg and Jo March, Dmitri and
Ivan Karamazov -- and you can think about those authors and their older siblings as well, and draw
any comparisons you like). And then there are plenty of examples of brilliant eldest siblings, but
given my own eldest status, I will refrain from citing any. (I told you this always gets personal.)

I.Q., though it does grab headlines, may shape family life less than personality and temperament.
``It's a part of a bigger picture that really involves family dynamics,'' Dr.~Sulloway said. ``Child
and family dynamics is like a chessboard; birth order is like a knight.''

Then there are all the other influences, from family size to socioeconomic status. ``Typically
firstborns tend to boss their younger siblings around, but what if you're a very shy person?''
Dr.~Sulloway said. ``Napoleon was a second-born and his older brother was a very shy guy, and he
usurped the older-sibling niche because his older sibling didn't occupy that niche.

``And why didn't he occupy that niche? Temperament.''

Now, of course birth order played into my patients' patterns, but so did gender and birth spacing
and, above all, temperament. That little boy was more even-tempered, more placid than either of his
sisters, easily soothed, and I think he would have shown that temperament no matter what.

But temperament also helped define his relationship to the four larger people in his immediate
circle. ``I wouldn't discount the impact of birth order,'' Dr.~Gorski told me. ``It sets up the
structure of one's place in relation to others from the beginning, as we learn how to react to
people of different ages and different relationships.''

Pediatricians are always being warned not to let a speech delay slip past because of parents'
beliefs that boys talk later or that youngest children talk later. I did eventually insist on a
hearing test and speech therapy for this little boy. As it turned out, his hearing was fine, and his
sisters drilled him over and over with ``use your words'' exercises until his speech improved. That
is one of the advantages of having hardworking older sisters.

\section{Leaving Google, Executive to Back Chinese Start-Ups}

Three days after announcing that he was resigning as the head of Google's Chinese operations, Kai-Fu
Lee said on Monday that he had raised \$115 million to create a company that would finance high-tech
start-ups in China.

The company, Innovation Works, will search for talented Chinese engineers and entrepreneurs and help
them develop the next generation of Internet and mobile computing technologies, Mr.~Lee said in an
interview by telephone.

``We're going to collect the best ideas, and we're going to hire the best engineers and
entrepreneurs,'' he said. ``After one year, we'll send the companies into the open. If they get
venture capital funding, great; if they don't, they won't live.''

Innovation Works is being backed by the YouTube co-founder Steve Chen as well as the Foxconn
Technology Group, Legend Holdings, the New Oriental Education and Technology Group and the WI Harper
Group. Foxconn is one of the world's biggest electronics companies, and it manufactures products
including the Apple iPhone and Hewlett-Packard computers. Legend is the parent company of Lenovo,
the Chinese computer maker.

Mr.~Lee, 47, who was born in Taiwan and grew up in the United States, said he was leaving Google
after four years because he had simply found a new challenge.

He dismissed as ``outrageously wrong'' the notion that his departure had anything to do with
Google's difficulties in complying with China's strict censorship rules or with fierce competition
from China's leading search engine, Baidu.

``In my career, I've never run away from anything -- I run toward things,'' he said. ``There's clear
evidence Google has doubled its market share in China under my leadership. It's on the upturn.''

In a conference call with Peter Liu, chairman of WI Harper, one of the lead investors, Mr.~Lee said
part of the logic behind forming the company was the lack of early-stage financing of high-tech
start-ups in China from so-called angel investors who could provide financing and coaching of young
companies.

Talented managers now working for technology companies in China are primed to leave those companies
to become great entrepreneurs, he said.

``Our experience is in company building,'' he said. ``I can coach them with my many years of
experience.''

China is already a hotbed of high-tech growth, with fast-growing Internet companies like Baidu,
Tencent, Shanda Interactive Entertainment and Alibaba Group already worth billions of dollars.

Venture capital funds are aggressively scouting out new companies that can capture share in what is
already the world's largest Internet market, with over 300 million users.

Innovation Works recently signed a lease to rent part of a building in Beijing's high-tech district,
next to Google's offices, and the company will house its start-ups there.

Despite the economic downturn, Mr.~Liu of WI Harper said he had no trouble raising the money for
Innovation Works.

``We raised the money in less than 30 days,'' he said, and called the hiring of Mr.~Lee ``a
miracle.''

``He can identify the next Google in China, or the next Baidu,'' Mr.~Liu said.

\section{At NBC's Site for Women, a True Makeover}

The Web site iVillage is getting another makeover. This time, its executives say they are convinced
they have the right look for the site for women that has struggled to fit into the NBC Universal
world since the company paid a hefty\cite{hefty} \$600 million for it three years ago.

Changes to the site, which already got a less-cluttered home page and new logo in April, are to
begin appearing Wednesday, with the introduction of a new entertainment section, and to continue for
six months.

Overhauls of areas devoted to food, health, home, family and beauty are also planned, as is a
makeover of Astrology.com, which NBC Universal owns and cross-links with iVillage. The technology
underlying the site is also being rebuilt.

Explaining the need for what the company is calling a ``total site relaunch,'' Lauren Zalaznick,
president of the NBC division that oversees it, said, ``We had a very passionate and loyal audience
base who had been underserved for a good couple years.''

The changes, she said, will ``serve them original content in a vibrant modern aesthetic.''

The number of unique visitors to iVillage in July was 21.4 million, up 23 percent from a year
earlier, according to comScore numbers provided by NBC Universal. The site, an Internet content
pioneer, was founded in 1995.

As part of the coming changes, the number of topic areas on the busy site will be pared down, with
``Love,'' ``Go Green'' and ``Weddings'' stripped away (and in some cases integrated into the other
sections).

``We are changing our philosophy: we're not everything to everybody anymore,'' said Jodi Kahn, an
executive vice president at iVillage who is overseeing the overhaul. She and Ms.~Zalaznick declined
to discuss how much the company was spending on the revamp\cite{revamp}.

The approach is to focus on newsy articles and emphasize input from readers. In an interview in
iVillage's offices, Ms.~Kahn described it as ``very elegantly weaving the topical newsy content,
wrapped around conversation and really usable tools that help women manage their day-to-day life.''

The editorial team is being led by Angela Matusik, who was recently hired from People magazine's Web
site.

The new entertainment section, in addition to featuring big photos and easily searched video, will
have a prominent area called ``For Kids'' to help parents decide on appropriate viewing for
children. Other tools include a user-friendly recipe search engine, in development, and, in the
proposed health section, a broadly simplified way to search for treatments by symptom and age group.

If users create a personalized profile for one section of the site -- say, by signaling that they
are pregnant or have specific health concerns -- the personalization will follow them through all
the sections, like a menu planner in the food section, Ms.~Kahn said.

Advertisers will be given prominent placement and integrated content, like coupons that print with a
recipe.

Content from other NBC Universal units, like Oxygen, Bravo and Universal films, is also prominently
displayed, but that is not a corporate requirement. ``We are agnostic,'' said Ms.~Kahn, adding, ``We
do not have a mandate; we have the access.''

Synergy with the rest of the NBC Universal company has been a challenge since the site was acquired.
Two attempts at a daily iVillage TV show failed, and the company earlier acknowledged that some of
the expected benefits -- like a spike in traffic if the site was mentioned on NBC's ``Today'' show
-- never materialized.

But with the purchase of the cable network Oxygen, the company has been able to put together a more
comprehensive strategy for women, and advertisers have appreciated having a ``digital anchor'' in
iVillage for ad buys across several of the company's properties, Ms.~Zalaznick said, adding, ``We
are starting to be able to deliver value back out to the marketplace.''

``Nobody ever looks backward in an acquisition,'' she said. ``My job is to make that purchase price
look like a bargain in retrospect, and we're on our way to doing that.''

\section{Apple iPod Updates Expected, Jobs Show Uncertain}

Apple Inc's media extravaganza next week is being greeted with unusual calm, as it is uncertain if
charismatic CEO Steve Jobs will appear and investors have nothing much to bet on apart from new
iPods with cameras.

Apple's stock has often swung wildly and rumors have swirled about game-changing products ahead of
such events, but this time the stock has held steady in the two weeks leading up to the September 9
event.

Few expect any surprises, but the potential is always there, given the host's track record.

Apple has traditionally held a September gathering where master showman Jobs showed off new iPods
for the holiday season and trumpeted new content for iTunes.

Revamped\cite{revamp} iPods are again on tap this year. But much of the focus has been on whether
Jobs will make his first public appearance since taking medical leave in January, and what new
partners or services Apple may bring into its booming online store.

Rumors also continue to rumble that the Beatles catalog may finally be coming to iTunes, but such
speculation is not new and many are doubtful.

Some analysts think Jobs is unlikely to show up, as other members of Apple's management team have
become more visible in his absence.

``It would be great if he (Jobs) does appear, but I'm not expecting it,'' said Kaufman Bros analyst
Shaw Wu. ``Most investors have moved beyond that.''

``It shouldn't have a material impact on the stock if he doesn't,'' he said, noting that it could
edge higher.

Apple's stock, which is up more than 90 percent so far this year, has been relatively
tame\cite{tame} this week as expectations have been subdued. But analysts don't foresee a bold new
product announcement at the event, such as the long-rumored tablet device, that would immediately
boost the stock.

REFRESH

The iPod nano and touch models are widely expected to be updated with digital cameras, while the
iTouch may also get a video camera.

``A non-event,'' predicted Broadpoint Amtech analyst Brian Marshall. ``We're going to see the focus
on iPods and iTunes and then leave it at that.''

Apple's invitation for the event was typically cryptic, featuring a dancing figure, an iPod and the
tagline, ``It's only rock and roll, but we like it.''

Apple has acknowledged that its dominant iPod business is changing as growth has slowed. IPods
accounted for less a than one-fifth of the company's revenue last quarter.

``The key here is what Apple ultimately sets for pricing and storage capacity for these new iPods,''
said Wu.

Apple has said it expects sales of its traditional iPods such as the shuffle, nano and classic to
decline as they are cannibalized by the newer iPod touch and the iPhone.

Beyond iPods, many are still talking about Jobs, who joked openly about his health issues at last
year's event.

Jobs' last appearance at an Apple event was in October of 2008. He returned from a 6-month medical
leave in June, during which time he underwent a liver transplant.

In his absence, marketing chief Phil Schiller, who has led other events this year, may again be
master of ceremonies.

However, FTN Equity Capital analyst Bill Fearnley Jr.~said it is ``likely'' that Jobs will make an
appearance.

``There is no other fall event that we are aware of that would present a keynote-type address or
presentation opportunity for Apple or its CEO,'' he wrote in a note.

While an appearance by Jobs might dominate media coverage of the event, an announcement that songs
from the Beatles -- the most high-profile iTunes holdout -- were headed to the online store would
also be headline-grabbing.

However Tim Bajarin, president of consulting group Creative Strategies, called the Beatles chatter
``wishful thinking.''

``This rumor does come around every year at this time and I am afraid that we will be disappointed
again,'' he said via email.

\section{Sleep May Be Nature's Time Management Tool}

If there is a society of expert sleepers out there, a cult\cite{cult} of smug\cite{smug} snoozers
satisfied that they're getting just the right number of restful hours a night, it must be a
secretive one. Most people seem insecure about their sleep and willing to say so: they would like to
get a little more; maybe they wish they could get by on less; they wonder if it's deep enough.

And they are pretty sure that being up at 2 a.m., pacing the TV room like a caged animal, cannot be
good. Can it?

In fact, no one really knows. Scientists aren't sure why sleep exists at all, which has made it hard
to explain the great diversity of sleeping habits and quirks in birds, fish and mammals of all
kinds, including humans.

Why should lions get 15 hours a night and giraffes\cite{giraffe} just 5 -- when it is the giraffes
who will be running for their lives come hunting time? How on earth do migrating birds, in flight
for days on end, sleep? Why is it that some people are early birds as young adults and night owls
when they're older?

The answer may boil down to time management, according to a new paper in the August issue of the
journal Nature Reviews Neuroscience. In the paper, Jerome Siegel, a professor of psychiatry at the
University of California, Los Angeles, argues that sleep evolved to optimize animals' use of time,
keeping them safe and hidden when the hunting, fishing or scavenging was scarce and perhaps risky.
In that view, differences in sleep quality, up to and including periods of insomnia\cite{insomnia},
need not be seen as problems but as adaptations to the demands of the environment.

``We spend a third of our life sleeping, and it seems so maladaptive\cite{maladaptive} -- 'the
biggest mistake nature has made,' scientists often call it,'' said Dr.~Siegel, who is also chief of
neurobiology at the V.A. Greater Los Angeles Healthcare system. ``But another way of looking at it
is this: unnecessary wakefulness is a bigger mistake.''

As a field of study, sleep research is anything but sleepy -- experts disagree strongly on almost
every theory offered, and this one is no exception. Among other objections, critics point out that
sleeping animals are less alert to predators\cite{predator} than they are while awake, and that
sleep appears to serve other essential functions. Some studies suggest that the brain consolidates
the day's memories during slumber; others suggest that it needs sleep to repair neural damage.

``My own theory, which is more consistent with the mainstream, is that neurons require sleep as part
of the long-term process'' of modification to support learning, wrote Dr.~Clifford Saper, a
neuroscientist at Harvard, in an e-mail message. But, he added, his theory and Dr.~Siegel's may not
be mutually exclusive.

For one thing, sleep is not nearly as vulnerable a state as it appears. Sleepers are highly
sensitive to some sounds, like a baby's whine or an unusual thump or voice. And as Dr.~Siegel put
it, sleepers are less vulnerable to harm than they would be if they were out on the street late at
night.

For another thing, the new paper argues, evidence from other animals strongly suggests that the need
for sleep drops sharply during the most important waking hours. Migrating\cite{migrate} killer
whales are alert and swimming for weeks on end, and seemingly just as alert as when well rested,
studies find. Recent research suggests that the same goes for white-crowned sparrows : they get far
less sleep than usual when migrating.

Consider the big brown bat, perhaps the longest-sleeping mammal of them all. It snoozes 20 hours a
day, and spends the other 4 hunting mosquitoes and moths in the dusk and early evening. ``Increased
waking time would seem to be highly maladaptive for this animal, since it would expend energy and be
exposed to predatory birds with better vision and better flight abilities,'' Dr.~Siegel writes.

In humans, it is well known that sleep quality changes with age, from the long, deep plunges of
early childhood to the much lighter, more frequently interrupted five or six hours that many elderly
people call a night's sleep. Doctors have long debated whether elderly people are sleep-deprived as
a result, or simply need less restful slumber\cite{slumber}. In Dr.~Siegel's view, it's a matter of
tradeoffs: older people no longer have a child's need to grow, which requires deep, long sleep and
may have more need and more ability to do things for themselves instead.

In short, when there is hay\cite{hay} to be made, animals tend to make it, whether the sun is
shining or not. Depending on the animal, a long period of waking may or may not be followed by a
long recovery sleep.

The theory also supports what people already suspect about early birds and night owls: they are most
alert when they are naturally most productive. And they can feel strung out if their work schedule
doesn't match.

None of which is to say that good sleep is unnecessary or that serious sleep problems do not exist.
It is and they do. But the theory does suggest that a stretch of insomnia may not be evidence of a
disorder. If sleep has evolved as the ultimate time manager, then being wired at 2 a.m. may mean
there is valuable work to be done. Time to turn off the ``South Park'' reruns and start doing it.

\section{Clinton's Health Defeat Sways\cite{sway} Obama's Tactics}

Before Congress's August break, the chief aides to Senate Democrats met in a nondescript Senate
conference room with three former advisers to President Bill Clinton. The topic: lessons learned the
last time a Democratic president tried, but failed disastrously, to overhaul the health care system.

With the aides as divided as their bosses on President Obama's signature initiative, their typically
tedious weekly session turned hotly spirited. So the Clinton White House veterans -- John
D.~Podesta, a former senior adviser; Steve Ricchetti, a Congressional lobbyist; and Chris Jennings,
a health policy aide -- homed in on their ultimate lesson of the failure 15 years ago, that there is
a political cost to doing nothing.

In 1994, Democrats' dysfunction over fulfilling a new president's campaign promise contributed to
the party's loss of its 40-year dominance of Congress. Now that memory is being revived, and it is
the message the White House and Congressional leaders will press when lawmakers return this week,
still divided and now spooked after the turbulent town-hall-style meetings, downbeat polls and
distortions of August.

Republicans early on united behind the lesson they took from the past struggle, that they stand to
gain politically in next year's elections if Democrats do nothing. But the Democrats' version
similarly resonates with all party factions, giving Mr.~Obama perhaps his best leverage to unify
them to do something. In now-familiar financial parlance, this one is ``too big to fail.''

``Certainly if you undercut your own leadership, that shortens the honeymoon and could possibly even
cripple the administration,'' said Representative Jim Cooper of Tennessee, a Democrat and member of
the fiscally conservative Blue Dog Coalition who opposed the Clinton plan and has criticized current
efforts. ``And no one here wants that.''

That 15-year-old lesson underscores how much the Clinton debacle has defined Mr.~Obama's drive for
his domestic priority from the beginning, providing a tip sheet for what not to do. Even Mr.~Obama's
decision to address a joint session of Congress on Wednesday night to jumpstart his health
initiative left some aides wary, given the inevitable parallels with Mr.~Clinton's September address
16 years ago to introduce his ill-fated plan.

Before Mr.~Obama was elected, advisers began debriefing Clinton veterans to draft ``lessons
learned'' memorandums. According to interviews with more than a dozen participants in the debates
then and now, those lessons have helped the president's proposals progress further through
Congress's committees than the plan advocated by Bill and Hillary Rodham Clinton did. All the while,
the administration has held the tentative\cite{tentative} support of powerful associations for
doctors, nurses, seniors, hospitals, drug makers and, as Mr.~Obama recently put it, ``even the
insurance companies,'' which did the most to defeat the Clintons.

But Mr.~Obama's performance has also raised questions about whether the administration has drawn too
much from some lessons and underestimated some hurdles unique to today's battle.

The losses of two important confidants -- former Senator Tom Daschle's withdrawal as Mr.~Obama's
choice for health and human services secretary amid a tax controversy, and Senator Edward
M.~Kennedy's long illness and death -- have been ``incalculable'' setbacks for pushing legislation
through Congress, said a top aide to Mr.~Obama. Neither the House nor the Senate met his deadline
for passing legislation before August.

Yet even if the administration did everything right, drafting legislation this complicated is never
going to be easy.

``That period of defining the issue and developing the pieces and the resources to actually
legislate is a relatively smooth river,'' said Charles Kahn, an insurance industry lobbyist in the
1990s who now represents for-profit hospitals. Once Congress starts filling in the details, he
continued, ``then you hit the rapids. And they've hit the rapids now.''

Mr.~Obama has two disadvantages that Mr.~Clinton did not. The deep recession has stoked concerns
about deficits, and moderate Republicans willing to cut deals are nearly extinct.

But Mr.~Obama also has advantages flowing from his election by a 53 percent majority -- the highest
number for a Democrat since 1964 and 10 percentage points more than Mr.~Clinton won. Congressional
Democrats, while hardly of one mind, are still more unified behind him than they were behind
Mr.~Clinton, committee leaders are more respectful and Mr.~Obama is richer in campaign money and
grass-roots support.

Add to that, Democrats say, the benefits of the lessons of the Clinton era:

\emph{Lesson 1: Failure Is Not an Option.}

As the Clinton-era veterans attest, and the Obama team is now arguing to Democrats, voting for a
health bill might be difficult politically, but doing nothing at this point would be worse.

``When a party fails to govern, it fails electorally,'' said Rahm Emanuel, a former Clinton aide who
now is Mr.~Obama's chief of staff.

\emph{Lesson 2: Know your audience -- insured taxpayers.}


Even as a candidate, Mr.~Obama showed that he was trying to avoid his predecessor's mistakes.

Mr.~Clinton had promised universal coverage, emphasizing care for the uninsured -- roughly 37
million then, now 46 million. That left the many more Americans with insurance, however flawed, to
wonder what the complex and initially costly changes would mean for them.

By the 2008 presidential campaign, Mr.~Obama's Democratic rivals, including Mrs.~Clinton, were again
promising universal coverage. Mr.~Obama never did, except for children. He emphasized that insured
Americans would see lower costs, more choices and better coverage.

Then, as now, he more often cited as potential beneficiaries not the mostly poor uninsured but the
working and middle classes, people like his mother, who had insurance but fought her carrier all the
while she was dying of cancer. When he spoke of covering the uninsured, Mr.~Obama argued that doing
so would also help the insured because hospitals, doctors and insurers would no longer have to pass
on unpaid expenses in higher premiums and prices to paying patients.

But for all his efforts, when Congress began writing legislation and its analysts priced the various
proposals this summer, the sticker shock drew taxpayers' attention -- just as in Mr.~Clinton's time
-- to the main expense, which was covering the uninsured. Democrats' plans would expand Medicaid for
the poor and subsidize both low-income workers buying insurance and small businesses seeking
coverage for employees.

So, just as in 1994, ``people are trying to figure out what they're getting aside from additional
costs'' for taxpayers, said Howard Paster, Mr.~Clinton's chief Congressional lobbyist. ``At the end
of the day, most people will complain about their insurance company, will grumble about costs, but
they're ultimately satisfied with the health care they're getting.''

Some Democrats compare the current moment to the early spring of 1994. The insurance lobby's ``Harry
and Louise'' commercials against the Clinton plan had come to embody the angst of the insured
middle-class. Democratic lawmakers were panicky. ``It became irrevocably lost,'' said Mr.~Daschle,
who as a Senate Democratic leader was at the center of that debate.

Mr.~Obama has more recently been spotlighting features that will appeal to the insured middle class,
the proposed ``consumer protections'' that will not cost taxpayers. Among them, companies must cover
anyone regardless of medical history, cannot drop policyholders who become sick and must cover
preventive care like mammograms.

\emph{Lesson 3: Move before the honeymoon ends.}

This is one of the lessons Mr.~Obama may have learned too well.

In 1993 Mr.~Clinton delayed his push for a health care bill until late that first year, by which
time he was weakened by other legislative battles and personal controversies. Mr.~Obama, by
contrast, moved quickly to exploit his post-inaugural momentum, amid the demands of the worst
economic crisis since the Depression.

But some say he moved too quickly, setting the August deadline for the House and the Senate to each
pass their bills, then appearing to take a loss when neither acted.

``The biggest mistake Obama made, and I want him to succeed, is trying to rush it,'' said former
Senator Bob Dole, who as Republican leader in the 1990s first negotiated with Mr.~Clinton and then
led the opposition. ``Why put some arbitrary deadline on a piece of legislation that's going to
affect every American?''

Meanwhile, Mr.~Obama, like Mr.~Clinton, has lost some leverage and luster as his poll numbers have
weakened.

\emph{Lesson 4: Leave the details to Congress.}

Mr.~Obama got a faster start than Mr.~Clinton by not repeating his mistake of trying to write the
law for the lawmakers. The Clintons' secretive labors on a 1,342-page bill cost nine months and
stoked resentment among Congress's proud Democratic committee barons, who felt left out.

Mr.~Obama went to the other extreme. He produced no plan, only fairly specific directives. He said
he wanted to create ``exchanges'' offering private insurance plans and a public option. He called
for insurance subsidies for individuals and small businesses. And he advocated changes in Medicare
and Medicaid payments to give the health industry incentives to control costs and improve care.

While Congressional Democrats welcomed the partnership, some now wonder if the president did not
``overlearn the lessons of 1994 by giving Congress too much leeway,'' as Mr.~Cooper of the Blue Dog
Coalition put it.

Administration officials counter that the president's initiative would not otherwise have gotten as
far as it has, with bills passed in four out of five committees, if he had not initially deferred to
Congress. But with the fifth and most crucial panel, the Senate Finance Committee, still struggling
for a bipartisan alternative, the president will most likely serve notice with Wednesday night's
speech that he is taking the lead, with what aides call a ``more prescriptive'' legislative
blueprint for Democrats to get behind.

``They're not going to get the ball over the finish line without his direct help and intervention,''
said Mr.~Podesta, the former Clinton lieutenant, who headed Mr.~Obama's post-election transition.

\emph{Lesson 5: Co-opt the opposition.}


In a lesson that holds some irony for Mr.~Obama given his campaign against special interests, he has
mostly rejected the Clintons' industry-bashing populism. That has helped keep powerful groups at the
table, to prevent their allying against him as they did against Mr.~Clinton.

The president privately reached early deals with the for-profit hospital group represented by
Mr.~Kahn, who led the ``Harry and Louise'' campaign, and with the drug manufacturers lobby. The
industries agreed to accept roughly \$230 billion in reduced Medicare and Medicaid payments over 10
years to help offset the cost of a health care bill, and the White House committed to support the
deals through the legislative process despite liberals' demands for bigger concessions.

For older Americans, the administration agreed with the advocacy group AARP that any bill would
eliminate the gap in Medicare coverage of prescription drugs.

``People have underestimated the strategic value of some of these alliances in terms of being able
to keep this thing going,'' said a former Clinton aide, who asked not to be identified because he
now lobbies for several industries.

The Obama team's outreach extended even to the insurance industry, until the administration began
going after it last month to shore up public support. White House officials remain divided over the
tactic, with some fearful of provoking the deep-pocketed industry's fury -- as the Clintons did.

\emph{Lesson 6: Take what you can get.}

What optimism remains among Democrats stems from their belief that Mr.~Obama, unlike the Clintons,
will take half a loaf and declare victory, and that most Democrats, mindful of 1994's election
debacle, will go along.

In his 1994 State of the Union address, Mr.~Clinton famously waved a pen and threatened to veto any
bill that did not ``guarantee every American'' private health insurance. Even an aide who
recommended that uncompromising signal, Paul Begala, now says it was a mistake. Others have said the
White House forfeited\cite{forfeit} a chance to compromise with Mr.~Dole and other Senate Republicans.

The question for Mr.~Obama is whether he will have any Republicans with whom to compromise. More
likely, he will have to mediate between the liberals and conservatives in his own party.

\section{China's Li Na Finds a Taste of Home in New York}

\lettrine{L}{i} Na stepped past the musicians plucking familiar string instruments and people doing
tai chi in the pavilion. She stood near tables of men huddled around games of checkers.

``It is exactly like China,'' Li said with a broad smile.

It was Columbus Park, the quiet, beating heart of Manhattan's bustling Chinatown. One difference
from her strolls in China is that no one seemed to recognize her Monday afternoon, seeing only a
young woman in jeans and a black cap, not China's best tennis player and a quarterfinalist at the
United States Open.

Li, 27, is not only a long way from home, but she is also as far into a Grand Slam tournament as she
has ever been. The day before her match against the resurgent, unretired Kim Clijsters, the
18th-seeded Li took a taxi from her Midtown hotel and spent a couple of hours strolling Chinatown
and having lunch.

``I love New York City,'' she said. ``People in China say: 'If you love your children, send them to
New York. If you hate your children, also send them to New York.' ''

She laughed at that lesser-known Chinese proverb\cite{proverb}. In news conferences, when asked to
address questions in English, Li seems unsure and introverted. Alone, on the street or at a
restaurant table, her English is near perfect and she continually punctuates jokes with a disarming
grin.

When she plays at Flushing Meadows, she is about a mile -- one stop on the 7 train -- from the
bustling Chinatown in Flushing, Queens, which some say has overtaken the Manhattan Chinatown in
recent years as the largest concentration of Chinese in the Western Hemisphere. Li has stayed at a
house in Flushing with other Chinese players and coaches the three previous times she played the
Open, but she wanted to get a close-up view of Manhattan this time.

She laughed about how people from outside China get confused about her name. Li is her family name,
she said. And Na is her first name, even if the Chinese place it at the end.

She has heard people see her name and call her ``Lina.'' Or ``Nali.''

``If you see the draw,'' she said of the tournament bracket, ``it's easy to find my name. It's so
short.''

It is increasingly easy to find her name on the bracket now that Li is among the final eight women.
The only other time she reached a Grand Slam quarterfinal was at Wimbledon in 2006. She lost to
Clijsters.

``I still can remember the second set,'' Li said. ``I had set point. I was 5-2 up, had set point,
40-30. I served, and she returned the ball so fast.''

Clijsters won that match, 6-4, 7-5. But Li won in three sets the last time the two met, in 2007 at
Miami.

At this Open, Li slid easily through the fourth round with a 6-2, 6-3 victory over Francesca
Schiavone. Afterward, to soothe a sore right knee, Li dunked her lower body into an ice bath for the
first time in her life.

``I thought I was going to die,'' she said. A decision was made not to practice Monday.

That is why Li sat in a booth at an unassuming restaurant down a covered alley between Elizabeth
Street and the Bowery. The place called itself ``Coluck'' on the menu, but the light-up sign, in
Chinese, reads, ``First Chinese Western Restaurant,'' Li said.

She ordered kimchi bokumbop, a traditional Korean fried-rice dish, and iced milk tea.

``They say it is not good for the body,'' Li said of the milk tea. ``But I like it a lot.'' (She
also likes beer, like Tsingtao, but not during the season. But if she wins the Open? ``For sure.'')

She rated the rice dish as ``not 100 percent.''

``Maybe 85 percent,'' she said.

From the huge central China city of Wuhan, Li was deemed a badminton player when she was 6. A coach
kept telling her that she looked as if she were playing tennis.

``In badminton, they use a lot from the wrist,'' Li said as she demonstrated the arm action. ``But I
use a lot from the shoulder.''

When Li was 8, the coach asked her parents -- her father died when she was 14 -- if she could switch
to tennis.

``They said, 'What's tennis?' '' Li said.

Li went on to become the first Chinese player to reach the top 30 (in 2006), then the top 20 (in
2007).

Problems with her right knee forced her to miss the second half of the 2007 season, but Li has since
made a habit of reaching the third and fourth rounds of important tournaments.

Her best moment came a little more than a year ago, at the Beijing Olympics, where Li reached the
semifinals by beating Svetlana Kuznetsova and Venus Williams. The stadium echoed with chants of ``Li
Na!'' But she then lost twice, missing out on a medal.

Smatterings of fans tend to make it to Li's matches at the Open, chanting encouragement in Mandarin.
It will be hard to drown out the affection certain to shower Clijsters.

Still, Li says she loves that she can find pieces of home so far from home.

``Sometimes, I think it's funny,'' she said as she weaved through pedestrians on Bayard Street.
``You're in America, but everyone is speaking Chinese.''

\section{AOL, Seeking to Revitalize Its Internet Role, Turns to a Candid Ex-Yahoo Executive }

\lettrine{B}{rad} Garlinghouse, a former Yahoo executive who was credited with highlighting the deep
problems at the Internet company, was named president of Internet and mobile communications at AOL,
the company announced early on Tuesday.

In his new role, Mr.~Garlinghouse will lead the effort to expand the reach of AOL's popular e-mail
and instant-messaging services. He will also head the company's Silicon Valley operations, which AOL
plans to expand, and lead the West Coast arm of AOL Ventures, a unit in charge of investing in
start-ups and spinning off businesses.

Mr.~Garlinghouse is the latest senior executive to join AOL, a unit of Time Warner. Tim Armstrong,
the new chief executive of AOL, is trying to revive the company's fortunes and burnish its tattered
image with consumers and investors.

Nearly a decade after it merged with Time Warner, AOL is struggling to overcome multiple problems,
including declining revenue, low morale after several rounds of layoffs and corporate infighting
with its parent. Time Warner has said it plans to spin off AOL into an independent company by the
end of the year.

Analysts say that Mr.~Armstrong, a well-regarded sales executive that Time Warner lured away from
Google in April, is fighting to save the company from irrelevance in an industry dominated by giants
like Google, Microsoft, Yahoo and Facebook. He has recruited other top executives from Google and
elsewhere and laid out a plan to focus on a handful of top priorities, including becoming a leader
in display advertising and a top online destination for original content.

In an interview on Monday, Mr.~Armstrong said that communications products, which include e-mail and
instant-messaging services, remained one of AOL's most important assets, keeping customers engaged
and helping drive traffic to other AOL properties. But in that area too, AOL faces challenges. Its
e-mail service was recently overtaken by Google's Gmail, which became the No.~3 e-mail service in
the United States, behind Yahoo Mail and Hotmail, owned by Microsoft.

At Yahoo, Mr.~Garlinghouse was responsible for communications products and helped to turn it into
the most used Web e-mail service in the United States.

``Brad is an 'all star' Internet executive,'' said Mr.~Armstrong, adding that he expected
Mr.~Garlinghouse not only to fulfill his specific duties, but also to play a broader role in helping
to turn around AOL.

At Yahoo, Mr.~Garlinghouse acquired some notoriety after a frank memo he wrote for senior executives
on internal disarray became public in late 2006. In the memo, which became known as the ``peanut
butter manifesto,'' Mr.~Garlinghouse complained that Yahoo was spreading itself too thin, like
peanut butter on toast. He said the company needed to focus on fewer priorities.

The memo became a rallying cry inside and outside the company and served as a roadmap while three
different chief executives struggled to turn Yahoo's fortunes around.

In an interview on Monday, Mr.~Garlinghouse acknowledged that AOL also faced significant challenges,
but suggested that it did not have the kinds of systemic problems that Yahoo was confronting in
2006.

``There is a clarity of vision and strategy,'' he said.

Mr.~Garlinghouse left Yahoo at the end of 2008 and served most recently as a senior adviser to
Silver Lake Partners, a private equity firm.

\section{Gas Explosion Kills 35 Miners And Traps 44, China Says}

\lettrine{T}{hree} days after one of China's senior leaders proclaimed coal mine safety to be a
priority for the government, a gas explosion at an unlicensed mine in central China killed at least
35 workers and trapped 44, the national work safety agency said Tuesday.

The explosion rocked a coal pit in the city of Pingdingshan, in Henan Province, about 1 a.m. on
Tuesday, a city spokesman told the official Xinhua news agency.

The spokesman said 93 miners were working in the pit when the blast occurred, and 14 managed to
escape. He also said the mine had been under repair and had not been cleared to resume operations.

An estimated 80 percent of the 16,000 mines operating in China are illegal, according to the State
Administration of Work Safety. Government figures show that about 3,200 people died in mining
accidents last year, a 15 percent decrease from 2007.

China's deputy prime minister, Zhang Dejiang, addressed the issue on Saturday at a coal industry
conference, calling for numerous improvements in mine safety to prevent, in particular, gas
explosions, Xinhua reported.

In February, in the country's deadliest coal mine accident in more than a year, at least 74 people
died in a gas explosion, and 114 were hospitalized with carbon monoxide poisoning. The blast at the
Tunlan Coal Mine in Shanxi Province, the coal mining heartland of China, also occurred just after
midnight.

In December 2007, a gas explosion in a mine in the city of Linfen, also in Shanxi, killed 105
miners, according to the work safety agency.

\section{Big U.S.~Bases Are Part of Iraq, but a World Apart}

\lettrine{I}{t} takes the masseuse\cite{masseuse}, Mila from Kyrgyzstan, an hour to
commute\cite{commute} to work by bus on this sprawling American base. Her massage parlor is one of
three on the base's 6,300 acres and sits next to a Subway sandwich shop in a trailer, surrounded by
blast walls, sand and rock.

At the Subway, workers from India and Bangladesh make sandwiches for American soldiers looking for a
taste of home. When the sandwich makers' shifts end, the journey home takes them past a power plant,
an ice-making plant, a sewage treatment center, a hospital and dozens of other facilities one would
expect to find in a small city.

And in more than six years, that is what Americans have created here: cities in the sand.

With American troops moved out of Iraq's cities and more than 100 bases across the country
continuing to close or to be turned over to Iraq, the 130,000 American troops here will increasingly
fall back to these larger bases.

While some are technically called camps or bases, they are commonly referred to as forward operating
bases, or F.O.B.'s. The F.O.B. is so ingrained\cite{ingrained} in the language of this war that
soldiers who stayed mainly on base were once derisively called Fobbits by those outside the wire.
But increasingly, the encampments are the way many Americans experience the war.

To be sure, thousands of Americans are with Iraqis at small bases, where they play an advisory role,
and thousands more are on the roads and highways providing the protection needed to carry out the
withdrawal.

But the F.O.B. has become an iconic part of the war, both for those fighting it and for the Iraqis,
who have been largely kept out of them during the war.

They are in some ways a world apart from Iraq, with working lights, proper
sanitation\cite{sanitation}, clean streets and strictly observed rules and codes of conduct. Some
bases have populations of more than 20,000, with thousands of contractors and third-country citizens
to keep them running.

But the bases are also part of the Iraqi landscape. Mortar shells still occasionally fall inside the
wire, and soldiers fall asleep to the constant sounds of helicopters, controlled detonations and
gunfire from firing ranges.

``It is definitely a strange place,'' said Capt.~Brian Neese, an Air Force physician. ``I've asked
the Civil Affairs guy if there is anything that I can do off base, and there just isn't anything for
me to do. What kills is not the difficulty of the job but the monotony\cite{monotony}.''

At the height of the war, more than 300 bases were scattered across Iraq. Over the next few months,
Americans hope to be at six huge bases, with 13 others being used for staging and preparing for a
complete withdrawal.

The first people you encounter when driving up to an American base are not actually American. They
are usually Ugandans, employed by a private security company, Triple Canopy, and those at Balad had
enough authority to delay for five hours an American Air Force captain escorting an American
reporter onto the base.

The Ugandans make up only one nationality of a diverse group of workers from developing nations who
sustain life on the F.O.B.'s for American soldiers. The largest contingents come from the
Philippines, Bangladesh and India. They live apart from both Western contractors and soldiers on
base, interacting with them only as much as their jobs demand.

``Everyone stays pretty much separate,'' said Mila, the massage therapist, whose last name could not
be used out of security concerns. She has been in Iraq a year, but she said other workers had been
here as long as six years, some never taking a break to go home. ``You miss nature, trees and
grass,'' she said.

The base has two power plants, and two water treatment plants that purify 1.9 million gallons of
water a day for showers and other uses. The water the soldiers drink comes from yet another plant,
run by a bottling company, which provides seven million bottles of water a month for those on base.

Fifteen bus routes crisscross the complex, with 80 to 100 buses on the roads at any given moment.
The Air Force officers who run the base have meetings to discuss road safety; with large, heavily
armored vehicles competing for space with sedans, there are bound to be collisions.

There are two fire stations as well, and because Balad has the single busiest landing strip in the
entire Defense Department, they can handle everything from an electrical fire in a trailer to a
burning airplane.

The Americans also installed two sewage treatment plants, given how deeply troubled Iraq's sewage
system remains.

The facilities, like much in Iraq, are run by KBR, a company based in Houston. But as Americans
prepare to turn bases over to Iraqis, they are working to bring in Iraqi companies to run some
facilities, a process that has been slow and complex largely because of safety concerns.

One of the few places Iraqis can be seen, in fact, is the ``Iraqi Free Zone,'' a fenced-in area
enclosed with barbed wire and blast walls. There, Iraqis sell pirated movies, discount cigarettes,
electronics and Iraqi tchotchkes.

Each large base in Iraq takes on its own distinct flavor. Most large American bases were once Iraqi
bases, but some, like Camp Bucca in southern Iraq near the Kuwaiti border, were created where there
was only sand.

An Iraqi interpreter at Bucca who was living in Texas with his family when the war started said that
when a contracting firm approached him, he asked where he would be working.

``They told me, 'You would be going to a place called Bucca,' '' said the interpreter, whose name
the military asked not to be printed for security reasons. ``I said, 'There is no city called
Bucca.' They showed me it on the map and I said, 'I am from Iraq and there is no city called Bucca.'
''

It turned out that the interpreter was correct. Bucca, which would house the largest American-run
prison in Iraq, was named after Ronald Bucca, a soldier with the 800th Military Police Brigade and a
fire marshal in New York who was killed in the Sept.~11 attacks.

Entertainers come to American bases here, the most frequent being N.F.L. cheerleaders. When the
Minnesota Vikings cheerleaders visited Bucca this spring, one could only wonder what the thousands
of detainees, among them Muslim extremists for whom the flash of an ankle is cause for severe
punishment, would have made of the spectacle less than a mile from their cells.

\section{Driver Texting Now an Issue In Back Seat}

\lettrine{A}{fter} decades of marriage, Terry and Debbie Buchen learned to work through various
marital\cite{marital} issues. Then something new came between them -- his cellphone.

Mr.~Buchen, 62, couldn't put it down while driving. The first time he sent e-mail messages from
behind the wheel, he drove his BMW S.U.V. into a ditch on a deserted stretch of road. He was alone
and driving slowly, and he wasn't injured. Still, the incident was ``very scary,'' his wife said.

Mr.~Buchen knew he had a to make a choice between his habit and marital bliss\cite{bliss}.

``I chose my wife,'' he said. But then Mr.~Buchen, an agronomist for golf courses, asked for a
compromise: he asked her to drive when they were together so he could stay connected with clients.
That didn't fly. ``If looks could kill,'' he said.

For all the conversations about distracted driving playing out in statehouses and on talk shows, the
most heated discussions, and the ones with the most lasting impact, may be happening between family
members and friends.

Such disputes are an extension of a longstanding source of tension -- sometimes light, other times
more antagonistic -- between drivers and their self-appointed watchdogs. It's just that now, the
back-seat driver is going after the BlackBerry.

These critics say such devices not only put lives at risk, but also steal attention from passengers
hoping for some quality catch-up time. The multitaskers counter that, in their view, social and work
demands override the increased crash risks.

Safety advocates who favor outlawing multitasking behind the wheel say the new generation of
back-seat hawks may be playing a crucial role in changing the culture -- much as they did in helping
enforce seat belt laws -- in a way these advocates say laws alone may not be able to.

In a survey conducted this year of 2,501 people, the AAA Foundation for Highway Safety found that 48
percent of people worry about a friend or family member driving unsafely. Of those people, 19
percent said the cause of their concern was multitasking behind the wheel.

Some drivers say that the second-guessing is unnecessary because of their ability to handle driving
and other tasks.

``I barely even look at my phone when I text,'' said Sarah Edwards, 29, a customer service
representative in Cleveland.

Her close friend, Amy Macauley, 32, isn't convinced. ``We've talked about this,'' Ms.~Macauley said.
``I've cited vague\cite{vague} statistics about how dangerous it is, but she doesn't pay
attention.''

She brings a particular sensitivity to the discussion. Her sister and the husband of a close friend
died from crash injuries, and her brother-in-law was seriously injured in a car accident.

``I know things can change in an instant when someone is behind the wheel,'' Ms.~Macauley said. To
her, talking on the phone or texting while driving, unless it's an emergency, is ``completely
gross.''

People like Ms.~Macauley are finding more ammunition to argue their side.

Studies show that people who talk on the phone while driving face four times as great a crash risk
as those who do not. The risk is considerably higher for motorists who text.

The federal government estimates that at any given time about 11 percent of drivers, or about two
million people, are talking on a cellphone.

Many of them are most likely doing so alone. Up to 85 percent of drivers have no passenger with them
in the car, estimates the Virginia Tech Transportation Institute.

But when an adult passenger is present, the Transportation Institute found, he or she can enhance
safety -- and reduce crash risk by up to 50 percent -- perhaps by acting as a second set of eyes on
the road, or encouraging safer behavior.

That means that a friend or family member cajoling a motorist to put down the phone can provide a
safety advantage, as long as the disagreement itself doesn't escalate to the point where it becomes
its own distraction.

``An emotional, heated argument probably isn't good, but that's just my opinion,'' said Charlie
Klauer, senior research associate at the Transportation Institute.

For all the evidence about the dangers of distracted driving, multitaskers say there are powerful
forces -- both social and economic -- that make it hard to put their devices down.

A 2008 poll taken by Nationwide Insurance of 1,500 motorists found that 48 percent of people who
multitasked behind the wheel did so because they felt an urgent need to address an issue pertaining
to school or work; 33 percent said they felt pressure to stay connected socially.

That pressure may affect some people more than others. An emerging body of science suggests that
there are some people who are more likely to be drawn to multitasking -- behind the wheel or
otherwise -- because of the way their brains are wired.

Clifford Nass, a co-director of an automotive research laboratory at Stanford University, said some
researchers had labeled such heavy multitaskers as ``explorers'' because they enjoyed the constant
hunt for information -- whether it pertained to work, entertainment or social life.

Vic Gideon, 48, a hospital executive from Cleveland, says he feels the call of his device for social
and work purposes, even when behind the wheel and even when his family, including his four children,
are in the car.

``Even if I'm going 60 miles an hour, I feel the need to check it,'' said Mr.~Gideon, referring to
his phone. ``It might be spam, a wrong number, whatever. But who cares? My cell vibrates. I
respond.''

Mr.~Gideon said it had created tension. ``My wife tells me I don't know how to drive,'' he said.
``And then, of course, not only do I not know how to drive, but I don't know how to drive safely. I
believe I'm being careful.''

The Nationwide Insurance poll found that both men and women were avid multitaskers. About 85 percent
of female drivers said they multitasked, compared with 78 percent of male drivers.

Grace Andrews, 49, a corporate consultant in Melrose, Mass., is the one taking heat in her family.
Her husband, Joe Nardone, 44, and her son, Colby Andrews, 12, despise her incessant use of the
phone, which she concedes can be over the top.

``I honestly do laugh at myself all the time,'' says Ms.~Andrews. ``Is it really possible that I am
talking on the phone, e-mailing and driving with my knees simultaneously?''

Her husband and son tell her that she cares more about the phone than she cares about them, and that
she's putting herself and others at risk.

``I could never imagine that we would get to this stage -- that this is the stuff we would fight
about,'' she said.

\section{Merger Would Create Mobile Giant in Britain}

\lettrine{D}{eutsche} Telekom and France T\'el\'ecom said Tuesday that they were planning to merge
their struggling mobile operations in Britain, creating the largest mobile operator there.

The 50-50 joint venture between the companies, the third- and fourth-largest operators in Britain,
would have 28.4 million customers and a 37 percent market share, according to Gartner, leapfrogging
the market leader, O2, with 27 percent, and Vodafone, with 25 percent. The companies said they
expected to sign an agreement by the end of next month.

One analyst, John Strand, a mobile industry consultant based in Copenhagen, said it was a sign of a
coming wave of mergers as operators struggled to survive in the saturated\cite{saturate} markets of
the European mobile industry.

``Many operators are having a difficult time competing in the current climate with the increased
regulation and the intense competition over the price of mobile broadband, which many are now
selling below cost,'' he said.

Deutsche Telekom and France T\'el\'ecom's plans need approval from each company's supervisory board
and from British and European regulators. They said they planned to maintain their brands in Britain
for 18 months before choosing a final brand, which could be the French one, Orange; the German one,
T-Mobile; or something else.

The companies said they planned to spend \textsterling600 million to \textsterling800 million, or
\$984 million to \$1.3 billion, through 2014 to eliminate unneeded mobile base stations, close
retail stores and streamline administrative operations. In a conference call, Tom Alexander, chief
executive of Orange in Britain, said the merger would lead to layoffs but declined to be more
specific.

Mr.~Alexander will serve as chief executive of the new venture. Richard Moat, T-Mobile's chief
executive in Britain, will become chief operating officer. The venture will be run by a board with
two members from each company.

The operators predicted the venture would eventually generate \texteuro4 billion, or \$5.7 billion,
in savings. ``We needed the scale to run a very efficient business,'' Gervais Pellissier, the France
T\'el\'ecom chief financial officer, said.

Both operators had reported declining profit and sales at their British businesses this year. In the
first quarter, Deutsche Telekom wrote down \$1.8 billion of the value of T-Mobile U.K.

``We evaluated all options over the last months and came to the conclusion that this would have the
most value for our shareholders,'' Timotheus Höttges, the Deutsche Telekom chief financial officer,
said. ``This will give us a stronger foothold in the U.K. market than we had before.''

Ernest Doku, an analyst with Omio, an online handset retailer in Britain, said T-Mobile and Orange
would be able to use their combined market strength to pry customers away from O2 and Vodafone.

``For these two networks to merge makes simple business sense, bringing operating costs down,''
Mr.~Doku said. ``They instantly turn from a bit part into a big player.''

T-Mobile and Orange have struggled to keep pace in Britain with O2, owned by Telef\'onica of Spain,
and Vodafone, the largest operator in Europe, which is based in Newbury, England.

The so-called mobile penetration rate in Britain, which measures the number of SIM card accounts
held by consumers divided by the total population, is expected to reach 128 percent by the end of
this year, up from 126 percent, according to Gartner. In other words, many consumers have more than
one mobile account.

The crowded mobile market in Britain, the largest in Europe after Germany, has five network
operators -- the other being the company 3 U.K. -- and multiple resellers. The British market has
been saturated for more than a decade, and competition has grown so intense that Orange is giving
away netbook computers to win long-term customers. Most Britons can buy mobile phone service through
no-strings-attached, pre-paid packages for as little as \textsterling10.

With similar saturation in other European countries, more mergers are expected, said Katja Ruud, an
analyst with Gartner in Stockholm. She said consolidation is especially possible in Spain and
Sweden.

``Consolidation in general is never 'done,' so to speak, and in markets where there are three or
more operators, you can expect this to happen,'' she said.

\section{Crux of Afghan Debate: Will More Troops Curb Terror?}

\lettrine{D}{oes} the United States need a large and growing ground force in Afghanistan to prevent
another major terrorist attack on American soil?

In deploying 68,000 American troops there by year's end, President Obama has called Afghanistan ``a
war of necessity'' to prevent the Taliban from recreating for Al Qaeda the sanctuary that it had in
the 1990s.

But nearly eight years after the American invasion drove Qaeda leaders from Afghanistan, the
political support for military action that followed the Sept.~11, 2001, attacks has faded. A war
that started as a swift counterattack against those responsible for the murder of 3,000 Americans, a
growing number of critics say, is in danger of becoming a quagmire with a muddled mission.

In interviews, most counterterrorism experts said they believed that the troops were needed to drive
Taliban fighters from territory they had steadily reclaimed. But critics on the right and the left
say that if the real goal is to prevent terrorist attacks on the United States, there may be
alternatives to a large ground force in Afghanistan. They say Al Qaeda can be held at bay using
intensive intelligence, Predator drones, cruise missiles, raids by Special Operations commandos and
even payments to warlords to deny haven to Al Qaeda.

After all, they point out, the Central Intelligence Agency has killed more than a dozen top Qaeda
leaders in the lawless Pakistani tribal areas, disrupting the terrorists' ability to plot and carry
out attacks against the United States and Europe.

Andrew J.~Bacevich, professor of international relations at Boston University, said the alternatives
would have at least as much chance of preventing attacks on the United States as a large-scale
counterinsurgency effort, which he said would last 5 to 10 years, require hundreds of billions of
dollars, sacrifice hundreds of American lives and have a ``slim likelihood of success.''

Despite the Obama administration's assertions that it has a new approach, ``the truth is they want
to try harder to do what we've been doing for the last eight years,'' Mr.~Bacevich said.

But most specialists on counterterrorism and counterinsurgency, inside and outside the government,
say terrorism cannot be confronted from a comfortable distance, such as by airstrikes or proxy
forces alone. It may take years to turn Afghanistan into a place that is hostile to Al Qaeda, they
say, but it may be the only way to keep the United States safe in the long term. Many agree with the
classified strategy for a troop buildup that Gen. Stanley A.~McChrystal, the top commander in
Afghanistan, has presented to Mr.~Obama and the Joint Chiefs of Staff in recent days.

They say a large American-led NATO ground force is needed to clear Taliban-held territory and hold
it while instructors train sufficient, competent Afghan soldiers and police officers to secure those
areas. The allied force, the argument goes, will buy time and space to help the Afghans build more
effective local, provincial and national governments, and create some semblance of an economy. Since
many polls in Afghanistan show little support for the Taliban, a stable, peaceful country would not
be likely to become a home for terrorists.

The administration's ``clear, hold, build'' strategy is meant to win over Afghans. It is based on
the counterinsurgency principle of protecting the population -- or, in this case, at least Afghan
population centers -- to win confidence and support, before isolating remaining insurgents to be
killed or captured. Defense Secretary Robert M.~Gates and Adm. Mike Mullen, chairman of the Joint
Chiefs of Staff, have maintained that the United States cannot fight terrorists from afar.

``The notion that you can conduct a purely counterterrorist kind of campaign and do it from a
distance simply does not accord with reality,'' Mr.~Gates told reporters last Thursday. ``The
reality is that even if you want to focus on counterterrorism, you cannot do that successfully
without local law enforcement, without internal security, without intelligence.''

Bruce Hoffman, a terrorism expert at Georgetown University, concurred, saying the argument that
terrorism can be prevented essentially by remote control was ``immensely seductive'' -- and
completely wrong.

``We tried to contain the terrorism problem in Afghanistan from a distance before 9/11,'' he said.
``Look how well that worked.''

Airstrikes risk killing civilians -- as shown by international concern over the possibility that
many civilians were among the scores killed in the NATO bombing of two fuel tankers in northern
Afghanistan last week -- and making enemies of the very people American commanders are trying to
sway. General McChrystal recently tightened rules on airstrikes to try to avert civilian deaths.

Mr.~Hoffman said the success of strikes from Predators in killing Qaeda suspects in Pakistan
depended on accurate information on terrorists' whereabouts from Pakistani intelligence. In
Afghanistan, without such sources, ``we'd be flying blind,'' he said.

Disengagement from Afghanistan could destabilize Pakistan and ``guarantee'' a future attack on the
United States from the region, Mr.~Hoffman said. For starters, a pullout could deny the United
States bases from which it carries out some Predator missions.

By all accounts, Pakistan remains the wild card that complicates all predictions on Afghanistan.
Proponents of the increased American force in Afghanistan say a withdrawal would reinforce
Pakistan's fears that the United States is not committed to security in the region, encouraging an
old Pakistani strategy of maintaining ties to Islamic militants.

But Daniel L.~Byman, director of Georgetown's Center for Peace and Security Studies, said
calculations on Pakistan's role were complicated. ``The more we escalate in Afghanistan, the more we
depend on Pakistan for logistics,'' he said, ``and the more Pakistan may feel it can resist our
pressure to go after the militants.''

As opposition to the war in Afghanistan builds, some of Mr.~Obama's strongest supporters say he must
do a better job of explaining how deploying a large American ground force there safeguards Americans
at home.

``He needs to reinforce that message more frequently and to the point,'' said Senator Bob Casey, a
Pennsylvania Democrat on the Foreign Relations Committee who just returned from Afghanistan and
Pakistan. ``He hasn't made it enough.'' 

\section{Despite Fears, Health Care Overhaul Is Moving Ahead}

\lettrine{T}{he} conventional wisdom, here and around the country, is that the centerpiece of
President Obama's domestic agenda -- remaking the health care system to cut costs and cover the
uninsured -- is on life support and that only a political miracle could revive it.

Here's why the conventional wisdom might be wrong:

While the month of August clearly knocked the White House back on its heels, as Congressional town
hall-style meetings exposed Americans' unease with an overhaul, the uproar does not seem to have
greatly altered public opinion or substantially weakened Democrats' resolve.

Critical players in the health-care industry remain at the negotiating table, meaning they are not
out whipping up public or legislative opposition.

Despite tensions between moderate and liberal Democrats, there is broad agreement within the party
over most of what a package would look like. Four of the five Congressional committees considering
health care legislation have already passed bills. Each would require all Americans to have
insurance and provide government subsidies for those who cannot afford it. Each would bar insurance
companies from refusing coverage for pre-existing conditions; imposing lifetime caps on coverage; or
dropping people when they get sick.

Getting a bill through the Senate remains a big challenge, but even there, the Obama administration
has a reasonable chance of corralling the 60 votes it would need to pass legislation more or less on
its terms. One wavering Democratic moderate, Senator Ben Nelson of Nebraska, signaled over the
weekend that he might be able to go along with one of the compromise proposals under discussion.
Senator Olympia J.~Snowe, the Maine Republican whose vote would be vital to Mr.~Obama, remains
deeply engaged in negotiations, and there are indications that one or two other Republicans, like
Senator George V.~Voinovich of Ohio, might be in play.

Politically, there is an imperative for Democrats to act; they remember well the disastrous
political fate that befell them in 1994, when they lost control of the House and Senate after
failing to pass a health bill under President Bill Clinton. Rahm Emanuel, the bare-knuckled
political operative and former Clinton aide who is now the White House chief of staff, has wasted
little time in reminding his fellow Democrats that, as he said in an interview Tuesday, ``the
inability to act here will have political consequences.''

None of this is to understate the magnitude of the task facing Mr.~Obama as he begins a final drive
for the legislation with a nationally televised address to Congress on Wednesday night. The size and
complexity of the legislation, the deep partisan divide, the undercurrent of concern among voters
about whether government is getting too big and intrusive, opposition from special interests -- all
create land mines that could still blow up the effort.

But even after weeks filled with seemingly ominous portents for Mr.~Obama's ambitions, there is
evidence that public opinion remains basically supportive of him. Despite intense controversy over
the ``public option,'' a government-backed insurance plan that would compete with the private
sector, a CBS poll at the end of August found that 60 percent of Americans still support the idea,
down from 66 percent in July. And half the respondents to the poll said Mr.~Obama had better ideas
on health care than Republicans, down from 55 percent.

Mr.~Obama likes to say that in the 100 years since President Theodore Roosevelt began advocating
universal health care, ``we've never had such broad agreement on what needs to be done.'' On Capitol
Hill, it is possible to see how a compromise could come together; Mr.~Nelson indicated over the
weekend that he could back a provision known as a ``trigger'' to create a public option if private
efforts to cover the uninsured failed.

And despite the fracas of August, the major stakeholders in the health care debate -- hospitals,
doctors, insurers and the pharmaceutical industry -- have not abandoned the negotiations. Ralph
G.~Neas, chief executive of the National Coalition on Health Care and a veteran of Washington
legislative fights, said this was especially significant.

``They're saying to themselves: 'We're going to get 30 to 40 or 50 million new customers. This is in
our economic self interest,' '' Mr.~Neas said. ``That, as much as anything else, could propel this
forward to a law that does provide quality health care for all.''

Mr.~Obama still clearly has not closed the deal, which is a major reason he will be making his case
directly to the American people and their elected representatives on Wednesday night. The CBS poll
found that 6 in 10 Americans say Mr.~Obama has not clearly explained what his plans for health
reform would mean.

That is a problem for the White House, though it also presents the president with an opportunity to
reframe the debate on his own terms. In his address on Wednesday, Mr.~Obama has promised to outline
what he wants to see in a bill; Republican leaders say the message from August is that Democrats and
the president need to start over.

``At this point, there really should be no doubt where the American people stand: the status quo is
not acceptable, but neither are any of the proposals we've seen from the White House or Democrats in
Congress,'' Senator Mitch McConnell of Kentucky, the Republican leader, said in a statement, adding:
``It should be clear by now that the problem isn't the sales pitch. The problem is what they're
selling.''

Yet Mark McClellan, who ran the Food and Drug Administration and later Medicare under President
George W.~Bush, said he saw the churning in August as a part of the public's education, a
``necessary step in the process'' and not a fatal blow.

Whether or not Mr.~Obama gets the kind of comprehensive bill he is hoping for, Dr.~McClellan said,
Congress is all but certain to take up health legislation by early next year to fix a measure that
would impose a draconian 21 percent cut in Medicare reimbursements to doctors. And once it is
tinkering with health care, he said, it is not that big a leap to imagine lawmakers using that bill
to take smaller steps toward expanding coverage and passing insurance market reforms.

``Everybody is talking about how the public is very concerned about some of the specifics that
they've heard,'' Dr.~McClellan said. ``But the public is also very concerned about some aspects of
the health care system, including the cost, including the security of their coverage. So depending
on how this plays politically, I think there is the foundation for building support for broader
legislation.''

\section{The House Trap}

\lettrine{E}{dward} and Maria Moller are worried about losing their house -- not now, but in 2013.

That is when the suburban San Diego schoolteachers will see their mortgage payments jump, most
likely beyond their ability to pay.

Like millions of buyers during the boom, the Mollers leveraged their way into a house they could not
otherwise afford by taking out a loan that required them to make only interest payments at first,
putting off payments on the principal for several years.

It was a ``buy now, pay later'' strategy on a grand scale, meant for a market where home prices went
only up, and now the bill is starting to come due.

With many of these homes under water -- worth less than the loans against them -- many interest-only
mortgages will soon become unaffordable, as the homeowners have to actually start paying principal.
Monthly payments can jump by as much as 75 percent.

The Mollers owe so much more than their house is worth, and have so few options, that they are
already anticipating doom.

``I'm praying for another boom,'' said Mr.~Moller, 34. ``Otherwise, we'll have to walk.''

Keith Gumbinger, an analyst with HSH Associates, said: ``This is going to be the source of
tomorrow's troubles. The borrowers might have thought these were safe loans, but it turns out they
bet the house.''

After three brutal years, evidence is growing that the housing market has turned a corner. Sales in
July were the highest in a year, and August gives signs of having been even better. In nearly all
major cities, home prices are now rising.

Celebration, however, might be premature. The plight of the Mollers and many others in a similar
position is likely to weigh on any possible recovery for years to come.

Experts predict a steady drumbeat of defaults over much of the next decade as these interest-only
loans mature. Auctioned off at low prices, those foreclosed houses could help brake any revival in
home prices.

Interest-only loans are not the only type of exotic mortgage hanging over the housing market.
Another big problem is homeowners with ``pay option'' loans; in many of these loans, principal
balances are actually increasing over time.

Still, interest-only loans represent an especially large problem. An analysis for The New York Times
by the real estate information company First American CoreLogic shows there are 2.8 million active
interest-only home loans worth a combined total of \$908 billion.

The interest-only periods, which put off the principal payments for five, seven or 10 years, are now
beginning to expire. In the next 12 months, \$71 billion of interest-only loans will reset. The year
after, another \$100 billion will reset. After mid-2011, another \$400 billion will reset.

John Karevoll, a longtime senior analyst for MDA DataQuick, sees the plight of interest-only owners
this way: ``You're heading straight for a big wall and you can't put the brakes on.''

The greater the length of the interest-only period, the more years the owners have to hope for a
recovery, government help, or a miracle. But a long interest-only period works against them, too. A
loan that is interest-only for its first 10 years means the entire house has to be paid off in the
next 20 rather than the more typical 30 years.

One possible solution: start paying extra each month now to pay down the principal before the loans
reset. But many homeowners took out the maximum they qualified for, and don't have the means to pay
more, or at least not enough to make a sizeable dent in the principal.

A decade ago, interest-only loans were rare. But as the boom heated up and desperate buyers sought
any leverage they could, these loans became ubiquitous. They were especially popular in Florida,
Nevada and above all California. In 2004, nearly half of all buyers in the state got one.

The Mollers bought in 2005, paying \$460,000 for their three-bedroom, thousand-square-foot house. A
quick refinance a few months later supplied cash to pay debts. Now the house is worth perhaps
\$310,000. After their interest-only period is up, they expect their monthly payments to increase 20
percent if not more.

``Everyone out here always preached\cite{preach} to me, 'Buy real estate. It's the best investment
you'll ever have,' '' said Mr.~Moller, who grew up in Iowa. ``Then all this stuff started crumbling
and I was like, 'You're kidding me.' ''

While default may be a long way off, the prospect is already dampening the couple's spending habits.
They are postponing the new windows the house needs. They recently bought a 2005 Nissan Murano
instead of a new car, and they have put off buying a flat-screen TV.

Mark Goldman, a San Diego mortgage broker, said many interest-only buyers thought they would be in
control when the loans reset. ``They expected to move or refinance,'' he said. ``But you can't do
either when you're under water.''

Among the people Mr.~Goldman put into interest-only loans was himself. He refinanced five years ago
to shrink his payments so he and his wife, Julie, could put their two sons through college. When the
interest-only period expired a few months ago, their payments went up by 40 percent.

The Goldmans have been in their house for 20 years, which means they still have some equity. Still,
they are unhappy to find themselves in ``a world different than we planned for,'' said Mr.~Goldman,
a lecturer in real estate finance at San Diego State. ``If you purchased your home with an
interest-only loan between 2003 and 2006, you're cooked.''

The federal government, through the finance company Fannie Mae, increased the scope of a program
this summer that might help some interest-only borrowers by letting them refinance. But it will not
help many in coastal California. Only loans owned by Fannie Mae are eligible, and during the boom
Fannie had a limit of \$417,000 -- not enough to buy a home at the peak in a middle-class community.

Dean Janis, a Southern California lawyer who bought a \$950,000 home in 2004, will see his
interest-only loan reset in December. He calculates that will send his payments up a minimum of 27
percent, to \$3,726. A rise in rates could eventually push it as high as \$6,700.

``I understand I took a risk,'' Mr.~Janis said. ``But I did not anticipate that the real estate
market would go down 30 percent.'' He talked with Wells Fargo about his options, and the lender said
he had none.

Homeowners with interest-only loans have a much greater likelihood of default, the First American
CoreLogic figures indicate. Nationally about 18 percent of prime interest-only loans are at least 60
days delinquent. In California, the level is even higher: 21 percent, a rate exceeded only in the
other bubble states of Florida and Nevada.

``The bailout is not trickling through to help many of us who have worked hard, under very difficult
circumstances, to keep paying our bills,'' Mr.~Janis said. ``I am stuck with nowhere to go -- absent
trashing my credit and defaulting.''

\section{Kraft Chief Has History of Success}

\lettrine{W}{hen} Irene B.~Rosenfeld, the chief executive of Kraft, was studying marketing at
Cornell in the late 1970s, she was determined to get as many people as possible to respond to a
survey she mailed out as part of her Ph.D. research on how consumers choose products.

So she attached a crisp \$1 bill to the surveys, to entice\cite{entice} people to fill them out and
mail them back.

Ms.~Rosenfeld now has a much bigger budget, but she is still working hard to convince people to go
along with her plans.

Kraft has just had its \$16.7 billion offer for Cadbury, the British candy company, rebuffed -- but
Ms.~Rosenfeld appears determined to push the deal forward.

She flew to London last month to lay out her merger plans to Cadbury's chairman, Roger Carr. After
the Cadbury board brushed off the overture, she decided to make the offer public on Monday in hopes
of forcing the company's hand.

In public comments she has insisted the deal would benefit the shareholders of both companies. And
when Cadbury insisted the offer undervalued the company, Ms.~Rosenfeld said that the chocolate and
gum maker had far less potential to grow on its own. .

``This announcement just shows how ambitious she is,'' said Erin Swanson, an equity analyst at
Morningstar. ``She's not content with the status quo. She's striving for continuous improvement and
looking to stir growth in what has been a more mature business.''

The events have set the stage for a bidding war, meaning Kraft might have to pay substantially more
if it wants to acquire Cadbury.

Anticipating that possibility, investors bid down Kraft shares, which dropped nearly 6 percent on
Tuesday, closing at \$26.45. Cadbury shares, as measured in depository receipts, soared almost 40
percent, closing at \$51.88, a gain of \$14.42.

Analysts said that several other companies could also bid for Cadbury, including Nestl\'e and
Hershey's.

The coming battle will certainly test Ms.~Rosenfeld's leadership of the company she took over in
2006, vowing to revive sluggish sales of its brands, which include Oreo cookies, Oscar Mayer lunch
meats and Velveeta cheese.

Analysts said that the company's results in the last two quarters were promising.

``There seems to have been some evidence that she has in fact started to turn things around,'' said
Matt Arnold, a consumer analyst with Edward Jones, a retail brokerage firm based in St. Louis.

Ms.~Rosenfeld has said she wanted to encourage innovation by giving managers at many different
levels of the company more power to make decisions and try new things. Mr.~Arnold said that appeared
to have made the company more nimble\cite{nimble} and quicker to introduce new or expanded product
lines.

He said Kraft had success with an expanded line of DiGiorno frozen pizzas, which the company now
sells in single-serve packages and in a premium variety. It has also revamped its Maxwell House
coffee blend and packaging and has begun selling Oreo Cakesters, a snack cake based on the cookie.

``They just found ways to cater to consumer wants a little more, and the bet paid off,'' Mr.~Arnold
said.

Ms.~Rosenfeld also turned back earlier efforts at cost-cutting that hurt the quality of the
company's products. For instance, she insisted on adding more cheese back into the mix for the
company's signature Macaroni and Cheese, after it had been cut back to save money.

While the Cadbury deal would be by far her biggest move to date, Ms.~Rosenfeld has not shied from
buying and selling. She sold the underperforming Post cereals division to Ralcorp for \$1.7 billion
in 2007, and in the same year paid \$7 billion for the cookie unit of the French company Groupe
Danone.

Several analysts said that a Cadbury acquisition could make sense for Kraft, which would greatly
expand its confectionery\cite{confectionery} business. Cadbury would also give it strong sales in
emerging markets, like India, where the chocolate maker has strong sales and has seen impressive
growth.

But they cautioned that those considerations could change if Kraft winds up paying significantly
more.

Ms.~Rosenfeld, who declined to be interviewed for this article, began working for Kraft in the early
1980s. She soon became a marketing manager with responsibility for Kool Aid and helped increase
sales, in part by changing television ads aimed at children that featured a rock and roll
soundtrack. She had similar success with Jell-O and was later made head of Kraft's Canadian and then
North American divisions.

But she left the company abruptly in 2003. In 2004 she was hired to be chief executive of PepsiCo's
Frito-Lay division. But two years later she returned to Kraft, this time as chief executive. She was
named chairman in 2007 and now holds both posts.

Ms.~Rosenfeld's thesis adviser at Cornell, Vithala R.~Rao, still teaches marketing there, and said
on Tuesday that he recalled the \$1 bills and the innovative way his former pupil conducted her
thesis survey nearly three decades ago (she received her Ph.D. in 1980).

Mr.~Rao said the trick worked, with the survey getting a higher response than might have been
anticipated without the incentive.

Mr.~Rao recalled that she ``always exhibited initiative in doing things,'' adding that ``she
obviously has done extremely well.''

\section{To Attract a Buyer, Make the Deal Irresistible}

\lettrine{G}{iven} the relentlessly dismal\cite{dismal} retail sales numbers, consumers seem to have
locked away their wallets and hid the key.

But some recent promotions show that shoppers will still chase a bargain, particularly one that
seems unlikely to be repeated anytime soon. Even big-ticket items costing hundreds or thousands of
dollars can be in high demand if they are priced right.

``The otherwise frugal\cite{frugal} consumer is willing to shell out for a blockbuster value,'' said
John D.~Morris, a managing director with BMO Capital Markets, ``whether it's cash for clunkers,
khakis or a cruise.''

In the last month, promotions like the government's cash-for-clunkers auto purchase plan and JetBlue
Airways' \$599 monthlong flying pass have won over otherwise wary consumers.

The pattern appears similar to what is playing out in the housing market, where sales are rebounding
from dismal levels a year ago as consumers take advantage of foreclosure sales and start to believe
that prices have hit bottom.

That has not been the case with retailers, where consumers have held off spending for months,
betting that a 40-percent-off sale was just a warm-up for an inevitable 50 percent sale. To overcome
the newfound patience of consumers, many retailers are trying to spur impromptu buying by putting
definitive time limits on sale items. The tactic has resonated\cite{resonate}, especially with
upper-middle-class to upper-class consumers.

``There are still shoppers that still have money,'' said Marie Driscoll, director of consumer
discretionary retail for Standard \& Poor's Equity Research. ``They're just much more
discriminating.''

Luxury-goods retailers are finding buyers, at the right price. At TheOutnet.com, a Web site selling
designer fashion at a discount, Alexander McQueen leather pants sold out at \$1,002 (reduced from
\$3,340).

Rabbit-fur coats by Roberto Cavalli also sold out for \$1,498 on Gilt.com, another high-end site,
marked down from \$5,030.

Hana Ben-Shabat, a partner in the retail practice of A.~T. Kearney, said the consumer buying luxury
items at a discount these days probably was earning \$200,000 to \$300,000 a year.

If that person's compensation was partly based on a bonus, overall pay has probably dropped
significantly. That means that such consumers may no longer be visiting the likes of Prada and Gucci
as often as they once did. But these buyers also know what such merchandise typically costs.

``They have a basis for comparison,'' Ms.~Ben-Shabat said. ``They see something that maybe last year
was \$2,500, and they see it below the \$2,000 mark and they think, 'I will buy it.' ''

Donald R.~Grimes, an economist and senior research associate at the Institute for Research on Labor,
Employment and the Economy at the University of Michigan, said many of these buyers sat out the
recession, because they were either nervous about the economy or because they did not want to appear
ostentatious\cite{ostentatious}.

But their confidence is slowly returning.

``Now, they feel, 'I've made it through,' '' Mr.~Grimes said. ``They're definitely looking for
deals.''

JetBlue's promotion drew in Sam Rosen, a Chicago Web site designer. He was planning to spend \$400
on airfare to attend a friend's bachelor party in Las Vegas this month. But when he saw a notice for
JetBlue's ``All You Can Jet'' pass, Mr.~Rosen jumped.

``I'm a smart enough guy to realize that for an extra \$200, I could go wherever I want,'' Mr.~Rosen
said.

JetBlue halted the program last month, earlier than it had planned, when it sold out of the passes
it designated under the plan. It has not said how many were available.

The first travelers begin using their passes on Tuesday; trips must be completed by Oct.~8. Along
with his Las Vegas trip, Mr.~Rosen will go to six cities before his month of flights is up.

He is visiting two places -- Boston and St.~Martin -- where he has never traveled, and will work
while on the road.

Another Chicago resident, Frank Earullo, could not resist the \$3,500 rebate he received when he
bought a Hyundai Santa Fe sport utility vehicle under cash for clunkers.

He was among the last of nearly 700,000 people who bought vehicles during the offer, which expired
Aug.~24. It helped a number of automakers post better August auto sales compared with dismal results
in 2008, although many analysts expect sales to fall sharply now that the offer no longer exists.

``They're pretty much forcing me to do this, it's such a good deal,'' he said of his new Santa Fe,
bought with a trade-in of his 1999 Ford Explorer.

Other promotions, though less ambitious than the clunkers rebates, are succeeding with rare price
cuts.

Zingerman's, a restaurant and food company in Ann Arbor, Mich., has for years offered a weeklong
``bake-cation'' seminar for \$1,000, where students spend five days learning to make everything from
croissants to pie.

This summer, it added a weekend version, for \$500, and the first session sold out in days (a second
has been added for October). Last month, Zingerman's also started a training course for food
professionals on cheeses, and priced the introductory session at \$975, rather than the \$1,200 it
plans to charge for subsequent sessions.

The offer attracted three participants from Dorothy Lane Market, a small chain of high-end specialty
food stores in the Dayton, Ohio, area. The store has sent employees around the world for training
classes.

Supervisors ``won't cut corners when it's something worth our time to do,'' said Dave Mader, manager
of specialty cheese at the chain's Springboro store.

Still, a broader recovery depends on the willingness of less-flush consumers to buy, and Professor
Grimes said he did not envision a return to the kind of impulsive spending that seemed commonplace
earlier this decade.

``People who are making \$80,000 or \$100,000 and who would go off to Disney World and blow \$2,000
on a weekend, they aren't going to do that,'' he said.

Still, the popularity of cash for clunkers was a boost for economic morale\cite{morale}, he said,
and could bode well for similar programs down the road.

``At least it gave some optimism, and got people doing something,'' he said.

\section{Palm Unveils a Smartphone for Younger Users}

\lettrine{T}{aking} the next step toward rebuilding its lineup of smartphones, Palm announced on
Wednesday a new cellphone called the Pixi.

Palm also announced it was dropping the price of its Pre smartphone by \$50 to \$150, with a
two-year service agreement and after a \$150 instant rebate and \$100 mail-in rebate.

Given the name, the Pixi is aimed at younger customers than those who bought the Pre. The company
said the phone would have a smaller touch screen and be slimmer than the Pre, but it would share
many of the same features, including eight gigabytes of storage and GPS navigation capabilities.
Like the Pre, the Pixi runs on the company's new operating system, called WebOS, which promises
speedy Web browsing.

``It's cute and it's a lot of fun,'' said Jonathan J.~Rubinstein, chief executive of Palm. ``Whereas
the Pre is a little more serious.''

For Palm, creating a family of smartphones is an important step toward establishing itself as a
contender in an increasingly competitive wireless market.

``This is a good step forward for Palm,'' said Michael Gartenberg, vice president for strategy and
analysis at Interpret, a market research firm based in Los Angeles. ``It shows that the Pre is not
just a one-hit wonder, but a family of devices that users can embrace.''

Mr.~Rubinstein said Palm had had success before with the strategy of taking a smartphone into the
mainstream and make it appeal to nonbusiness users. In 2007, it introduced the Centro smartphone, a
more stylish version of the Palm Treo, to attract everyone but the business user. Palm says it
believes that a gap in the market between business smartphones and general-use cellphones still
exists. ``There are a tremendous number of people who are tired of traditional feature phones and
want a smartphone,'' he said. ``A lot of our Centro customers came from exactly that, and so we see
this a much larger opportunity for that market.''

To capture the fancy of younger users, the Pixi will also have a removable backplate that can be
swapped with a series of interchangeable panels, the first of which the company plans to commission
from illustrators, graphic artists and fashion designers. Instead of having a slide-out keyboard
like the Pre, the Pixi has a full, exposed keyboard, which Palm executives hope will attract a
younger buyer who favors features like text-messaging.

The company said the Pixi would be distributed exclusively by Sprint and be available in time for
the holiday shopping season. Although the company did not disclose the price of the smartphone,
Mr.~Rubinstein said it would be less expensive than the Palm Pre, which was initially priced at
\$200 after a rebate, with a two-year contract.

Almost equally crucial to Palm gaining headway against competitors like Apple and Research in
Motion, Mr.~Gartenberg said, was cultivating a system of developers eager to write and publish
programs that could run on its devices.

``The more devices the better, the more apps the better,'' said Mr.~Gartenberg.

Currently, the company is only accepting applications into its app store on a limited basis. It said
that tens of thousands of developers had downloaded the WebOS software development kit, the main set
of tools needed to write apps, since the company made it available in July. This fall, Palm plans to
begin reviewing apps submitted by developers for sale in its app store.

``We've been very methodically developing an incredible App Catalog experience,'' said
Mr.~Rubinstein. ``What we have accomplished is really terrific, but we are just beginning.''

\section{11th-Hour Filings Oppose Google's Book Settlement}

\lettrine{A}{fter} a flurry\cite{flurry} of last-minute filings on Tuesday, a federal judge must now
begin untangling the mountain of competing claims about how a legal settlement granting Google the
right to create the world's largest digital library and bookstore would affect competition, authors'
rights and readers' privacy.

The \$125 million proposed settlement among Google, the Authors Guild and the Association of
American publishers, which is awaiting review by Judge Denny Chin of the Federal District Court for
the Southern District of New York, has prompted dozens of opposing filings from individuals, rival
companies like Amazon and Microsoft, advocacy organizations, groups representing authors and
publishers and even some foreign governments.

It has also received the support of companies like Sony, civil rights groups and some antitrust and
economics experts in academia.

Legal scholars say that Judge Chin will have to address not only whether the settlement is fair to
the authors, publishers and rights holders covered by it, but also whether it benefits the public at
large.

``The number and quality of opposition filings is very unusual,'' said Jay Tidmarsh, a professor of
law at Notre Dame Law School. ``The court is going to have to look at the public interest in the
settlement.''

The agreement, which would bring millions of rarely seen books online, has clear benefits to readers
and authors. But scholars say the judge is likely to weigh those benefits against arguments that the
settlement would limit competition. Opponents say it would give Google a quasi-exclusive license to
profit from millions of out-of-print books and create a consortium that would have power to set
prices for digital books. Google, the Authors Guild and the Association of American Publishers have
vigorously disputed those claims, but the claims are being investigated by the Justice Department.

Opponents have raised other issues including contending that the settlement tramples on the rights
of some authors and that it does not protect the privacy of readers. The court has the power to
either approve or strike down the settlement, an option that would revive the lawsuits filed in 2005
by the authors and publishers against Google over its plan to digitize millions of books from
libraries without authorization from rights holders. But if Judge Chin finds problems with the
settlement, he could also offer the parties a road map for overcoming them.

``If the judge has some significant concerns, it is much more likely that he would invite the
parties to address those concerns rather than reject the agreement,'' said Andrew I.~Gavil, a law
professor at Howard University. Professor Gavil said that Judge Chin was likely to give special
consideration to the opinion of the Justice Department, which has until Sept.~18 to make its views
known. A hearing on the settlement is scheduled for Oct.~7.

On Tuesday, several groups filed briefs in opposition, including Microsoft and Yahoo, and a
coalition representing those companies and others. The coalition, which calls itself the Open Book
Alliance, opposes the agreement on antitrust grounds. The group is co-led by Gary Reback, an
antitrust attorney in Silicon Valley who in the 1990s helped persuade the Justice Department to file
its landmark antitrust case against Microsoft. He said the court could address some of the antitrust
objections by forcing Google to license its database of digital books to others.

``Google should be ordered to license the database with all attendant rights to a number of
competitors, under the supervision of the Justice Department,'' Mr.~Reback wrote in the brief. He
traced the birth of Silicon Valley to a similar ``compulsory license'' mandated by the Justice
Department. ``Silicon Valley exists precisely because the Antitrust Division ordered AT\&T to
license its key invention, the transistor, for nominal payments,'' he wrote.

Defenders of the agreement say the antitrust concerns are unfounded, and argue that others besides
Google could obtain similar licenses without any mandates from the court.

``We have never said that the same kinds of outcomes would not be available to Microsoft or Amazon
or anyone else who is willing to make the same investments,'' said Richard Sarnoff, former chairman
of the Association of American Publishers and co-chairman of the American unit of Bertelsmann, the
parent company of Random House. ``We have a road map to do it now.''

\section{Presidential Pep Talk Kicks Off Year for Students}

\lettrine{M}{illions} of American schoolchildren, oblivious to the uproar that preceded a
back-to-school speech by President Obama, heard him exhort them to greatness on Tuesday, watching,
applauding and in some classrooms cheering a nationally broadcast address that urged them to set
high goals, knuckle down in their studies and persevere through failure.

``Don't ever give up on yourself, because when you give up on yourself you give up on your
country,'' Mr.~Obama told students packed into a high school gymnasium in a Washington suburb.

Several school districts in Maryland, Texas, Virginia and other states, where clamor by
conservatives accusing the White House of partisan motives was loudest, decided not to show the
speech. Some school officials said schedules were too packed to accommodate a presidential
interruption, while others said they had taped the speech to show later this week.

And in thousands of schools that did show it, some children opted out after they or their parents
decided they should not watch. At a 3,700-student high school in suburban Chicago, that meant 250
students sat out the speech in a cafeteria and elsewhere. At an elementary in Greensboro, N.C., a
fourth-grade girl sat alone, eating her lunch as scores of her classmates watched the president in
an adjoining classroom.

But certainly a large segment of the nation's 50 million public school students watched the speech,
and it seemed to resonate with many, including Mariah Key, 10, who watched with fifth-grade
classmates at M.~Agnes Jones Elementary School in Atlanta. Mariah said that her grandparents had
recently divorced, and that money was tight at home.

``President Obama was saying to put everything else aside and focus on what's important -- your
education,'' Mariah said.

The speech was the first of its kind to be nationally broadcast by a president since the elder
President George Bush delivered a speech in 1991 from a Washington school urging students not to use
drugs.

By the time Mr.~Obama took the microphone at Wakefield High School in Arlington, Va., just before
noon, most of the sizzle and smoke was gone from the conservative firestorm leading up to the
address.

Right-wing radio commentators had moved on to other topics, and the infuriated parents who besieged
school offices in at least a dozen states from Virginia to California had mostly stopped calling.

``I'm still puzzled over the controversy,'' said Benny Gooden, superintendent of the Fort Smith
Public Schools in western Arkansas. ``We got word that the president would admonish students to do a
good job, and I said, it sounds like he's going to give my speech.''

Still, ``multiple hundreds'' of parents called Mr.~Gooden's district offices. The district gave
parents a ``nonparticipatory option,'' but, he said, some parents nonetheless kept their children
home.

``Now that's a fine solution -- keep them home from school,'' he said. ``This is not a conservative
or liberal or libertarian issue, this is an American issue, of wanting every young person to know
that if you apply yourself you can have a bright future. The president gave that message eloquently,
and I hope our students took it to heart.''

Millions of students watched, but not all with rapt attention. At Bolingbrook High School southwest
of Chicago, about 200 students viewed the address on a large screen in an auditorium, and a few of
them looked downright drowsy.

``It made me want to fall asleep,'' said Sarah Vogt, a sophomore. One of Sarah's classmates seemed
to take Mr.~Obama's study-hard message a bit too literally: calculator in hand, she worked on
homework throughout the presidential address.

Andrew Quick, a 15-year-old student at Oviedo High School near Orlando, Fla., said he considered the
speech to be a potentially disruptive interruption of his school day, so decided not to watch it.
``I'm a Republican,'' he said, ``and I really don't like Obama all that much.''

The uproar over the speech erupted most vigorously, perhaps, in Texas, and some districts there
reversed the opt-out procedure: the only students they allowed to watch the president were those
whose parents called or e-mailed with explicit permission.

J.~W. Williams Middle School in Rockwall, Tex., used that system. The result was that 100 of the
school's 900 students filed into a library to see the speech projected on a big screen, among them
Haley Welch, 13, a blue-eyed eighth grader with braces, who wore a T-shirt saying ``Obama said knock
you out!''

``He made me think I can do anything I want to,'' Haley said after the speech. ``Even if I have
problems at home, school is a different place where I can leave that behind.''

About 25 miles to the southwest at Booker T.~Washington High School in downtown Dallas, Ron Kirk,
the United States trade representative, watched the speech with students. Interviewed by reporters
there, Mr.~Kirk, a former Dallas mayor, said of the controversy over the speech, ``There are few
moments in my life when I'm embarrassed to say I'm from Texas.''

``This is one of them,'' he added.

Mr.~Kirk was one of 23 cabinet members and other top administration officials who fanned out to
visit schools across the country to hammer home Mr.~Obama's message. In New York, Secretary of State
Hillary Rodham Clinton visited the Manhattan Charter School.

At Jericho High School on Long Island, Jasleen Kaur, a 16-year-old senior, said the speech sounded a
lot like pep talks she had heard many times from her father, whom she described as a staunch
Republican businessman.

``They both want to motivate us because honestly we are the future,'' Ms.~Kaur said.

And as Mr.~Obama ended his 25-minute pep talk, a dozen students in a classroom at Cameron Elementary
School in West Covina, Calif., sat up and clapped.

``I liked the speech,'' said Christian Cristin-Sanchez, summing up views of several classmates. ``He
said you have to work hard, to always listen to others and have respect.''

At the Hill Campus of Arts and Sciences middle school in Denver, seventh graders in Sue Hoopfer's
medieval world history class discussed the speech after it concluded. Ms.~Hoopfer asked what
resonated most, and a flurry of hands shot up.

``That someone always has it worse than you,'' Noah Adelstein said.

``Asking a question is not a sign of weakness, it's a sign of strength,'' Ella Zeiler added.

Naomi Kamidate expressed another thought: ``I think it was pretty cool to have the leader of the
country speaking to us.''

\section{First Stem Cell Drug Fails 2 Late-Stage Clinical Trials}

\lettrine{W}{hat} might become the first drug derived from human stem cells failed in two late-stage
clinical trials, dealing a setback to the drug's developer and to the stem cell field.

The developer, Osiris Therapeutics, said Tuesday that its drug, Prochymal, was no more effective,
over all, than a placebo in treating a life-threatening complication of bone marrow transplantation,
though certain patients seemed to be helped.

Shares of Osiris, which is based in Columbia, Md., fell 34 percent, to \$8.03.

Prochymal is a preparation of mesenchymal stem cells, which are obtained from the bone marrow of
healthy young adults.

Because the cells are derived from adults, they sidestep the ethical issues stemming from the
destruction of human embryos needed to make embryonic stem cells. Unlike most other types of adult
stem cells, mesenchymal cells grow well in culture, so thousands of doses can be produced from a
single donation.

Stem cells, particularly in the form of bone marrow transplants, are already used in medicine.
Osiris is hoping that Prochymal will become the first stem cell product approved by the Food and
Drug Administration and sold as a mass-produced pharmaceutical product.

But the failure in the two trials could make it hard to reach that goal. Both trials tested
Prochymal as a treatment for graft-versus-host disease, which occurs when immune cells in donated
marrow attack the recipient's organs as foreign tissue.

In one trial, in which Prochymal was used along with steroids, 45 percent of patients responded to
Prochymal and steroids compared with 46 percent who had a response to steroid and a placebo.

In a second trial, in which Prochymal was tested in patients who were not benefiting from steroids,
35 percent of those getting the drug had a resolution of graft-versus-host disease for at least 28
days, compared with 30 percent getting the placebo. The difference was not statistically
significant.

Osiris said, however, that in the second trial, the drug did provide a statistically meaningful
benefit in patients having graft-versus-host disease that specifically affected their livers or
their gastrointestinal tracts.

C. Randal Mills, the company's chief executive, said he hoped the drug could be approved for those
patients, given the seriousness of the disease. ``Prochymal is having a clear effect in the liver,
and that is the most underserved patient population,'' he said in an interview.

But the F.D.A. is usually reluctant to approve a drug based on its working in only a subset of
patients.

In March, Osiris stopped enrollment in a trial testing Prochymal as a treatment for Crohn's disease,
saying it was unlikely the drug would be better than a placebo, because of a high positive response
to the placebo. Genzyme has the right to sell Prochymal outside the United States and Canada.

Dr.~Darwin J.~Prockop, an expert on mesenchymal stem cells who was not involved in the trials, said
there were still a lot of unknowns about how the cells work. ``Understanding it well enough to
translate to the clinic -- that's the hurdle we're at,'' said Dr.~Prockop, who is director of the
Institute for Regenerative Medicine at Texas A\&M Health Science Center.

\section{You Think Life Is Mild North of the Border?}

\lettrine{A}{ny} lingering illusion that Canada is a milder, blander\cite{bland} version of the
United States is dispelled by ``Durham County,'' a Canadian-made crime series that begins on Monday
on the Ion network.

It takes a while to realize that this scary, well-made thriller was not made in America. ``Durham
County'' has all the signposts of a Hollywood production: lush theme music and stark cinematography,
a brooding homicide detective, barbaric rapes and murders, philanderers, sociopaths, yoga moms and
alienated teenagers.

But the landscape, chosen for the forests of electrical power lines that forebodingly crosshatch the
horizon, is hard to place, and some of the accents have a slightly unusual lilt. Mostly, though,
it's the violence that comes with an unfamiliar inflection.

Nothing is too gory for American television. Even the more lighthearted procedurals like ``NCIS''
and ``Bones'' zoom in on dismembered limbs, maggot-covered corpses and swollen, rotting bellies.
Most American crime dramas punch up brutality with musical scores and lurid sound effects: the
squish of a knife to the stomach, the crack of a fist on a jaw, the blast of a gun to the temple.
This show, on the other hand, turns down the volume at those moments, and at times even averts the
camera, which only amplifies the horror; viewers are stripped of a preconceived and familiar
catharsis and left with savage acts that look all the more real for not being artistically enhanced.

And nothing on ``CSI'' or ``Southland'' is as disturbing as this program's opening picnic scene. A
pedophile hosts a fête champêtre for two teenagers in schoolgirl uniforms that turns into tragedy,
all of it seen and uninterrupted by a voyeur hidden in the bushes.

``Durham County,'' in short, is very, very creepy and unsettling, and entirely addictive, a modern
murder mystery with a touch of Patricia Highsmith misanthropy. Even bystanders seem a little
sinister, including the young children who play quietly while wearing masks of Japanese anime
characters.

It's an odd but promising choice for the little-known Ion, which only a few years ago was the Pax
network, provider of wholesome family shows and reruns of ``Growing Pains.'' Ion is trying to appeal
to more adult tastes with reruns of ``Boston Legal'' and ``Criminal Minds.'' This series appeals to
grownups who prefer their crimes served fresh and with film noir understatement.

Mike Sweeney (Hugh Dillon, ``Flashpoint'') is a Toronto homicide detective who has relocated to the
suburbs to start a new life after his partner was killed and his wife, Audrey (H\'el\`ene Joy),
developed breast cancer. Audrey is in recovery, still wearing scarves and wigs, and hoping to
rekindle her relationship with her husband and two daughters, particularly the estranged eldest,
Sadie (Laurence Leboeuf), a teenager who hates their shabby housing development and dreams of
becoming an investigator like her father. (As a hobby Sadie creates crime scenes in her dollhouse,
molding tiny clay figures into bloodied murder victims.)

Sweeney thinks he is making a fresh start, but old troubles crop up wherever he goes, including in
the form of his new neighbor, Ray Prager (Justin Louis), a onetime promising hockey star who went to
high school with Sweeney and still bears a grudge. Prager is a narcissistic charmer who bullies his
peppy, pretty wife and their shy, bookish son.

The six-episode series has already had two seasons in Canada, but it never gives viewers a
reassuring sense that the hero will prevail. Nor does it amp suspense by withholding the killer's
identity; the viewer knows who did it, and the mystery is in how the detective will solve the case.
A little like Helen Mirren in the British series ``Prime Suspect,'' Sweeney has enemies, but he is
by far his own worst one.

``Durham County'' evokes a world etched in uncertainty and suspicion, where a menacing atmosphere is
darkened by a toxic environment and a hostile community. There is nothing gentle or folksy about it.
It's not what a lot of people think of when they think of Canada, but it's a lot more intriguing.

\section{The Rise of the Female Anchor}

\lettrine{O}{NE} female network TV anchor is a breakthrough. Two become a catfight.

That equation is almost inevitable no matter who the women are who make it to the top of television
news. In the case of Diane Sawyer and Katie Couric, who in January will resume their former
morning-show rivalry on the evening news, it's already printed up in the program of public
perception.

So before the Betty \& Veronica comparisons flare up again, it's worth noting what doesn't fit
classic clich\'es of female advancement on television. The breakthrough that Ms.~Sawyer accomplished
isn't that she became the first solo female anchor. Ms.~Couric took that trophy -- and the attendant
triumph and flack -- in 2006. Ms.~Sawyer's accomplishment is more subtle -- she is a gorgeous,
glamorous television personality who got the top job by waiting around.

Particularly when it comes to its female stars, the television news business favors ``All About
Eve'' coups and annexations -- plucking a fresher face from the chorus line in the case of Deborah
Norville or raiding enemy territory as CBS did when it wooed Ms.~Couric away from NBC. Ms.~Sawyer
has spent more than 10 years as a co-host of ABC's ``Good Morning America,'' and didn't flinch or
change networks when ABC experimented with a team of younger anchors after Peter Jennings died, and
still later chose her former ``G.M.A.'' co-host, Charles Gibson.

Patience is not normally a virtue in the news business, but Ms.~Sawyer made it her ally, letting
time smooth bumps in her r\'esum\'e that at one time seemed insurmountable. She moved directly from
working with former President Richard Nixon on his memoirs to CBS News back when the line between
journalism and government was virtually inviolate -- until, that is, Tim Russert and George
Stephanopoulos came along and changed the rules. Early on, her golden allure and teen beauty queen
status were almost suspect -- she was accused of being too fetching to be a ``serious'' journalist.
Time has not altered her appearance -- at 63 she is almost absurdly good looking -- but mores have
loosened. In 2003, Meredith Vieira did a stint in the Broadway show ``Thoroughly Modern Millie.''
When Ms.~Couric sat in for Jay Leno in 2003, workers cut through her desk to expose her gams.

Women who let their ambition show too openly are usually punished for it. Ms.~Sawyer, who is as
relentless and driven as any of her peers, makes an art of coy deflection. The morning after ABC
made the announcement, Ms.~Sawyer's response was almost geisha-like in its deference to Mr.~Gibson.
``It's just such an honor to be a part of any team he's captained,'' she said softly.

Waiting, of course, has at least one drawback. Ms.~Sawyer and Ms.~Couric will outnumber their male
counterpart 2 to 1 in an era when networks are losing their primacy and even their creative
advantage over cable, and network news, in particular, is sinking in relevance and prestige. (Back
in 1976, when Barbara Walters was lured by ABC, the network very likely considered the anchor job
too important to entrust to a woman and paired Ms.~Walters with the more authoritative-seeming Harry
Reasoner.) As in other fields, women seem to break through the glass ceiling just as the
air-conditioning is being turned off in the penthouse office suites. Women anchors may turn out to
be what women doctors once were in the Soviet Union, a majority without status or financial
advantage.

Nowadays, viewers tend to treat network evening news shows less as a source of information than as a
weather vane -- if all three networks go live to a presidential speech about health care, it must be
important; if none of them carry a dead statesman's funeral, he must not have mattered so much after
all. Network news may not be as important as it once was, but, oddly, the role of anchor hasn't
evolved much at all since the days of Mr.~Jennings and Tom Brokaw. It is one that suits Ms.~Sawyer's
velvety voice and regal demeanor. Ms.~Couric, who tried to tinker with the part and inject it with
some of her personality and morning-show pizazz, failed, and has since returned to the more
traditional, starched style favored by the NBC anchor Brian Williams, who, perhaps not
coincidentally, has a commanding lead in the ratings over ABC and CBS.

Ms.~Sawyer is not likely to have difficulty adopting a more neutral, impassive persona for the new
job. She is at her best, and most natural, when avoiding human emotion. Her worst on-camera moments
have been those awkward, almost mawkish pantomimes of empathy or outrage -- her 2000 ``Good Morning
America'' interview with the 6-year-old Cuban refugee Eli\'an Gonz\'alez, during which she crayoned
with the boy and stood on her head, or the way she scolded the Bush-bashing Dixie Chicks in 2003 for
what some saw as a lack of patriotism. She's better, or more believable, in hard news interviews and
trips to North Korea.

It's an odd bit of a role reversal: Ms.~Couric is a morning-show natural who had to tone down her
peppy cadence and casual style to suit the formality of evening news programs, while Ms.~Sawyer,
born to be an evening anchor, spent much of her career twisting her natural elegance into the shape
of slap-happy morning television. And Mr.~Williams, who ascended to the position of NBC anchor on
the shoulders of an old boys' club, now has to reposition himself as a member of a persecuted
minority, the white male anchorman.

\section{Pentagon Keeps Wary Watch as Troops Blog}

\lettrine{O}{ver} the course of 10 months in eastern Afghanistan, an Army specialist nicknamed Mud
Puppy maintained a blog irreverently chronicling\cite{chronicle} life at the front, from the terror
of roadside bombs to the tyrannies of master sergeants.

Often funny and always profane\cite{profane}, the blog, Embrace the Suck (military slang for making
the best of a bad situation), flies under the Army's radar. Not officially approved, it is hidden
behind a password-protected wall because the reservist does not want his superiors censoring it.

``Some officer would be reviewing all my writing,'' the 31-year-old soldier, who insisted that his
name not be used, said in an e-mail message. ``And sooner or later he would find something to nail
me with.''

There are two sides to the military's foray\cite{foray} into the freewheeling world of the
interactive Web. At the highest echelons of the Pentagon, civilian officials and four-star generals
are newly hailing the power of social networking to make members of the American military more
empathetic, entice recruits and shape public opinion on the war.

Gen.~Ray Odierno, commander of American forces in Iraq, is on Facebook. The chairman of the Joint
Chiefs of Staff, Adm. Mike Mullen, has a YouTube channel and posts Twitter updates almost daily.

The Army is encouraging personnel of all ranks to go online and collaboratively rewrite seven of its
field manuals. And on Aug.~17, the Department of Defense unveiled a Web site promoting links to its
blogs and its Flickr, Facebook, Twitter and YouTube sites.

The Web, however, is a big place. And the many thousands of troops who use blogs, Facebook, Twitter
and other social media sites to communicate with the outside world are not always in tune with the
Pentagon's official voice. Policing their daily flood of posts, videos and photographs is virtually
impossible -- but that has not stopped some in the military from trying.

The Department of Defense, citing growing concerns about cybersecurity, plans to issue a new policy
in the coming weeks that is widely expected to set departmentwide restrictions on access to social
networking sites from military computers. People involved with the department's review say the new
policy may limit access to social media sites to those who can demonstrate a clear work need, like
public information officers or family counselors.

If that is the case, many officials say, it will significantly set back efforts to expand and
modernize the military's use of the Web just as those efforts are gaining momentum. And while the
new policy would not apply to troops who use private Internet providers, a large number of military
personnel on bases and ships across the world depend on their work computers to gain access to the
Internet.

To many analysts and officers, the debate reflects a broader clash of cultures: between the
anarchic, unfiltered, bottom-up nature of the Web and the hierarchical, tightly controlled, top-down
tradition of the military.

``We as an institution still haven't come to grips with how we want to use blogging'' and other
social media, said Lt.~Gen.~William B.~Caldwell IV, the commander of the Army Combined Arms Center
at Fort Leavenworth, Kan.

One of the Army's leading advocates for more open access to the Web, General Caldwell argues that
social networking allows interaction among enlisted soldiers, junior officers and generals in a way
that was unthinkable a decade ago.

He requires students at the Command and General Staff College at Fort Leavenworth to blog, and the
college now sponsors 40 publicly available blogs, including his own, where policies are freely
debated.

But getting approval for those blogs, as well as for YouTube and Facebook access at the college, was
a struggle. ``At every corner, someone cited a regulation,'' General Caldwell said. In recent
months, however, ``the Army has made quantum leaps'' in embracing the Web, he added.

Noah Shachtman, editor of Wired.com's national security blog, Danger Room, which has reported
extensively on the new policy review, said he recently asked students at West Point whether they
would allow soldiers to blog. Almost every cadet said no.

``Then I asked, 'How many of you think you can stop the flow of information from your soldiers?' ''
Mr.~Shachtman recalled. ``Everybody agreed there is no way to stop this information from going out
anyway. So there is this sort of dual-headedness.''

Skeptics of the Pentagon review say it is motivated partly by a desire among certain officials to
exert control over the voices of troops on the Web.

Since the advent of military blogging during the Iraq war, some commanders have remained
uncomfortable with the art form, citing concerns about both security and decorum.

Over the years, blogs have been censored or shut down, and several years ago the Army instituted
requirements that bloggers register with their commanding officers and submit posts for review. As a
result, some bloggers say, blogs have become tamer -- or, as in the case of Mud Puppy's blog, gone
underground.

Officials knowledgeable about the review say it is a result of growing concerns at the United States
Strategic Command, which oversees the military's use of the Internet, that social networking sites
make military computers vulnerable to viruses, hackers, identity thieves, terrorists and even
hostile governments. (Those concerns are not focused on the military's secure system for classified
material, which does not use the public Internet.)

The review may already be having a chilling effect. The Marine Corps recently restated a ban on
using any social media on its network. And the Army, which in June gave some bases access to
Facebook, Twitter and other networking sites, recently urged units to avoid creating new social
media pages until the final department policy was issued.

Still, even as they consider restricting the troops' access to social media, the most senior
Pentagon officials have clearly come to view Facebook, Twitter, YouTube and blogging as crucial
elements of their public information operations.

``This department, I think, is way behind our curve'' in using social media, Secretary of Defense
Robert M.~Gates said in July as he extolled the use of Twitter by Iranian dissidents.

To critics, the Pentagon's social media sites are goofy at best, propagandistic at worst. ``It's
like your parents' using modern slang and failing miserably,'' said Sgt. Selena Coppa, who writes a
blog, Active Duty Patriot, which frequently criticizes the Iraq war and, she says, has gotten her
into trouble with her superiors.

But to many troops, the deeper question is whether the military will allow personnel in the field to
use the sites the Pentagon itself wants to exploit. For a generation raised on the Web, any
restrictions will damage morale, those people say.

``What comes out of my blog is the experiences of a soldier right in the middle of all of this,''
Mud Puppy (a nickname for military police), who recently returned home to Illinois, wrote in a
recent e-mail message. ``I think that people need to hear from us, more than they need to hear from
the big whigs. War has a cost, and that cost is paid by soldiers.''

\section{China to Issue Yuan-Denominated Bonds in Hong Kong}

\lettrine{T}{he} Chinese Ministry of Finance said Tuesday that it would issue 6 billion yuan worth
of government bonds in Hong Kong, a major step to internationalize its currency at a time of concern
about the dollar.

The yuan bond issue, the equivalent of \$879 million, will ``promote the yuan in neighboring
countries and improve the yuan's international status,'' the ministry said on its Web site.

``The first step toward internationalization is regionalization,'' Shi Lei, a currency analyst at
Bank of China in Beijing, said during an interview. ``China wants to develop the offshore market in
Hong Kong.''

While domestic banks like Bank of China and the Export-Import Bank of China have issued
yuan-denominated bonds in Hong Kong for a couple of years at the encouragement of Beijing, this is
the first time that government bonds, comparable to U.S.~Treasury securities, are to be issued. The
sale is set for Sept.~28.

In July, the People's Bank of China, the country's central bank, started a program for local
companies to settle trade in yuan, but it has so far spurred little trade. Zhi Ming Zhang, an
analyst at HSBC in Hong Kong, said the government bond issue might show foreign investors they could
rely on the yuan.

``If I'm doing trade with China, where am I going to park this money?'' Mr.~Zhi asked, referring to
the yuan. The yuan-bond market needs security and liquidity to make such settlements attractive, he
said, and government bonds will provide security and a pricing benchmark. The next step, he added,
would be to increase the number of Chinese issuers and investors in the yuan.

Experts estimate that China holds about 75 percent of its \$2 trillion in foreign reserves in
dollar-denominated assets, but since the global financial crisis began, that position has made
Beijing uneasy. Since the beginning of 2007, the dollar has slipped more than 20 percent against the
yen, and more than 12 percent against the yuan, and investors are concerned as the United States
continues to pile up debt to finance its huge stimulus package.

In March, China's prime minister, Wen Jiabao, expressed concern about the dollar's slide and
encouraged the United States to ensure its stability.

While the bond issue announced Tuesday is a step toward making the yuan a global currency, the size
of the sale is small compared with those of U.S.~Treasury securities, and the time it will take to
establish the yuan internationally remains uncertain.

``There is no timetable,'' said Mr.~Shi, the Bank of China analyst, adding that developing the
market would take ``at least three to five years.''

In another move to make the yuan accessible to investors, BOC Suisse Fund Management, Bank of
China's asset-management arm based in Geneva, said Friday that it had received approval from the
Swiss financial regulator to create a new set of funds, nearly half of them denominated in the
Chinese currency.

\section{Times Reporter Is Freed in Afghan Raid That Kills Aide}

\lettrine{B}{ritain} ordered a predawn commando raid in northern Afghanistan on Wednesday to rescue
a British reporter for The New York Times and his Afghan interpreter after Afghan agents learned
that the Taliban was planning to move the hostages into Pakistan, a senior Afghan official said
Wednesday.

The raid by British Special Forces and Afghan soldiers freed the reporter, Stephen Farrell, but the
interpreter, Sultan Munadi, and a British paratrooper were killed in a fierce firefight, as were
least one Afghan civilian and dozens of Taliban fighters, officials said.

A senior Afghan official and Mr.~Farrell described a situation where after two days in captivity,
the hostages' situation turned more menacing\cite{menace}. They said it seemed likely that Taliban
leaders from outside the immediate district in Kunduz Province were planning to move the captives
across the border into neighboring Pakistan, largely outside the reach of NATO forces.

While Mr.~Farrell said he was treated well -- given food, water and blankets and never harmed -- the
militants increasingly taunted Mr.~Munadi. At one point one of the Taliban reminded Mr.~Munadi of a
case two years ago in which an Italian journalist taken hostage in Helmand Province was freed while
his Afghan translator was beheaded.

``I did not think they were going to kill me,'' Mr.~Farrell said Wednesday in a telephone interview
from the British Embassy in Kabul. ``I did think they were going to kill him.''

The Taliban captors talked freely on their telephones, increasing the chances that NATO
eavesdroppers also picked up on the change in mood and believed time was running short to act. ``My
sense is they probably just got them in the right place and needed to get on to it,'' said Peter
Gilchrist, a retired British major general and a former senior commander in Afghanistan.

A senior American military official in Washington said that the United States provided intelligence
assistance and had helicopters and attack aircraft at the ready, if needed, but that the operation
was planned and carried out by British commanders and civilian officials.

A senior NATO spokesman in Afghanistan, Rear Adm. Gregory J.~Smith, said that Gen. Stanley
A.~McChrystal, the top commander in the country, was informed of the raid after it took place.

The Times of London reported on its Web site on Wednesday that Prime Minister Gordon Brown of
Britain personally approved the raid. But other British officials said Mr.~Brown gave only general
authorization to the military to rescue Mr.~Farrell, who has dual British and Irish citizenship, and
Mr.~Munadi.

In a statement, Mr.~Brown said the raid was a British operation supported by the Afghan authorities
and NATO allies, including the United States. He praised the heroism of the British commandos and
confirmed ``with very deep sadness'' the death of one of them.

``This operation was carried out after extensive planning and consideration,'' Mr.~Brown said,
adding that whenever ``British nationals are kidnapped, we and our allies will do everything in our
power to free them.

``Sadly, we were unable to rescue Stephen's Afghan interpreter, Sultan Munadi, and we send his
family our condolences,'' he said.

President Hamid Karzai ``strongly condemned the killing of an experienced Afghan journalist,'' his
office said in a statement on Wednesday, according to Agence France-Presse. The statement said
Mr.~Munadi ``was killed mercilessly by the enemies of Afghanistan'' -- shorthand for Taliban
insurgents -- but did not give further details.

Neither The Times nor Mr.~Farrell's family knew that the military operation was taking place. But
they had discussed with the military under what conditions they might attempt a rescue. Basically
the answer was they might act if they had intelligence that they could act on quickly and with a
high probability of success.

Gunmen seized Mr.~Farrell and Mr.~Munadi on Saturday while they were working in a village near
Kunduz. They were reporting on the aftermath of NATO airstrikes on Friday that blew up two fuel
tankers hijacked by Taliban militants and killed scores of people, including an uncertain number of
civilians.

Mr.~Farrell, speaking to colleagues at The Times, said that he and Mr.~Munadi were moved several
times a day amid the cornfields, rice plantations and mud-brick villages in the Char Dara district
over their four days of captivity. In the first two days, he said, they felt optimistic that they
would be released.

``They drove us around Char Dara, almost always in the same old Toyota Corolla, sometimes with
masked and turbaned motorcycle outriders, rocket-propelled grenades sticking out of backpacks in
full daylight, just a few miles from the main Kabul to Kunduz road,'' Mr.~Farrell said.

Mr.~Farrell said his captors delighted in showing off, at one point driving to within about 1,500
feet of Afghan government and NATO watch towers, gleeful at their daring.

The area of Kunduz Province where they were held is ethnically mixed, with Tajiks and Pashtuns
living in ethnically distinct villages. But there was no doubt who was in charge. ``At no point did
we see a single NATO soldier, Afghan policeman, soldier or any check to the Taliban's ability to
move at will,'' Mr.~Farrell said.

Six to eight guards took turns watching the hostages, and they were hugely unpredictable. ``One
became enraged when I urinated standing up, deeming it an offense to local families,'' Mr.~Farrell
said. ``He then calmed down and asked me to teach him how to count to 10 in English.''

But on the third day, the situation became more threatening. It became harder for the captors to
find safe houses. They would get lost driving down ever-narrower and ever more obscure country
lanes.

``When they finally found a house with electricity two youngsters produced a tape recorder and began
blaring hours of religious sermons, praising Osama bin Laden, the mujahadeen of Chechnya, Somalia,
Helmand, Kandahar and anyone fighting the Americans,'' Mr.~Farrell said.

Some new Taliban figures, evidently more senior, arrived on the scene. There was much discussion of
moving the captives outside the Kunduz area, including to Baghlan, a medium-size city not far from
the sanctuaries of the Taliban and Al Qaeda in Waziristan in Pakistan.

Mr.~Munadi became worried about his own fate, and told Mr.~Farrell, ``I think you're going to be
O.K., but they've got it in for me.''

Late on the third night, loud explosions nearby and the drone of aircraft overhead sent the captors
and their hostages into the darkness, racing across open fields. It was a false alarm, but with a
waning moon Mr.~Farrell warned Mr.~Munadi that a rescue attempt could be coming soon. Indeed, it
would.

At about 2:30 a.m. local time Wednesday, the allied helicopters descended on the hideout and
disgorged the British commandos. British and NATO officials have refused to discuss details of the
operations, but most of the nations with sizable forces in Afghanistan keep a special military
hostage-rescue team at the ready for occasions just like this.

With surveillance drones and helicopters overhead, the captors scattered, Mr.~Farrell said. The two
men initially stayed put, fearing that they would be caught in any cross-fire. Then one of the
captors came back and tipped his gun toward them, he said, but left without firing.

``We absolutely expected them to cut us down as they ran,'' Mr.~Farrell said. ``We were crouching
targets in a long, narrow room devoid of anything but walls and matting, which felt like a death
trap.''

The two men waited a bit, then made their way out of the room into a courtyard. The lost each other
in the darkness for a moment, before linking back up. With Mr.~Munadi leading, they scuttled along a
narrow ledge along the outer wall of the compound. ``We could see nothing more than a few feet in
front of us,'' Mr.~Farrell said. ``We had no idea who was where, and there were bullets flying
through the air.''

After crouching and running for some 60 feet, the two men got to a corner. Mr.~Munadi was about two
feet ahead of Mr.~Farrell, and walked out into the clearing saying in an accent, ``Journaliste,
journaliste.'' It was not clear whether he was assuring commandos that he was not a Talib, or
assuring the Taliban that he was not with the commandos. There was a hail of bullets -- unclear
whether from friend or foe -- and Mr.~Munadi fell.

Mr.~Farrell said he reared back from the gunfire and dived into a ditch. He waited a couple of
minutes until he was clear which direction the British voices were coming from, then shouted,
``British hostage! British hostage!'' A few seconds later with hands raised high, he walked to the
British troops and safety.

On Wednesday, Mr.~Farrell blamed himself for Mr.~Munadi's death, though he said the two of them
discussed while in captivity the possibility that they might not survive.

Mr.~Farrell said of his colleague, ``He was trying to protect me up to the last minute.'' As they
left the room under commando siege, ``he moved out in front of me.''

``He was three seconds away from safety,'' Mr.~Farrell said. ``I thought we were safe. He just
walked into a hail of bullets.''

\section{Karzai Praises Election Panel}

\lettrine{P}{resident} Hamid Karzai stressed his support for Afghanistan's election commission on
Wednesday as international trust in the group appeared to be eroding after allegations of widespread
fraud in last month's presidential vote.

Mr.~Karzai's praise for the Independent Election Commission came a day after it announced
preliminary results giving him more than 50 percent of the vote, the threshold for avoiding a runoff
against his top challenger, Abdullah Abdullah, a former foreign minister.

The president said he ``applauds the elections commission for its efforts in pursuing the process in
an impartial and faithful national spirit.'' He also said he ``welcomed the partial election
results.''

The Aug.~20 election was marred by attacks by Taliban insurgents and has been tainted by allegations
of widespread fraud, including ballot-box stuffing and fake polling places.

A separate United Nations-backed body that is the final arbiter of the vote ordered a partial
recount on Tuesday because of what it called ``clear and convincing evidence of fraud.'' One Western
official said this week that had the Independent Election Commission not decided on Monday to undo a
decision it made the day before to enforce stricter safeguards, Mr.~Karzai's vote total would still
be under 50 percent, forcing him into a second election against Mr.~Abdullah.

The commission had moved on Sunday to carry out precautions intended to catch a number of voting
irregularities. But the official said that as it became clear that those safeguards would prevent
Mr.~Karzai from surpassing 50 percent, the decision was reversed Monday.


\section{Jobs in Fighting Form After Liver Transplant}

\lettrine{A}{fter} a prolonged absence from the public eye, a thin but energetic Steven P.~Jobs
reappeared Wednesday at an Apple news conference and addressed his illness in personal terms for the
first time.

``I feel great. I probably need to gain about 30 pounds, but I feel really good,'' said Mr.~Jobs,
Apple's chief executive, in an interview after the event. ``I'm eating like crazy. A lot of ice
cream.''

At the music-themed event, Apple introduced new iTunes software and a new line of iPod Nano music
players with video cameras. But the product announcements were eclipsed by the public return of
Mr.~Jobs, five months after he received a liver transplant at a hospital in Tennessee. Mr.~Jobs
needed the operation after suffering from complications from a rare form of pancreatic cancer.

``I'm vertical, I'm back at Apple and I'm loving every day of it,'' he said on stage, after a
standing ovation\cite{ovation} from the crowd, which included Apple employees and journalists
gathered for the news conference.

The applause ``was an endearing statement for someone who is clearly still recovering,'' said Gerry
Purdy, a wireless analyst at Frost \& Sullivan who has covered the mobile computing industry for 20
years. ``You just feel like this is an opportunity to share in the greatness of someone who is an
icon in the industry.''

Mr.~Jobs, who appeared as gaunt as he was when he went on a leave of absence in January, immediately
addressed the matter of his health.

``I'm very happy to be here, and thank you all,'' he said. ``As some of you know, five months ago I
had a liver transplant. I now have the liver of a mid-20s person who died in a car crash and was
generous enough to donate their organs. I wouldn't be here without such generosity, so I hope all of
us can be as generous and elect to be organ donors.''

But after that gentle opening, Apple's chief executive, who wore his characteristic outfit of jeans
and a black turtleneck, was as tough as ever. He used the news conference to take jabs at competing
products from Microsoft, Sony, Nintendo, Dell and Cisco Systems.

Mr.~Jobs's public return will reassure many Apple fans and shareholders, who view the charismatic
co-founder of the company as crucial to Apple's success.

``Apple users and investors think the world of Steve Jobs and the world of Apple products, and that
he was here today is a statement that he is back in charge,'' said Gene Munster, an analyst at Piper
Jaffray. ``The wizard of Oz is back in Oz.''

The actual products announced Wednesday were not nearly as magical as recent hits like the iPhone.

Apple announced a new version of its iTunes software that will allow users to better organize their
applications for the iPhone and the iPod Touch and to share music and videos among up to five
computers in the home.

The company also introduced a line of iPod Nanos that have a video camera, microphone, speaker, FM
radio tuner and pedometer and start at \$149.

Mr.~Jobs drew a sharp comparison between the Nanos, which are the size of a five-pack of chewing
gum, and Cisco's popular Flip video cameras, which are the size of a pack of playing cards. The new
Nano is an effort to turn around sales of the iPod, which have been sluggish for the first time and
which actually declined 7 percent in the last quarter from a year earlier.

``They need something to reinvigorate the line and give people another reason to have an iPod,''
said Charles S.~Golvin, an analyst focusing on mobile and consumer electronics at Forrester
Research.

Jonathan Kaplan, senior vice president for consumer products at Cisco, said his company was
flattered by the attention from Mr.~Jobs but said that the Flip camera records video in high
definition, while the Nano represents ``last generation's technology.''

``This will help consumers realize it's not all about photography, that video can make a difference
in the world,'' Mr.~Kaplan said. ``It's nice that Apple is finally realizing that video is a place
they need to be.''

Apple also pitched a cheaper, \$199 version of the iPod Touch with more memory but did not announce,
as some had expected, a new version with a video camera.

In the interview, Mr.~Jobs said that Apple had been originally unsure how to market the iPod Touch,
but settled on selling it as a video game device that can take advantage of the many free or cheap
games available in the company's App Store. ``We don't need to add new stuff,'' he said. ``We need
to get the price down where everyone can afford it.''

Despite much speculation by Apple fans, the company made no announcements about a tablet-size
computer or reading device.

But in the interview, Mr.~Jobs said general-purpose devices are more appealing than specialized
devices like Amazon.com's Kindle e-book reader.

``I think people just probably aren't willing to pay for a dedicated device,'' he said. ``You notice
Amazon never says how much they sell; usually if they sell a lot of something, you want to tell
everybody.''

Apple's shares, which rose in the days leading to the news conference, fell about 1 percent in
Wednesday's trading, closing at \$171.14.

\section{Obama Vows to 'Deliver on Health Care'}

\lettrine{P}{resident} Obama sought to reframe the contentious debate over health care on Wednesday,
asking a critical Congress and a skeptical nation to reach consensus on legislation to expand health
coverage to millions of Americans and lower medical costs through an ambitious overhaul that has
eluded lawmakers for generations.

``I am not the first president to take up this cause, but I am determined to be the last,''
Mr.~Obama said to a standing ovation from members of Congress. He added, ``Our collective failure to
meet this challenge – year after year, decade after decade – has led us to a breaking point.''

In a speech to a joint session of Congress, the president attempted to regain his political footing
on his signature priority of remaking the nation's health care system. He presented his most
detailed outline yet of a plan that he said would guarantee all Americans coverage, regardless of
pre-existing medical conditions, while letting people keep their own insurance if they wanted.

The president placed a price tag on the plan of about \$900 billion over 10 years, which he added
was ``less than we have spent on the Iraq and Afghanistan wars.'' But he devoted much of his address
to making the case for why such a plan is necessary.

``As soon as I sign this bill, it will be against the law for insurance companies to drop your
coverage when you get sick or water it down when you need it most,'' Mr.~Obama said. ``They will no
longer be able to place some arbitrary cap on the amount of coverage you can receive in a given year
or a lifetime.''

The president called for ``making a not-for-profit public option available'' to consumers. But he
suggested he was amenable to refining that idea, which has become a lightning rod for criticism.

``We should remain open to other ideas that accomplish our ultimate goal'' of insuring all
Americans, he said.

The speech was the president's second address before a joint session of Congress. But the political
backdrop on Wednesday was far different than Mr.~Obama's appearance in the House chamber on the 36th
day of his term, when an optimistic wave of momentum was at his back and his Republican rivals were
dispirited and in disarray.

``What we have also seen in these last months is the same partisan spectacle that only hardens the
disdain many Americans have toward their own government. Instead of honest debate, we have seen
scare tactics,'' Mr.~Obama said. ``Some have dug into unyielding ideological camps that offer no
hope of compromise.

``Too many have used this as an opportunity to score short-term political points, even if it robs
the country of our opportunity to solve a long-term challenge,'' he said.

He added, ``And out of this blizzard of charges and counter-charges, confusion has reigned.''

The president paused to remember Senator Edward M.~Kennedy, Democrat of Massachusetts, who died last
month of cancer before his goal of a thoroughhealth care overhaul could be realized. The senator's
widow, Victoria Kennedy, was sitting in the first lady's box in the House chamber alongside several
ordinary Americans selected by the White House to help bolster the president's case that the health
care system was in crisis.

``That is why we cannot fail,'' Mr.~Obama said. ``Because there are too many Americans counting on
us to succeed – the ones who suffer silently, and the ones who shared their stories with us at town
hall meetings, in emails, and in letters.''

The president said that he had not closed the door on reaching a bipartisan compromise on the health
care legislation. He gave a nod to Senator John McCain, Republican of Arizona, and embraced his
proposal to create a high risk pool to help cover people with pre-existing conditions against
catastrophic expenses.

But Mr.~Obama also signaled a more defiant tone, saying that he would not allow a misinformation
campaign to flourish as the White House believes that it did during the August Congressional recess
when angry voters flooded town meetings across the nation.

``I will not waste time with those who have made the calculation that it's better politics to kill
this plan than improve it,'' Mr.~Obama said. ``I will not stand by while the special interests use
the same old tactics to keep things exactly the way they are. If you misrepresent what's in the
plan, we will call you out.''

The president said ``there are arguments to be made'' from both the left and right side of the
ideological spectrum on health care. But throughout his speech, he sought to reassure Americans that
he did not support ``a radical shift that would disrupt the health care most people currently
have.''

``Since health care represents one-sixth of our economy,'' he said, ``I believe it makes more sense
to build on what works and fix what doesn't, rather than try to build an entirely new system from
scratch.''

Mr.~Obama took a sharp tone against the insurance industry.

``Now, I have no interest in putting insurance companies out of business,'' Mr.~Obama said. ``They
provide a legitimate service, and employ a lot of our friends and neighbors. I just want to hold
them accountable.''

The speech, which Mr.~Obama and his aides were putting the finishing touches on until shortly before
he arrived at the Capitol, was intended to restart the debate on Capitol Hill on the health care
legislation. The president said he remained firm in his deadline of signing some type of measure
before year's end.

Senator Mitch McConnell of Kentucky, the Republican leader, signaled before the address that
Republicans still would not be supportive of the president's retooled health care approach.

``Americans don't understand how they'll be able to keep the health plans they have if government is
allowed to undermine the private market,'' Mr.~McConnell said on the Senate floor earlier Wednesday.
``And they don't understand why the administration doesn't seem to be listening to these and many
other concerns.''

A few hours before the speech, Senator Max Baucus, chairman of the Finance Committee, said that his
panel would take up sweeping legislation and start voting on it in two weeks, with or without the
support of Republicans.

``The time has come for action, and we will act,'' Mr.~Baucus said. He added, ``Irrespective of
whether there are any Republicans, I will move forward.''

\section{Study Says Ghostwriting Rife in Medical Journals}

\lettrine{S}{ix} of the top medical journals published a significant number of articles in 2008 that
were written by ghostwriters financed by drug companies, according to a study released Thursday by
editors of The Journal of the American Medical Association.

Among authors of 630 articles who responded anonymously to an online questionnaire created for the
study, 7.8 percent acknowledged contributions to their articles by people whose work should have
qualified them to be named as authors on the papers but who were not listed.

In the scientific literature, ghostwriting usually refers to medical writers, often sponsored by a
drug or medical device company, who make major research or writing contributions to articles
published under the names of academic authors.

The concern, the researchers said, is that the work of industry-sponsored writers has the potential
to introduce bias, affecting treatment decisions by doctors and, ultimately, patient care.

According to the study, responding authors reported a 10.9 percent rate of ghostwriting in The New
England Journal of Medicine, the highest rate among the journals.

Editors of the Boston-based journal said Thursday that they were ``puzzled'' and ``skeptical'' of
the findings.

The study also reported a ghostwriting rate of 7.9 percent in JAMA, 7.6 percent in The Lancet, 7.6
percent in PLoS Medicine, 4.9 percent in The Annals of Internal Medicine, and 2 percent in Nature
Medicine.

``These journals are the top of the medical field,'' Joseph S.~Wislar, a survey research specialist
and lead author of the study, said in a phone interview. He recommended that they take more action
to require that all contributors be listed in acknowledgments if they are not named as authors.

Three JAMA editors, Annette Flanagin, Phil B.~Fontanarosa and Catherine D.~DeAngelis, joined
Mr.~Wislar in the study.

The new study, which has not yet been peer-reviewed or published in a medical journal, was made
public Thursday morning at an international meeting of journal editors in Vancouver.

``It was very compelling, and I find it quite shocking, to be honest,'' Ginny Barbour, chief editor
of PLoS Medicine, the journal of the Public Library of Science, said after the meeting. ``We are a
journal that has very tough policies, very explicit policies on ghostwriting and contributorship,
and I feel that we've basically been lied to by authors.''

Some of the same researchers (though not Mr.~Wislar) also sent out a questionnaire to authors of
articles published in 1996 in three of the same publications. That study reported ghost authorship
rates of 16.2 percent in The New England Journal of Medicine, 15.3 percent in The Annals of Internal
Medicine, and 7.1 percent in JAMA.

Comparisons between the studies may not be valid because they relied on different methodologies and
covered different authors. The older study involved a mail-in questionnaire sent to authors based in
the United States while the new study, involving an online questionnaire, solicited responses from
authors based both inside and outside the United States. In both cases, the studies have the
potential for reporting bias because they did not choose respondents randomly but relied on authors
to elect to answer the questions; moreover, authors were asked to disclose their own behavior, with
the potential for them to underreport the use of a ghostwriter, which is considered an academic
crime akin to plagiarism.

Finally, the response rates from authors of articles varied widely, ranging from 58.3 percent for
one journal to 85.9 percent for another journal, the researchers said.

Karen P.~Buckley, spokeswoman for The New England Journal of Medicine, said she was ``completely
shocked'' at the high rate of ghostwriting reported by its authors. She said the journal was
continually strengthening its safeguards.

Editors of the journal released a statement through Ms.~Buckley saying the JAMA study used an
improperly broad definition of ghostwriting. But Annette Flanagin, a JAMA editor and co-author of
the new report, responded that it was the standard definition of the term.

\section{Remembering a Future That Many Feared}

\lettrine{T}{he} day dawned different and stayed that way. Traffic was thin and sidewalks quiet. The
stock exchange didn't open, nor the airports, the schools, Broadway. People loaded up on bottled
water, batteries, canoes. The law enforcement presence was intense: men with machine guns, gunboats
circling the harbor.

Downtown, fires burned, smoke plumed. The odor stood.

It was a city humbled and scared, where the possibilities of destruction had been recalibrated. It
was Sept.~12, 2001. The day after.

So much has been said and written about what happened on 9/11. The following day is forgotten, just
another dulled interlude in the aftermath of an incoherent morning.

But New Yorkers were introduced that day to irreducible presumptions about their wounded city that
many believed would harden and become chiseled into the event's enduring legacy.

New York would become a fortress city, choked by apprehension and resignation, forever patrolled by
soldiers and submarines. Another attack was coming. And soon.

Tourists? Well, who would ever come again? Work in one of the city's skyscrapers? Not likely. The
Fire Department, gutted by 343 deaths, could never recuperate.

If a crippled downtown Manhattan were to have any chance of regeneration, ground zero had to be
rebuilt quickly, a bricks and mortar nose-thumbing to terror.

Eight years later, those presumptions are cobwebbed memories that never came to pass. Indeed,
glimpses into a few aspects of the city help measure the gap between what was predicted and what
actually came to be.

You could start at one downtown street corner. The wisdom of the day after was that New York would
never again bunch together important institutional nerve centers, binding them together in
vulnerability.

On Sept.~11, American Express had its headquarters at the southwest corner of West and Vesey
Streets. It is still there. Since then, Verizon has settled its headquarters into the northeast
corner. Goldman Sachs has assumed the northwest. All that's missing is the southeast corner. That
will be filled by the tallest building in America.

The Souvenir Man

David Cohen pointed out what the tourists like: replica taxicabs, ``I Love New York'' T-shirts and
thimbles -- any gewgaw inscribed with New York. ``See this digital picture postcard?'' he said.
``Nice little item.''

Mr.~Cohen, 83, is the patriarch of Grand Slam, a family-run novelty and baseball clothing store on
Broadway between 46th and 47th Streets, in the heart of Times Square. Eight years ago, he could not
have imagined the heaving commerce, the new big buildings, and especially not the complacent scene
outside his doors. People basked in the balmy weather at tables and chairs, under sheltering patio
umbrellas, spread across Broadway. If they worried about anything, it was sunburn.

How about that? People, at the behest of the mayor himself, flocking to Times Square to relax!

When fear engulfed the city on Sept.~12, many wrote off Times Square. Chemical bombs were sure to
explode there. A suicide bomber strapped with explosives was destined to blow himself up at lunch
hour.

``It was creepy,'' Mr.~Cohen said. ``It was, 'Oh my God, what's next?' I thought this would be the
next hit.''

Business was slow for months. Souvenirs didn't seem to mean the same anymore. ``Yeah, it took a
dive,'' Mr.~Cohen said. He shortened the store's hours.

But he did not leave. ``You can't live in fear,'' he said. ``Things happen and then they don't
happen.''

Now the weak economy squeezes sales, but pedestrian traffic in Times Square is far higher than it
was before Sept.~11. Vastly enhanced security has been put in place, and even when incidents defy
it, like the small bomb that exploded at the military recruiting station in March 2008, people shrug
it off, keep coming.

``This is the best spot in New York,'' Mr.~Cohen said. ``Listen, the Square is the place.''

The Garage Manager

The fires wouldn't go out. The smell persisted. What company would ever open its doors in Lower
Manhattan? Who would live there? Who could feel secure?

The police stopped and searched trucks. Only a few cars were allowed below 14th Street.

Still, Wilson Ortega, 34, came to work. He managed the parking garage at 56 North Moore Street in
TriBeCa.

On Sept.~12, business was, as he put it, ``off 100 percent.'' But cars were still in there, and
maybe people wanted them.

As streets reopened, car pooling into Manhattan was mandated during rush hours. Bombs preyed on
peoples' minds. Many garages throughout the city began checking trunks and jabbing mirrors on the
ends of poles beneath cars. Some still do, but in large part the practices are additional relics of
the times.

``Yeah, I checked,'' Mr.~Ortega said.

Every trunk was searched. He acknowledged that he had no training in explosives, didn't know exactly
what he was looking for, but he did every car for several months, then those he didn't recognize,
the nonregulars, for nearly a year. Some people were insulted, wouldn't pop the trunk, and he turned
them away. He never found a thing.

The trade center site remains a conflicted construction project. But on North Moore it is cars in,
cars out, just as before.

``I never thought things would be the same again,'' he said. ``But, man, I was wrong. We came back
strong.''

The Firefighter

The number was 343. Back in those awful days, Chief Charlie Williams, 9th Battalion, Manhattan,
thumbed down the death list looking for the firefighters he could have said hello to by name: ``Hi
Tom, hi Joe, hi Ray.'' After about 40, he stopped. It was enough.

The loss of life to the Fire Department was staggering. Many asked, who would put out the fires of
tomorrow?

In addition to the deaths, there was a stampede of retirements. The wives didn't want to join the
widows. And the expansive opportunity for overtime pay afforded a tantalizing opportunity for
firefighters to retire at bulgier pensions.

There were 11,339 uniformed members of the Fire Department on Sept.~10, 2001. By Jan.~28, 2003, the
ranks had declined to 10,630.

Chief Williams asked himself: ``Do I want to go back and do this job?'' His wife would have liked
him to walk away. But he wasn't done.

Fresh recruits were rushed in. There was a long, difficult period. Even now, the experience level is
not the same. But there are 11,415 uniformed personnel, more than before.

``The bell rings and the men put out the fires,'' Chief Williams said. ``The city is well served.''

In the aftermath of Sept.~11, the firefighters were elevated to superhuman status. People flocked to
the firehouses, wanting to shake hands with firefighters, snap their pictures, just say thanks.
Chief Williams obliged, though he allowed how it got overbearing at times; he had to shut himself in
his office to do his work.

The bravery was always real. But the mythology -- well, that, too, wasn't going to last. In the
ensuing years, there were embarrassing incidents: the firefighters who had sex with a woman at one
Bronx firehouse, a drunken brawl at another in Staten Island, on-duty drinking and drug use.

``The worship was definitely an inflated thing,'' Chief Williams said. ``You couldn't sustain
that.''

His own lungs went bad on him, traced back to the trade center, and he retired last year. He chose
the date: Sept.~11.

The Flag Printer

People bought them from hardware stores and Wal-Mart and street vendors and unfurled them outside
their homes and on the antennas of their cars. They billowed down the Henry Hudson and the F.D.R.

Flags.

People wore their patriotism and defiance openly. A new cohesiveness, a oneness, was going to remold
the character of American citizenry.

Christopher Gravagna didn't feel right that people had to buy their patriotism. ``That was
ridiculous,'' he said. ``Why should people capitalize on flags at that time?''

He had a printing business in Long Island City, Queens, doing work for clubs and concerts. On
Sept.~12, demand for his services essentially stopped and didn't resume for weeks. So he decided to
print paper American flags with the motto ``United We Stand'' and give them away. He and his
employees handed out more than 100,000.

He saw them everywhere.

``It helped feed this feeling that we have to be one, we have to be together on this,'' Mr.~Gravagna
said. ``We're a strong country. We're strong New Yorkers.''

The flags -- cloth and paper -- are mostly gone. Some come out, as they always did, on Memorial Day,
on the Fourth of July, and on Sept.~11, but that is it.

That special mood? ``It's definitely diminished a lot,'' Mr.~Gravagna said. ``Did I expect it?
No.~But as a New Yorker, I understand it. I guess part of it has to do with capitalism. In America,
we have issues. And time passes. It just passes.''

No one, perhaps, displayed as many flags as Mr.~Gravagna himself. He taped them to the windows of
his Queens apartment and in his Nissan Sentra. They festooned his offices.

After a while, they came down. The last one he possessed he had framed. He hung it on his office
wall. Four years ago, someone stole it.

The Skyscraper Dentist

``The windows here open,'' Dr.~Charles Weiss said.

He unlatched one. The view south was dazzling, as only a 1,000-foot-high view can be. There was the
Empire State Building and, way off, the Statue of Liberty and Ellis Island, as well as a spot where
two duplicate towers once stood.

On Sept.~12, it seemed no one would choose to work in a skyscraper again. Especially those with the
emblematic names, the ones everyone knew about, high-rise terrorist bounty.

Workers stuffed parachutes under their desks, were given particle masks, acquainted themselves with
Geiger counters.

On Sept.~11, Dr.~Weiss, a dentist, repaired teeth on the 69th floor of the Chrysler Building, at
42nd Street and Lexington Avenue. He still does.

Capitulation was not his style. He recalled a book, ``The Last Angry Man,'' in which a pugnacious
Brooklyn doctor refuses to yield to the bums he calls ``galoots.'' Dr.~Weiss thought, as an
assertion of faith, ``I'm not going to let the galoots get me.''

On Sept.~12, the Chrysler Building was essentially closed, but he got in. He called patients to
reschedule them. Some wanted some time before readdressing their cavities. He didn't see anyone
until the following Monday.

As far as he knows, they all came back. The patients. The people who worked for him. His colleagues
who minded the other dental chairs on the floor.

There are always some squeamish patients who fear heights. Dr.~Weiss, now 82, dispatches a nurse
down to the lobby to ride the elevator up with them. That happened before Sept.~11, too.

Waiting patients now flipped through magazines as the drills sang. ``There's a tremendous drive of
human beings to make the most of life,'' Dr.~Weiss said. ``We're not hermits. We rise up and move
on.''

Dr.~Weiss drank in the view some more, watched the ant cars crawling across the ever-clogged city.
``I never get tired of that view,'' he said. ``Never.''

\section{Immigrant Finds Path Out of Maze of Detention}

\lettrine{H}{olding} tight to her sister's hand in the bustling streets of New York's Chinatown last
week, Xiu Ping Jiang looked a little dazed, like someone who has stepped from a dark, windowless
place into a sunny afternoon.

In a sense, she has. For a year and a half Ms.~Jiang, a waitress with no criminal record and a
history of attempted suicide, was locked away in an immigration jail in Florida. Often in
solitary\cite{solitary} confinement\cite{confinement}, she sank ever deeper into mental illness,
relatives say, not eating for days, or vomiting after meals for fear of being poisoned.

With no lawyer to plead for asylum\cite{asylum} on her behalf, she had been ordered to be deported
to her native China, from which her family says she fled in 1995 after being forcibly sterilized at
age 20. Too ill to obtain the travel documents needed for the deportation to take place, she was
trapped in an immigration limbo: a fate that detainee advocates say is common in a system that has
no rules for determining mental competency and no obligation to provide anyone with legal
representation.

Then, through a fluke, her case came to the attention of The New York Times, which published an
article on May 4 about her ordeal\cite{ordeal} and the efforts of her sisters in New York to help
her. An immigration judge in Florida reopened the case.

Now Ms.~Jiang, 36, is free on bail\cite{bail}, living in Brooklyn with her older sister, Yun, a
United States citizen, and receiving the medical and psychiatric help she needs while awaiting a
fresh immigration hearing close to home -- this time with a lawyer. And her case is being held up as
an example of the system's worst and best approaches toward the mentally ill, as advocates press the
Obama administration for change.

If immigration courts were required to offer the same basic protections for the mentally disabled as
any other court, advocates say, Ms.~Jiang's prolonged detention -- which cost hundreds of thousands
of dollars and put her life at risk -- could have been avoided.

Hers is one of several cases cited in a 15-page letter to Attorney General Eric H.~Holder Jr.~that
asks for such protections, including the appointment of counsel to anyone with a mental disability
in deportation proceedings, and the appointment of guardians ad litem to speak on behalf of those
found mentally incompetent. The letter, signed by 77 mental health experts, civil rights lawyers and
immigration advocates around the country, was sent July 24, but not made public until Thursday.

The Justice Department has made no formal response, but Matthew A.~Miller, a spokesman, said the
department provided specialized training on handling the mentally ill at the annual conference of
immigration judges last month.

``Persons with mental disabilities face significant challenges in removal proceedings, and these
cases also present substantial challenges for immigration courts,'' he said in an e-mail message.
``We recognize our obligation\cite{obligation} in this area, and we will continue to review how we
can improve our efforts to provide due process and reasonable accommodation\cite{accommodation}.''

Sunita Patel, a lawyer at the Center for Constitutional Rights who signed the letter, said such
training was a positive step, but insufficient.

``As Ms.~Jiang's case demonstrates,'' she said, ``these difficult cases require all of the proposed
protections. And change should happen fast, before more people with mental disabilities disappear
into the detention maze.''

Ms.~Jiang was held at the Glades County Detention Center in December 2007, after immigration agents
stopped her at a Greyhound bus station in West Palm Beach, Fla., where she was traveling to a
restaurant job. She had no immigration papers.

At her first brief hearing, the immigration judge, Rex. J.~Ford, issued an order of deportation
saying that she had failed to show up because she kept answering his question, ``What is your
name?'' without waiting for the interpreter's translation into Mandarin.

He sent her back to the jail, where she was deprived of proper treatment and her mental state
rapidly deteriorated\cite{deteriorate}, according to court petitions on her behalf. She was
periodically confined to an isolation cell on suicide watch, records show.

``Immigration detainees with mental-health issues often receive little, if any, treatment,'' the
Florida Immigrant Advocacy Center wrote this year in a report on detention. ``In many cases, their
conditions worsen or they destabilize while in detention. They are misdiagnosed, improperly
medicated, cruelly treated or denied psychiatric care altogether. Worse, many of those immigrants
should not have been detained in the first place.''

Ms.~Jiang did not speak at all at a second hearing, before another immigration judge, Scott
G.~Alexander. He, too, ordered her deported, but reopened the case after The Times article was
published, and in June 2009, transferred her from the jail in Glades County to a hospital detention
center in Columbia, S.C., for a new assessment of her mental state. There, Ms.~Jiang jumped out of a
second-floor window, her lawyer says, only to enter an immigration agent's vehicle and wait inside,
apparently thinking it was a taxi.

Her physical condition had deteriorated, too. An obvious lump in her neck, ignored at the jail, was
biopsied at the South Carolina center and found to be thyroid\cite{thyroid} cancer; she underwent
surgery, said Theodore N.~Cox, an immigration lawyer in New York who took the case without fee.

In July, Mr.~Cox persuaded the immigration judge in Charlotte, N.C., Theresa Holmes-Simmons, to free
Ms.~Jiang on \$3,500 bail, despite opposition from Immigration and Customs Enforcement lawyers who
argued that her leap from a window showed she was a flight risk.

On July 17, she went home to Brooklyn with her sister Yun, 38. ``But she was too scared,'' the
sister recalled. ``I had to bring her to the hospital'' -- Bellevue Hospital Center, which has
treated immigrants regardless of their ability to pay since 1736.

There, after 17 days as a voluntary psychiatric inpatient, she was visibly on the mend when a
reporter visited on Aug.~3. Her cheeks were round.

Only months before, when her older sister had visited her in the Florida jail, Ms.~Jiang,
emaciated\cite{emaciate} and incoherent, had scratched her face and refused to speak to her, the
sister said. Now, she added, laughing, ``She calls me on the phone, and says, 'Visit and bring me
some food!' ''

That afternoon, the older sister had brought a feast culled from fruit stands and Chinese
restaurants: melon, bags of peaches, sticky rice wrapped in banana leaves, noodles with duck, even
the sweet black-bean-paste treats of her childhood in Fujian Province. Ms.~Jiang, in blue hospital
pajamas, dug happily into the yellow flesh of a melon with a plastic spoon, then paused to answer a
question about her detention.

``When they locked me up, I was very cold,'' she said softly through an interpreter. ``I asked for a
blanket, but they didn't give me one. It was dark. I got scared, so I think I went crazy.''

Her reason for jumping from a detention center window? ``I wanted to stay in the United States,''
she said. ``I like the United States. China had the sterilization -- nothing is good.''

Her sister eyed her anxiously. ``Whenever she talks about the past, her hands start to tremble and
she gets this faraway look in her eyes,'' she said.

Since her discharge from Bellevue in August, Ms.~Jiang has been receiving weekly treatment. Doctors
recommended delaying a second operation on her neck until her emotional state was more stable, her
sister said.

Meanwhile, she does simple chores in the restaurant where her sister is a cashier: cutting broccoli,
watching the employees' children.

``I'm feeling O.K.,'' she said last week with a timid smile, fingering the scar at her throat. Her
sister smoothed her hair, and guided her onto the subway home.

\section{G.M. Agrees to Sell Opel to Magna, With Strings}

\lettrine{G}{eneral} Motors told the German government Thursday that it had decided to sell a
majority stake in its unprofitable European operations as agreed last spring -- on condition that
the government and labor unions pledged to help with the ``necessary'' restructuring.

Chancellor Angela Merkel of Germany, who is up for re-election in less than three weeks, was quick
to claim credit for clinching the potentially job-saving deal even though she had, at times, been
skeptical of government rescues during the past year's financial and economic crisis.

``This has shown that the patience and single-mindedness of the federal government has paid off,''
Mrs.~Merkel said at a news conference.

But G.M. stressed that it still needs union agreement -- in writing -- for the restructuring of the
four Opel plants in Germany, which employ almost half of G.M. Europe's 55,000-strong European work
force. The powerful IG Metall union said that it expected the negotiations to be ``tough.''

And the German government has not yet said how much public money it will commit. It had originally
pledged over \texteuro4.5 billion, or \$6.5 billion, in loan guarantees. Members of the trust
overseeing the sale of Opel were sharply divided over the accord, with some criticizing the amount
that the German taxpayer would have to provide in order to save Opel.

Under the deal, Magna International, a Canadian-Austrian auto parts and engineering company, and
Sberbank, its Russian partner, would purchase a 55 percent stake in Opel, which is based in Germany
and is the main G.M. operation in Europe, as well as the smaller Vauxhall, which is based in
Britain.

G.M., which acquired Opel in 1929, would keep a 35 percent stake, while Opel's employees would be
provided 10 percent.

The accord was initially reached in late May, but it had been thrown into doubt in recent weeks
after a revitalized G.M. emerged from bankruptcy protection in the United States.

The new G.M. board entertained a rival bid from a private equity group, and was considering holding
on to the unit itself. But Mrs.~Merkel, in an uncharacteristically risky move, threw all her weight
behind Magna, which had promised to preserve the most jobs in Germany.

Her government lobbied hard in Washington this summer as well, although U.S.~officials said that
President Barack Obama did not want to be dragged into the negotiations, despite the
U.S.~government's 60 percent stake in the bailed-out automaker.

Prime Minister Vladimir Putin of Russia, who has long regarded the development of Russia's lagging
auto industry as a strategic goal, also backed Magna and Sberbank from the start. The Opel
production in Russia would be taken over by GAZ, a car manufacturer owned by the Russian oligarch
Oleg Deripaska, who is close to the Kremlin.

G.M. had reservations about the Russian connection, fearing that intellectual property rights and
other licenses would not be respected. But John Smith, G.M. chief negotiator, said Thursday at a
news conference in Berlin that those issues had been resolved.

Mr.~Smith said that the deal could be closed ``in November or December.'' That would be well after
the German elections, in which rising unemployment is a major issue.

``At the moment, I think it is fair to say that in some shape or form all four German facilities
will carry on,'' he said, referring to the production plants in Rüsselsheim, Eisenach, Bochum and
Kaiserslautern.

But he added that the entire European automobile industry was struggling with overcapacity, and
suggested that an Opel factory in Antwerp, Belgium, which employs 2,000 workers, would be idled.

In its statement, G.M. said that ``several key issues'' remained to be settled over the next few
weeks, ``including the written support of the labor unions to support the deal with the necessary
cost restructuring for viability, and the completion of a definitive financing package from the
German government.''

Under the original terms reached in May, the German government was to provide \texteuro4.5 billion
in loan guarantees. In return, Magna has said it would put up \texteuro500 million for restructuring
costs and \texteuro350 million in private equity.

Fred Irwin, chairman of the trust established by the German government last May to oversee Opel's
assets during the negotiations, would not say how much Magna would be paying in the end. But he said
the deal was ``in the best interests of Opel and its employees.''

In addition to saving jobs in Germany, government officials saw the Magna bid as giving Opel greater
access to Russia, an expanding market.

Still, the twist with the Russian deal, analysts said, is that G.M. must now find the right way of
cooperating with the new Opel business in Russia, or risk the wrath of Mr.~Putin, whose prot\'eg\'e,
Dmitri A.~Medvedev, also backed the deal.

G.M. already has stakes in three assembly plants in Russia, and G.M.'s Chevrolet is Russia's No.~2
selling brand, after Lada. G.M. sold 98,800 Opels and Vauxhalls in Russia last year, compared to
235,466 Chevrolets, for a total market share of 11.2 percent.

``Russia is still a market where protectionism and the state lurks,'' said Ferdinand Dudenhöffer,
director of the Center for Automotive Studies at the University of Duisburg-Essen. ``It's always a
good thing not to surprise the Russians.''

Mr.~Dudenhöffer estimated Opel could grab 15 percent of the Russian market by 2015, equivalent to
750,000 vehicles.

In a conference call with reporters, Mr.~Smith emphasized that Magna and Sberbank would still need
G.M., despite the American giant's searing experiences of the last year.

``Magna has never designed a car from scratch, nor have they sold one,'' Mr.~Smith said.

He said that Opel and Chevrolet would likely cooperate further on distribution, a good way to ``save
some money in a way that the customer does not care about.''

Mr.~Smith predicted that ``New Opel,'' as he called it, would be profitable by 2011 and generate
enough cash to pay back the German government-backed loans by 2014.

For the first three years, the new management would have eight top executives, four appointed by
G.M., with Magna and Sberbank appointing the other four, including the chief executive. After that,
only Magna can choose the chief executive, Mr.~Smith said.

The agreement places a series of restrictions on Opel in four major car markets: the United States,
Canada, China and South Korea. The restrictions aim to head off open competition between the Opel
brand and brands fully under G.M. control.

\section{Motorola Phone Focuses on Social Networks}

\lettrine{M}{otorola} introduced the first of a new generation of smartphones Thursday that it hopes
will reverse its plummeting cellphone sales.

The phone, called the Cliq, is meant for young people obsessed with social networks. Instead of the
traditional menu of features, the Cliq's home screen is an ever-changing mosaic of e-mail, Twitter
tweets and status updates, superimposed over photos of the people sending those messages.

``It's alive,'' said Sanjay Jha, the co-chief executive of Motorola, who was hired a year ago from
Qualcomm to revive its cellphone business. ``Think of it like the text bubbles in cartoons, with new
information pushed to you all the time.''

The Cliq will be available through T-Mobile in the United States in the fall. The price has not been
announced, but analysts expect a \$100 price tag. The phone has a 3.1-inch touch screen and a
slide-out keyboard. Motorola is expected to introduce a second, more expensive smartphone in a few
weeks that will work on the Verizon network.

Both phones are based on Google's Android operating system, but Motorola has substantially modified
the Google software to make its phone stand out from the dozen or so other Android phones that are
expected to be introduced before the holiday sales season. While phone users will be able to
download and use any of the growing list of applications for the Android system, the look of
Motorola's phones and the way they operate will be different from other handsets using the operating
system.

``I get the question all the time, 'How are you going to differentiate on Android?' '' Mr.~Jha said.
``Android is open and flexible, so we can build on top of it.''

Indeed, much of what Mr.~Jha said makes the Cliq special is not in the phone at all, but in an
Internet-based service called Motoblur that integrates all of a user's e-mail and social networking
accounts.

Mr.~Jha said that the Cliq represented the turning point for Motorola. It is not meant to be like
the iPhone, one handset with a few variations that will lure tens of millions of buyers. Rather it
is the first of dozens of handsets, built from the same underlying Android software and Motoblur
service, that will have different features to appeal to all sorts of customer groups around the
world.

``This is one product, but we have changed the way we do business,'' Mr.~Jha said.

Users can simply enter their account information for these services, and the phone numbers, photos
and other information from their friends will be automatically downloaded into the phone's address
book. Motoblur also keeps copies of any other contact information that the user enters, so that it
can be moved automatically onto a new phone, provided the phone is made by Motorola.

Some of the early Android handsets were tied instead to Google's e-mail and other online services.

Tero Kuittinen, an analyst with MKM Partners, said that most other smartphone makers were also
building services to integrate messages from different services. The first was the Palm Pre.

``Message integration hasn't been enough to turn the Pre into a big hit when it was the only phone
that had it,'' he said.

Mr.~Kuittinen said that Motorola faced very stiff competition in the smartphone market, with dozens
of models available from most every manufacturer.

While the top end of the market is dominated by Research In Motion's BlackBerry and Apple's iPhone,
Motorola is competing in the very crowded lower-price tier.

``The challenge for Motorola is to win back customers they have lost in the last two years,'' he
said.

\section{Ford Backs Ban on Text Messaging by Drivers}

\lettrine{T}{he} Ford Motor Company on Thursday became the first automaker to endorse a federal ban
on sending text messages while driving.

Ford issued a statement in support of legislation proposed by Senator Charles E.~Schumer, Democrat
of New York, that would cut by 25 percent the federal highway financing given to states that did not
comply with a text-messaging ban. Ford also said it supported a similar proposal in the House by
Carolyn McCarthy, Democrat of New York.

``The most complete and most recent research shows that activity that draws drivers' eyes away from
the road for an extended period while driving, such as text messaging, substantially increases the
risk of accidents,'' Susan M.~Cischke, Ford's group vice president for sustainability, environment
and safety engineering, said in the statement.

``Ford believes hands-free, voice-activated technology substantially reduces that risk by allowing
drivers to keep their hands on the wheel and eyes on the road. Ford supports a ban on hand-held text
messaging while driving.''

The ban would not affect use of Ford's in-car communications and entertainment system, called Ford
Sync, which allows most mobile phones to be used hands-free. The system can also read text messages
aloud to the driver.

Ford Sync is standard equipment on many models and is available on other vehicles for about \$400.

Mr.~Schumer praised Ford for its support of a ban.

``Ford deserves credit for stepping up as the first car company to endorse a ban on this dangerous
habit,'' he said in a statement. ``We are gathering a critical mass of support for this bill, which
will give us the momentum we need to get it passed.''

Other major automakers have not yet taken a position on the bill.

Text-messaging bans have already been enacted in 14 states and the District of Columbia.

The Governors Highway Safety Association, which represents state highway safety officials, said last
month that it favored a nationwide text-messaging ban. Earlier, it had been opposed to such a law,
saying that enforcement would be too difficult.

The association said it based its new position on a study by the Virginia Tech Transportation
Institute concluding that text messaging by drivers makes them 23 times more likely to crash or
narrowly avoid a crash.

\section{Obama Factor Plays to Senator's Advantage}

\lettrine{M}{ere} survival was never the issue. Those familiar with Louisiana and its Republican
senator, David Vitter, knew he would survive the summer of 2007, after he showed up on the client
list of a Washington prostitution ring and then refused to address the matter outside of admitting
to a ``very serious sin'' at a brief news conference.

What came as a surprise to many here is how he became a strong early favorite going into his 2010
re-election race. That turnabout is largely due to one person. ``Along comes Obama,'' said Elliott
Stonecipher, a political analyst and demographer based in Shreveport, ``and it changed everything.''

It is difficult to overstate President Obama's unpopularity in most of Louisiana. He lost handily to
Senator John McCain in November here, picking up only 14 percent of the white vote (the state is
roughly two-thirds white). His health care plan is unpopular. His cap-and-trade plan to reduce
greenhouse gases, in a state so dependent on oil and gas, is anathema.

In fact, throughout the South -- which largely voted against Mr.~Obama last year -- the anger at the
president's policies has been palpable, as demonstrated by the outburst of a South Carolina
Republican congressman, Joe Wilson, during the president's health care address to Congress on
Wednesday evening.

Mr.~Vitter, an intense and competitive politician, has taken this displeasure and run with it. Well
before his Democratic opponent, Representative Charlie Melancon, announced his Senate candidacy on
Aug.~27, Mr.~Vitter's campaign had been attacking him for his support of Mr.~Obama.

In Mr.~Vitter's advertisements and on his Web sites -- including charliemelanconforsenate.com, which
was secured by the Vitter campaign -- a newspaper headline is featured: ``Melancon backs Obama
presidency.'' That apparently says enough.

Though nearly 22 percent of the state's adult residents have no health insurance -- one of the
highest rates in the nation -- pollsters and political experts say voters in the state are
overwhelmingly against Mr.~Obama's health care proposals.

The strategy of running against the White House got a trial run recently in a special State Senate
election in Mr.~Melancon's own district in southeastern Louisiana. It began as a three-way race
among Brent Callais, a Republican former parish councilman; Damon J.~Baldone, a Democratic state
representative; and Norby Chabert, a Democrat whose father and brother have held the seat in years
past.

Days before the primary on Aug.~1, according to Mr.~Chabert, Mr.~Baldone was the front-runner. But
then the Republican Party sent out mailers highlighting that the two Democrats had voted for
Mr.~Obama. (``You might be a liberal if \ldots you voted for Barack Obama,'' one read, next to a
picture of a smiling hippie.) Mr.~Callais came out on top, and Mr.~Baldone ended up in last place.

``We knew that the president was going to be a cross we had to bear,'' Mr.~Chabert said. But the
relentlessness of the anti-Obama strategy caught him by surprise. ``Of course it's effective,'' he
said. ``How could it not be?''

Mr.~Chabert fought back and beat Mr.~Callais in the Aug.~29 runoff, 54 percent to 46 percent, though
Republicans say they at least kept the race closer than it would have been.

So given Louisiana's increasingly reddish hue, the prevailing political wisdom is that a real threat
to Mr.~Vitter would come from his right. So far, that threat has not shown up.

Lt.~Gen. Russel L.~Honor\'e, retired, who was appointed by President George W.~Bush to lead military
relief efforts after Hurricane Katrina, said recently that he was not running, ``as of this time.''
No Republican officeholders have stepped up.

Until two weeks ago, Mr.~Vitter's most prominent potential opponent was Stormy Daniels, a writer,
director and actress in films of the X-rated variety.

Ms.~Daniels, who is to announce whether she will run later this month, was charged with a
misdemeanor in Florida in July after her husband told the police that she had hit him repeatedly in
the head during an argument over laundry and bills. (The charge was dropped last week.)

But now there is Mr.~Melancon (pronounced meh-LAW-sawn), a former sugar industry lobbyist who is
known around southern Louisiana as Charlie Boy. Mr.~Melancon is, understandably, emphasizing his
political independence. He is a true Blue Dog, he insists to voters, a ``pro-life, pro-gun Southern
Democrat'' who is opposed to Mr.~Obama's public health insurance option. A spokesman for the state
Democratic Party had to think to come up with major policy issues on which Mr.~Melancon and
Mr.~Vitter actually disagreed.

Mr.~Vitter himself, in a telephone interview, had no such trouble. ``No Louisiana conservative would
have voted for the bailout,'' he said. ``No Louisiana conservative would have voted for the
stimulus. No Louisiana conservative who cares about spending and debt would have voted for the Obama
budget.''

Another hurdle for Mr.~Melancon is that New Orleans, the state's Democratic bastion, has only
three-fourths of the population it had before Katrina. And the Cajun country of southwestern
Louisiana, once dependably Democratic, is becoming steadily more Republican.

``If Charlie Boy can't cut the mustard in Acadiana, he's done,'' said Paul Hardy, a former state
lieutenant governor and secretary of state, referring to the swampy Cajun region.

But there is always the matter of the prostitute.

The Democrats, who accuse Mr.~Vitter of hypocrisy for his family values rhetoric and his call on
President Bill Clinton to step down after the Monica Lewinsky affair, are not about to let voters
forget about it.

``He should not expect to be pampered by the Louisiana Democratic Party,'' said Kevin Franck, a
state party spokesman.

Mr.~Vitter is aware of this, having seen his favorability ratings fall after the admission. ``I
addressed it very directly in 2007 and said it was a very serious sin and apologized very
directly,'' he said. ``I've said that, and I'm sure I'll repeat that.''

Influential religious figures in the state, while acknowledging their disappointment with
Mr.~Vitter, say worries about Democratic policies may overshadow everything else.

``Most assuredly Senator Vitter has his baggage,'' said T.~F. Tenney, bishop emeritus of the
Pentecostals of Louisiana. ``But on the other hand, the representative is identified with the party
that has a lot more liberal baggage. Now how are people going to weigh that out -- are they going to
vote for a party, or are they going to vote for the man?''

Mr.~Hardy and others say that Mr.~Vitter may have a problem with women, who tend to be a little less
forgiving about these kinds of things.

Then again, it is Louisiana, which has its own brand of moral philosophy. Take Acadiana, where this
election will likely be won or lost.

``The Cajun mentality has never admired someone who is untrue to their spouse,'' said Morgan
Goudeau, a Democrat who was the district attorney of St.~Landry Parish for 24 years. ``But if it's
going to be done, it would be better done with a prostitute than with a neighbor's wife.''

\section{Obama Keeps Up Health Care Push, Citing Uninsured}

\lettrine{P}{resident} Obama, seeking to buttress his case for the kind of comprehensive health care
overhaul that has eluded Washington for decades, said Thursday that the number of uninsured
Americans rose by nearly 6 million as the recession intensified during the last 12 months.

On the morning after his blunt address on health care to a joint session of Congress, Mr.~Obama
addressed a group of nurses on the White House campus -- and received the endorsement of their
professional association, administration officials said. He used his brief appearance to reinforce
Wednesday night's message that his plan will bring ``security and stability'' to those who have
insurance, and coverage to those who do not.

``Now is the time to act,'' Mr.~Obama said, ``and I will not permit reform to be postponed or
imperiled by the usual ideological diversions.'' The president drew applause from the nurses, in
sharp contrast to the cool, even hostile, reactions from Republicans to his Wednesday night address.

The White House is in what Obama aides regard as the final, crucial phase of the health care debate,
in which the president will move aggressively to try to close the deal with lawmakers and the
American public. While Democratic leaders acknowledged that the speech probably did not guarantee
any new support, Mr.~Obama nonetheless invited a group of centrist Democrats -- essential to passing
the legislation -- to meet with him at the White House on Thursday afternoon.

And if Mr.~Obama picked up any Republican supporters beyond perhaps Senator Olympia Snowe of Maine,
they were invisible and silent on Thursday.

``I must say I've never heard a more partisan speech by a president in that House chamber, and I've
listened to five presidents now as a member of the House and the Senate,'' said Senator Jon Kyl of
Arizona, the Republican whip.

Recalling Mr.~Obama's statement that he would not ``waste time with those who have made the
calculation that it's better politics to kill this plan than to improve it,'' Mr.~Kyl asked, ``How
about those of us who simply have legitimate differences of opinion with him?''

Mr.~Kyl's Arizona colleague, Senator John McCain, said of Mr.~Obama speech, ``The math doesn't add
up and the record doesn't add up.''

``There is very little if anything in this package that calls for real spending reduction, and \$1
trillion is basically what it's going to cost,'' Mr.~McCain said on NBC's ``Today'' show.

House Republican leaders were also cool. Representative John A.~Boehner of Ohio, the minority
leader, said several of the president's statements were debatable, including Mr.~Obama's insistence
that his outline for change would not run up the deficit and would not lead to too much government
involvement.

``We appreciated having the president here last night,'' Mr.~Boehner said. ``Unfortunately, what the
American people got wasn't a new health care plan. It was just another lecture.''

And Representative Eric Cantor of Virginia, the House minority whip, said it was less important to
enact health care change quickly than it was ``we get it done right.''

``There are many things we can work together on,'' he said.

On Saturday, the president will return to his bully pulpit, with a campaign-style rally in
Minneapolis.

``Most Americans do have insurance and have never had less security and stability than they do right
now,'' Mr.~Obama told the nurses, ``because they're subject to the whims of health insurance
companies.''

The White House said Mr.~Obama's assertion that there are 6 million more unemployed Americans came
from a new Gallup Survey that tracked changes in the number of uninsured between September 2008 and
today.

The president also took the opportunity on Thursday to again denounce those who he said have been
guilty of exaggerations, or outright untruths, through ``chatter and the noise on radio and TV.''

But it was an outburst by Representative Joe Wilson, a South Carolina Republican, on Wednesday night
during the president's address that was still a hot topic at the Capitol on Thursday.

``You lie!'' Mr.~Wilson shouted as Mr.~Obama said illegal immigrants would not be covered by any of
the proposals making their way through Congress. Mr.~Wilson apologized for his outburst, and the
president said on Thursday that he had accepted Mr.~Wilson's expression of contrition.

``We all make mistakes,'' Mr.~Obama said.

Mr.~Boehner said Mr.~Wilson's behavior had been ``inappropriate,'' and that he was glad Mr.~Wilson
apologized.

Senator Mitch McConnell of Kentucky, the Senate Republican leader, agreed. ``I think we ought to
treat the president with respect,'' he said, adding that anything else is inappropriate.

Mr.~Wilson's outburst also created trouble for Representative Charlie Wilson, an Ohio Democrat, who
said on Thursday that he had received messages mistakenly criticizing him for attacking the
president. The Wilsons are not related.

Mr.~Obama vowed on Wednesday night to aggressively go after those who misrepresent what health care
change is all about, and he said he would ``not waste time'' with those who have made a political
calculation to oppose him. But he left the door open to working with Republicans to cut health costs
and expand coverage to millions of Americans.

The White House offensive comes after a rocky August, in which many lawmakers held public meetings
that deteriorated into shouting matches over health care.

The president placed a price tag on the plan of about \$900 billion over 10 years, which he said was
``less than we have spent on the Iraq and Afghanistan wars.'' And he sought to reassure the elderly
and the Americans who already have insurance that they would not be worse off.

As expected, Mr.~Obama repeated his support for a government insurance plan to compete with the
private sector, though he said he would consider alternatives to the ``public option.''

He sketched out a vision for a plan in which it would be illegal for insurers to drop sick people or
deny them coverage for pre-existing conditions, and in which every American would be required to
carry health coverage, just as drivers must carry auto insurance.

Mr.~Obama did embrace some fresh proposals. He announced a new initiative to create pilot projects
intended to curb medical malpractice lawsuits, a cause important to physicians and Republicans.

He endorsed a plan, included in a proposal by Senator Max Baucus, the Montana Democrat who is
chairman of the Senate Finance Committee, to help pay for expanding coverage by taxing insurance
companies that offer expensive, so-called gold-plated insurance plans. And the president promised to
include a provision that ``requires us to come forward with more spending cuts'' if the savings he
envisions do not materialize.

As Mr.~Boehner and Mr.~Cantor made clear on Thursday, Republicans are clearly primed for a fight;
many, like Senator Charles E.~Grassley, the Iowa Republican who has been deeply involved in health
negotiations, released statements about the speech even before it began. Mr.~Grassley called on Mr.
Obama to ``start building the kind of legislation that could win the support of 70 to 80 senators,''
a goal Mr.~Grassley said could not be achieved if the bill contained a new government plan.

\section{In Toronto, Directing Is Clearly Women's Work}

\lettrine{F}{ROM} the tattoo on Diablo Cody's bicep to Lone Scherfig's leopard-spot pumps, it was
impossible not to notice: The 34th Toronto International Film Festival opened on Thursday with the
women in charge.

While still struggling to find their place in the movie industry at large -- the number of directors
at American studios remains well over 90 percent male -- female filmmakers have managed to occupy
some of this 10-day festival's most valuable slots: those showcase screenings and press conferences
in the first couple of days, when everyone is still paying attention.

Thursday's most raucous event was almost certain to be the 11:59 p.m. red-carpet debut of 20th
Century Fox's ``Jennifer's Body,'' directed by Karyn Kusama (``Girlfight'') from a script by
Ms.~Cody (``Juno''), in which Megan Fox plays a high school sex bomb who, quite literally, turns
into a man-eater.

According to Natalie Johnson, a spokeswoman for Fox, tickets to the midnight show at the landmark
Ryerson Theater, which seats more than 1,200, were gone within two hours of going on sale last week.
(``Hell is a teenage girl,'' runs a theme-setting line from the film.)

``Jennifer's Body,'' which opens in commercial theaters next Friday, got its first festival
screening at noon on Thursday. Several hundred press and film industry types, normally a jaded
bunch, were lined up for a look at the Kusama-Cody-Fox combination's take on female vengeance.

The audience was laughing in all the right places, a good sign for the film, which is walking a fine
line between comedy, horror and a postpunk feminism that is telegraphed by the title's cute pink
script in the opening credits.

But the deeper question is whether the Toronto festival's first couple of days might help propel a
clutch of female directors to the front of Hollywood's award race.

Something like that happened in 2003, when a Toronto screening of ``Lost in Translation'' put Sofia
Coppola on the path to a best-director Oscar nomination. She is one of only three women ever to earn
that distinction, the others being Lina Wertmüller for ``Seven Beauties'' and Jane Campion for ``The
Piano.''

None of them won the directing Oscar. But 2003 became known as a good year for women, as Niki Caro,
directing ``Whale Rider''; Catherine Hardwicke, directing ``Thirteen''; Patty Jenkins, with
``Monster''; and Shari Springer Berman, with ``American Splendor,'' all joined Ms.~Coppola in making
a strong impression.

Ms.~Campion is back in contention for prizes this year with ``Bright Star,'' a romance about the
poet John Keats and his muse Fanny Brawne, from the new film company Apparition. The film began
screening here Thursday, as Ms.~Campion and her team, including the actors Ben Whishaw and Abbie
Cornish, gathered in advance of a Friday night presentation ahead of its commercial opening next
week.

By Thursday morning Ms.~Scherfig, a Danish director, was already in motion. Preparing for a 6 p.m.
screening of ``An Education,'' her offbeat romance from Sony Pictures Classics, Ms.~Scherfig was at
the Four Seasons Hotel, doing the occasional press interview and getting ready for a reunion with
her cast members, who include Peter Sarsgaard and Carey Mulligan.

``For me it was never an issue to project anything that had to do with gender,'' Ms.~Scherfig said
of her own take on filmmaker demographics. ``All my films have had men in their late 30s in the
lead.''

This stop was not Ms.~Scherfig's first -- she and the film had just dropped in from a festival in
Telluride, Colo., and had made a splash at the Sundance festival in January in Park City, Utah --
nor the last, as she was planning to head for yet another festival appearance in London.

Ms.~Caro was also expected in Toronto with her latest film, ``The Vintner's Luck,'' which was to be
toasted at a New Zealand film cocktail party Friday, ahead of a weekend screening. Other women with
films in the Toronto festival include Rebecca Miller (``The Private Lives of Pippa Lee''), Samantha
Morton (``The Unloved'') and Leanne Pooley (``The Topp Twins'').

Drew Barrymore, like Ms.~Morton an actress making her directorial debut, has already felt like a
presence with ``Whip It,'' a roller-derby romp starring the 2007 awards season darling Ellen Page
(``Juno''), even though the film was not scheduled to screen until Sunday. Ms.~Barrymore's pose on
the cover of an unofficial festival guide is almost as sassy as the e-mail promotions for ``All Girl
Roller Derby Action,'' a live exhibition of the sport with appearances by the ``Whip It'' cast, in a
public square on Sunday evening.

All that, and Kathryn Bigelow's ``Hurt Locker,'' which was shown last year at Toronto, has been
generating Oscar talk since its release earlier this summer by Summit Entertainment.

Anyone who is looking can read the signs, including those ubiquitous photographs of Ms.~Cody
flashing her bicep tattoo of a bikini-clad beauty: Women have staked a claim on the season.

\section{Looking for a (Long) Leg Up}

\lettrine{T}{EN} years ago, when Amy Astley, then the beauty director of Vogue, began working on a
prototype of a spinoff magazine for teenagers, the question she was most commonly asked by potential
readers was this: ``How can I become a fashion model?''

``It was really depressing,'' said Ms.~Astley, now the editor of Teen Vogue.

They all wanted to be models, and of course they couldn't, so Ms.~Astley found herself offering
teenagers a harsh dose of reality, rather than something positive or inspiring.

``That was the most visible career in fashion at the time,'' she said. ``But I've seen a really
profound change. They don't ask me anymore how to become a model, or if I've met Britney or Justin
Timberlake.''

Now they ask how they can get her job.

What changed, Ms.~Astley said, is that teenagers around the world have become interested in all
sorts of careers in fashion as a result of the industry's increasingly outsize place in popular
culture. ``Project Runway,'' the designer competition originally set at Parsons the New School for
Design, has alone been credited with causing a spike in applications to fashion schools. At Parsons,
applications have gone up 41 percent over the last five years. At Pratt Institute, they have gone up
20 percent.

``We have tour buses stopping outside so tourists can take pictures of us,'' said Simon Collins, the
dean of fashion at Parsons. ``They're looking for Tim Gunn or Heidi Klum or something.''

And given the intense scrutiny directed at ``The September Issue,'' the behind-the-scenes
documentary of a single issue of Vogue, and the ever-escalating coverage of Fashion Week, which
begins Thursday in Bryant Park, it stands to reason that there will be even more interest to come.
Feeding on the demand, Teen Vogue has its own book coming out next month called ``The Teen Vogue
Handbook,'' a how-to guide for students dreaming of jobs as a designer, stylist, photographer or
editor.

For much of America's youth, fashion is where it's at. But this wave of Anna Wintours and Michael
Korses in training is coming at a moment when the industry is shrinking; retailers are collapsing;
several magazines within Teen Vogue's parent company, Cond\'e Nast, have closed; and jobs, of any
sort, are scarce. A report last month from the NPD Group estimated that 12 percent of fashion
companies will not survive the recession.

The situation is not entirely grim for new fashion graduates, even though the National Association
of Colleges and Employers said in March that employers expected to hire 22 percent fewer seniors
graduating in 2009 for entry level positions. Normally about 90 percent of Parsons seniors find
jobs, but that figure dipped by only 10 percent. At Pratt and at the Fashion Institute of
Technology, they have remained about the same.

``But we are seeing a trend toward some jobs disappearing into unpaid internships, which is a little
troubling,'' said Judy Nylen, the director of career services at Pratt.

So what is a young person trying to break into fashion supposed to do?

DON'T EXPECT TO START AT THE TOP

Let us take the example of Sang A Im-Propp, who was a pop star in Korea before she decided, while on
a business trip to New York, that she wanted to be in fashion. This was nearly a decade ago, and
Ms.~Im-Propp's command of English was tenuous, but she enrolled at Parsons and in short order found
herself an internship with Victoria Bartlett, a noted stylist and designer whom she admired and
hoped would introduce her to the glamorous world of design. Instead, Ms.~Im-Propp found it difficult
to understand Ms.~Bartlett's heavy British accent, and at first she thought she had misunderstood
just what Ms.~Bartlett was asking her to do. Get cupcakes?

Not just any cupcakes, but the glossy butter-cream confections from the Cupcake Cafe, which is a
four-block crosstown walk from Ms.~Bartlett's studio through the dodgy garment district, and it was
freezing outside.

``It made me cry a lot,'' Ms.~Im-Propp said. ``Vicky is an amazing artist, but she can be
difficult.''

But Ms.~Im-Propp persisted, and after many cupcake runs, was entrusted with the research projects,
location scouting and shopping collections Ms.~Bartlett did not have time to see. When she decided
in 2006 to start her own collection of handbags, under the label Sang A, Ms.~Bartlett personally
recommended her to a showroom.

Ms.~Bartlett, reminded of the cupcake episodes, said she had been a deliberate taskmaster with
interns, having encountered a few who did not live up to expectations. There was one who, without
asking, borrowed her car on several occasions, and another who ran off to Los Angeles with the
company computer.

``You can't be a princess in this business,'' she said. ``People see fashion from the end result,
which is kind of a false facade. They only see this beautiful, glamorous world, but I don't think
they realize it is one of the hardest careers out there.''

BE PREPARED TO SUFFER

This is a point that is made repeatedly in ``The Teen Vogue Handbook,'' which compiles stories of
how many now-famous designers and editors got their starts. The common themes in the histories of
Marc Jacobs, Ms.~Wintour, Gucci Westman and Mario Testino are that hard work and persistence lead to
opportunities to be seized -- with a caveat, of course, that it'll be a lot harder than that for
anyone else. Karl Lagerfeld says that rather archly: ``Are you ready to accept injustice? The idea
of the fashion industry may look better from the outside. It can look like the world of dream jobs
-- for a very happy few.''

Still interested?

It's a wonder, given the portrayal of fashion in ``The Devil Wears Prada'' and the counterpoint of
``The September Issue,'' that anyone would think this is a glamorous business. (There is the specter
of Grace Coddington in the new film directing a colleague, ``Don't be too nice, not even to me,
because you'll lose.'')

In the real world, it's just as tough. Prabal Gurung, the hot newcomer, said that his first job was
working for a cuckoo designer in India who staged a fashion show with models eating chicken wings.
Kelly Cutrone, the salty publicist and reality-show fixture, often starts a job interview by telling
the applicant, ``You'll be very lucky if you start and end your career liking clothes.''

``The truth is,'' said the designer Phillip Lim, ``a lot of doors are shut right now, and no one is
going to open them.'' But Mr.~Lim cited his own start as a reason not to give up hope. As a young
salesman at Barneys in Los Angeles, he was so na\"ive that he simply picked up the phone and called
the office of Katayone Adeli and asked for a job, which brings us to our next point.

NAÏVITÉ IS BLISS

``Don't listen to other people,'' Ms.~Astley said. ``If you want to work in fashion, you should do
it. You should move to New York.''

Among the current crop of interns at Teen Vogue, there is little fear that the future of fashion
will happen without their participation. They tell stories about 12-hour days of sorting through
piles of shopping bags looking for a single skirt; or blisters from running garment bags around the
city; or the bummer of being sent to a famous designer's showroom and glimpsing only the messenger
center. Or the thrill, in the case of Media Brecher, who is 20 and a student at Barnard College, of
seeing a headline she suggested for a denim story, ``Bleach Streak,'' appear in the August issue.

(It does help to have connections, as Ms.~Brecher is the daughter of The Wall Street Journal wine
columnists Dorothy J.~Gaiter and John Brecher; and Hope Brimelow, another intern, met Teen Vogue
editors in Paris while her mother, a producer on ``60 Minutes,'' was filming a profile of Anna
Wintour. But Ms.~Astley said that connections are not required at a publication that employs as many
as 40 interns at a time.)

``Anyone can contribute to fashion now because of how accessible it has become,'' said Vanessa
Stylianos, 20, a New York University communications major who started working in the fashion closet
last month.

ON THE OTHER HAND \ldots

Ms.~Astley recalled a recent job applicant who was clearly unqualified to work at her magazine.

``I interviewed someone who hadn't seen 'Twilight,' '' she said. ``You can't work at Teen Vogue if
you haven't seen 'Twilight.' ''

With designers, too, it helps to be aware of their work, even wearing an item of their clothing to a
job interview, but goodness, not head to toe. This is a blunder seen by Karen Harvey, a fashion
headhunter, who said that it looks as if you are trying too hard.

``You want to demonstrate your personal style and then integrate something of theirs,'' she said.
``They want to see how your styles will connect.''

ASK LOTS OF QUESTIONS

This can be tricky. Most employers value a go-getter who doesn't need to be told what to do, but
they also appreciate a display of interest. A suggestion from Kenneth Wyse, the president of
licensing at Phillips-Van Heusen, who also is president of the Fashion Scholarship Fund, an industry
group that helps mentor young people with internships: Submit a list of questions once a week.

``One of my favorite stories was that after working here for a summer, a young lady said to me,
'Gee, I had no idea that a collection business could lose money,' '' Mr.~Wyse recalled.

BE PERSISTENT, AND IT HELPS TO BE CHARMING

It was March 1996 in the suburbs of Detroit, where Jamie Rubin, age 10, was doing her homework,
writing a speech about what she wanted to be when she grew up. Ms.~Rubin had no idea, but her
favorite dress was a really simple gray jumper by Nicole Miller, and so, she said, ``I just decided
that's what I wanted to do, to give little girls their favorite dress.''

As part of the assignment, Ms.~Rubin filmed a video in which she charmingly announced, ``When I grow
up, I'm going to be the next Nicole Miller!'' She was so convincing that her parents sent the video
to the designer's showroom in New York, where it became an office favorite. Ms.~Miller was so taken
by Ms.~Rubin that she invited her for a tour of her showroom and, eight years later, when Ms.~Rubin
enrolled at Parsons, offered her an internship for a summer, working right by her side.

``I felt like I was meeting Oz,'' Ms.~Rubin said. ``I couldn't even speak when I first met her.''

If her story sounds too easy, Ms.~Rubin, now 23, will point out that things were not so simple once
she graduated. Even though she had also interned for Women's Wear Daily and for Proenza Schouler,
she had little response from the hundreds of r\'esum\'es she sent out, except for one she sent to a
showroom where she was dying to work, as she noted in a cover letter that she accidentally addressed
to one of its competitors. ``I got an e-mail back saying, 'That's nice, if you want to a job there,
you should send them an application,' '' Ms.~Rubin said.

She did land an interview at Dolce \& Gabbana and bought a blouse for the occasion, but it was
loose-fitting and the reason Ms.~Rubin suspects she never got a call back. After much persistence,
she was offered a job at a creative agency that represents Tod's, Hogan and other luxury brands.

``The door will close on you 900 times,'' she said. ``So you've got to keep your skin tough and your
goals very focused. I walk into work every day and know I'm going to be challenged and inspired, and
that's the recipe for happiness in any job.''

\section{Record Numbers of Chinese Wed on 09/09/09}

\lettrine{C}{ities} across China have reported record numbers of marriages after tens of thousands
of couples decided that 9/9/09 would be an auspicious date to wed.

The number nine in Mandarin, jiu, sounds the same as the word for ``perpetual'' and ``forever.''

In Beijing on Wednesday, 18,979 couples married or received licenses, according to the official
Chinese news agency Xinhua. That figure shattered the capital's previous daily record of more than
15,000, set by couples who chose Aug.~8 last year, 8/8/08, because eight is a lucky number in China
and is pronounced ba, which sounds like the word for becoming rich. Aug.~8 also was the opening day
of the Olympic Summer Games in Beijing.

Huge starry-eyed turnouts were reported in cities throughout China on Wednesday, Xinhua said, and it
appeared likely that the day's total would also break the national record of more than 314,000 set
last Aug.~8.

The matrimonial import of the day went beyond China.

Jimmy Bell, a manager of 99 Cents Only in Hollywood, said nine couples were married at the Sunset
Boulevard store on Wednesday. Each couple paid 99 cents -- at Register No.~9 -- and received
receipts for the transaction.

``No, we haven't had any refunds,'' Mr.~Bell said Thursday, laughing. ``Not yet.''

\section{China Weighs Tariffs on Some U.S.~Exports as Tensions Rise}

\lettrine{C}{hina} unexpectedly increased pressure Sunday on the United States in a widening trade
dispute, taking the first steps toward imposing tariffs on American exports of automotive products
and chicken meat in retaliation for President Obama's decision late Friday to levy\cite{levy}
tariffs on tires from China.

The Chinese government's strong countermove followed a weekend of nationalistic
vitriol\cite{vitriol} against the United States on Chinese Web sites in response to the tire tariff.
``The U.S.~is shameless!'' said one posting, while another called on the Chinese government to sell
all of its huge holdings of Treasury bonds.

The impact of the dispute extends well beyond tires, chickens and cars. Both governments are facing
domestic pressure to take a tougher stand against the other on economic issues. But the trade battle
increases political tensions between the two nations even as they try to work together to revive the
global economy and combat mutual security threats, like the nuclear ambitions of Iran and North
Korea.

Mr.~Obama's decision to impose a tariff of up to 35 percent on Chinese tires is a signal that he
plans to deliver on his promise to labor unions that he would more strictly enforce trade laws,
especially against China, which has become the world's factory while the United States has lost
millions of manufacturing jobs. The trade deficit with China was a record \$268 billion in 2008.

China had initially issued a fairly formulaic criticism of the tire dispute Saturday. But rising
nationalism in China is making it harder for Chinese officials to gloss over American criticism.

``All kinds of policymaking, not just trade policy, is increasingly reactive to Internet opinion,''
said Victor Shih, a Northwestern University specialist in economic policy formulation.

Eswar Prasad, a former China division chief at the International Monetary Fund, said that rising
trade tensions between the United States and China could become hard to control. They could cloud
the Group of 20 meeting of leaders of industrialized and fast-growing emerging nations in Pittsburgh
on Sept.~24 and 25, and perhaps affect Mr.~Obama's visit to Beijing in November.

``This spat\cite{spat} about tires and chickens could turn ugly very quickly,'' Mr.~Prasad said.

China exported \$1.3 billion in tires to the United States in the first seven months of this year,
while the United States shipped about \$800 million in automotive products and \$376 million in
chicken meat to China, according to data from Global Trade Information Services in Columbia, S.C.

For many years, American politicians have been able to take credit domestically for standing up to
China by taking largely symbolic measures against Chinese exports in narrowly defined categories. In
the last five years, the Commerce Department has restricted Chinese imports of goods as diverse as
bras and oil well equipment.

For the most part, Chinese officials have grumbled but done little, preferring to preserve a trade
relationship in which the United States buys \$4.46 worth of Chinese goods for every \$1 worth of
American goods sold to China.

Now, the delicate equilibrium is being disturbed.

China's commerce ministry announced Sunday that it would investigate ``certain imported automotive
products and certain imported chicken meat products originating from the United States'' to
determine if they were being subsidized or ``dumped'' below cost in the Chinese market. A finding of
subsidies or dumping would allow China to impose tariffs on these imports.

The ministry did not mention the tire dispute in its announcement, portraying the investigations as
``based on the laws of our country and on World Trade Organization rules.''

But the timing of the announcement -- on a weekend and just after the tire decision in Washington --
sent an unmistakable message of retaliation\cite{retaliation}. The official Xinhua news agency Web
site prominently linked its reports on the tire dispute and the Chinese investigations.

The commerce ministry statement, posted on its Web site, also hinted obliquely\cite{obliquely} at
the harm that a trade war could do while Western nations and Japan struggle to emerge from a severe
economic downturn. ``China is willing to continue efforts with various countries to make sure that
the world economy recovers as quickly as possible,'' the statement said.

The Chinese government sometimes organizes blog postings to defend its own policies. But some
postings on the tire decision have been implicitly critical of the Chinese government, making it
unlikely that they are part of an orchestrated effort.

``Why did our government purchase so much U.S.~government debt?'' said one posting signed by a
``Group of Angry Youths.'' The item continued, ``We should get rid of all such U.S.~investments.''

China has accumulated \$2 trillion in foreign reserves, mostly in Treasury bonds and other
dollar-denominated assets, and held down the value of its currency, which has kept Chinese goods
quite inexpensive in foreign markets. China's exports have soared -- China surpassed Germany in the
first half of this year as the world's largest exporter -- while China's imports have lagged, except
for commodities like iron ore and oil that China lacks.

Worries that China might sell Treasury bonds -- or even slow down its purchases of them -- have been
a concern for the Bush and Obama administrations as they have tried to figure out how to address
China's trade and currency policies.

At the same time, the Chinese economy relies heavily on exports to the United States, while the
American economy is much less dependent on exports in the other direction. Exports to the United
States, at 6 percent of China's entire economic output, account for 13 times as large a share of the
Chinese economy as exports to China represent for the United States economy.

Products involved in trade disputes between the United States and China together make up only a
minuscule sliver of the two countries' trade relationship.

The bigger risk for China, economists and corporate executives have periodically warned, is that
trade frictions could cause multinationals to rethink their heavy reliance on Chinese factories in
their supply chains. The Chinese targeting of autos and chickens affects two industries that may
have the political muscle in the United States to dissuade the Obama administration from
aggressively challenging China's policies.

General Motors sees much of its growth coming from its China subsidiary, the second-largest auto
company in China after Volkswagen. And the farm lobby in the United States has long pressed for
maximum access to a market of 1.3 billion mouths.

But spotlighting automotive trade may be risky for China. G.M. and Ford both rely mostly on local
production to supply the Chinese market, while China is rapidly increasing auto parts shipments to
the United States.

\section{Freshly Minted Opel Deal Draws Critics}

\lettrine{G}{erman} Chancellor Angela Merkel is having to defend the sale of the troubled carmaker
Opel to a Canadian-Russian consortium, just days after hailing the deal.

Criticism grew sharper over the weekend on the terms of the deal, which was announced on Thursday in
Berlin, beginning with unions that said more jobs would be cut than expected.

Among the most vocal\cite{vocal} critics were government representatives of the Opel Trust, which
was established in May to oversee the search for a buyer of 65 percent of Opel. Members said the
German government was putting politics before business considerations in its deal with Magna
International, a Canadian auto parts dealer, and Sberbank, the Russian partner.

On Sunday, the Belgian foreign affairs minister, Yves Leterme, concerned that Germany might have
sought to protect its own factories at the cost of others, said he would discuss the deal with
European trade ministers at an informal meeting.

And the European Commission, the executive arm of the European Union, warned of protectionism in the
sale of Opel, which is based in Germany and is G.M.'s main operation in Europe.

``If something happens against the rules, action will be taken,'' said a spokesman for Neelie Kroes,
the E.U.'s competition commissioner.

Mrs.~Merkel, who is facing elections on Sept.~27, denied that the government was using taxpayer
money to bail out Opel.

``We are not rescuing corporations, but are maintaining a chance for them in the financial crisis,''
Mrs.~Merkel told the newspaper Süddeutsche Zeitung in comments published on Saturday. ``

Details of the sale are not final with regard to job cuts, restructuring plans and government
subsidies.

The German government, along with four of the states that have Opel plants, had earlier promised 4.5
billion euros, or \$6.6 billion, in government guarantees. But they have not yet confirmed those
guarantees.

In return, Magna has said it would put up 500 million euros for restructuring costs. But Fred Irwin,
chairman of the Opel Trust, has not confirmed that either.

Berthold Huber, chairman of IG Metall, the influential trade union, told the mass circulation
newspaper Bild am Sonntag in comments published on Sunday that it was time that ``Magna was serious
about its intentions'' with regard to job cuts.

G.M. had stressed last week that it still needed union agreement, in writing, for restructuring the
four Opel plants in Germany. They employ almost half of G.M. Europe's 55,000-strong European work
force. IG Metall said in a statement that it expected the negotiations to be ``tough'' and feared
that more jobs would be cut than expected because Magna had not specified that in writing.

The toughest criticism may have come from the German government's own representatives on the Opel
Trust, which approved the deal last week at the last minute, apparently under immense pressure.

Dirk Pfeil, a member of the trust, said the sale to Magna ``was a serious mistake.''

``Had I known that it was going to be a purely political decision, and that business management
considerations would be left in the shadows,'' Mr.~Pfeil told German Radio, he would have come to a
different conclusion.

``It is unique in my professional life that the party which has been acting as a guarantor also
decides on the buyer,'' Mr.~Pfeil said. ``That is a bit odd, and it does not fit into the way things
are done.''

Another government representative, Manfred Wennemer, voiced his concern the moment the deal was
announced by Mr.~Irwin and John Smith, G.M.'s chief negotiator.

Mr.~Wennemer, who is the former head of the auto supplier Continental, said he doubted whether Magna
could make Opel competitive and said the deal was risky for German taxpayers. He even raised the
possibility of insolvency for Opel in the next year or two.

Under the deal, Magna and Sberbank would purchase a 55 percent stake in Opel and the smaller
Vauxhall, which is based in Britain.

G.M., which acquired Opel in 1929, would keep a 35 percent stake, while Opel's employees would be
provided 10 percent.

\section{In Shanghai, Police Detain Former Worker at Coca-Cola}

\lettrine{P}{olice} here have detained a former employee of a Coca-Cola bottling plant, whom they
accused of corruption and bribery\cite{bribery}.

The detention of the employee from the Shanghai Shen-Mei Beverage and Food Company, a bottling plant
partly owned by Coca-Cola, was reported over the weekend by China's state-run news media, which said
the employee took about \$1.5 million in bribes.

A spokesman for Coca-Cola, Kenth Kaerhoeg, confirmed on Sunday that a female middle manager at the
plant was detained by Shanghai police this year and then dismissed by the bottling company.

Mr.~Kaerhoeg declined to give further details about the case, but he said Coca-Cola was cooperating
with the investigation.

Calls to Shen-Mei, in which Coca-Cola has a minority stake, went unanswered on Sunday and Shanghai
police officials could not be reached for comment.

The detention is the second prominent bribery case this year involving a global company operating in
China.

In July, four Shanghai-based employees of the British-Australian mining giant Rio Tinto, including
an Australian citizen, were detained and later formally charged with corruption and bribery.

That case emerged after tense iron ore negotiations between Rio Tinto and Chinese steel makers and
after Rio Tinto scrapped a deal with a big Chinese company, leading some analysts to speculate that
it was partly politically motivated.

Rio Tinto has strongly denied its Shanghai employees were involved in bribery or corruption. Chinese
officials have insisted the case is not politically motivated.

China recently stepped up an anticorruption drive that has ensnared numerous high-ranking government
officials and even Huang Guangyu, a flamboyant entrepreneur listed by Forbes as the country's
second-richest man.

That corruption drive, coupled with the Rio Tinto case -- which initially involved allegations of
stealing state secrets -- has unnerved many foreign business people working in China.

But even though corruption is pervasive in China, very few executives working for global companies
have been detained or arrested by Chinese police. Legal experts say that is partly because
investigators focus on the government officials who accept the bribes; but they also say the
government is reluctant to go after foreign companies.

There have been indications in China's state-run news media, though, that investigators could begin
to target multinational corporations.

And many foreign companies operating here are warning employees to use extra caution because the
United States Justice Department has stepped up its enforcement of the Foreign Corrupt Practices
Act, which forbids the paying of bribes to government officials to win or influence business deals.

In July, Control Components, a company based in California, pleaded guilty to foreign bribery
charges and agreed to pay an \$18.2 million fine.

The company had been accused of paying bribes to officials in more than 30 countries over many
years, including to Chinese state-owned companies.

In the Coca-Cola case, the National Business Daily, a Chinese state-controlled publication, reported
over the weekend that the bottling plant middle manager had worked in the marketing department of
the bottling company.

Coca-Cola has huge operations in China, including about 38 bottling plants, and about 30,000
employees in China, and is pushing to expand aggressively here.

This year, the company's move to pay \$2.4 billion to acquire one of China's biggest juice makers,
the Huiyuan Juice Group, was blocked by the Chinese government, which cited antitrust concerns.

\section{To Hopeful Makers, the Electric Car's Time Is Here}

\lettrine{T}{he} electric car is at the starting line, and the gun is about to sound. Now automakers
must prove that the technology -- and the market are ready.

After years of talk and prototypes, some automobile makers believe the electric vehicle is about to
become more than just a science experiment.

The French company Renault will unveil a lineup at the Frankfurt motor show this week that includes
a purely electric sedan, without a backup internal combustion engine. Renault says the vehicle will
be in showrooms by 2011.

General Motors, which recently emerged from bankruptcy protection, is offering the Chevrolet Volt as
one of its comeback cars. The Volt is an electric vehicle with a backup gasoline-run generator for
longer trips.

BMW will be unveiling the Vision EfficientDynamics, a plug-in diesel\cite{diesel}-electric concept
car. And Volkswagen is adding an electric version to its Up! concept car.

``This is not a false dawn,'' said Paul Scott, vice president and a founder of Plug In America, a
group that has long accused automakers of moving too slowly on electric cars. ``This is the real
thing.''

Still, questions about relatively limited range and high costs vex electric vehicles. Two prominent
skeptics are Audi, the German luxury carmaker, and Toyota, maker of the Prius Hybrid. These firms
are developing fully electric vehicles of their own, but for the moment, are emphasizing other
technologies. Audi is pushing ``clean diesel'' engines, and Toyota is working on a hydrogen-based
fuel cell technology that would obviate the need for heavy batteries.

Other companies argue that public fascination with electric cars can be made profitable -- soon.

``Once we put the cars on the road and can have customers test-drive the cars, we are quite
convinced that we will transform this interest into a clear market,'' said Thierry Koskas, the
director of Renault's electric vehicle program.

Renault's French rival, PSA Peugeot Citro\"en, sold 10,000 electric cars from 1995 to 2005, mainly
in the form of the Peugeot 106, a compact car. Waning consumer interest and a European phase-out of
cadmium, a crucial ingredient in the car's batteries, ended that venture. But Peugeot is back,
having signed a deal with Mitsubishi last week to sell electric cars in Europe by the end of 2010.

Electric vehicles might capture the public imagination in much the way that hybrids -- the Toyota
Prius being the most important --

have done in the last few years. But if ambitious rollout schedules give way to delays, automakers
could end up feeding consumer cynicism\cite{cynicism}.

``It could be that some manufacturers are awakening expectations that they cannot fulfill in a
reasonable time frame,'' said Willi Diez, director of the Institute for Automotive Research in
Nürtingen-Geislingen, near Stuttgart.

Advocates of traditional engines have by no means ceded\cite{cede} the field.

In an interview with Lawrence Ulrich, an auto blogger for the Web portal MSN, Johan de Nysschen,
president of Audi of America , called the Chevrolet Volt as ``a car for idiots.'' He argued that the
\$40,000 base price at which G.M. had hinted smacked of fantasy or, at the least, a narrow appeal to
the most environmentally committed.

``No one is going to pay a \$15,000 premium for a car that competes with'' conventional sedans that
cost around \$25,000, he said. ``So there are not enough idiots who will buy it.''

Renault, however, says its cars are a ``rational choice'' for getting around town. ``We don't just
want people to buy an electric vehicle because it's fashionable,'' Mr.~Koskas said.

In Frankfurt, Renault is introducing a lineup of all-electric cars that do without a ``range
extender,'' which is a combustion engine that turns an electricity generator when the battery's
charge is depleted.

(The Volt will have a gasoline engine to charge the batteries on long trips. In the Prius, the range
extender is an entire gasoline powertrain alongside the electrics.)

Since most Europeans drive less than 100 kilometers, or 62 miles, a day, or less than the range of a
single battery charge, Renault says it can start selling cars before an elaborate infrastructure for
charging them is in place. Its electric cars are scheduled to go on sale by 2011 in some countries
in Europe and possibly Asia, but not in the United States, Mr.~Koskas said.

But Nissan, the Japanese automaker that Renault controls, is planning to introduce the Leaf, an
electric sedan, in late 2010 in Japan, the United States and Europe. (It will share many components
with Renault vehicles.) Renault is not disclosing the cost of its vehicles.

For skeptics like Mr.~Nysschen, though, government support for electric vehicles tilts the playing
field. The Volt and others, he said in a follow-up statement, exist by the grace of
``taxpayer-funded subsidies.''

The Volt's attractiveness will rest in part on a \$7,500 United States government tax credit for
purchasers of electric vehicles. And France is offering almost the same amount.

Mr.~Koskas counters that the subsidies should be considered part of the start-up phase and should
fade away as scale is achieved.

``We will not need this to last forever,'' he said. ``We are in an industry that depends very much
on volumes.''

Indeed, some critics are hedging\cite{hedge} their bets.

Mr.~Nysschen's company, Audi, created a Web site, electricityuntamed.com, which is counting down to
the opening of the Frankfurt auto show Tuesday, where it will unveil a prototype electric car. And
Toyota is showing its own Prius Plug-In, scheduled for full release in 2012.

\section{Ad Shift Throws Blogs a Business Lifeline}

\lettrine{L}{isa} Sugar began blogging about celebrity\cite{celebrity} gossip in her spare time four
years ago. Now she and her husband, Brian, have a little media empire called, sensibly enough, Sugar
Inc., with 12 blogs, 11 million readers a month and advertisers like Chanel and Sony.

The dream of quitting the day job and making a living from blog revenue has proved to be far-fetched
for most bloggers. But a few entrepreneurs, like the Sugars, have found success in blog networks.

Such networks put blogs on various topics under some form of central control, like a digital-era
Cond\'e Nast. Though they do not command nearly the same ad rates that glossy magazines do, they are
attracting ad dollars while magazines are losing them.

Sugar's ad revenue increased 20 percent in the first half of the year, and the company is on track
to double its revenue and turn a profit this year, said Mr.~Sugar, the company's chief executive.
Gawker Media, one of the earliest and biggest blog networks, reported that ad revenue was up 45
percent in the first half of this year.

Both companies are private, and neither would disclose more specific figures, but by some estimates
the larger networks have annual revenue in the low tens of millions of dollars.

Meanwhile, advertising revenue for magazines dropped 21 percent in the first half, and the number of
ad pages sold dropped 28 percent, according to the Publishers Information Bureau.

The blog networks that have survived the downturn in advertising and the explosion of competing
content on the Web credit their obsessive coverage of narrow topics, along with business models that
reach beyond advertising.

``It's actually really hard creating compelling content that brings an audience,'' Mr.~Sugar said.

Hobbyist bloggers first got the idea that their online entries could be profitable in 2002 when Nick
Denton, a former reporter for The Financial Times, started what would become Gawker Media. The catty
network quickly grew into a powerhouse, and it now has eight blogs, 20 million monthly readers and
more than 150 full- and part-time employees. When Jason Calacanis sold his blog network, Weblogs
Inc., to AOL for a reported \$25 million in 2005, the notion that blogging could be a business was
cemented.

Today, blog networks range from big, like Gawker and Sugar, to smaller and more focused, like The
Business Insider, VentureBeat and the GigaOm Network, which cover business and technology. There are
also large networks of blogs that share ads but not editorial control, like BlogHer and Glam Media.

Sugar's blogs -- with names like PopSugar (celebrity gossip), BellaSugar (beauty) and LilSugar
(mothering) -- are all edited and designed with 28-year-old women in mind. The writers and sales
staff are fanatical about going after that ideal visitor, Mr.~Sugar said. The posts are short, light
and sarcasm-free, with big photos and headlines like ``The Jolie-Pitts Bring the Twins to the Golden
Arches'' and ``10 Sexy Bedhead Hairstyles to Try Today.''

Other blog networks take a similarly narrow focus. Gawker Media's business cards used to read
``Unhealthily obsessed,'' said Lockhart Steele, the company's former managing editor. He follows the
same motto at Curbed, a company he founded that publishes blogs on real estate, dining, shopping and
travel in New York, Los Angeles and San Francisco.

``We are completely obsessed with every square inch of the areas we cover,'' Mr.~Steele said. ``It's
the classic blogger model: give people everything they could possibly want on a topic.'' For the
Curbed sites, that sometimes means two dozen posts a day just about New York real estate.

That all-encompassing coverage is meant to keep readers coming back many times a day, in turn
attracting advertisers.

Blogs, with their unpredictable and sometimes edgy\cite{edgy} content, can be frightening to
advertisers, but blog networks are less risky, said Shenan Reed, a founder of Morpheus Media, a
digital marketing agency that represents brands like Louis Vuitton and L'Or\'eal.

``When you're dealing with a company where the editorial control is living under one roof, you feel
like there's a consistency in the message, which is what makes Sugar, Gawker and Curbed
fantastically interesting to us,'' Ms.~Reed said.

Blog networks also make it easier for advertisers, who do not have time to sift through the millions
of blogs on the Web, to reach a big blog audience. For example, Morpheus ran an ad campaign for
Neiman Marcus on the Sugar network in which it dressed each of the Sugar blogs' cartoon mascots in
outfits from the store, and readers could click on the clothes to buy them.

Yet for Sugar, the dream of growing a large and profitable media business on advertising alone did
not come true.

``Outside of a few examples, I think people way overestimated the amount of ad revenue that was
there for the taking for blogging as a business,'' said Scott Rosenberg, author of the recent book
``Say Everything: How Blogging Began, What It's Becoming, and Why It Matters.''

In 2007, Sugar, which is backed by Sequoia Capital and has 105 employees, acquired ShopStyle, an
e-commerce site. Today it brings in half of Sugar's revenue. At ShopStyle, shoppers can browse
online retailers' selections, and Sugar gets paid when they click through to a retailer or make a
purchase.

Sugar would not have been able to grow as it has without both sources of revenue, said Mr.~Sugar,
who was formerly vice president for e-commerce at J.~Crew. The company plans to eventually make
money selling virtual goods, too, like stilettos in a fashion game.

Blog networks have also had to change the type of content they publish. Though many started by
publishing links to and commentary on other Web sites' material, there is now so much material
online that they need to offer original reporting to stand out, Mr.~Steele said.

Each of Curbed's sites has a full-time editor and a few freelancers, and they have started competing
with local papers for scoops. Curbed often depends on readers' tips and might publish 10 updates to
an article in response to them.

When Gawker blogs get exclusive content, like leaked videos, they can get more than a hundred times
their typical traffic. ``Scoops pay,'' Mr.~Denton said.

Sugar gets press passes to the Oscars and recently started producing videos about shopping and
celebrities. It also developed blogging software called OnSugar. When bloggers use the software to
showcase clothes and readers click to buy them, the bloggers and Sugar get a cut. Sugar plans to
eventually sell ads on these blogs.

``Perpetual\cite{perpetual} movement is the essence of survival and prosperity online,'' said
Michael Moritz, the Sequoia investor who backed Google, Yahoo and Sugar. ``If online media and
entertainment companies don't improve every day, they will just wind up as the newfangled version of
Reader's Digest -- bankrupt.''

\section{Facebook Is Going for Some Twitter Sensibility}

\lettrine{L}{ike} a balding hipster\cite{hipster} who imitates a young trendsetter's style, Facebook
is updating itself to look a lot more like Twitter.

Unlike Facebook, where friends mutually agree to let one another into their online lives, Twitter
lets people share updates and links with anyone who cares to read them.

That has turned Twitter into a tool for people to peer into the collective mind and see what people
are talking about in real time. It is also a tool for businesses to reach customers and monitor what
their customers are saying about them.

Facebook seems to be very interested in those features. Since last fall, when Facebook tried and
failed to acquire Twitter, it has been slowly introducing features that mimic\cite{mimic} Twitter.

Last week, Facebook added two new, Twitter-like features. Users can now ``tag'' friends or companies
that they mention in status updates, and they can use a pared-down version of the site called
Facebook Lite that resembles Twitter's stream of status updates.

Meredith Chin, a Facebook spokeswoman, played down the changes. ``We've been making iterations to
our product over time to reflect the rapid evolution of how people share information online,'' she
said.

But others see another force at work. ``Twitter envy: Facebook has it, absolutely,'' said Jeremiah
Owyang, a social media consultant at the Altimeter Group, which advises businesses on using new
technologies. ``Facebook absolutely recognizes that Twitter is a threat, and they're doing what they
can to replicate the features before Twitter gets mainstream adoption.''

To tag another Facebook member in a status update, users type the $@$ symbol before the friend's
name. The $@$ symbol is a convention that Twitter users started. In response, Twitter added a
section on its site where people can see any tweets that mention them. The mentions are hyperlinked
so others could click on them to see the subject's profile page.

Andrew Huang, a product manager at Facebook, said it is ``a common Internet mechanism,'' and he
expects Facebook members to use it more for storytelling than Twitter users do.

Facebook has long allowed people to tag friends in pictures, but until now, not in status updates.
When people are tagged, they get notified by e-mail, the update appears on their profile pages and
their names are hyperlinked to their pages.

Mr.~Huang, who developed the new tagging feature, said it would enable users ``to talk about their
real-world connections'' and ``interact with each other more.'' It will also enable people and
businesses to monitor what others are saying about them on the site, which was previously much
harder to do. That has been one of Twitter's vital selling points to businesses.

Adoption by businesses is a revenue-generating opportunity for both companies. Twitter, which does
not yet have any significant revenue, has said that it will soon introduce features that help
businesses interact with customers. Facebook offers businesses special pages and the option to buy
ads to show to users who like similar companies.

Luna Park, a chain of three restaurants in San Francisco and Los Angeles, uses both Facebook and
Twitter to send out promotions. Chuck Meyer, Luna Park's general manager in San Francisco, said
Facebook is more useful because the restaurant can post photos and longer updates. But he said
customers use Twitter more because they think of it as a place to follow businesses and Facebook as
a place to chat with friends. Luna Park has about 1,580 fans on Facebook and 2,350 followers on
Twitter.

Mr.~Meyer is pleased that Facebook added a tagging feature similar to Twitter's because when people
mention Luna Park, their friends can go to Luna Park's profile page with a single click and Luna
Park will get an alert. The new feature will also lure people to the site with e-mail notifications
that they have been tagged.

``A lot of companies are envious of Twitter because people spend a lot more time there, and this
allows Facebook to do the same thing -- it gives them another opportunity to get people to come back
to the site,'' said Jason Keath, a social media consultant in Charlotte, N.C.

The second new feature, Facebook Lite, is meant for people with very slow Internet connections or
new users who want an introduction to the core features of the site, Ms.~Chin said. But it might
also appeal to veteran Facebook users who like the simplicity of Twitter.

Facebook Lite is essentially a stream of updates, like Twitter. It includes photos and comments,
which are not available on Twitter, but disposes of other distracting sections that clutter the
traditional Facebook homepage.

Facebook has made other Twitter-like changes. In March, it updated users' homepages to show the full
stream of updates from all friends in real time, instead of just the updates selected by an
algorithm. And last month, Facebook allowed brands and celebrities to send status updates directly
to Twitter without visiting Twitter's site. Twitter users can send tweets to Facebook.

Twitter says it is happy to share with Facebook. ``Twitter continues to reduce friction between many
services,'' said Biz Stone, a founder. ``Our services are complementary to mobile networks, social
networks, search engines, software platforms, television networks and maybe a few other areas we
haven't thought of yet.''

No matter how many features they share, it is unlikely that Facebook will make Twitter unnecessary
for its users -- or vice versa, Mr.~Keath said. ``I don't think that divide is going to close soon.
There's going to be certain aspects where Facebook can compete or maybe take over Twitter, but over
all, they are safe in their niches\cite{niche}.''

\section{Miley Cyrus Is on the Line Just for You}

\lettrine{T}{hat} hoariest of social networking devices -- the telephone -- is making a comeback
among Hollywood marketers.

Facebook, MySpace and Twitter are now considered essential parts of the entertainment factory's
marketing arsenal. Add another service to the list: SayNow, a tiny Silicon Valley company whose
low-key approach -- connecting stars and their fans through voice mail -- is gaining traction,
particularly among teenage audiences.

SayNow has been quietly amassing huge numbers of callers and winning celebrity converts. The free
service says it averages 10 million fan calls a month, or about one million voice minutes each day.
Use has doubled over the last year, according to Nikhyl Singhal, a company founder. About 65 percent
of callers are teenage girls.

``We effectively create a mobile army, allowing artists to directly motivate and direct their
fans,'' said Mr.~Singhal, who also plans to court athletes and politicians as participants.

SayNow, founded in 2005 with \$6 million in venture capital, is essentially a telephone company.
Stars are given a personal phone number to post on their Web sites or mention in television
interviews. Fans call and listen to recorded messages.

Callers can also leave voice messages of their own, many of which are then made public on
SayNow.com. Callers can sign up to get text alerts when stars like Selena Gomez, of ``Wizards of
Waverly Place,'' or the Jonas Brothers (22 million SayNow calls from fans and counting) record a new
message.

A newly introduced ``live'' function allows celebrities to hold what are essentially giant
conference calls, chattering away as thousands of fans listen on mute.

Artists can also return calls for one-on-one conversations that are broadcast to the group,
something that the pop singer Jesse McCartney recently did. ``When I come down to Texas, I expect
you in the front row,'' he recently told a weeping Catherine from Fort Worth. Her response: ``I will
definitely be there!''

The marketing power of SayNow extends beyond verbal pitches. Studios like Summit Entertainment,
which used the service to help promote the recent movie ``Bandslam,'' can see where fan interest is
hot and where it lags by looking at which area codes are responding.

SayNow also connects callers to ticket sellers -- ``Did you know Miley Cyrus has a new movie coming
out? Press one to find a theater near you'' -- and sells ring tones. The company makes money by
selling 15-second ads that are played after messages. It shares a small portion of revenue with the
stars or studios that use the service.

Hollywood has long tried to use the phone as a way to deepen the connection between consumers and
entertainment. Will Smith, back when he was known as the Fresh Prince in the 1980s, had a 1-900
number that fans could call to hear him rap (with D.~J.~Jazzy Jeff). Dora the Explorer has been
calling children at home (parents sign up online) for the better part of a decade.

But studios and stars are turning more heavily to the phone as a marketing tool as new services like
SayNow pop up; cellphones become more ubiquitous; and evidence grows that their most crucial
customers -- teenagers -- are not all that enamored\cite{enamor} with the likes of Twitter. A voice
message is more intimate and authentic than a status update.

A phone number, for instance, figured heavily into the success of ``District 9,'' the
science-fiction movie that became a surprise hit this summer. Billboards and bus ads encouraged
people to call a toll-free number to report nonhuman activity, resulting in thousands of messages.

Teenagers typically look at marketing by text message as intrusive, but the simplicity of SayNow
could be ``incredibly powerful,'' said Matt Britton, the managing partner of Mr.~Youth, a New York
consultancy.

``This has an authentic, unfettered access feel,'' Mr.~Britton said. ``Right now, it's all about
removing the third-party barriers -- the studio, the publicist -- that come between consumers and
the star. Teens want to feel closely connected.''

SayNow, of course, has its challenges. The biggest is cooperation from celebrities. ``If I'm a
celebrity and I blog or I tweet, nobody believes it's really me -- you can't fake a voice,''
Mr.~Singhal said. ``And that's a problem for us: we need Miley Cyrus herself.''

The company recently hired full-time ``artist liaisons'' in New York and Los Angeles to persuade
more agents and publicists that their clients should participate.

Because the service is aimed at children, there are also risks involving privacy and protection.
SayNow says it never shares phone numbers and its system keeps the identity of every caller private.
A new SayNow feature, however, allows youngsters to connect with strangers to talk about similar
interests, something that unsavory adults could abuse.

Mr.~Singhal said recording made that feature safe. ``Predators don't want a voice record of their
activities,'' he said.

Katie Hobson, 11, does not have a cellphone, but SayNow came to her attention recently after she
read in J-14 Magazine that some of her favorite celebrities used it. ``I'd personally love to get a
call from Taylor Lautner,'' said Katie, a sixth grader in Los Angeles.

Mr.~Lautner, an actor who plays a heartthrob werewolf in ``New Moon,'' the coming ``Twilight''
sequel, does not have an account. But Tiffany Thornton, who plays the diva Tawni Hart on Disney
Channel's ``Sonny With a Chance,'' joined SayNow two weeks ago. Within hours, she had dozens of
messages.

``I have Twitter to keep in touch with my fans, but it's not very personal,'' she said. So far,
Ms.~Thornton has used the service just to say hello (``I hope you're having a blessed day'') and to
promote an appearance at the New York State Fair.

``Some of the kids cry in the messages,'' she said. ``It's so touching and amazing.''

\section{Public Option Fades From Debate Over Health Care}

\lettrine{I}{t} was just one line in a campaign manifesto\cite{manifesto}, and it hardly seemed the
most significant or contentious. As a presidential candidate, Barack Obama said he would ``establish
a new public insurance program'' alongside private health care plans.

That proposal took on a life of its own, but it now appears to be dying, a victim of an ineffectual
White House strategy, the president's failure to argue passionately for the ``public option'' and
all-out opposition by the insurance industry and much of the health care industry.

In the campaign, Mr.~Obama said the public plan would compete with private insurers on the price and
quality of care, thus benefiting consumers. What Mr.~Obama did not foresee is that, to some people
on the right and the left, it would become the most important issue in the debate over health care,
touching off a battle over the role of government in one of the nation's biggest, fastest-growing
industries.

Once in office Mr.~Obama and his advisers have sent conflicting signals about how critical a
government-run health plan would be. He prefers a public plan but is open to other ideas.

Dancing around the issue for eight months, Mr.~Obama has seemed, at various times,
pragmatic\cite{pragmatic}, flexible or indecisive.

``I just want to figure out what works,'' Mr.~Obama said in March at a White House forum. If he
could drive down health costs and expand coverage ``entirely through the market,'' he said, ``I'd be
happy to do it that way.'' And ``if there was a way of doing it that involved more government
regulation and involvement, I'm happy to do it that way, as well,'' he added.

Champions of the public plan said it could save money by using Medicare rates and fee schedules to
pay hospitals and doctors. In a book last year, one of Mr.~Obama's top advisers, former Senator Tom
Daschle, said consumers should have the option of enrolling in ``a government-run insurance program
modeled after Medicare, a proven and popular program.''

That is exactly what worries health care providers, who say Medicare pays them less than market
rates paid by private insurers. And they have pressed their concerns on Capitol Hill with a small
army of lobbyists.

Conservatives have another concern. They see the public option as a step toward a single-payer
system in which the government would pay most of the nation's health care bill and could supplant
private insurers.

``A public plan is essentially a stalking horse for a single-payer plan,'' said Senator Judd Gregg,
Republican of New Hampshire. ``It is more than the camel's nose under the tent. It is the camel's
neck, and probably front legs, under the tent. There is no way the private sector will be able to
compete.''

In trying to answer this charge, Democrats feel torn. Mr.~Obama and many Democrats deny that they
want to drive private insurers from the market. But others embrace the ultimate goal of ``Medicare
for all.''

The American Federation of State, County and Municipal Employees, which supports Mr.~Obama and the
public option, has long supported a single-payer system.

In a memorandum to union leaders last year, Gerald W.~McEntee, president of the 1.6-million-member
federation, said a public plan would ``create a competitive check on the private market and build
both public support and the infrastructure for a single-payer system.''

Then in April, Representative Jan Schakowsky, Democrat of Illinois, said insurers were right to fear
that a public plan could ``put the private insurance industry out of business.''

That might happen because of ``the superiority of the public health care option,'' said
Ms.~Schakowsky, one of 86 co-sponsors of a bill to establish a single-payer system.

Such comments provided new ammunition to Republicans already worried about the costly commitments
undertaken by the federal government to stave off an economic collapse.

Senator Jon Kyl of Arizona, the No.~2 Republican in the Senate, voiced this sense of bailout fatigue
in June when he said, ``We have the takeover of the auto companies and banks and A.I.G. and student
loans -- and now health care.''

In battering the proposal for a public option, Republicans have made effective use of estimates by
the Lewin Group, a consulting concern, which said that more than 100 million people might sign up
for the government-run insurance plan.

By contrast, the Congressional Budget Office has estimated that 11 million to 12 million people
might enroll.

Mr.~Obama cited the lower estimate in a speech to a joint session of Congress on Wednesday, to
buttress\cite{buttress} his assertion that fears of a public plan were overblown. ``We believe that
less than 5 percent of Americans would sign up,'' Mr.~Obama said.

But a leader of the Congressional Progressive Caucus, Representative Lynn Woolsey, Democrat of
California, predicted that over time ``more and more people will select the public option.''

If only 5 percent of people will enroll, she asked, ``Why are the private insurance companies so
worried?''

The different estimates by Lewin and the Congressional Budget Office are based on different
assumptions about who would be permitted to enroll in the public plan -- workers in all companies or
just those in smaller businesses.

Many people who cite studies by the Lewin Group do not know it is a unit of Ingenix, a wholly owned
subsidiary of UnitedHealth, one of the nation's largest insurers. John F.~Sheils, a vice president
of the Lewin Group, said the parent company had no influence over its research.

Momentum for the public option has waned, in part, because senators have been focusing on an
alternative: nonprofit member-owned insurance cooperatives.

Apart from the question of whether co-ops would be workable or effective, they provide a politically
convenient middle ground for centrists. With no immediate prospect of getting the votes for a public
option in the Senate, some liberals have said they too are willing to consider the idea -- if it
enables them to pass a bill, and if the co-ops are strong enough to put competitive pressure on
insurance companies.

Senator Kent Conrad, Democrat of North Dakota, who floated the idea in early June, said co-ops would
accomplish ``much of what those who want a public option are calling for -- something to compete
with private for-profit insurance companies.''

At the same time, Mr.~Conrad said, co-ops address Republican concerns because they are not
controlled by the government.

Liberal Democrats are not giving up. Senator Tom Harkin, Democrat of Iowa, said the president and
Senate Democratic leaders had not made a serious effort to round up votes for a public option. If
they did, he said, it could pass.

While the White House has struggled to define its position, insurance companies have never wavered.
Starting two weeks after the 2008 election, they have said they would accept greater federal
regulation of their market practices if Congress also required everyone to have health insurance.

These may have been tactical concessions, to abate public wrath, but they were well received in
Congress. While making these offers, the industry conserved its resources for the bigger battle over
a public option.

\section{Clean Water Laws Are Neglected, at a Cost in Suffering}

\lettrine{J}{ennifer} Hall-Massey knows not to drink the tap water in her home near Charleston,
W.Va.

In fact, her entire family tries to avoid any contact with the water. Her youngest son has scabs on
his arms, legs and chest where the bathwater -- polluted with lead, nickel and other heavy metals --
caused painful rashes. Many of his brother's teeth were capped to replace enamel that was eaten
away.

Neighbors apply special lotions\cite{lotion} after showering because their skin burns. Tests show
that their tap water contains arsenic\cite{arsenic}, barium\cite{barium}, lead,
manganese\cite{manganese} and other chemicals at concentrations federal regulators say could
contribute to cancer and damage the kidneys and nervous system.

``How can we get digital cable and Internet in our homes, but not clean water?'' said
Mrs.~Hall-Massey, a senior accountant at one of the state's largest banks.

She and her husband, Charles, do not live in some remote corner of Appalachia. Charleston, the state
capital, is less than 17 miles from her home.

``How is this still happening today?'' she asked.

When Mrs.~Hall-Massey and 264 neighbors sued nine nearby coal companies, accusing them of putting
dangerous waste into local water supplies, their lawyer did not have to look far for evidence. As
required by state law, some of the companies had disclosed in reports to regulators that they were
pumping into the ground illegal concentrations of chemicals -- the same pollutants that flowed from
residents' taps.

But state regulators never fined or punished those companies for breaking those pollution laws.

This pattern is not limited to West Virginia. Almost four decades ago, Congress passed the Clean
Water Act to force polluters to disclose the toxins they dump into waterways and to give regulators
the power to fine or jail offenders. States have passed pollution statutes of their own. But in
recent years, violations of the Clean Water Act have risen steadily across the nation, an extensive
review of water pollution records by The New York Times found.

In the last five years alone, chemical factories, manufacturing plants and other workplaces have
violated water pollution laws more than half a million times. The violations range from failing to
report emissions to dumping toxins at concentrations regulators say might contribute to cancer,
birth defects and other illnesses.

However, the vast majority of those polluters have escaped punishment. State officials have
repeatedly ignored obvious illegal dumping, and the Environmental Protection Agency, which can
prosecute polluters when states fail to act, has often declined to intervene\cite{intervene}.

Because it is difficult to determine what causes diseases like cancer, it is impossible to know how
many illnesses are the result of water pollution, or contaminants' role in the health problems of
specific individuals.

But concerns over these toxins are great enough that Congress and the E.P.A. regulate more than 100
pollutants through the Clean Water Act and strictly limit 91 chemicals or
contaminants\cite{contaminant} in tap water through the Safe Drinking Water Act.

Regulators themselves acknowledge lapses\cite{lapse}. The new E.P.A. administrator, Lisa P.~Jackson,
said in an interview that despite many successes since the Clean Water Act was passed in 1972, today
the nation's water does not meet public health goals, and enforcement of water pollution laws is
unacceptably low. She added that strengthening water protections is among her top priorities. State
regulators say they are doing their best with insufficient resources.

The Times obtained hundreds of thousands of water pollution records through Freedom of Information
Act requests to every state and the E.P.A., and compiled a national database of water pollution
violations that is more comprehensive than those maintained by states or the E.P.A. (For an
interactive version, which can show violations in any community, visit www.nytimes.com/toxicwaters.)

In addition, The Times interviewed more than 250 state and federal regulators, water-system
managers, environmental advocates and scientists.

That research shows that an estimated one in 10 Americans have been exposed to drinking water that
contains dangerous chemicals or fails to meet a federal health benchmark in other ways.

Those exposures include carcinogens\cite{carcinogen} in the tap water of major American cities and
unsafe chemicals in drinking-water wells. Wells, which are not typically regulated by the Safe
Drinking Water Act, are more likely to contain contaminants than municipal water systems.

Because most of today's water pollution has no scent or taste, many people who consume dangerous
chemicals do not realize it, even after they become sick, researchers say.

But an estimated 19.5 million Americans fall ill each year from drinking water contaminated with
parasites\cite{parasite}, bacteria or viruses, according to a study published last year in the
scientific journal Reviews of Environmental Contamination and Toxicology. That figure does not
include illnesses caused by other chemicals and toxins.

In the nation's largest dairy states, like Wisconsin and California, farmers have sprayed liquefied
animal feces onto fields, where it has seeped into wells, causing severe infections. Tap water in
parts of the Farm Belt, including cities in Illinois, Kansas, Missouri and Indiana, has contained
pesticides at concentrations that some scientists have linked to birth defects and fertility
problems.

In parts of New York, Rhode Island, Ohio, California and other states where sewer systems cannot
accommodate heavy rains, untreated human waste has flowed into rivers and washed onto beaches.
Drinking water in parts of New Jersey, New York, Arizona and Massachusetts shows some of the highest
concentrations of tetrachloroethylene, a dry cleaning solvent that has been linked to kidney damage
and cancer. (Specific types of water pollution across the United States will be examined in future
Times articles.)

The Times's research also shows that last year, 40 percent of the nation's community water systems
violated the Safe Drinking Water Act at least once, according to an analysis of E.P.A. data. Those
violations ranged from failing to maintain proper paperwork to allowing carcinogens into tap water.
More than 23 million people received drinking water from municipal systems that violated a
health-based standard.

In some cases, people got sick right away. In other situations, pollutants like chemicals, inorganic
toxins and heavy metals can accumulate in the body for years or decades before they cause problems.
Some of the most frequently detected contaminants have been linked to cancer, birth defects and
neurological disorders.

Records analyzed by The Times indicate that the Clean Water Act has been violated more than 506,000
times since 2004, by more than 23,000 companies and other facilities, according to reports submitted
by polluters themselves. Companies sometimes test what they are dumping only once a quarter, so the
actual number of days when they broke the law is often far higher. And some companies illegally
avoid reporting their emissions, say officials, so infractions go unrecorded.

Environmental groups say the number of Clean Water Act violations has increased significantly in the
last decade. Comprehensive data go back only five years but show that the number of facilities
violating the Clean Water Act grew more than 16 percent from 2004 to 2007, the most recent year with
complete data.

Polluters include small companies, like gas stations, dry cleaners, shopping malls and the Friendly
Acres Mobile Home Park in Laporte, Ind., which acknowledged to regulators that it had dumped human
waste into a nearby river for three years.

They also include large operations, like chemical factories, power plants, sewage treatment centers
and one of the biggest zinc smelters, the Horsehead Corporation of Pennsylvania, which has dumped
illegal concentrations of copper, lead, zinc, chlorine and selenium into the Ohio River. Those
chemicals can contribute to mental retardation and cancer.

Some violations are relatively minor. But about 60 percent of the polluters were deemed in
``significant noncompliance'' -- meaning their violations were the most serious kind, like dumping
cancer-causing chemicals or failing to measure or report when they pollute.

Finally, the Times's research shows that fewer than 3 percent of Clean Water Act violations resulted
in fines or other significant punishments by state officials. And the E.P.A. has often declined to
prosecute polluters or force states to strengthen their enforcement by threatening to withhold
federal money or take away powers the agency has delegated to state officials.

Neither Friendly Acres Mobile Home Park nor Horsehead, for instance, was fined for Clean Water Act
violations in the last eight years. A representative of Friendly Acres declined to comment. Indiana
officials say they are investigating the mobile home park. A representative of Horsehead said the
company had taken steps to control pollution and was negotiating with regulators to clean up its
emissions.

Numerous state and federal lawmakers said they were unaware that pollution was so widespread.

``I don't think anyone realized how bad things have become,'' said Representative James L.~Oberstar,
a Minnesota Democrat, when told of The Times's findings. Mr.~Oberstar is chairman of the House
Transportation and Infrastructure Committee, which has jurisdiction\cite{jurisdiction} over many
water-quality issues.

``The E.P.A. and states have completely dropped the ball,'' he said. ``Without oversight and
enforcement, companies will use our lakes and rivers as dumping grounds -- and that's exactly what
is apparently going on.''

The E.P.A. administrator, Ms.~Jackson, whose appointment was confirmed in January, said in an
interview that she intended to strengthen enforcement of the Clean Water Act and pressure states to
apply the law.

``I've been saying since Day One I want to work on these water issues pretty broadly across the
country,'' she said. On Friday, the E.P.A. said that it was reviewing dozens of coal-mining permits
in West Virginia and three other states to make sure they would not violate the Clean Water Act.

After E.P.A. officials received detailed questions from The New York Times in June, Ms.~Jackson sent
a memo to her enforcement deputy noting that the E.P.A. is ``falling short of this administration's
expectations for the effectiveness of our clean water enforcement programs. Data available to E.P.A.
shows that, in many parts of the country, the level of significant noncompliance with permitting
requirements is unacceptably high and the level of enforcement activity is unacceptably low.''

State officials, for their part, attribute rising pollution rates to increased workloads and
dwindling resources. In 46 states, local regulators have primary responsibility for crucial aspects
of the Clean Water Act. Though the number of regulated facilities has more than doubled in the last
10 years, many state enforcement budgets have remained essentially flat when adjusted for inflation.
In New York, for example, the number of regulated polluters has almost doubled to 19,000 in the last
decade, but the number of inspections each year has remained about the same.

But stretched resources are only part of the reason polluters escape punishment. The Times's
investigation shows that in West Virginia and other states, powerful industries have often
successfully lobbied to undermine effective regulation.

State officials also argue that water pollution statistics include minor infractions, like failing
to file reports, which do not pose risks to human health, and that records collected by The Times
failed to examine informal enforcement methods, like sending warning letters.

``We work enormously hard inspecting our coal mines, analyzing water samples, notifying companies of
violations when we detect them,'' said Randy Huffman, head of West Virginia's Department of
Environmental Protection. ``When I look at how far we've come in protecting the state's waters since
we took responsibility for the Clean Water Act, I think we have a lot to be proud of.''

But unchecked pollution remains a problem in many states. West Virginia offers a revealing example
of why so many companies escape punishment.

One Community's Plight

The mountains surrounding the home of Mrs.~Hall-Massey's family and West Virginia's nearby capital
have long been mined for coal. And for years, the area enjoyed clean well water.

But starting about a decade ago, awful smells began coming from local taps. The water was sometimes
gray, cloudy and oily. Bathtubs and washers developed rust-colored rings that scrubbing could not
remove. When Mrs.~Hall-Massey's husband installed industrial water filters, they quickly turned
black. Tests showed that their water contained toxic amounts of lead, manganese, barium and other
metals that can contribute to organ failure or developmental problems.

Around that time, nearby coal companies had begun pumping industrial waste into the ground.

Mining companies often wash their coal to remove impurities. The leftover liquid -- a black fluid
containing dissolved minerals and chemicals, known as sludge\cite{sludge} or slurry -- is often
disposed of in vast lagoons or through injection into abandoned mines. The liquid in those lagoons
and shafts can flow through cracks in the earth into water supplies. Companies must regularly send
samples of the injected liquid to labs, which provide reports that are forwarded to state
regulators.

In the eight miles surrounding Mrs.~Hall-Massey's home, coal companies have injected more than 1.9
billion gallons of coal slurry and sludge into the ground since 2004, according to a review of
thousands of state records. Millions more gallons have been dumped into lagoons.

These underground injections have contained chemicals at concentrations that pose serious health
risks, and thousands of injections have violated state regulations and the Safe Drinking Water Act,
according to reports sent to the state by companies themselves.

For instance, three coal companies -- Loadout, Remington Coal and Pine Ridge, a subsidiary of
Peabody Energy, one of the largest coal companies in the world -- reported to state officials that
93 percent of the waste they injected near this community had illegal concentrations of chemicals
including arsenic, lead, chromium, beryllium or nickel.

Sometimes those concentrations exceeded legal limits by as much as 1,000 percent. Those chemicals
have been shown to contribute to cancer, organ failures and other diseases.

But those companies were never fined or punished for those illegal injections, according to state
records. They were never even warned that their activities had been noticed.

Remington Coal declined to comment. A representative of Loadout's parent said the company had
assigned its permit to another company, which ceased injecting in 2006. Peabody Energy, which spun
off Pine Ridge in 2007, said that some data sent to regulators was inaccurate and that the company's
actions reflected best industry practices.

West Virginia officials, when asked about these violations, said regulators had accidentally
overlooked many pollution records the companies submitted until after the statute of limitations had
passed, so no action was taken. They also said their studies indicated that those injections could
not have affected drinking water in the area and that other injections also had no detectable
effect.

State officials noted that they had cited more than 4,200 water pollution violations at mine sites
around the state since 2000, as well as conducted thousands of investigations. The state has
initiated research about how mining affects water quality. After receiving questions from The Times,
officials announced a statewide moratorium\cite{moratorium} on issuing injection permits and told
some companies that regulators were investigating their injections.

``Many of the issues you are examining are several years old, and many have been addressed,'' West
Virginia officials wrote in a statement. The state's pollution program ``has had its share of
issues,'' regulators wrote. However, ``it is important to note that if the close scrutiny given to
our state had been given to others, it is likely that similar issues would have been found.''

More than 350 other companies and facilities in West Virginia have also violated the Clean Water Act
in recent years, records show. Those infractions include releasing illegal concentrations of iron,
manganese, aluminum and other chemicals into lakes and rivers.

As the water in Mrs.~Hall-Massey's community continued to worsen, residents began complaining of
increased health problems. Gall bladder diseases, fertility problems, miscarriages and kidney and
thyroid issues became common, according to interviews.

When Mrs.~Hall-Massey's family left on vacation, her sons' rashes\cite{rash} cleared up. When they
returned, the rashes reappeared. Her dentist told her that chemicals appeared to be damaging her
teeth and her son's, she said. As the quality of her water worsened, Mrs.~Hall-Massey's once-healthy
teeth needed many crowns. Her son brushed his teeth often, used a fluoride rinse twice a day and was
not allowed to eat sweets. Even so, he continued getting cavities until the family stopped using tap
water. By the time his younger brother's teeth started coming in, the family was using bottled water
to brush. He has not had dental problems.

Medical professionals in the area say residents show unusually high rates of health problems. A
survey of more than 100 residents conducted by a nurse hired by Mrs.~Hall-Massey's lawyer indicated
that as many as 30 percent of people in this area have had their gallbladders removed, and as many
as half the residents have significant tooth enamel\cite{enamel} damage, chronic stomach problems
and other illnesses. That research was confirmed through interviews with residents.

It is difficult to determine which companies, if any, are responsible for the contamination that
made its way into tap water or to conclude which specific chemicals, if any, are responsible for
particular health problems. Many coal companies say they did not pollute the area's drinking water
and chose injection sites that flowed away from nearby homes.

An independent study by a university researcher challenges some of those claims.

``I don't know what else could be polluting these wells,'' said Ben Stout, a biology professor at
Wheeling Jesuit University who tested the water in this community and elsewhere in West Virginia.
``The chemicals coming out of people's taps are identical to the chemicals the coal companies are
pumping into the ground.''

One night, Mrs.~Hall-Massey's 6-year-old son, Clay, asked to play in the tub. When he got out, his
bright red rashes hurt so much he could not fall asleep. Soon, Mrs.~Hall-Massey began complaining to
state officials. They told her they did not know why her water was bad, she recalls, but doubted
coal companies had done anything wrong. The family put their house on the market, but because of the
water, buyers were not interested.

In December, Mrs.~Hall-Massey and neighbors sued in county court, seeking compensation. That suit is
pending. To resolve a related lawsuit filed about the same time, the community today gets regular
deliveries of clean drinking water, stored in coolers or large blue barrels outside most homes.
Construction began in August on a pipeline bringing fresh water to the community.

But for now most residents still use polluted water to bathe, shower and wash dishes.

``A parent's only real job is to protect our children,'' Mrs.~Hall-Massey said. ``But where was the
government when we needed them to protect us from this stuff?''

Regulators 'Overwhelmed'

Matthew Crum, a 43-year-old lawyer, wanted to protect people like Mrs.~Hall-Massey. That is why he
joined West Virginia's environmental protection agency in 2001, when it became clear that the
state's and nation's streams and rivers were becoming more polluted.

But he said he quickly learned that good intentions could not compete with intimidating politicians
and a fearful bureaucracy.

Mr.~Crum grew up during a golden age of environmental activism. He was in elementary school when
Congress passed the Clean Water Act of 1972 in response to environmental disasters, including a fire
on the polluted Cuyahoga River in Cleveland. The act's goal was to eliminate most water pollution by
1985 and prohibit the ``discharge of toxic pollutants in toxic amounts.''

``There were a bunch of us that were raised with the example of the Clean Water Act as
inspiration,'' he said. ``I wanted to be part of that fight.''

In the two decades after the act's passage, the nation's waters grew much healthier. The Cuyahoga
River, West Virginia's Kanawha River and hundreds of other beaches, streams and ponds were
revitalized.

But in the late 1990s, some states' enforcement of pollution laws began tapering off, according to
regulators and environmentalists. Soon the E.P.A. started reporting that the nation's rivers, lakes
and estuaries were becoming dirtier again. Mr.~Crum, after a stint\cite{stint} in Washington with
the Justice Department and the birth of his first child, joined West Virginia's Department of
Environmental Protection, where new leadership was committed to revitalizing the Clean Water Act.

He said his idealism was tested within two weeks, when he was called to a huge coal spill into a
stream.

``I met our inspector at the spill site, and we had this really awkward conversation,'' Mr.~Crum
recalled. ``I said we should shut down the mine until everything was cleaned up. The inspector
agreed, but he said if he issued that order, he was scared of getting demoted or transferred to the
middle of nowhere. Everyone was terrified of doing their job.''

Mr.~Crum temporarily shut the mine.

In the next two years, he shut many polluting mines until they changed their ways. His tough
approach raised his profile around the state.

Mining companies, worried about attracting Mr.~Crum's attention, began improving their waste
disposal practices, executives from that period said. But they also began complaining to their
friends in the state's legislature, they recalled in interviews, and started a whisper campaign
accusing Mr.~Crum of vendettas\cite{vendetta} against particular companies -- though those same
executives now admit they had no evidence for those claims.

In 2003, a new director, Stephanie Timmermeyer, was nominated to run the Department of Environmental
Protection. One of West Virginia's most powerful state lawmakers, Eustace Frederick, said she would
be confirmed, but only if she agreed to fire Mr.~Crum, according to several people who said they
witnessed the conversation.

She was given the job and soon summoned\cite{summon} Mr.~Crum to her office. He was dismissed two
weeks after his second child's birth.

Ms.~Timmermeyer, who resigned in 2008, did not return calls. Mr.~Frederick died last year.

Since then, hundreds of workplaces in West Virginia have violated pollution laws without paying
fines. A half-dozen current and former employees, in interviews, said their enforcement efforts had
been undermined by bureaucratic disorganization, a departmental preference to let polluters escape
punishment if they promise to try harder, and a revolving door of regulators who leave for
higher-paying jobs at the companies they once policed.

``We are outmanned and overwhelmed, and that's exactly how industry wants us,'' said one employee
who requested anonymity for fear of being fired. ``It's been obvious for decades that we're not on
top of things, and coal companies have earned billions relying on that.''

In June, four environmental groups petitioned the E.P.A. to take over much of West Virginia's
handling of the Clean Water Act, citing a ``nearly complete breakdown'' in the state. The E.P.A. has
asked state officials to respond and said it is investigating the petition.

Similar problems exist in other states, where critics say regulators have often turned a blind eye
to polluters. Regulators in five other states, in interviews, said they had been pressured by
industry-friendly politicians to drop continuing pollution investigations.

``Unless the E.P.A. is pushing state regulators, a culture of transgression and apathy\cite{apathy}
sets in,'' said William K.~Reilly, who led the E.P.A. under President George H.~W. Bush.

In response, many state officials defend their efforts. A spokeswoman for West Virginia's Department
of Environmental Protection, for instance, said that between 2006 and 2008, the number of
cease-operation orders issued by regulators was 10 percent higher than during Mr.~Crum's two-year
tenure.

Mr.~Huffman, the department's head, said there is no political interference with current
investigations. Department officials say they continue to improve the agency's procedures, and note
that regulators have assessed \$14.7 million in state fines against more than 70 mining companies
since 2006.

However, that is about equal to the revenue those businesses' parent companies collect every 10
hours, according to financial reports. (To find out about every state's enforcement record and read
comments from regulators, visit www.nytimes.com/waterdata.)

``The real test is, is our water clean?'' said Mr.~Huffman. ``When the Clean Water Act was passed,
this river that flows through our capital was very dirty. Thirty years later, it's much cleaner
because we've chosen priorities carefully.''

Some regulators admit that polluters have fallen through the cracks. To genuinely improve
enforcement, they say, the E.P.A. needs to lead.

``If you don't have vigorous oversight by the feds, then everything just goes limp,'' said Mr.~Crum.
``Regulators can't afford to have some backbone unless they know Washington or the governor's office
will back them up.''

It took Mr.~Crum a while to recover from his firing. He moved to Virginia to work at the Nature
Conservancy, an environmental conservation group. Today, he is in private practice and works on the
occasional environmental lawsuit.

``We're moving backwards,'' he said, ``and it's heartbreaking.''

Shortcomings of the E.P.A.

The memos are marked ``DO NOT DISTRIBUTE.''

They were written this year by E.P.A. staff, the culmination\cite{culmination} of a five-year
investigation of states' enforcement of federal pollution laws. And in bland, bureaucratic terms,
they describe a regulatory system -- at the E.P.A. and among state agencies -- that in many ways
simply does not work.

For years, according to one memo, federal regulators knew that more than 30 states had major
problems documenting which companies were violating pollution laws. Another notes that states'
``personnel lack direction, ability or training'' to levy fines large enough to deter polluters.

But often, the memos say, the E.P.A. never corrected those problems even though they were widely
acknowledged. The E.P.A. ``may hesitate to push the states'' out of ``fear of risking their
relationships,'' one report reads. Another notes that E.P.A. offices lack ``a consistent national
oversight strategy.''

Some of those memos, part of an effort known as the State Review Framework, were obtained from
agency employees who asked for anonymity, and others through Freedom of Information Act requests.

Enforcement lapses\cite{lapse} were particularly bad under the administration of President George
W.~Bush, employees say. ``For the last eight years, my hands have been tied,'' said one E.P.A.
official who requested anonymity for fear of retribution. ``We were told to take our clean water and
clean air cases, put them in a box, and lock it shut. Everyone knew polluters were getting away with
murder. But these polluters are some of the biggest campaign contributors in town, so no one really
cared if they were dumping poisons into streams.''

The E.P.A. administrators during the last eight years -- Christine Todd Whitman, Michael O.~Leavitt
and Stephen L.~Johnson -- all declined to comment.

When President Obama chose Ms.~Jackson to head the E.P.A., many environmentalists and agency
employees were encouraged. During his campaign, Mr.~Obama promised to ``reinvigorate the drinking
water standards that have been weakened under the Bush administration and update them to address new
threats.'' He pledged to regulate water pollution from livestock operations and push for amendments
to the Clean Water Act.

But some worry those promises will not be kept. Water issues have taken a back seat to other
environmental concerns, like carbon emissions.

In an interview, Ms.~Jackson noted that many of the nation's waters were healthier today than when
the Clean Water Act was passed and said she intended to enforce the law more vigorously. After
receiving detailed questions from The Times, she put many of the State Review Framework documents on
the agency's Web site, and ordered more disclosure of the agency's handling of water issues,
increased enforcement and revamped technology so that facilities' environmental records are more
accessible.

``Do critics have a good and valid point when they say improvements need to be made? Absolutely,''
Ms.~Jackson said. ``But I think we need to be careful not to do that by scaring the bejesus out of
people into thinking that, boy, are things horrible. What it requires is attention, and I'm going to
give it that attention.''

In statements, E.P.A. officials noted that from 2006 to 2008, the agency conducted 11,000 Clean
Water Act and 21,000 Safe Drinking Water Act inspections, and referred 146 cases to the Department
of Justice. During the 2007 to 2008 period, officials wrote, 92 percent of the population served by
community water systems received water that had no reported health-based violations.

The Times's reporting, the statements added, ``does not distinguish between significant violations
and minor violations,'' and ``as a result, the conclusions may present an unduly alarming picture.''
They wrote that ``much of the country's water quality problems are caused by discharges from
nonpoint sources of pollution, such as agricultural runoff, which cannot be corrected solely through
enforcement.''

Ultimately, lawmakers and environmental activists say, the best solution is for Congress to hold the
E.P.A. and states accountable for their failures.

The Clean Water Act, they add, should be expanded to police other types of pollution -- like farm
and livestock runoff -- that are largely unregulated. And they say Congress should give state
agencies more resources, in the same way that federal dollars helped overhaul the nation's sewage
systems in the 1970s.

Some say changes will not occur without public outrage.

``When we started regulating water pollution in the 1970s, there was a huge public outcry because
you could see raw sewage flowing into the rivers,'' said William D.~Ruckelshaus, who served as the
first head of the Environmental Protection Agency under President Richard M.~Nixon, and then again
under President Ronald Reagan.

``Today the violations are much more subtle -- pesticides and chemicals you can't see or smell that
are even more dangerous,'' he added. ``And so a lot of the public pressure on regulatory agencies
has ebbed away.''

\section{Wall Street's Math Wizards Forgot a Few Variables}

\lettrine{I}{N} the aftermath of the great meltdown of 2008, Wall Street's quants have been cast as
the financial engineers of profit-driven innovation run amok. They, after all, invented the exotic
securities that proved so troublesome.

But the real failure, according to finance experts and economists, was in the quants' mathematical
models of risk that suggested the arcane stuff was safe.

The risk models proved myopic, they say, because they were too simple-minded. They focused mainly on
figures like the expected returns and the default risk of financial instruments. What they didn't
sufficiently take into account was human behavior, specifically the potential for widespread panic.
When lots of investors got too scared to buy or sell, markets seized up and the models failed.

That failure suggests new frontiers for financial engineering and risk management, including trying
to model the mechanics of panic and the patterns of human behavior.

``What wasn't recognized was the importance of a different species of risk -- liquidity risk,'' said
Stephen Figlewski, a professor of finance at the Leonard N.~Stern School of Business at New York
University. ``When trust in counterparties is lost, and markets freeze up so there are no prices,''
he said, it ``really showed how different the real world was from our models.''

In the future, experts say, models need to be opened up to accommodate more variables and more
dimensions of uncertainty.

The drive to measure, model and perhaps even predict waves of group behavior is an emerging field of
research that can be applied in fields well beyond finance.

Much of the early work has been done tracking online behavior. The Web provides researchers with
vast data sets for tracking the spread of all manner of things -- news stories, ideas, videos,
music, slang and popular fads -- through social networks. That research has potential applications
in politics, public health, online advertising and Internet commerce. And it is being done by
academics and researchers at Google, Microsoft, Yahoo and Facebook.

Financial markets, like online communities, are social networks. Researchers are looking at whether
the mechanisms and models being developed to explore collective behavior on the Web can be applied
to financial markets. A team of six economists, finance experts and computer scientists at Cornell
was recently awarded a grant from the National Science Foundation to pursue that goal.

``The hope is to take this understanding of contagion and use it as a perspective on how rapid
changes of behavior can spread through complex networks at work in financial markets,'' explained
Jon M.~Kleinberg, a computer scientist and social network researcher at Cornell.

At the Massachusetts Institute of Technology, Andrew W.~Lo, director of the Laboratory for Financial
Engineering, is taking a different approach to incorporating human behavior into finance. His
research focuses on applying insights from disciplines, including evolutionary biology and cognitive
neuroscience, to create a new perspective on how financial markets work, which Mr.~Lo calls ``the
adaptive-markets hypothesis.'' It is a departure from the ``efficient-market'' theory, which asserts
that financial markets always get asset prices right given the available information and that people
always behave rationally.

Efficient-market theory, of course, has dominated finance and econometric modeling for decades,
though it is being sharply questioned in the wake of the financial crisis. ``It is not that
efficient market theory is wrong, but it's a very incomplete model,'' Mr.~Lo said.

Mr.~Lo is confident that his adaptive-markets approach can help model and quantify liquidity crises
in a way traditional models, with their narrow focus on expected returns and volatility, cannot.
``We're going to see three-dimensional financial modeling and eventually N-dimensional modeling,''
he said.

J.~Doyne Farmer, a former physicist at Los Alamos National Laboratory and a founder of a
quantitative trading firm, finds the behavioral research intriguing but awfully ambitious,
especially to build into usable models. Instead, Mr.~Farmer, a professor at the interdisciplinary
Sante Fe Institute, is doing research on models of markets, institutions and their complex
interactions, applying a hybrid discipline called econophysics.

To explain, Mr.~Farmer points to the huge buildup of the credit-default-swap market, to a peak of
\$60 trillion. And in 2006, the average leverage on mortgage securities increased to 16 to 1 (it is
now 1.5 to 1). Put the two together, he said, and you have a serious problem.

``You don't need a model of human psychology to see that there was a danger of impending disaster,''
Mr.~Farmer observed. ``But economists have failed to make models that accurately model such
phenomena and adequately address their couplings.''

When a bridge over a river collapses, the engineers who built the bridge have to take
responsibility. But typically, critics call for improvement and smarter, better-trained engineers --
not fewer of them. The same pattern seems to apply to financial engineers. At M.I.T., the Sloan
School of Management is starting a one-year master's in finance this fall because the field has
become too complex to be adequately covered as part of a traditional M.B.A. program, and because of
student demand. The new finance program, Mr.~Lo noted, had 179 applicants for 25 places.

In the aftermath of the economic crisis, financial engineers, experts say, will probably shift more
to risk management and econometric analysis and concentrate less on devising exotic new instruments.
Still, the recent efforts by investment banks to create a trading market for ``life settlements,''
life insurance policies that the ill or elderly sell for cash, suggest that inventive sales people
are browsing for new asset classes to securitize, bundle and trade.

``Good or bad, moral or immoral, people are going to make markets and trade via computers, and this
is a natural area of financial engineers,'' says Emanuel Derman, a professor at Columbia University
and a former Wall Street quant.

\section{Turning to Windmills, but Resistance Lingers}

\lettrine{W}{endie} Howland grows her own food and heats her water with rooftop solar panels. She
drives a Prius with a bumper sticker that boasts ``One Less S.U.V.''

But when Mrs.~Howland tried to take the next step in green living -- installing a 132-foot windmill
in her backyard that would generate enough electricity to power her home -- she hit a wall. The
planning board in this pastoral Cape Cod town twice rejected the project citing safety concerns and
predicting ``an adverse effect on the character of the neighborhood.''

Mrs.~Howland's defeat was sealed by a Superior Court ruling in July that backed the planning board's
decision, underscoring the steep odds that residential windmill plans face nationwide. After
investing some \$40,000 in a 10-kilowatt turbine and legal fees, Mrs.~Howland and her husband,
Francis, are giving up their two-year fight.

``It's ludicrous,'' said Mrs.~Howland, 58, a health care consultant. ``We were trying to make our
bills smaller as we got older, in a clean and responsible fashion, and it boggles my mind that
ordinary people like us aren't allowed to do that.''

The decision is likely to be scrutinized by towns across the region and even the nation as they
grapple with how to regulate windmills on residential property. In wind-rich regions, clashes like
Mrs.~Howland's are increasingly common as conservation-minded people seek to install small wind
turbines on their property.

Battles over the height and noise level of residential windmills, and even over the shadows cast by
their blades, are springing up from Maine to California, even as the Obama administration promotes
renewable energy and the federal stimulus package provides 30 percent tax credits for homeowners who
install wind turbines.

Many towns still enforce old laws that prohibit anything taller than 30 feet or 40 feet on
residential land -- a height too low for sufficient wind power generation, experts say. Wind
turbines need to be at least 30 feet higher than anything within 500 feet, including trees, which
often means a tower of 80 feet or more. The Howlands' windmill would have been more than three times
the height of an average utility pole, to ensure that the surrounding white pines did not interfere.

``Everyone recoils at that,'' said Jonathan D.~Fitch, the Howlands' lawyer. ``It reminds me of the
litigation involving cell towers in the beginning -- a lot of neighborhood hostility back then, but
today you hardly notice them.''

Even before the Howland case, Cape Cod was a battleground in the windmill debate: a proposal to
build the nation's first offshore wind farm has met with stiff opposition from those who say the
project would spoil the natural beauty of Nantucket Sound.

While residential turbines remain a tiny fraction of the wind energy market, they are popping up
often enough for many communities, especially in New England, the Midwest and the West, to start
regulating them. Nearly 2,700 wind units with capacities of 10 kilowatts or less, the size used for
residences, were sold nationwide last year, up from 1,167 in 2007, according to the American Wind
Energy Association, a trade group.

But challenges persist even in communities that have passed ordinances on windmills, like Bourne,
where residents overwhelmingly approved a bylaw regulating windmills at a town meeting in 2007. The
wind energy association estimates that one-third of small wind projects are thwarted by vague or
overly strict local laws, or by outdated zoning rules that preclude them.

``It's mostly that they will make the rules so prohibitive that they essentially ban
installations,'' said Ron Stimmel, the association's small wind advocate. ``There are humongous
hurdles that no other type of installation has been faced with.''

In Waukee, Iowa, the City Council is considering a law that would ``allow and encourage'' the use of
small wind-energy systems but would ban them on single-family homes. In Islip, N.Y., a new provision
allows residents to install windmills up to 45 feet tall -- probably not high enough to make the
investment worthwhile, Mr.~Stimmel said.

And in Bourne, a town of about 19,000, the 2007 bylaw limits residential windmills to 75 feet high
unless the homeowner can prove that a taller windmill is needed and would not ``increase any adverse
impacts.''

In July, during the Howlands' trial, Brian Wall, a lawyer for Bourne, told Judge Christopher Muse of
Barnstable County Superior Court that the Howlands had not provided enough evidence that their
windmill would be quiet and safe.

``If this 1,600-pound thing falls or something falls off it in the area around it,'' he said, citing
the possibility of ice flying off the blades in winter, ``it needs to be safe.''

Mr.~Wall said the town needed an engineer's guarantee that in heavy winds the windmill would not
crash onto neighboring homes -- the closest is 105 feet from its proposed base -- even though its
pole is designed to fold in such weather, not break apart.

He added that no qualified engineer had certified that the windmill would be quiet or that its
``shadow flicker'' -- the strobe of sunlight passing through rotating blades -- would not bother the
Howlands' neighbors.

In his ruling, Judge Muse wrote that while he did not believe the windmill would pose a safety
threat, the planning board's decision fell within its discretion and was not ``unreasonable,
whimsical, capricious or arbitrary.''

Megan Amsel of the Cape and Islands Self Reliance Corporation, a nonprofit group that promotes
alternative energy in Cape Cod and southeastern Massachusetts, said she had seen some disastrous
wind projects -- not because they proved unsafe, but because they did not generate enough power.

``It's really hard to find a good installer,'' she said, adding that there are no certification
requirements. ``I've seen some real disasters, and it can give this emerging industry a black eye.''

The wind needs to blow at least 12 miles an hour for a turbine to generate electricity -- a
requirement that rules out many sites -- and the initial cost is steep. Mr.~Stimmel said that the
average cost of buying and installing a residential turbine was \$30,000, and that it took 6 years
to 30 years to recover that cost through energy savings.

The total cost for the Howlands' turbine, including installation, would be \$72,000, they said.

To Mrs.~Howland, the town's claim that a windmill would hurt the character of her neighborhood was
especially galling. None of her closest neighbors objected, she said.

``There is no aesthetic clause in the bylaw because aesthetics are subjective,'' she said. ``For
every person I know who says, 'I don't like it,' I hear another say, 'I think they're beautiful.' ''


\section{Foreign Languages Fall as Schools Look for Cuts}

\lettrine{I}{N} Edgemont, a high-performing Westchester school district, children as young as 7
could recite colors and days of the week in Spanish, but few if any learned to really converse, read
or write. So this fall, the district canceled the Spanish lessons offered twice weekly at its two
elementary schools since 2003, deciding the time and resources -- an estimated \$175,000 a year --
could be better spent on other subjects.

Class consolidation in Yonkers resulted in the loss of four foreign-language teaching positions, and
budget cuts have cost Arlington, N.Y., its seventh-grade German program, and Danbury, Conn., several
sections of middle school French and Spanish.

And in New Jersey, the Ridgewood district is replacing its three elementary school Spanish teachers
with Rosetta Stone, an interactive computer program that cost \$70,000, less than half their
combined salaries.

``There's never a replacement for a teacher in the classroom,'' said Debra Anderson, a Ridgewood
spokeswoman. ``But this was a good solution in view of the financial constraints.''

After years of expanding language offerings, suburban districts across the New York region are now
cutting back on staff and instructional time, phasing out less popular languages, and rethinking
whether they can really afford to introduce foreign tongues to their youngest students while under
constant pressure to downsize budgets and raise achievement in English and other core subjects.

But such cuts have dismayed and frustrated some educators and parents, who say that children need
more, not fewer, foreign language skills to compete in a global marketplace.

``In many cultures, a lot of business does not get done around the business table, it gets done in
side conversations and social situations,'' said Marty Abbott, director of education for the
American Council on the Teaching of Foreign Languages and a former high school Spanish and Latin
teacher. ``If you can't participate in those discussions, you get left out.''

Foreign languages play an increasingly prominent role in urban schools that serve diverse ethnic
communities. For instance, the New York City schools offered courses in Haitian-Creole, Vietnamese,
and Portuguese last year, and opened the first public school dedicated to Arabic language and
culture in Brooklyn in 2007. Last week, the city's first Hebrew-language charter school opened, also
in Brooklyn.

Advocates for foreign-language instruction would like to see the lessons integrated into the core
curriculum rather than treated as electives easy to ax at budget time. They also say that
instruction should begin as early as possible -- ideally in preschool -- because academic research
shows that younger children are more accepting of other cultures and better able to master the
pronunciation and intonation of foreign words. Some even contend that learning a foreign language
can foster\cite{foster} cognitive\cite{cognitive} skills that lead to higher standardized test
scores in other subjects.

On Long Island, more than 200 residents of Long Beach signed petitions over the summer opposing the
district's decision to phase out a dual-language English-Spanish program at Lido Elementary School
that had served as a model for other districts.

``I think it's a terrible shame,'' said Sebastian Arengo, a software engineer, whose 6-year-old twin
daughters are in the program. ``It's at the right age for kids to speak both languages, and it's
also a great way to bring the English-speaking and Spanish-speaking communities together here in
Long Beach.''

Robert Greenberg, superintendent\cite{superintendent} of the 4,000-student Long Beach district, said
that the program had been created mainly to help Spanish-speaking students learn English through
bilingual classes, but that it has evolved into an immersion\cite{immersion} program for those who
want to learn Spanish. ``I have Latino families wanting me to teach their children Spanish, but
that's not the intent of the program,'' he said.

Separate from the dual-language program, the district has provided 90 minutes a week of Spanish
instruction to all kindergarten and first-grade students since 2007, and plans to expand that
program by one grade every year. ``We made an instructional decision that we're teaching all
children Spanish rather than a few,'' he said.

Many superintendents say they remain committed to teaching languages, but simply cannot afford to do
more at this time. In Rockland County, the 9,400-student Clarkstown district spent about \$60,000
last year to hire a full-time Spanish teacher for one of its 10 elementary schools but postponed
plans to do the same at the other schools this fall ``until we determine the economy is getting
better,'' said Meg Keller-Cogan, the superintendent.

In Connecticut, the New Hartford district cut its one foreign language teacher at Ann Antolini
Elementary School from full time to three days a week to save \$35,000. Fewer hours mean that
Spanish will no longer be taught to third and fourth graders. ``It was just for budget reasons and
it was a very painful decision,'' said Philip O'Reilly, the superintendent, adding that other staff
members had hours reduced -- and two were laid off -- to cut costs.

Some educators said they were re-evaluating foreign-language programs not just because of finances
but to update them and incorporate new technology. The Ridgewood district, which started
twice-weekly Spanish lessons at its elementary schools in 2005, said its interactive software, made
by Rosetta Stone, allows students to learn at their own pace.

The 10,400-student Arlington district decided to phase out German -- leaving Spanish, French and
Italian -- because it was the least popular choice among students; last year, 44 seventh graders
enrolled in introductory German compared to more than 300 in Spanish. The district also phased out
Russian more than a decade ago.

``It was a low-fill, high-cost area of instruction, and if that wasn't taken, something else would
have been,'' said Frank V.~Pepe Jr., the superintendent.

The district will continue to require every student to study a foreign language in seventh and
eighth grade; nearly a quarter of all students study a language through their junior or senior year,
according to district officials.

``I'm not pleased we eliminated German,'' Mr.~Pepe said. ``I'm not pleased at all.''

\section{China Watched for Sign of New Leader}

\lettrine{C}{hina}'s governing Communist Party will convene its annual policy meeting on Tuesday
with a sober\cite{sober}, if not soporific\cite{soporific}, mandate\cite{mandate} to root out
government corruption and make the party adapt to changing times.

But lurking\cite{lurking} in the background is a more compelling topic: Who will become China's next
ruler in 2012?

Analysts will watch the meeting, the annual plenary session of the party's 17th Central Committee,
to see whether Vice President Xi Jinping is given the additional title of vice chairman of the
Central Military Commission.

Such an appointment would be seen as a confirmation that Mr.~Xi, 56, is set to succeed President Hu
Jintao when Mr.~Hu's second term ends in 2012. Any Chinese leader must have experience in leading
the military, which is under party control. Mr.~Hu was awarded the same post in 1999, three years
before he became the party's general secretary in 2002.

Yet Chinese politics are so opaque\cite{opaque} that no outsider can say for certain that Mr.~Xi,
the presumed heir, will win the position -- or that there will be a mark against him should he not.

``There is no foregone conclusion these days,'' said a political analyst at a Beijing institution
tied to the Communist Party.

Whether that is true is a central question hanging over the meeting this week. Since the founding of
the People's Republic 60 years ago, the Communist Party has governed both the Chinese people and
itself strictly from the top down, with all important actions approved by a handful of party leaders
united by power and personal relationships.

Officially, at least, the 2,000 or more Central Committee members meeting this week have been given
an agenda to shake up that model. The members are supposed to prepare plans to bring democracy to
the party's inner deliberations, choosing new leaders by consensus, not by the dictates of those at
the top.

``A new crop of leaders who grew up after the reform and opening up started are going to step into
new leadership roles'' in 2012, Zhen Xiaoying, a professor at the Communist Party's central party
school, stated in a recent article in the state-run newspaper People's Daily. He was referring to
the period of economic reform that began in 1978.

``The era of relying on authority and personal charm to run the party is over,'' he said.

Mr.~Xi and Mr.~Hu epitomize that shift. Mr.~Hu, 66, joined the party in 1964, two years before Mao's
Cultural Revolution brought China a decade of social and political chaos. Mr.~Xi joined in 1974, two
years after President Richard M.~Nixon first visited Beijing and China began to reconnect to the
outside world.

Mr.~Hu was the party's designated successor to Jiang Zemin, who ruled a battened-down China after
the bloody suppression of the 1989 Tiananmen Square protests.

The process of political succession in China's one-party system is always shrouded in intrigue. The
party elite elevated Mr.~Xi to the ruling Politburo Standing Committee in 2007 and gave him the
highest rank of any leader of his age group, signaling that he had been chosen to succeed Mr.~Hu
when the latter's second five-year stint as top leader ends in 2012. But the party's internal
deliberations on such matters are in the highest order of state secret, and there has been no public
confirmation of Mr.~Xi's status.

Whatever changes the plenum\cite{plenum} orders are unlikely to resemble democracy as Westerners
know it. China has long shunned Western democracy, branding it anarchy, and embraced what it calls
``democratic centralism'' -- essentially, passing carefully reviewed suggestions from lower-level
party organs to leaders at the top.

Xinhua, the official Chinese news agency, recently quoted Mr.~Hu as saying that democratic
centralism would remain China's guiding version of democracy. One liberal political analyst who has
called for a more open Chinese society, Liu Junning, argued in a telephone interview that prospects
for genuine changes this week were dim.

``I think it is important in China first to strengthen formal institutions such as the legislature
and the court system, rather than informal structures such as the ruling party,'' he said. ``Let's
see if there are any open factions within the party -- any open opposition, any open minorities.''

China's governing elite, like any group, has factions, but they are tightly cloaked. Mr.~Xi, for
example, is widely believed to be the favorite of Mr.~Jiang, who still has considerable sway in
retirement.

After Mao wreaked havoc with the party hierarchy by designating and then toppling multiple
successors, the party's elite clawed back the power to oversee political succession. Mr.~Hu was
effectively designated China's future top leader in 1992, leaving Mr.~Jiang, then the new No.~1
official, little choice in the matter. Likewise, Mr.~Hu's apparent favorite, Deputy Prime Minister
Li Keqiang, was not selected as his future successor, though Mr.~Li is now considered likely to be
the next prime minister.

The plenum will be closely watched for any signs that internal politicking has kept the succession
contest alive.

The analyst at the Communist Party institution, who spoke on condition of anonymity because he was
not authorized to speak to the press, said he believed that Mr.~Xi might not win the military post
this week. ``If he doesn't,'' he said, ``it would show that there's more of a balance of power. But
it would not mean that Xi lost the opportunity.''

The plenum is also scheduled to take up anticorruption measures that could include a requirement
that some party officials disclose their holdings of property or financial instruments.

\section{Google Releases News-Reading Service}

\lettrine{G}{oogle}, long seen as an enemy by many in the news industry, is making a bold attempt to
be seen as a friend with a new service it hopes will make it easier for readers to read newspaper
and magazine articles.

On Monday, the company introduced an experimental news hub called Fast Flip that allows users to
view news articles from dozens of major publishers and flip through them as quickly as they would
the pages of a magazine. Google will place ads around the news articles and share resulting revenue
with publishers.

Fast Flip, which is based on Google News, tries to address what Google considers a major problem
with news sites: they often are slow to load, and so they turn off many readers. Google, the leader
in Web search services and advertising, argues that if reading news online was closer to the
experience of scanning through physical newspapers or magazines, people would read more.

``Browsing news on the Web is much slower than it is in print,'' said Krishna Bharat, a
distinguished researcher at Google who developed Google News in 2002. ``When it is fast, people will
look at more news and more ads, and that's something that publishers want to see.''

The service is being initiated with the cooperation of about three dozen publishers, including major
news outlets like BBC News, The New York Times, The Washington Post and Newsweek; magazines like
Cosmopolitan, The Atlantic, Esquire and Good Housekeeping; and Web-only publications like
TechCrunch, Salon.com and Slate.

Some of the publishers said they viewed the experiment with caution, adding that no single solution
could address the industry's main problem: plunging advertising revenues.

``I don't look at this as the solution to the future of journalism,'' said Richard Gingras, the
chief executive of Salon Media Group, who previously worked as an adviser to Google executives on
media strategy. ``But who knows? We will learn from it.''

Fast Flip, which is available at \href{http://fastflip.googlelabs.com}{fastflip.googlelabs.com},
first appears as a collection of images of news articles that Google has culled from the sites of
its partners. The articles are displayed side by side in various horizontal rails that show them
arranged by popularity; sections like politics, business, travel or entertainment; topics like
tennis, Iran or the Beatles; and news source.

Flipping from one article to the next, or from one rail to the next, is quick. The articles, which
are images of Web pages that have been stripped of ads and other items that slow them down, load
with scant delay. Readers can zoom into a specific section, publication or article. They can often
read the majority of an article directly on Google, although if they click on it, they will be taken
to the publisher's Web site.

``We are helping people immerse themselves in the content,'' Mr.~Bharat said. ``If they don't like
something, they can move on quickly until they find something they want to read.''

Microsoft is also trying to make Web content more visually appealing with the new Visual Search
feature introduced Monday on its Bing search engine. Although not directed at news, it displays some
categories of search results using thumbnail images instead of text.

Google plans to place display ads alongside the stories and share the resulting revenue with
publishers. Mr.~Bharat declined to discuss what percentage of the revenue will be kept by Google but
said publishers would receive the majority.

Some publishers acknowledged that, if successful, Fast Flip would compete with already beleaguered
news Web sites. But they were persuaded the experiment could be useful.

``Of course there is a concern,'' said Martin A.~Nisenholtz, senior vice president for digital
operations for The New York Times Company. ``That doesn't mean you don't participate.''

Some news publishers have long complained that Google has unfairly profited at their expense by
selling ads on Google.com and Google News alongside newspaper content. Google has countered that its
services help news publishers by driving traffic to them, and that any publisher is free to block
Google from indexing its content.

``The interesting thing about this service, when compared to search, is that there is a revenue
model for us on Google,'' said Scott Havens, vice president of digital strategy and operations for
The Atlantic.

Ken Doctor, an analyst with Outsell, said that the decision by Google to begin paying publishers for
news content on its site is a significant change. ``It is a chink\cite{Chink} in Google's armor,''
Mr.~Doctor said. ``It could be a path to peace and rationalization of the relationship.''

While the experiment includes some major publishers, several of the top newspaper chains, including
the News Corporation, Gannett, McClatchy, Tribune and MediaNews, are not part of it.

Mr.~Bharat said that while Fast Flip tries to recreate some of the experience of reading news
offline, the service will also incorporate many Web features. For instance, it will rely partly on
Google's algorithms and partly on user behavior to rank articles. Those that are clicked on or
e-mailed the most often will rise in the rankings. And when users save an article they like, friends
may be automatically notified.

There have been other attempts to make reading electronic news more efficient. They include Times
Reader, an application that allows users to download The Times onto their computers, and versions of
newspapers available for e-book readers like Amazon.com's Kindle. But those efforts have gained only
limited acceptance.

Mr.~Bharat said that Google will offer a version of Fast Flip for some mobile phones, and may allow
news publishers to use Fast Flip's underlying technology directly on their sites, he said.

\section{Candidates' Rapport\cite{rapport} at Debate Is a Lullaby\cite{lullaby} for Germany}

\lettrine{T}{he} debate in Germany between Chancellor Angela Merkel and Foreign Minister
Frank-Walter Steinmeier on Sunday was billed as a duel, but it will go down in history as ``the
duet'' for the harmonious way the two candidates agreed on issues ranging from the financial crisis
to Afghanistan.

More than 14 million people got less than they bargained for when they tuned in to the only debate
before the Sept.~27 parliamentary election that will decide whether Mrs.~Merkel remains as
chancellor, a toothless tussle roundly panned in the German news media as a bore. The front page of
Bild, the country's largest newspaper and a tabloid\cite{tabloid} not known for its subtlety, rhymed
the famous Obama slogan with the words ``Yes we g\"ahn,'' which means ``yawn.''

After Mrs.~Merkel praised their common work on the financial crisis, one of the moderators made the
comment, so widely cited in Germany on Monday it seemed all but certain to go down as the unofficial
name of the debate, saying, ``But Mrs.~Merkel, that sounds now more like a duet\cite{duet} and less
like a duel.''

The problem is that Mrs.~Merkel and Mr.~Steinmeier are paired as chancellor and vice chancellor in
the current government, an alliance of Mrs.~Merkel's conservative Christian Democrats and
Mr.~Steinmeier's left-leaning Social Democrats. And it is pretty hard to attack an opponent's record
when it is the same as one's own, especially when they could well end up also governing together
after the election.

The smaller opposition parties, whose candidates were shut out of the debate, claimed that their
absence all but ruled out any pointed discussion of the issues. Claudia Roth, leader of the Green
Party, called it ``a political air guitar championship.'' Guido Westerwelle, leader of the
pro-business Free Democrats, who hope to supplant the Social Democrats as the junior party in a
government with Mrs.~Merkel, said that ``a television duel without the opposition is like a football
game without the second team.''

While there were just two candidates, the television stations that broadcast the event only whittled
down the number of moderators to four. An editorial in the newspaper Süddeutsche Zeitung was one of
several suggesting that there had been more clashes between the journalists asking the questions
than the candidates answering them.

There were a few points of contention. Mr.~Steinmeier spoke out most forcefully on the need for a
minimum wage, while Mrs.~Merkel said she would cut taxes to promote economic growth and create jobs.
Mrs.~Merkel also said she would consider extending operation of nuclear power plants, which by law
are supposed to be phased out in Germany by 2021.

Polls showed that the 90-minute debate ended in essentially a draw, though Mr.~Steinmeier, who is
foreign minister as well as vice chancellor, surpassed viewers' expectations. Mrs.~Merkel seemed
cranky.

Voter enthusiasm has been so low throughout the campaign season that Bild last week offered the
chance to win one million euros, or roughly \$1.5 million, to someone who correctly guesses the
outcome of the election and can prove that he or she actually voted.

Both parties claimed momentum coming out of the debate.

\section{A Slight Uptick and a Cool Reaction to Obama's Speech}

\lettrine{P}{resident} Obama's speech on overhauling the financial system was interrupted once by
spontaneous\cite{spontaneous} -- if milquetoast\cite{milquetoast} -- applause. That seemed to
reflect Wall Street's reaction.

Stock markets meandered\cite{meander} on Monday afternoon as the crowds ebbed\cite{ebb} and the
security and pomp accompanying Mr.~Obama's visit to Lower Manhattan were dismantled. Stocks were
slightly higher after Mr.~Obama urged bank executives to accept new federal regulations and change
the way they do business.

``We will not go back to the days of reckless behavior and unchecked excess that was at the heart of
this crisis, where too many were motivated only by the appetite for quick kills and bloated
bonuses,'' Mr.~Obama said in his speech. ``Those on Wall Street cannot resume taking risks without
regard for consequences, and expect that next time, American taxpayers will be there to break their
fall.''

Shares wavered as he spoke to the polite, if unenthusiastic, crowd of financial executives, public
officials and consumer advocates. Stocks edged up in afternoon trading.

At the close, the Dow Jones industrial average was 21.39 points, or 0.2 percent, higher at 9,626.80,
while the broader Standard \& Poor's 500-stock index was 6.61 points, or 0.63 percent, higher at
1,049.34. The Nasdaq was up 10.88 points, or 0.52 percent, at 2,091.78.

With September -- historically Wall Street's worst month -- now half over, some analysts are cheered
to see that stocks are up since the end of August.

``There's still a lot of strategists, a lot of money managers waiting for a pullback in the
market,'' said Owen Fitzpatrick, head of equity trading at Deutsche Bank Private Wealth Management.
``It will come, but it may not be as violent as we've seen in the past, just given the amount of
cash that's sitting on the sidelines waiting to come in.''

After a week containing few significant economic reports, investors this week will receive a
passel\cite{passel} of benchmarks on the economy's progress, with new figures on retail sales,
wholesale and retail prices, new-home construction and industrial production.

The day's lift suggested that investors had smoothed over concerns about a rift between the United
States and China over tires, chickens and trade barriers. The Nikkei index in Japan fell 2 percent,
and Hong Kong's Hang Seng fell 1 percent, and shares were mixed across Europe.

But in New York, shares of the tire makers Goodyear and Cooper Tire and Rubber bounded higher, days
after the Obama administration announced tariffs of as much as 35 percent on Chinese tires.

The Chinese government criticized the move and suggested it would fight back by imposing trade
restrictions on American exports of chickens and automotive products. Shares of chicken producers
including Tyson Foods and Pilgrim's Pride each fell more than 2 percent in New York.

Shares of Bank of America were mostly unchanged after a federal judge in New York threw out a
settlement between the bank and the Securities and Exchange Commission over bonuses paid to Merrill
Lynch executives as Bank of America was preparing to take over Merrill.

Other financial shares were mostly higher.

Meanwhile, investors unwound some of their recent bets in other corners of the market.

Gold prices, which soared to \$1,006 an ounce last week, fell on Monday as traders cashed in some
gains. Inflation-wary investors have been piling their money into gold recently, despite few signs
that inflation is an imminent problem.

And the dollar regained some ground, after sliding to its lowest levels since September. As prices
of gold and copper, silver and platinum have risen, the dollar has skidded lower against those
commodities and other major currencies.

And crude oil prices fell 43 cents, to \$68.86 a barrel, continuing their dance near the benchmark
of \$70.

Interest rates were higher. The Treasury's benchmark 10-year note fell 21/32, to 101 22/32 and the
yield rose to 3.42 percent from 3.35 percent late Friday.

Following are the results of Monday's Treasury auction of three- and six-month bills:

\section{New Clues to Sex Anomalies\cite{anomaly} in How Y Chromosomes\cite{chromosome} Are Copied}

\lettrine{T}{he} first words ever spoken, so fable\cite{fable} holds, were a
palindrome\cite{palindrome} and an introduction: ``Madam, I'm Adam.''

A few years ago palindromes -- phrases that read the same backward as forward -- turned out to be an
essential protective feature of Adam's Y, the male-determining chromosome that all living men have
inherited from a single individual who lived some 60,000 years ago. Each man carries a Y from his
father and an X chromosome from his mother. Women have two X chromosomes, one from each parent.

The new twist in the story is the discovery that the palindrome system has a simple weakness, one
that explains a wide range of sex anomalies from feminization to sex reversal similar to Turner's
syndrome, the condition of women who carry only one X chromosome.

The palindromes were discovered in 2003 when the Y chromosome's sequence of bases, represented by
the familiar letters G, C, T and A, was first worked out by David C.~Page of the Whitehead Institute
in Cambridge, Mass., and colleagues at the DNA sequencing center at Washington University School of
Medicine in St.~Louis.

They came as a total surprise but one that immediately explained a serious evolutionary puzzle, that
of how the genes on the Y chromosome are protected from crippling mutations.

Unlike the other chromosomes, which can repair one another because they come in pairs, one from each
parent, the Y has no evident backup system. Nature has prevented it from recombining with its
partner, the X, except at its very tips, lest its male-determining gene should sneak into the X and
cause genetic chaos.

Discovery of the palindromes explained how the Y chromosome has managed over evolutionary time to
discard bad genes: it recombines with itself. Its essential genes are embedded in a series of eight
giant palindromes, some up to three million DNA units in length. Each palindrome readily folds like
a hairpin, bringing its two arms together. The cell's DNA control machinery detects any difference
between the two arms and can convert a mutation back to the correct sequence, saving the Y's genes
from mutational decay.

After Dr.~Page discovered the palindromes, he wondered whether the system had weaknesses that might
explain the male sex chromosome anomalies that are a major object of his studies. In the current
issue of Cell, with Julian Lange and others, he describes what they call the ``Achilles' heel'' of
the Y chromosome and the wide variety of sexual disorders that it leads to.

The danger of the palindrome protection system occurs when a cell has duplicated all its chromosomes
prior to cell division, and each pair is held together at a site called the centromere. Soon, the
centromere will split, with each half and its chromosome tugged to opposite sides of the dividing
cell.

Before the split, however, a serious error can occur. Palindromes on one Y chromosome can
occasionally reach over and form a fatal attraction with the counterpart palindrome on its neighbor.
The two Y's fuse at the point of joining, and everything from the juncture to the end of the
chromosome is lost

The double-Y's so generated come in a range of lengths, depending on which of the palindromes makes
the unintended liaison. Like other chromosomes, the Y has a left arm and a right arm with the
centromere in between. The male-determining gene lies close to the end of the left arm. If the
palindromes at the very end of the right arm make the join, a very long double-Y results in which
the two centromeres are widely separated. But if the joining palindromes are just to the right of
the centromere, a short double-Y is formed in which the two centromeres lie close together.

Dr.~Page detected among his patients both short and long double-Y's and those of all lengths in
between. He and his colleagues then noticed a surprising difference in the patients' sexual
appearance that depended on the length between the centromeres of their double-Y's.

The patients in whom the distance between the Y's two centromeres is short are males. But the
greater the distance between the centromeres, the more likely the patients are to be anatomically
feminized. A few of the patients were so feminized that they had the symptoms of Turner's syndrome,
a condition in which women are born with a single X chromosome.

The explanation for this spectrum of results, in Dr.~Page's view, lies in how the double-Y's are
treated in dividing cells and in the consequences for determining the sex of the fetus.

When the centromeres are close together, they are seen as one and dragged to one side of the
dividing cell. As long as the Y's male-determining gene is active in the cells of the fetal sex
tissue, or gonad, the gonads will turn into testes whose hormones will masculinize the rest of the
body.

But when the centromeres lie far apart, chromosomal chaos results. During cell division, both
centromeres are recognized by the cell division machinery, and in the tug of war the double-Y
chromosome may sometimes survive and sometimes be broken and lost to the cell.

Such individuals can carry a mixture of cells, some of which carry a double-Y and some of which
carry no Y chromosome. In the fetal gonads, that mixture of cells produces people of intermediate
sex. In many of these cases the patients had been raised as female but had testicular tissue on one
side of the body and ovarian tissue on the other.

In the extreme version of this process, the distribution of cells may be such that none of the fetal
gonad cells possess a Y chromosome, even though other cells in the body may do so. Dr.~Page and his
colleagues found five of the feminized patients had symptoms typical of Turner's syndrome. The
patients had been brought to Dr.~Page's attention because their blood cells contained Y chromosomes.
Evidently by the luck of the draw, the blood cell lineage had retained Y chromosomes but the all
important fetal gonad cells had been denied them.

In 75 percent of women with Turner's syndrome, the single X comes from the mother. ``Since they are
females, everyone imagines it's Dad's X that is missing,'' Dr.~Page said. ``But it could easily be
Dad's Y.''

That the degree of feminization parallels the distance between the two centromeres of the double Y
chromosome is ``a fantastic experiment of nature,'' Dr.~Page said. Despite having studied the Y
chromosome for nearly 30 years, he has learned that it is always full of surprises.

``I continue to see the Y as an infinitely rich national park where we go to see unusual things, and
we are never disappointed,'' he said.

Dr.~Cynthia Morton, editor of the American Journal of Human Genetics, said the new explanation of
Turner's syndrome was plausible. ``It's another beautiful David Page contribution to the science of
genetics,'' Dr.~Morton said.

\section{An Organ of Many Talents, at the Root of Serious Ills}

\lettrine{S}{hould} anybody in the reliably pestilent\cite{pestilent} health care debate be casting
about for a mascot\cite{mascot} organ to represent some of the biggest medical crises that we
Americans face, allow me to nominate a nonobvious candidate: the pancreas\cite{pancreas}.

It may lie in the hidden depths of the abdominal cavity, and its appearance, size and purpose may be
obscure\cite{obscure} to the average person. Yet the pancreas turns out to be a linchpin in two
epidemics that are all too familiar.

As the organ entrusted with the manufacture of insulin and other hormones that help control blood
sugar, the pancreas gone awry is a source of diabetes, which afflicts more than 23 million people in
this country, including the newest member of the Supreme Court.

And as the tireless brewer of digestive juices that help shear\cite{shear} apart the amalgamated
foodstuffs that we consume each day, the pancreas is at the frontlines of our expanding waistlines,
the mass outbreak of fatness that has already claimed 60 percent of Americans and shows no sign of
slackening\cite{slacken}.

Researchers are discovering that the pancreas helps mediate much of the appetite-related cross talk
between the brain and the gastrointestinal tract, the streams of chemical signals that say, I'm
starving down here, how about some dinner, or, enough already, step away from that dessert cart and
no one will be hurt. By better understanding the precise role of the pancreas in conveying
sensations of hunger or satiety, suggested Rodger A.~Liddle of Duke University Medical School, we
may find new ways to combat obesity.

Other researchers are intrigued by the pancreas's ability to shield itself from harm, to churn out
huge quantities of enzymes that can rapidly reduce a cheeseburger and fries to so many particles of
amino acids, carbohydrates and fats, without digesting its own tissue in the process. They suspect
that the organ's set of self-protective mechanisms has a terrible downside, and helps explain why
pancreatic cancer can be so difficult to treat -- insights that are just beginning to offer hope in
the fight against one of the most lethal of all malignancies.

If you've been remiss in appreciating your pancreas, don't feel bad: the ancients were, too. Early
anatomists were certainly aware of the pancreas but mostly ignored it, and the organ's name reflects
that ho-hum attitude. Pancreas is Greek for all-meat, a reference to its seeming homogeneity from
one end to the other.

Much of the neglect may well have been practical. ``The pancreas has always been difficult to
study,'' Dr.~Liddle said. For one thing, it's hidden. Measuring some six to eight inches in length,
and slippery and lobular to the touch, the yellowish-brown pancreas is located deep in the abdomen,
wedged between the stomach and the spinal cord and extending horizontally right above the waistline.
Think of it as an oblong eel, the tail tucked into the stomach, the head butting up against the
curve of the duodenum, of the small intestine.

Add to that inaccessibility a prima donna sensitivity. ``If you do anything to the pancreas, you
initiate an inflammatory response,'' Dr.~Liddle said. ``It tends to become inflamed\cite{inflamed}
more easily than other organs.'' In fact, inflammation of the pancreas, or pancreatitis, is a
relatively common and often debilitating condition, brought on by excess alcohol, drug reactions,
gallstones, genetic predisposition or other reasons. Unfortunately, said Dr.~Anthony Kalloo, a
professor of medicine at the Johns Hopkins University School of Medicine, the symptoms of
pancreatitis, like chronic abdominal pain radiating into the back, could be misdiagnosed or
dismissed as a hypochondriac's lament. As a result, Dr.~Kalloo said, patients do not always receive
the right pain medications, the optimal diet, surgery when necessary.

For all the difficulty of studying the pancreas, researchers eventually came to appreciate the organ
as a gland of many talents, serving both an exocrine role -- secreting its products through ducts,
as the breast secretes milk and the sweat glands perspiration -- and an endocrine role, fabricating
hormones and squirting them into the bloodstream, as the ovaries and testes dispense sex hormones
and the thyroid thyroxine.

Roughly 90 percent of the pancreas is devoted to its exocrinic role of generating digestive enzymes
and funneling them into the small intestine, a burbling broth that flows forth from the pancreas at
a rate of perhaps a quart a day.

The other 10 percent of pancreatic tissue consists of so-called islet cells, the endocrine players
that synthesize insulin and glucagon to manipulate and titrate blood sugar, the body's energy
currency, as needed. In people with Type 1, or juvenile-onset, diabetes -- among them Justice Sonia
Sotomayor of the Supreme Court -- an autoimmune reaction ends up destroying many of these islet
cells, resulting in the need for lifelong insulin injections. Among sufferers of Type 2, or
adult-onset, diabetes, the reasons for insulin imbalance are more varied, and the condition can
often be treated through diet and exercise alone.

Imagine the pancreas as a tree, Dr.~Liddle suggested. The trunk and branches are the ducts that
deliver digestive juices, the leaves the factories that make digestive enzymes, and the islet cells
birds' nests scattered throughout -- in the tree but not of it.

When cancer strikes, it generally arises in the ductal tissue of the pancreas, the woody parts of
our metaphoric tree, and intriguingly, they feel the part. ``These tumors are rock-hard masses,''
said Peter Olson, an oncology researcher at the University of California, San Francisco. ``They're
white on dissection, very tough and fibrous.''

Pancreatic cancer is almost impossible to cure. About 34,000 Americans will be diagnosed with it
this year, and nearly as many will die of it. As doctors have long known, some of that lethality is
positional: there is no easy way to screen the deep-set pancreas for early signs of malignancy, and
by the time symptoms arise, the cancer has already spread to other organs.

Another reason for the ferocity, however, might be the nature of the tumors themselves. Most cancers
are thought to spur the growth of new blood vessels to supply them with the extra oxygen and
nutrients necessary for frenzied cell division, but pancreatic tumors are markedly devascularized.
``The number of blood vessels in a pancreatic tumor is 10 percent what it is in normal tissue, of
the pancreas or anywhere else,'' said David A.~Tuveson of the Cambridge Research Institute in
England. The results are devastating. In the anoxic microenvironment beneath the fibrous, bloodless
capsule, any malignant cells that survive become increasingly unstable and virulent, like
superroaches proliferating in the wake of a pesticide bomb. Moreover, without blood vessels, nothing
can get into the tumor to kill the renegade cells, so chemotherapy is almost useless.

Reporting recently in the journal Science on results with genetically engineered mice, Dr.~Tuveson
and his colleagues described a new approach to treating pancreatic cancer, in which the tumors were
revascularized and thus made sensitive to cancer drugs. Clinical trials are now under way to test
the basic strategy in people, and with all due caveats, Dr.~Tuveson said, ``I am cautiously
optimistic.''

\section{In China's Alleys, Shouting Vendors Sow Echoes of the Past}

\lettrine{N}{ot} long after daybreak, before the city begins its full-throated roar, the shouts and
calls can be heard up and down the old alleyways and deep within the walled courtyards that form the
crowded heart of the Chinese capital.

``Goat meat, goat meat!''

``Eggs, rice, eggs, rice!''

``Scrap, household scrap!''

With more emphasis on song than lyric, they are the marketing jingles of itinerant fruit vendors,
sellers of roasted duck, and stooped men who have mastered the art of resuscitating blunt kitchen
knives. Like the familiar whine of cicadas in August, their garbled calls are the soundtrack of the
Beijing summer, and many residents look forward to the return of the hawkers' glutinous rice cakes,
mismatched crockery and pet grasshoppers that sing.

Even more numerous than the hawkers are the recyclers, sun-scorched migrants from the countryside
who survive by collecting yesterday's newspapers, spent computers or tattered cotton blankets that
will be spun into next winter's comforters.

``If you can't yell loudly, you'll starve,'' said Chen Lin, 37, a bony, animated man who earns about
\$5 a day salvaging dead appliances and anything else containing metal.

``No one really knows what I'm yelling,'' he said, ``but they remember my song and this brings them
out of their house.''

The singing hawkers and recyclers are reminders of the days when Beijing was a thickly populated
maze of hutongs, or alleys, that crept outward from the grandiose imperial quarters occupied by
China's emperors and the officials and artisans who served them.

Cao Huiping, 45, a taxi driver whose childhood compound was stuffed with 17 unrelated families,
recalls when vendors filled the air with a cacophony of competing tunes.

``One minute it would be someone selling sugar, then as soon as their song faded it would be the
flour dealer, then the fabric salesman,'' said Mr.~Cao, whose home has since been replaced by an
upscale mall. ``Now I live in a building where people don't even know each other and everyone shops
at the supermarket.''

Gated apartment complexes are the hawker's enemy. So, too, are the air-conditioners that drown out
sales calls and keep residents inside.

The city authorities are no friends of the street vendors either. Stringent laws and urban
management officials, known as chengguan, keep them on the run with fines and harassment. ``The best
time to be out is lunchtime, when the chengguan are on break,'' said Meng Xiandong, 54, a vendor of
dried sweet potatoes, as he nervously scanned the crowds.

A good place to get a taste of old Beijing is Qianmen, a poor but colorful quarter south of
Tiananmen Square that is a jumble of twisting hutongs and ramshackle houses. On most days, one can
find peddlers selling meticulously skinned pineapples, a man offering two kinds of honey -- plain
and medicinal -- and an ornery cobbler who can resole a pair of shoes in as much time as it takes to
down a steaming bowl of hand-shaved noodles.

Cradling a brass teapot and watching over three pairs of caged lovebirds, Wu Xiulong, 76, sat in
front of the doorway of his courtyard and reminisced about the vendors whose arrival he used to
await as a child: the bean-cake man, the corncob seller, the baker who produced the flakiest
flatbread. ``Oh, back then they were baked on both sides, so crunchy, with sesame seeds,'' he said.
``It was so delicious. Now they're all gone.''

Zhao Cai, a 66-year-old knife sharpener, is one of the old-timers who can still be found wandering
around with a beaten-up toolbox that doubles as a bench. His call is bracing but melodious, although
once he sets to work on a blade, the noise of grindstone on metal brings out the old women with
their beloved worn-out cleavers.

``I hate stainless steel,'' he said as he pedaled the grindstone. ``No one makes knives like they
used to.''

Unsentimental and gruff, his accent betraying his hometown in China's far northeast, Mr.~Zhao has
been plying his trade for more than 30 years.

``When you're good at sharpening knives, you get to know everyone,'' he said.

How good is he? Customers sometimes foolishly test his handiwork by touching the sharpened edge.
``I've had ladies draw blood and swear they didn't feel a thing,'' he said.

As he spoke, there was a loud crash behind him, followed by a choking plume of dust. Workers in
orange vests were tearing signs off nearby buildings, part of a government campaign to make the
neighborhood more attractive to tourists who like a bit less visual chaos in their Old Beijing
experience. ``I don't recognize some of the streets around here anymore,'' he said before fleeing
the advance of the demolition crew.

One man who needs no vocal announcement is Li Hailun, a grasshopper salesman whose wares, hundreds
of wingless insects imprisoned in round, woven enclosures, produce a deafening, high-pitched
symphony. From July to October, Mr.~Li, 28, bikes around the city with his chirping quarry, each of
which sells for 50 cents to a dollar, depending on the quality of the song and the gullibility of
the buyer. Add a dollar if the critter comes in a graceful wooden cage.

Much of Mr.~Li's village, about a two-hour drive from the capital, is engaged in the grasshopper
trade: women weave the cages, boys catch the insects and the men pedal them to nostalgic city
dwellers. When sales lag, he heads to the gates of the nearest children's hospital. Subtracting the
occasional fine for vending without a license, Mr.~Li pockets \$200 a week, a tidy sum for a sorghum
farmer biding his time between planting and harvest.

The bugs draw a crowd wherever Mr.~Li goes. On a recent day, passers-by debated whether to feed them
carrots, scallions or rice. A woman said that toddlers raised alongside a trilling insect were not
easily startled by noise.

``When the grasshopper guy comes out, you know summer has arrived,'' said a man who was seeking to
replace the one, tethered to his rear-view mirror, that was on its last legs.

Mr.~Li, ever the salesman, added his own poetic pitch. He declared that the Chinese had been raising
grasshoppers for hundreds of years. Even Qianlong, a Qing dynasty emperor, was a connoisseur of the
fighting variety.

``Everyone loves grasshoppers,'' Mr.~Li said. ``When they sing, you can't help but feel happy.''

\section{For Speech-Impaired, Insurance Fights Remedy}

\lettrine{K}{ara} Lynn has amyotrophic lateral sclerosis, or A.L.S., which has attacked the muscles
around her mouth and throat, removing her ability to speak. A couple of years ago, she spent more
than \$8,000 to buy a computer, approved by Medicare, that turns typed words into speech that her
family, friends and doctors can hear.

Under government insurance requirements, the maker of the PC, which ran ordinary Microsoft Windows
software, had to block any nonspeech functions, like sending e-mail or browsing the Web.

Dismayed by the PC's limitations and clunky design, Ms.~Lynn turned to a \$300 iPhone 3G from Apple
running \$150 text-to-speech software. Ms.~Lynn, who is 48 and lives in Poughkeepsie, N.Y., said it
worked better and let her ``wear her voice'' around her neck while snuggling with her 5-year-old
son, Aiden, who has Down syndrome.

Medicare and private health insurers decline to cover cheap devices like iPhones and netbook PCs
that can help the speech-impaired, despite their usefulness and lower cost.

Instead, public and private insurers insist that, if Ms.~Lynn and others like her want insurance to
pay, they must spend 10 to 20 times as much for dedicated, proprietary devices that can do far less.

The logic: Insurance is supposed to cover medical devices, and smartphones or PCs can be used for
nonmedical purposes, like playing video games or Web browsing.

``We would not cover the iPhones and netbooks with speech-generating software capabilities because
they are useful in the absence of an illness or injury,'' said Peter Ashkenaz, a spokesman for the
federal Centers for Medicare and Medicaid Services. Private insurers tend to follow the government's
lead in matters of coverage. Two years ago, iPhones and netbooks barely existed, so it may not be
surprising that the industry has yet to consider their role as medical devices.

But the health care system has long had trouble keeping up with Moore's Law, the principle that
computing power rapidly increases even as costs fall sharply.

Doctors must still bring a patient into their offices instead of, say, inspecting an e-mailed photo
of a rash if they want most insurers to pay for the consultation. Digitizing medical records is such
a vast undertaking that the government is now spending billions of dollars to jump-start it.

In the case of A.L.S., also called Lou Gehrig's disease, advocates spent years fighting to have any
speech-specific devices covered by insurance, finally succeeding in 2001.

For the millions of Americans with A.L.S., Down syndrome, autism, strokes and other speech-impairing
conditions, the insurance industry's aversion to covering mainstream devices adds to the challenges
they face. Advocates say using an everyday device to communicate can ease the stigma and fear of
making the adjustment.

At the same time, current policies mean that the government and private insurers may be spending
unnecessary dollars on specialty machines.

Dr.~Stanley E.~Harris, who helps set device coverage policies for Horizon Blue Cross Blue Shield of
New Jersey, said that if enough patients requested new types of devices, the insurer would study
their usefulness. ``We're looking for evidence-based data to support the effectiveness of whatever
is being requested,'' he said.

In the meantime, people with speech disabilities have a choice: pay for a cheaper product from their
own pockets, try to borrow one from a private assistance group or spend their insurer's money on a
specialty device from a company like DynaVox Mayer-Johnson or Prentke Romich.

DynaVox, a leading maker of devices for the speech-impaired, has computers that start at \$8,000 and
run Windows, just like 90 percent of all PCs. To meet insurance rules, DynaVox disables the general
computing tools. After the insurer pays, customers can pay \$50 to DynaVox to reactivate the full
functions.

The proprietary devices have some special qualities. They are sturdier than typical computers and
have better speakers and links to support services.

But the prices may seem hard to justify based on components alone. One \$5,000 DynaVox product is
essentially the speech software bundled with a two-pound keyboard that has a six-inch screen. And
the manufacturers mark up standard accessories by as much as 2,000 percent. Prentke Romich, for
example, charges \$250 for a Bluetooth wireless adapter similar to those that cost \$20 in stores.

Jim Shea, vice president for marketing at DynaVox, says his company's prices run high because it
must do a lot of custom work and research to serve a niche that mainstream companies ignore. ``We
are not riding the wave of consumer electronics in terms of cost,'' he said. ``We're building the
devices here in Pittsburgh from scratch.''

In addition, the do-it-yourself approach isn't for everybody, he said. ``You have to be somewhat
savvy, get the software and set it up,'' he said.

Disease experts say companies like DynaVox and Prentke Romich make many sophisticated, helpful
products. Still, advocates argue, advances in computing and easy-to use speech software have opened
doors to use cheap mainstream alternatives. Indeed, the price drops have made it possible for A.L.S.
assistance groups to buy dozens of netbooks, install specialized software like Proloquo2Go and lend
them to clients.

Betsy Caporale, a speech language pathologist in Danville, Calif., has tested various devices and
software with children who have Down syndrome and autism.

``The iPhone has been a runaway success with these kids,'' she said. ``It takes them about 10
minutes to learn how to use the iPhone, and there is this cool factor for them.''

Ms.~Lynn, from Poughkeepsie, would like to see insurers loosen their rules to accommodate
general-purpose devices and give people like her more financial flexibility. Since insurers will
typically cover only one device every five years, people with degenerative conditions like A.L.S.
often hold off any claims until their condition worsens, and they really need an expensive specialty
product that can track their eye and head movements.

Perhaps the government could set a certain dollar limit and then let patients find the products that
fit their needs, Ms.~Lynn suggested. ``I really would like to see Medicare do away with the
dedicated-device rule and the one-device limit,'' she said by e-mail.

But so far, government and private insurers are not swayed. ``We look at determining the
effectiveness of the technology -- and not the cost -- first,'' Mr.~Harris said.

For Ms.~Lynn, the iPhone, with the special software, is cheaper, more effective and essential.
``Technology has become as important to me as air, food, water,'' she wrote.

\section{Military Chief Challenged on Troop Levels in Afghan War}

\lettrine{T}{he} chairman of the Joint Chiefs of Staff, Adm. Mike Mullen, told Congress on Tuesday
that success in Afghanistan would probably require more troops and certainly much more time, a
position seconded by a top Republican but challenged by a leading Democrat.

The intense dialogue, at a morning hearing of the Senate Armed Services Committee, previewed the
sharp national debate expected over coming weeks as the White House considers how best to pursue its
new strategy in Afghanistan in the face of growing skepticism from members of President Obama's own
party.

The committee chairman, Senator Carl Levin of Michigan, insisted that accelerated efforts to train
and equip Afghan security forces should precede any deployment of American troops beyond those
already committed by the Obama administration. Mr.~Levin's stance is expected to have great sway, as
he is the most powerful Democrat in Congress on military matters.

But the committee's ranking Republican, Senator John McCain of Arizona, his party's most experienced
voice on military affairs, countered by declaring that more troops were ``vitally needed'' in
Afghanistan -- and that any delay in ordering more combat forces to the fight put American lives at
risk.

Admiral Mullen said that no specific troop request had yet been received from Gen. Stanley
A.~McChrystal, the senior American and NATO commander in Afghanistan.

``But I do believe that -- having heard his views and having great confidence in his leadership -- a
properly resourced counterinsurgency probably means more forces, and, without question, more time
and more commitment to the protection of the Afghan people and to the development of good
governance,'' Admiral Mullen said.

``We will need resources matched to the strategy,'' he added.

Broad as they were, Admiral Mullen's comments were his most specific to date in a public setting on
whether more troops would have to be sent to Afghanistan, and they and seem certain to frame the
debate facing the White House, Congress and the nation in coming weeks.

The hearing officially was called to consider Admiral Mullen's nomination to serve a second term as
chairman, but it immediately turned into an analysis of the administration's broader policy for
Afghanistan and Pakistan, in particular whether more American combat forces should be sent rapidly
or whether it would be wiser to immediately begin shifting the bulk of the fighting to local forces.

A range of officials have said that the White House hopes to have several weeks at least before
being faced with dealing with any requests for more forces for Afghanistan -- and the political
implications of such a request here at home.

Admiral Mullen acknowledged the importance of the training effort as advocated by Mr.~Levin, but he
pointed out that training could not quickly provide sufficient levels of security required by the
new counterinsurgency strategy.

``I share your view that larger and more capable Afghan National Security forces remain vital to
that nation's viability,'' Admiral Mullen said. ``I share your view -- and have stated publicly --
that the path to achieving the president's goal is through our training efforts there. We must
rapidly build the Afghan Army and police.''

But he cautioned that ``sending more trainers more quickly may give us a jump-start, but only
that.''

``Quality training takes time and patience,'' he said. ``Private trust by the Afghans -- so vital to
our purpose -- is not fostered in a public hurry.''

Mr.~Levin, who met with commanders and troops in Afghanistan during the congressional Labor Day
recess, said that training Afghan Army and police units ``would demonstrate our commitment to the
success of a mission that is in our national security interest, while avoiding the risks associated
with a larger U.S.~footprint.''

And he said emphatically that ``these steps should be urgently implemented before we consider a
further increase in U.S.~ground combat troops, beyond what is already planned to be deployed by the
end of the year.''

Mr.~Levin said new goals should be established for Afghan security forces: The Afghan Army, he said,
should grow to 250,000 by the end of 2012, and the police to 160,000 by that date. The current
target is 134,000 army personnel members and 96,000 police officers by the end of next year.

Although Mr.~Levin and Mr.~McCain have a cordial working relationship on the committee, they were
blunt in expressing opposing views on whether more combat troops should be deployed to Afghanistan.

Mr.~McCain recalled that initial attempts in Iraq to shift the security burden to local forces
instead of American forces failed dismally. ``I've seen that movie before,'' he said.

He lauded\cite{laud} Mr.~Obama's decision earlier this year to send 21,000 additional troops to
Afghanistan. Those extra combat troops, Mr.~McCain said, were the correct priority -- and he said
more troops were ``even more necessary now.''

Additional combat troops ``are vitally needed,'' Mr.~McCain said, and warned that each day of delay
``puts lives in danger.''

``We will need more U.S.~combat forces in Afghanistan, not less or the same amount we have today,''
Mr.~McCain said.

The final brigade\cite{brigade} of fresh troops ordered to Afghanistan -- a group of about 4,000
trainers -- is scheduled to land by November, bringing the American troop level there to 68,000.

Other members of the committee said the civilian agencies of the United States government needed to
accelerate their assistance for rebuilding Afghanistan.

And Senator Jack Reed, a Rhode Island Democrat who recently visited Afghanistan, warned that such
civilian efforts were urgently needed to fix a government in Kabul that he described as
``dysfunctional'' and that was not viewed as legitimate\cite{legitimate} by much of the population.

\section{Chewy Chicken Feet May Quash\cite{quash} a Trade War}

\lettrine{C}{hina} is threatening to cut off imports of American chicken, but poultry experts have
at least one reason to suspect it may be an empty threat: Many Chinese consumers would miss the
scrumptious\cite{scrumptious} chicken feet they get from this country.

``We have these jumbo, juicy paws the Chinese really love,'' said Paul W.~Aho, a poultry economist
and consultant, ``so I don't think they are going to cut us off.''

Chicken exports were thrust to the forefront of American-Chinese trade tensions on Sunday when China
took steps to retaliate for President Obama's decision to levy tariffs on Chinese tires. The Chinese
announced that they were considering import taxes on automotive products and chicken meat, a
development that some trade experts feared could escalate.

American executives expressed concern about losing what recently has become the largest export
market for their chickens, one that is expanding rapidly as the Chinese population grows more
prosperous. But the executives also expressed relief that, so far, Chinese importers have told them
to keep the feet and wings coming.

``We were told by our customers in China to continue to pack and ship product,'' said Michael
D.~Cockrell, chief financial officer of Sanderson Farms, a major poultry producer based in
Mississippi. ``It gives us a little bit of optimism that we will get over this.''

At a time when feed prices are high and domestic chicken sales to restaurants are down because of
the recession, the Chinese market is important to the industry. Exports of American poultry totaled
\$4.34 billion last year. Of that amount, \$854.3 million worth of chicken meat (less than 2 percent
of total revenue by the American chicken industry) was exported to China and Hong Kong. But industry
executives said the exports to China were particularly profitable.

About half of the chicken parts sold to China are wings and feet, which are worth only a few cents a
pound in the United States. As delicacies in China, they fetch 60 cents to 80 cents a pound, a price
that no other foreign market comes close to matching, according to industry experts.

Mr.~Aho said the big chicken feet result from the American preference for white chicken meat. A bird
bred for big breasts is necessarily bred to have big, strong feet and legs, he said. The United
States is by far the world's leading supplier of king-size chicken feet.

Despite China's fondness for American chicken, the trade has been rife with problems since 2004,
when the countries banned each other's poultry products after an outbreak of bird flu. China quickly
lifted its ban, but the United States did not, because of continuing concerns about the safety of
Chinese chicken.

The Agriculture Department partly rescinded the import ban in 2006 by ruling that China could export
cooked poultry meat to the United States as long as it first imported the raw chicken meat from the
United States or Canada. But Congress quickly inserted a provision in an appropriations bill that
effectively prohibited the import of chickens processed in China, with lawmakers citing unclean
conditions.

Rosa L.~DeLauro, Democrat of Connecticut Democrat who leads opposition in the House to the imports,
said the ban had nothing to do with trade policy. ``For me it's about health,'' she said in an
interview.

China appeared to be ready to cut off imports of American chicken products in July, and American
poultry producers said the issuance of import permits slowed for a time. But sales have since
returned to normal levels.

In an effort to assuage\cite{assuage} Beijing, American poultry producers have made it clear that
they have nothing to do with the Congressional import ban and say they do not fear competing with
Chinese canned or frozen chickens.

``We believe in free and open trade and we feel our industry has a lot more to lose by being an
obstructionist in trade than in supporting China's position,'' said James H.~Sumner, president of
the U.S.A.~Poultry and Egg Export Council. ``If the product is fully cooked, then that would destroy
any possible pathogens plus the product would be subject to further inspection when it enters the
United States.''

Two weeks ago, Mr.~Sumner's group and the National Chicken Council joined other American food
organizations in sending a letter to Ron Kirk, the United State trade representative, cautioning
that action against Chinese tires could lead to retaliation. ``For some, the Chinese market is the
difference between profitability and possible bankruptcy,'' the letter warned.

Now that the Chinese are threatening retaliation, industry officials say they can only hope Chinese
taste buds outweigh protectionist impulses.

``It complicates the issue for the Chinese'' because of their consumer demand for American chicken
parts, said Daniel Griswold, a trade expert at the Cato Institute in Washington. On the other hand,
he said the American poultry industry also has a lot to lose, adding, ``If we are playing a game of
chicken with China we are going to be big losers.''

\section{New Standard Links Mileage and Gas Emissions}

\lettrine{T}{he} Obama administration proposed rules on Tuesday that would impose the first
nationwide limits on greenhouse gas emissions from vehicles and require a nearly 40 percent increase
in fuel efficiency for cars and light trucks sold in the United States by 2016.

Officials said the rules offered concrete proof of the administration's commitment to address global
warming just days before a series of international meetings on climate change open in Washington and
at the United Nations. Broader legislation to deal with climate change is haltingly moving through
Congress.

The government projects that the vehicle regulations will raise new car and truck prices by an
average of \$1,100, but that drivers will save \$3,000 over the life of a vehicle in fuel bills.
Officials also said the new program, which is to take effect in 2012, would reduce carbon dioxide
emissions by nearly a billion tons and cut oil consumption by 1.8 billion barrels from 2012 to 2016.

The regulations, 1,227 pages, will go through a 60-day public comment period before they become
final early next year.

The emissions program was first announced in May by President Obama as a way of resolving legal and
regulatory conflicts among several federal agencies and a group of states, led by California, that
want to impose stricter mileage and emissions standards than those set by Congress and a succession
of presidents.

Automakers had complained that they faced a thicket of varying rules that were almost impossible to
meet. The Obama compromise was endorsed by all the major auto companies, state officials and most
environmental advocates.

Addressing autoworkers at a General Motors plant in Lordstown, Ohio, Mr.~Obama said on Tuesday that
the new rules would benefit manufacturers, workers and consumers.

``For too long,'' Mr.~Obama said, ``our auto companies faced uncertain and conflicting fuel economy
standards. That made it difficult for you to plan down the road. That's why, today, we are launching
-- for the first time in history -- a new national standard aimed at both increasing gas mileage and
decreasing greenhouse gas pollution for all new cars and trucks sold in America.

``This action will give our auto companies some long-overdue clarity, stability and
predictability.''

Lisa P.~Jackson, the administrator of the Environmental Protection Agency, which along with the
Department of Transportation will enforce the new rules, said at a White House briefing, ``The new
standards are aggressive and achievable, and ensure that consumers have a full range of vehicle
choices.''

In addition to providing domestic and foreign auto manufacturers with a single national efficiency
standard, the proposed rules allow them to continue to build and import all classes of vehicles,
from the smallest gas-electric hybrids to monster S.U.V.'s. The mileage standard varies by vehicle
size, but companies will have to achieve a fleet average of 35.5 miles a gallon in combined city and
highway driving.

The regulations say that car companies can meet the new standards largely by using existing
technologies.

``These include improvements to engines, transmissions and tires; increased use of start-stop
technology; improvements in air-conditioning systems (to the extent currently allowed by law);
increased use of hybrid and other advanced technologies; and the initial commercialization of
electric vehicles and plug-in hybrids,'' the proposed regulations state.

Manufacturers can also claim credits toward the standards by selling so-called flexible-fuel
vehicles capable of running on a combination of gasoline and ethanol, and by selling more efficient
cars in California and in other states that planned to adopt its stringent rules. If all those
tactics are fully employed, the standard comes down by 1 to 1.5 miles a gallon by 2016, according to
analysts for environmental groups.

Automakers that sell fewer than 400,000 vehicles in the United States would be allowed to meet a
weaker standard to keep per-unit costs down. This group includes Mercedes-Benz, BMW, Volkswagen and
Subaru.

The United States Chamber of Commerce and a group of automobile dealers have already indicated their
intent to challenge the rules in court, saying the E.P.A. does not have authority to allow
California to set its own emissions standards for vehicles. (The national program essentially
ratifies one approved by California in 2004.) But the Alliance of Automobile Manufacturers, the
trade group representing most domestic and foreign car makers, endorsed it.

\section{Obama Follows Bush on Salmon\cite{salmon} Recovery}

\lettrine{I}{n} its first major effort to address the plight of endangered salmon in the Pacific
Northwest, the Obama administration on Tuesday affirmed basic elements of a recovery plan set forth
last year by the Bush administration.

The announcement angered critics of federal conservation policies who said the Bush plan did not go
far enough in improving fish habitats in the Columbia River basin or water levels in rivers for
migrating fish, and did not take immediate action to explore whether to remove four dams on the
lower Snake River.

Thirteen species of salmon are listed as endangered or threatened, and critics say the new Obama
plan, like the Bush one, is too ready to accept only slight gains in their populations, a potential
violation of the Endangered Species Act.

Obama administration officials said that while the plan affirmed the scientific and legal basis of
the Bush approach, it included revisions that would hasten and expand efforts to improve habitats,
monitor any effects of climate change and put in place contingency plans should fish populations
``decline significantly.''

The Obama plan leaves open as ``a last resort'' the possibility of removing dams if certain fish
populations decline to historic lows, but even then, critics say, the decision would depend on a
multiyear study of whether removing dams would improve salmon populations.

The issue of dam removal has become more complicated as the Obama administration seeks ways to
produce clean energy. The dams help provide low-cost hydroelectric power to the region but block
salmon and steelhead trout from reaching their historic spawning areas.

``It's clear that dams provide good clean energy,'' said Dr.~Jane Lubchenco, the administrator of
the National Oceanic and Atmospheric Administration, which oversaw the review of the Bush plan.
``They allow integration of wind into the grid. It's not clear what impact their removal would have
on salmon, and we believe that removal of them is not necessary in the short term. We want to give
these other actions a chance to work.''

Judge James A.~Redden of Federal District Court in Oregon is presiding over a legal challenge to
federal recovery policies brought by environmentalists, fishermen, the Nez Perce Indian tribe and
the State of Oregon.

Judge Redden has rejected two federal plans for restoring salmon in the Columbia basin, one by the
Clinton administration and an earlier plan by the Bush administration. He is expected to decide
whether to accept the Obama plan within the next several weeks.

Even as some criticized the Obama plan as not going far enough, others said it went too far.

``The extremists who brought this lawsuit may be critical about this plan because dam removal wasn't
delivered on a silver platter with promises of wrecking balls arriving next week, but they got what
they wanted from the Obama administration, and they'll try and convince Judge Redden to give them
even more,'' said Representative Doc Hastings, a Republican who represents part of eastern
Washington.

Mike Carrier, the natural resources policy director for Gov. Theodore R.~Kulongoski of Oregon, a
Democrat, said the Obama administration had ``wisely'' chosen to reverse some of the Bush
administration's environmental policies but in the case of salmon recovery was ``fundamentally still
embracing'' the Bush approach.

Nicole Cordan, the policy and legal director for Save Our Wild Salmon, a coalition that includes
many of the plaintiffs in the case, said, ``Yes, dam breaching is on the table, but the table is
over the river and through the woods and 1,000 miles away.''

\section{U.S.~Embassy Under Fire During Biden's Iraq Visit}

\lettrine{V}{ice} President Joseph R.~Biden Jr.~arrived here on Tuesday for two days of talks with
Iraqi leaders, and shortly after saying he was headed to bed for the night at the American Embassy
four mortars landed nearby around the heavily fortified international zone.

Mr.~Biden, making his second unannounced visit as vice president in little more than three months,
was never in imminent\cite{imminent} danger, but the attack underscored the violence that remains a
backdrop to halting American efforts to promote political reconciliation.

``I'm here to listen,'' Mr.~Biden said of his plans to meet with Prime Minister Nuri Kamal al-Maliki
and an array of other senior leaders, many of whom remain at odds.

At least two Iraqis were killed and five wounded in the attacks, according to the Ministry of the
Interior. Two mortar rounds landed in an apartment complex only a few hundred feet from the embassy
compound, according to witnesses.

It was not clear whether the attacks -- a recurring, if less frequent occurrence in the area known
as the Green Zone -- were aimed at Mr.~Biden's entourage. But word of his arrival at Baghdad
International Airport at 4:20 on a hot, dusty afternoon had been widely broadcast by international
and Iraqi news media.

The attacks began shortly after 7 p.m. Flashes typical of mortar fire could be seen in the Karada
neighborhood across the Tigris River from the embassy.

Inside the embassy's sprawling compound, a piercing ``duck and cover'' alarm began moments after the
American military commander, Gen. Raymond T.~Odierno, told reporters traveling with Mr.~Biden that
security remained at its lowest levels since the war began -- despite major bomb attacks like the
ones on Aug.~19 that badly damaged two government ministries and killed at least 132 people.

American helicopters swiftly appeared over the neighborhood where the mortars originated, while
ambulances and fire trucks moved toward the apartment complex where two of the mortars landed. A
third mortar apparently landed harmlessly on the banks of the Tigris; the fourth also landed nearby.

Shortly after the attacks, Iraqi police opened fire on a car speeding down Palestine Street in
eastern Baghdad. That in turn prompted an American patrol to open fire on the Iraqi police in the
confusion. ``We had to hide behind blast walls,'' an Iraqi officer said, speaking on condition of
anonymity because he was not authorized to speak to the media.

No one was killed in the cross-fire. The driver of the car was arrested after the police found
materials that could be used in a mortar attack, the officer said.

Mr.~Biden, appointed by President Obama to oversee the administration's Iraq policy as American
troops begin to withdraw, did not come to discuss security as much as politics. On Wednesday he is
to meet not only Mr.~Maliki, but also the Iraqi president, both vice-presidents, the deputy prime
minister, the speaker of Parliament and Kurdish leaders, as well as the new United Nations
representative.

Mr.~Biden, who visited Iraq in January as senator and vice president-elect, told reporters that he
had built ``some degree of trust'' with Mr.~Maliki and other Iraqi leaders and had acted as ``an
interlocutor on their behalf'' in his personal meetings and what he described as ``basically
weekly'' telephone calls.

Iraq's leaders are often at odds, and they have openly quarreled in recent weeks over security and
diplomacy in the aftermath of the Aug.~19 bombings, trading accusations of responsibility.
Mr.~Maliki's government has accused Syria of providing haven to extremists who might have been
involved in those attacks, even as others, including a potential rival, Vice Presdent Adel Abdul
Mahdi, said such accusations were premature.

Despite what seems to be waning American influence in Iraq, Mr.~Biden said he hoped to help the
country's leaders resolve ``outstanding political issues,'' including the legislation necessary to
hold parliamentary elections in January and another to distribute oil revenues.

The elections, now scheduled for Jan.~16, have colored almost everything the government has done --
or not done -- in recent weeks, paralyzing Parliament as leaders and their parties jockey for
position. ``I think there has been progress,'' Mr.~Biden said when asked why no significant progress
has occurred since his last visit in July.

At the same time, he acknowledged that there was little hope of resolving the most pressing issues
until after the elections. ``A successful election is the necessary condition for these outstanding
political issues to be resolved,'' he said.

Another focus of Mr.~Biden's visit will be tensions between Arabs and Kurds over 15 different
disputed territories along the sectarian fault line in northern Iraq that runs from Kirkuk to Mosul.

General Odierno -- at the request of Mr.~Maliki, according to American commanders -- has begun
planning joint patrols in the region that would include American, Iraqi and Kurdish regional
soldiers. Shortly after Mr.~Biden's visit in July, Mr.~Maliki and the president of the
semi-autonomous Kurdish region, Massoud Barzani, met in Irbil after barely speaking to each other
for more than a year and began discussing the proposal.

The proposal, however, has prompted public protests from Arab and Turkmen leaders in the north, and
so far no joint patrols have begun.

General Odierno said on Tuesday night that the intention was to prevent extremists from exploiting
the gaps along the disputed territories where Iraq's national forces face off again Kurdish forces.

``What we're trying to eliminate are these seams,'' he said.

In July, Mr.~Biden pointedly warned Iraq that the United States would be unlikely to remain engaged
in the country if its leaders reverted to sectarian violence, but his latest visit suggested the
administration was deepening its efforts here after what many Iraqis perceived as an initial period
of drift.

``I'll be back again,'' Mr.~Biden said.

\section{Many Faces, and Phases, of Steel in Cars}

\lettrine{F}{or} all the talk of reinvention in the auto industry, of car companies using high-tech
materials to make lighter and stronger vehicles that are safer and more fuel-efficient, the mainstay
of the mass-produced automobile is the same as it has always been: steel. The modern car still
contains more of it than anything else, about 60 percent by weight.

But automotive steel has changed quite a bit since the first Model T rolled off the assembly line.
Metallurgists and manufacturers have learned to manipulate steel's microstructure through precise
control of processing to create sheet steels of increasing strength. Prompted by crash-worthiness
requirements and the need to make cars lighter to improve gas mileage, automakers are replacing
conventional steels with advanced high-strength ones.

Where once a single grade of steel might have sufficed, the typical ``body in white,'' as automakers
call a car's basic skeleton, might now be a patchwork of a dozen or more steels of different types
and strengths, tailored through computer modeling to handle the stress and strain of normal driving
-- and of severe crashes.

``The day of the mild steel part at the Ford Motor Company is dead,'' said Ford's chief safety
engineer, Steve Kozak. ``The majority of steels we're using now are high strength or ultra high
strength.''

The advanced steels go by names like D.P. (for dual phase) and TRIP (for transformation-induced
plasticity). The strongest ones are used in parts like door beams, where the aim is to stop a
foreign object (like another car's bumper) from entering the passenger compartment, and windshield
pillars, where the goal is to prevent the roof from flattening like a pancake in a rollover.

In the front and rear of the vehicle, where there is more room to absorb the energy of a collision,
steels that deform more easily, and get stronger as they do, might be used. Even body panels, which
are usually made from milder steel, are bake hardenable, getting stronger as they are heated during
paint curing to resist denting better.

North American carmakers draw a lot of criticism for overall design and reliability of their cars,
but they generally receive high marks for their use of advanced steels. ``I'd say North America
doesn't have to take a back seat to anyone,'' said Richard A.~Schultz of Ducker Worldwide, a
consulting firm that tracks materials use by automakers.

The steel industry, Mr.~Schultz said, has worked to develop new products, to keep steel in cars and
to fend off the increased use of aluminum and other materials. They have been developing
higher-strength products for decades, beginning with high-strength low-alloy, or H.S.L.A., steels.
But most were costly to produce and difficult to make into parts. ``They were looked at as more for
aerospace applications,'' Mr.~Schultz said.

Only in recent years have researchers and manufacturers figured out how to make economical
high-strength steels that are pliant enough to be stamped and formed and that can be welded or
otherwise joined to other parts in the complex auto assembly process, where time is money.

``At the end of the day we have to manufacture everything that's there,'' said James G.~Schroth, lab
group manager in the Materials and Processes Lab of General Motors' Research and Development Center.
``The real changes in technology have been the higher-strength materials we can still make into
parts.''

Using higher-strength steel, a part can be downgauged -- made from thinner stock -- to help improve
fuel economy. But some parts, notably exterior panels, are already so thin -- less than a
millimeter, or one twenty-fifth of an inch, in many cases -- that stiffness becomes an issue. ``You
reach a thickness where you can't go even thinner,'' Mr.~Schroth said.

Like other steels, high-strength varieties begin as molten iron produced from ore and coke in a
blast furnace. ``Despite literally hundreds of years of development of the process, it's still the
best and most economical way of making high-purity iron,'' said Blake K.~Zuidema, director of
automotive product applications for the steelmaker ArcelorMittal.

The molten iron, usually with scrap added, is made into steel in a basic oxygen furnace, in which
oxygen is blown into the melt (``basic'' in the name refers to the alkaline materials used to line
the furnace). The gas oxidizes the iron and removes some impurities and most of the carbon.
(High-carbon steel is very hard but difficult to work.)

For most automotive purposes, the steel is cast into a slab. At this point, Dr.~Zuidema said, ``you
have a product that is not too much different than conventional steels,'' except for the addition of
small amounts of other metals -- alloys.

Some of the strength of steel is gained through alloying, said David K.~Matlock, a metallurgist at
the Colorado School of Mines and director of the Advanced Steel Processing and Products Research
Center, which is supported by auto and steel companies.

Small amounts of metals like niobium or vanadium help keep the steel hard during hot rolling, where
the slab is squeezed into a long sheet about two to four millimeters thick. But depending on the
metal and how much of it is used, alloying can be costly. ``For automotive steels, you want them to
be inexpensive,'' Dr.~Matlock said. ``The alloying that goes into them has to be low.''

The next step in transforming the steel relies not on adding metals, but on using temperature to
change the arrangement of the atoms in steel. Like most metals, steel has a crystalline structure,
with iron and other atoms arranged in orderly fashion. But steel is not a single continuous crystal,
like, say, the silicon that is used for semiconductor production. Rather, it consists of
crystallites, or grains, that can vary in size and orientation.

To change the crystal structure, manufacturers use techniques like cold rolling, annealing and
quenching. The steel sheet is rolled again, at room temperature, and heated to a high temperature
for a specific time before being rapidly cooled, or quenched, by water. The process is fast,
continuous and precisely controlled.

``Very fast cooling gives us the ability to get unique structures with far less alloying,''
Dr.~Zuidema said.

The process results in phase transformations, or changes in the steel's crystalline structure. When
the sheet is heated above about 1,300 degrees Fahrenheit, some of the basic cube-shaped crystals,
known as ferrite, transform to another cube-shaped phase, austenite. The rapid cooling then converts
the austenite to a third phase, martensite, which has an elongated shape and traps more carbon
atoms. ``Martensite is a very refined structure,'' Dr.~Matlock said, and adds strength.

The result is a steel with dual phases of ferrite and martensite, and by precisely controlling
temperature and timing, the amounts of the two can be varied to produce steels of differing
strengths and formability characteristics. Steelmakers can, in effect, dial up the strength and
ductility desired. ``You get the final structure you want,'' Dr.~Matlock said, ``and create the
final properties you need for the specific design.''

The so-called martensitic transformation can be thought of as a slight repositioning of the atoms
that reshapes the crystal. The change is so rapid that it travels through the steel at about the
speed of sound. Such transformations are not limited to steels, or even metals; the protein crystals
in the tail sheath of some viruses, for example, undergo a martensitic transformation, changing
their structure rapidly before the tail invades a bacterium.

In some high-strength steels, all of the ferrite ends up being converted to martensite; in others,
some austenite remains. High-martensite steels, including those with small amounts of boron added,
are difficult to form, so parts are usually made by hot stamping, in which the steel is heated
first, stamped and then cooled to allow phase transformation and hardening. Steels with retained
austenite are easier to form and get harder during the forming process, or even when crumpling in a
crash, as the deformations convert the austenite to martensite.

For the car manufacturer, there is a complex calculus involving strength, weight requirements and
costs that go into choosing what steel goes where. A part made of stronger steel may save some
weight, but it may require more work to produce the necessary tools and dies. ``Tool development
costs a lot of money,'' said Mr.~Schultz, of Ducker Worldwide. Other steels could require different
welding methods, so it might not be practical or cost-effective to introduce them into an assembly
line.

What is certain, though, is that cars are going to have to continue becoming lighter and stronger.
New federal regulations, for example, call for the roof to be able to withstand three times the
vehicle's weight in a rollover, twice the previous standard. And Mr.~Schultz has calculated that to
meet fuel economy goals for 2020, carmakers will have to keep replacing mild steels, and even some
early-generation high-strength grades, with higher-strength steels (and in some cases with aluminum,
plastics and other materials).

For steelmakers, that means continuing development of more advanced products that are stronger and
yet remain formable. ``It's a systematic engineering process,'' said Dr.~Zuidema, of ArcelorMittal.
``We have a fairly comprehensive knowledge of how each of the potential phases in the iron-carbon
system affects the mechanical properties.'' So if they know what properties are desired, they can
determine the structure needed, and thus the combination of alloying and processing to achieve that
structure.

For the steel companies, even incremental improvements will help the car companies meet crash and
fuel standards -- and help the steelmakers keep other materials manufacturers at bay.

``You're fighting for ounces here,'' Mr.~Schultz said. ``You're turning over every rock. I've told
my steel friends, fasten your seatbelts.''

\section{Training Apps That Help You Sweat the Details}

\lettrine{T}{o} prepare for his first marathon last year, Michael Nolan trained for six months with
the New York Road Runners Club, running up to 20 miles a day five days a week.

Nevertheless, he finished with an average speed of about 11 minutes a mile, a full 60 seconds a mile
off his target pace. This year, he vowed to be faster. ``I didn't want to take that long again,'' he
said.

So he got a personal trainer. Now as Mr.~Nolan prepares for this year's New York marathon, he is
leaner, stronger and ``easily'' averages eight-and-a-half-minute miles on training runs, he said.

Mr.~Nolan's new workouts are not coached by a running guru, but by iPhone applications that show
video workout instructions and tabulate every set of burpees, a full-body exercise for strength
training, and step-ups.

The sports and health industries are just beginning to tap the computing power of smartphones.
Applications range from simple calorie counters to heart-rate monitors that use complex metabolic
calculations.

These apps can help an athlete achieve a personal best, but some doctors say that more important is
their ability to produce no-fail routines for the sedentary\cite{sedentary} and obese\cite{obese},
which could improve health and drive down medical costs. Here are some of the popular fitness apps
out there.

WEIGHT LOSS In theory, losing weight is simple: Just burn more calories than you eat. Martin
Gramckow, an avid\cite{avid} bicyclist who lives in California, had considered that fact since a
cyclist he met on a ride bragged about losing 50 pounds. ``I'm huffing and puffing trying to keep
with him, while he leisurely pedals along and tells me how he did it,'' said Mr.~Gramckow.

The answer was calorie-counting. But Mr.~Gramckow thought logging every morsel\cite{morsel} that
passed his lips ``was always too much work.'' Then he saw Calorie Tracker for the iPhone, \$2.99
from LiveStrong.com, an affiliate of the Lance Armstrong Foundation.

As with similar apps, Mr.~Gramckow can search for foods by name or meals from restaurant chains and
the app calculates and stores all of the nutritional information.

``Sure enough I find a couple of things in my meal plan that are out of whack\cite{whack},'' said
Mr.~Gramckow, who trimmed 10 pounds in less than a month. ``I'm not far away from being the fittest
I've been in a long time,'' he said.

Calorie Tracker, which is also available for the BlackBerry (\$2.99), won't give you a breakdown of
proteins, carbohydrates and fats (unless you sync it and view the result on your computer), but a
free iPhone app called Lose It! will. The app also allows you to enter the ingredients of your own
recipes and store a list of meals you commonly eat.

My Food Diary may have the most complete food database, and it can be accessed from any mobile
browser, but it has no phone-specific app. Phone access is free with a \$9 monthly membership to the
Web site. Make sure you point your browser to the site formatted for the phone,
mobile.myfooddiary.com.

INDOOR FITNESS: FitDeck Mobile (\$4.99 for iPhone, \$14.99 for BlackBerry) was created by a former
Navy Seal instructor, Phil Black, as a simple workout that required no equipment. Illustrations show
exercises like jumping jacks and push-ups.

Fitsync is a Web site with a collection of workout routines that can be loaded onto phones using
Android, Palm or Windows Mobile software as well as an iPhone. The company claims a library of 1,600
exercises organized into more than 400 workouts like ``Rock Hard Challenge'' and ``Bikini Body
Cardio.'' Scheduled workouts from the Web site can be sent to your phone on workout days. An annual
subscription, normally \$39.95, is being offered at \$2.99 for Android phones.

Mr.~Nolan, the marathoner, used apps from PumpOne, which has 20 workout apps for the iPhone,
including ones for specific sports like golf or specific goals like burning fat. The newest app from
the site, Fitness Builder 2 Plus, has 5,000 images and videos of exercises, and 500 ready-made
workout programs for a \$4.99 monthly subscription or a \$99.99 onetime fee. Created by an exercise
physiologist, Declan Condron, the programs can be researched by muscle group, effort level, type of
equipment or goal (like speed or strength). It also lets you enter the number of reps and the weight
used.

Smheart Link (\$124.95) makes a heart monitor that links to an iPhone. Smheart Link works with four
apps that allow gym rats to arrange a display screen to show what they want, like calories burned,
average heart rate or elapsed time. It also links to sensors on indoor and outdoor bikes that
measure cadence and estimated speed and distance. Hard-core riders can attach a power meter that
measures the watts a cyclist generates, a measure bikers often use in competitive training.

For precise data, you can visit a New Leaf-trained technician, who employs metabolic testing
equipment to measure your heart and lung efficiency in a grueling aerobic test (average cost
\$175-\$200). Using that data, New Leaf sends custom exercise programs to your phone and works like
a virtual personal trainer, telling you how high to push your heart rate, when, and for how long.

It also tracks results through the Smheart Link heart monitor. New Leaf's detailed reports show how
many fat and carb calories you have burned, and how much time is spent in each of five heart rate
zones.

OUTDOOR FITNESS: One of the beauties of a mobile device is that you don't have to be wired to a
machine in a gym. Several apps are designed with hiking, biking and running in mind.

MyTracks, a free app for Android phones, uses GPS signals to track your time, distance, speed and
elevation as you hit the trail. The data can be loaded into a spreadsheet on Google Docs to
determine whether you are getting faster, or you can put the map and statistics in MyMaps to share
with friends. You can also send an e-mail message of your route to running buddies from the app. If
you want heart-rate data or calories burned, you will need a second device.

The Nike + iPod Sport Kit is made especially for runners (it can be used in a gym for cardio
workouts as well); it tracks time, distance, pace and estimates calories burned. It uses a \$29
sensor compatible with a Nike+ shoe and beams information to an iPod or iPhone. It does take some
effort to calibrate. You will need to run a known course to set it up for highest accuracy.

The cycling app iMapMyRide turns the iPhone into a GPS cyclometer, recording time, distance, speed,
altitude and estimated calories burned. You can see your position on a map, and it has a button to
pause your ride. Your maps and statistics are loaded to the MapMyRide Web site, where you can
measure your progress, and, if you like, share your routes with other riders.

MapMyRide has a free app if you don't mind seeing ads, or an ad-free \$4.99 version that also stores
more rides and statistics. There is a separate running app as well. The Web site requires a separate
subscription, but there is a free option as well as an unlimited-use \$99 annual membership.

There is one respect, however, in which these apps don't go the distance. Powering the screen while
also using GPS or a heart monitor (and maybe listening to music as well) will leave your batteries
wheezing like a tubby, two-pack-a-day smoker. If you're really into getting fit, your smartphone may
hit the wall before you do.

\section{Early Warning for a Deadly Kidney Disease}

\lettrine{C}{arol} Johnson, an otherwise healthy woman in her 60s, was at her wits' end. Despite a
good diet and regular exercise, she was gaining weight -- 20-plus pounds. Her blood pressure was too
high, even though she was taking three drugs for it. And she didn't feel well.

Since 2005, two doctors had told her she had a high blood level of creatinine, a product of muscle
breakdown that can suggest abnormal kidney function. Yet neither doctor could figure out why. Even
repeated kidney infections failed to alert her current doctor to the real problem.

Then, in March 2008, Mrs.~Johnson, a retired special-education teacher living in Independence, Mo.,
noticed an advertisement for a free screening test offered by the National Kidney Foundation.

The test, part of a nationwide program called KEEP (for Kidney Early Evaluation Program), can serve
as an early warning sign for a disease that often shows no symptoms until patients are on the verge
of kidney failure. It uses a formula to convert the creatinine reading to a better measure of how
well the kidneys are filtering wastes from the blood -- the EGFR, or estimated glomerular filtration
rate. (The precise rate, called the GFR, is measured by a much more involved test.)

Mrs.~Johnson's EGFR was 49, well below the normal reading of 90 or above and a sign that she had
chronic\cite{chronic} kidney disease. Suspecting that she was taking more medication than her
kidneys could handle, her doctor stopped two of them.

And using information she found on the Web, Mrs.~Johnson made some changes on her own. She cut out
red meat, colas and sweets. She started eating still more vegetables and whole grains. And she kept
to her daily exercise program, alternating among a gym workout, water aerobics and brisk walking.

Within six months Mrs.~Johnson, now 67, had lost the extra weight, and her blood pressure and EGFR
were normal. ``If not for going to KEEP and finding out what was really wrong,'' she said in an
interview, ``I would have accepted the fact that nothing more could be done.''

Multibillion-Dollar Problem

Like many other Americans with chronic kidney disease, Mrs.~Johnson might have ended
up with a heart attack or stroke, bone disease or kidney failure. Once kidneys fail, patients must
undergo regular dialysis or have a kidney transplant, procedures that add billions of dollars to the
nation's health care costs.

The kidney foundation undertook the KEEP program nine years ago because family doctors and
internists often fail to order an EGFR when performing routine blood tests on patients at risk for
kidney disease -- even, as happened to Mrs.~Johnson, when there are hints of kidney malfunction.

Fully half of the first 11,000 people whose blood and urine were tested in the KEEP program had
evidence of kidney disease, but only 3 percent of them knew that before being screened.

Belgian nephrologists writing in The Mayo Clinic Proceedings in 2006 reported that physicians were
often unaware of the severity\cite{severity} of kidney malfunction because they relied solely on a
creatinine measure. They noted that as many as four out of five patients with chronic kidney disease
were referred to a specialist for treatment late in the disease, within six months of needing
dialysis or a transplant.

According to the latest national surveys, 26 million Americans have chronic kidney disease, and the
numbers rise daily as more people become overweight and develop diabetes or high blood pressure, the
two leading causes of kidney malfunction.

The KEEP program, which has so far examined more than 130,000 participants, screens adults with the
most common risk factors for chronic kidney disease -- diabetes, high blood pressure or a family
history of kidney disease.

To find nearby upcoming screenings, go to www.kidney.org and click on the ``Get Tested'' box, or
call the foundation at (800) 622-9010.

Follow-up surveys of 72,000 people who have taken advantage of the free community-based screenings
indicate that the test has prompted many to take action, like seeing a physician, adjusting their
diets and doing more to control their blood pressure and diabetes.

Constellation of Risks

Chronic kidney disease is a life-threatening condition even for the young and middle-aged. Men under
55 and women under 65 identified through the KEEP program as having chronic kidney disease had twice
the risk of heart attack or stroke, and four times the risk of dying, as those free of the disease,
according to a study led by Dr.~Peter A.~McCullough of the William Beaumont Hospital in Royal Oak,
Mich.

Dr.~McCullough explained that kidney damage caused biological changes that accelerated injury to
blood vessels throughout the body. He urged that all adults at risk for kidney disease, including
smokers, be routinely checked for protein in their urine and reduced kidney function. Those found to
have a kidney problem should also be tested for abnormalities that can result in bone disease or
calcium deposits in blood vessels, he said.

Dr.~Joseph A.~Vassalotti, chief medical officer of the kidney foundation, said that in addition to
the leading risk factors, many other conditions could increase the risk for chronic kidney disease.
Among them are heart disease, lupus, amyloidosis (a buildup of amyloid proteins in the body),
chronic kidney infection, polycystic kidney disease, hepatitis B or C, and multiple myeloma.

In addition, kidney damage can result from overuse of corticosteroids or of nonsteroidal
anti-inflammatory drugs like ibuprofen and naproxen, from excessive cola consumption and from
exposure to certain environmental toxins and radiological materials like barium.

The risk of developing kidney failure is also higher among Hispanics, African-Americans and American
Indians than other ethnic groups, partly because of higher rates of high blood pressure and diabetes
and poorer access to medical care.

As with Mrs.~Johnson, if kidney problems are caught early and their underlying cause is corrected or
properly treated, kidney damage and other complications can be slowed or halted, and sometimes even
reversed.

As Sally Burrows-Hudson, a registered nurse in Sunnyvale, Calif., and a specialist in kidney
disease, wrote in The American Journal of Nursing, ``When patients at risk for chronic kidney
disease are identified early and treated aggressively, the disease progression can be slowed or
stopped, complications can be prevented or controlled, and clinical outcomes can be favorable.''

\section{Japan's New Prime Minister Takes Office, Ending an Era}

\lettrine{Y}{ukio} Hatoyama, who led his party to a landmark victory in elections last month, took
office as prime minister and named a cabinet of loyal allies on Wednesday, promising to bring change
to a country mired in stagnation after a half-century of virtually uninterrupted, one-party rule.

Mr.~Hatoyama has said the Democratic Party will reverse Japan's long economic malaise, increasing
social benefits and aligning policies more closely with the public's needs rather than those of big
business or the country's bureaucrats. He has also spoken of redefining Japan's relationship with
the United States, its closest ally.

``I am trembling\cite{tremble} with deep emotion over this moment of historical change, while at the
same time I know I have taken on an immense responsibility,'' Mr.~Hatoyama said at a news
conference. ``We are entering the realm of the unknown.''

Mr.~Hatoyama, 62, is a management professor with a doctorate in engineering from Stanford; a critic
of globalization; a political scion who cast out Japan's postwar political order with a decisive
victory over the long-ruling Liberal Democratic Party.

He inherits an economy that is emerging from its deepest recession in decades, brought on by a
collapse in the country's mainstay exports amid the global economic crisis. He must also deal with
debt approaching twice the size of its gross domestic product, a legacy of years of government
spending on public works projects that fueled the politics of the departing Liberal Democrats.

Unemployment is at a record high of 5.7 percent, while the cost of supporting a rapidly aging
population is threatening the country's public finances.

Blaming years of mismanagement by bureaucrats for Japan's woes\cite{woe}, Mr.~Hatoyama has made
reining in their power a major goal. His rhetoric struck a chord with voters demoralized after
decades of insider-driven politics and wasteful spending.

The Democrats also promise to redistribute more funds directly to the country's struggling
households, and build a stronger social safety net. They hope that more economic security will spur
consumer spending, and wean Japan off its dependence on exports for economic growth.

``First and foremost, we will engage in policies that will boost household incomes,'' Mr.~Hatoyama
said. ``We will realize a world where politics does not rely on bureaucrats.''

The cabinet Mr.~Hatoyama picked reflects those goals.

Hirohisa Fujii, Japan's new finance minister, is a Finance Ministry bureaucrat-turned-politician
likely to play an important role in efforts to control the bureaucracy and rein in spending. He
served as finance minister for a coalition government in the mid-1990s, when the Liberal Democrats
briefly lost power.

Mr.~Fujii will be joined in his task by the deputy prime minister, Naoto Kan, 62, a founder of the
Democratic Party who will lead a new agency called the National Strategy Bureau. Mr.~Kan will take
on tasks once carried out by the bureaucracy, including drawing up national budgets and filling top
bureaucratic posts.

Mr.~Kan, a former activist, lawyer and lawmaker, has a record of battling bureaucrats, most famously
over a government cover-up of tainted blood products that caused hundreds of hemophilia patients to
contract the AIDS virus.

Critics say the Democrats, who swept into office with 308 out of 490 Parliament seats, lack a
coherent growth strategy. ``The Democratic Party is still trying to figure out how best it should
engage with the global economy,'' said Izuru Makihara, a politics professor at Tohoku University.
``They say they will overhaul current policy, but it's not clear what they will replace it with.''

In his cabinet lineup, Mr.~Hatoyama also accommodated the Democrats' two coalition partners -- the
Social Democratic Party and the People's New Party, a conservative splinter group.

Those parties, with anti-free-market platforms, are expected to push the Democrats to reverse many
of the pro-market reforms championed by the former prime minister, Junichiro Koizumi.

The Democrats' election platform also called for a re-examination of Tokyo's ties with Washington,
including the presence of 50,000 American service members in Japan. But Mr.~Hatoyama's choice of
foreign minister, the 56-year-old moderate Katsuya Okada, signals thatJapan's pro-American foreign
policy will not drastically change.

\section{China Says It Disrupted Bomb Plot in Tense Area}

\lettrine{C}{hina} announced the breakup of a bomb-making plot in the volatile western region of
Xinjiang on Wednesday, apparently an indication that authorities had not only failed to suppress the
ethnic hatreds there but also that the weapons of choice in the feud\cite{feud} could be getting
more lethal\cite{lethal}.

The announcement, from the Public Security Ministry, said the authorities had arrested six people
who had established three bomb-making workshops and assembled about 20 explosive devices in a town
430 miles outside the region's capital, Urumqi.

The ministry said that the suspects had planned to plant the explosives on cars, motorcycles and
people, but that they were foiled by what it called timely arrests. The police did not identify the
ethnicity of the suspects. But the names of those alleged to be the ringleaders suggested that they
were Uighurs, an ethnic group that dominated the province before an influx\cite{influx} in recent
decades of Han, China's major ethnic group.

Tensions between the Uighurs and Han in Xinjiang reached a boiling point in July when Uighurs rioted
in Urumqi. At least 197 people, mostly Han, died in the worst incidence of ethnic violence in China
in decades.

The weapons used in that mayhem\cite{mayhem} and subsequent clashes were largely confined to
hatchets, sticks, cudgels\cite{cudgel}, knives, needles, toothpicks and pins. But the Wednesday
announcement was the first suggestion that antagonists with scores yet to settle were making bombs.

Uighurs say the government has hidden the true tally\cite{tally} of Uighurs killed by security
forces or in revenge attacks by Han after the July riots. Han accuse the authorities of failing to
protect them from the Uighurs. Earlier this month, Han residents took to the streets of Urumqi
demanding better security.

On Tuesday, the police announced that they had detained a total of 75 suspects accused of
needle-stabbing attacks in the Xinjiang region. A public security official said two of those
detained confessed they stabbed a resident in a supermarket ``to create panic in society.''

Others admitted that they had organized needle attacks to inflame ethnic hatred, according to the
state-run news agency, Xinhua. On Saturday, two men and a woman were sentenced to 7 to 15 years in
prison for stabbing or threatening victims with syringes.

Nicholas Bequelin, a researcher for Human Rights Watch, questioned whether the police had uncovered
hard evidence or were merely trying to pacify\cite{pacify} the Han by accusing Uighurs of unproved
plots.

``I remain perplexed by these allegations,'' he said in a telephone interview from Bangkok. ``There
are certainly security issues in Xinjiang. But this seems like a major distraction, as if the
government wants to make a point.''

``Definitely, authorities are trying to reassure people that their concerns about security are being
taken seriously and that the government is actively removing this threat,'' he said.

Authorities also appear to be trying to reassure Xinjiang residents that the alleged syringe attacks
are not lethal. On Sunday, Xinhua reported that no trace of toxic or viral substances were found in
blood samples taken from victims.

\section{Profit Rises at Oracle, but Sales Decline}

\lettrine{T}{he} Oracle Corporation's profit rose 4 percent in the latest quarter, matching Wall
Street's forecasts, despite a drop in sales that revealed that businesses are still being
tight-fisted about buying new software.

The sales figure was short of analysts' expectations, and Oracle's shares fell in after-hours
trading.

The company's results, reported after the market closed Wednesday, reflect a familiar pattern that
has emerged for Oracle during the recession.

Oracle's sales of new software licenses fell 17 percent to \$1 billion, while revenue from software
updates and technical support contracts climbed 6 percent to \$3.1 billion. While many businesses
are still reluctant to pay for new software, existing Oracle customers usually pay the company to do
the follow-up work on software they have already bought, which is why the numbers sometimes go in
different directions.

The rise in support work helped lift Oracle's net income for the quarter to \$1.12 billion, or 22
cents a share. That compared with \$1.08 billion, or 21 cents a share, in the same quarter a year
ago.

Excluding one-time items, profit was 30 cents a share, matching the average estimate of analysts
polled by Thomson Reuters.

Sales fell 5 percent to \$5.05 billion, short of expectations for \$5.25 billion.

Guidance for the September-November quarter, Oracle's fiscal second quarter, were in line with
analysts' projections.

Oracle expects profit of 35 to 36 cents a share, excluding one-time items. Analysts predicted 36
cents a share, on average. Revenue is expected to be flat to up 3 percent over last year, which
translates to a range of \$5.6 billion to \$5.77 billion. Analysts expected \$5.72 billion.

Oracle sees revenue from new software licenses coming in flat to down 10 percent over last year.

Oracle is the world's top seller of database software, which companies use to archive and retrieve
data like payroll or customer information. Oracle is also a major player in the ''middleware''
market, which refers to software that allows computing applications to talk to each other.

The company, run by Lawrence Ellison, is trying to branch out by buying the struggling computer
server maker Sun Microsystems for \$7.4 billion, a deal that would thrust Oracle into the hardware
market, a new area for the company.

The deal is being held up by European Union antitrust regulators, who are worried about Oracle's
plans for Sun's MySQL open-source database, which is popular among Web companies and competes
against Oracle's proprietary database. The underlying programming code for open-source software is
distributed free on the Internet; companies make money off it by selling support contracts for
products built from that code.

There is still uncertainty about how the transaction will play out, but Oracle is moving ahead with
plans for Oracle-Sun products.

\section{A High School Football Team is Beaten, but Not Defeated}

\lettrine{A}{s} halftime of the Meadowcreek Mustangs' opener mercifully arrived, Coach David Ward
summoned\cite{summon} his players, some of them in midseason form with chins lowered to
clavicles\cite{clavicle}.

``Y'all run to the lockah room,'' Ward said, his accent a marriage of Massachusetts upbringing and
job jumping around the South.

They jogged, as ordered, their blue-jerseyed backs turned to a malfunctioning scoreboard that did
not convey down, distance and yard line. It did, alas, show Home (Meadowcreek) trailing Visitor
(Luella), 33-0.

The Mustangs\cite{mustang} played with purpose in the second half but were ultimately subjected to a
rule all too familiar to them -- a continuously running clock when the score becomes lopsided -- and
lost, 47-0. In his postgame address, Ward invoked St.~Peter and channeled ``Rocky,'' fists raised in
a boxer's pose while saying in a punch-drunk voice, ``Ref, ya stop this fight and I'll kill ya.''

His message: Keep slugging\cite{slug} away, no matter how beaten down you are. It is one that must
be pounded into the Mustangs.

As the high school football season begins, everyone who smears on eye black and buckles a helmet has
big hopes. But in the yin and yang of football, for each team that tears through its schedule
unscathed, there is one that rarely basks\cite{bask} in victory's glow.

Then there is Meadowcreek, to whom the adage\cite{adage} ``win some, lose some'' holds little
meaning. The Mustangs halted a 43-game losing streak last season and have won once in their last 55
games, falling twice more since opening night.

This is not so much a bad team as one in a bad situation. With an enrollment of 2,400, Meadowcreek
belongs to Georgia's large-school classification. Its classrooms and hallways offer a
tapestry\cite{tapestry} of ethnic backgrounds, from Afghan to Zambian, speaking 79 languages. Many
students were reared in cultures in which football more likely connotes goalkeepers and corner
kicks. Although every other pupil\cite{pupil} is Hispanic, only a few dot the 45-man roster.

The student body is poor, with 88 percent qualifying for reduced-price or free meals. Students
struggle academically, shrinking the pool of eligible players.

Further, the Mustangs play in a powerhouse region that includes the state's second-ranked Class 5A
team, Grayson, and two well-financed programs with rich legacies.

Parkview claimed three straight state titles -- the last in 2002, the first two engineered by
defensive back Jeff Francoeur, now the Mets' starting right fielder. In 2005, Parkview was lauded as
one of the nation's 10 foremost athletic schools by Sports Illustrated.

Brookwood, the state runner-up in 2002 and 2005, is supported by a Touchdown Club, whose members
have raised more than \$300,000 in some seasons. Meadowcreek, with barely a dozen registered
boosters, spends in the hundreds.

Yet there is a whiff\cite{whiff} of hope.

After failing for two years, Meadowcreek recently attained the Annual Yearly Progress goals, a state
designation for schools that reach certain standards in math and English/language arts. (It must do
so for three straight years to escape the state's Needs Improvement List.) While its college
entrance exam scores are among the lowest in Gwinnett County, the most populous school system in
Georgia, its graduation rate rose 7 percent from 2008.

Among 60 applicants for head football coach this year, Meadowcreek chose Ward, a 52-year-old dynamo
and self-proclaimed Mr.~Fix-It.

On his r\'esum\'e, Ward noted that he had inherited four teams with a combined seven wins in eight
previous seasons and managed a collective 13 victories in his first year.

Ward, a former ironworker, said, ``Having been in construction, I like to build something from the
ground up.'' He also likes to refurbish\cite{refurbish}, having bought fixer-upper or foreclosed
houses in various towns where he coached as far away as Louisiana and Nebraska while his wife tended
to their home in Lawrenceville, Ga.

Ward's approach sounds simple: ``Make the kids feel like this is a club.''

He is prone to beginning sentences with the words, ``There's an old saying \ldots''

A favorite: ``It's not about the X's and O's, it's about the Jims and Joes.''

Another: ``You can't learn to win unless you learn to lose.''

Here, learning to lose can be a lesson never completed. Mustangs players say previous seasons have
been plagued by infighting, mainly verbal, as frustrations bubbled over.

``It's been a challenge,'' linebacker David Smalls said. ``You've got the history. Eighty-five
percent of the school doesn't support us.''

Asked how he stays motivated, Smalls placed his hands on the sides of his face, mimicking a
racehorse wearing blinders.

``We actually got a team this year,'' he said, ``instead of a bunch of individuals.''

Running back Martrez Ogletree discovered a ``much more positive'' group attitude than expected after
transferring from Brookwood. ``We can win games,'' he insisted.

He hears what teammates hear: some students ridiculing them. ``Not to my face,'' said Ogletree,
whose pate has sprouted traces of a Mohawk. ``I guess they're scared of me.''

Players want to win but, to brace themselves for disappointment, rely on other stimuli for
participating. For Ogletree and another Brookwood transfer, tight end Lavon Hooks, it is showcasing
themselves for colleges. They embrace their new surroundings, even if pregame meals now consist of
sandwiches instead of steaks.

For quarterback Tydrell Wright, an African-American who persuaded four Hispanic freshmen to try out,
it is the adrenaline rush. For lineman Joe Pardo, it's to prove he can master a sport other than
wrestling.

For kicker/safety Abel Ordonez, it's the love of the game. For lineman Wifred Irie, a native of the
Ivory Coast: ``It's not really about wins and losses. It's about having fun and playing to the best
of your ability.''

For years, Chris Ibawuchi, a strapping 6 feet 2 inches and 215 pounds, declined coaches' invitations
to play. A senior, he's a rookie lineman, playing ``just for fun.''

``If we go zero and 10, like we do other years,'' said Ibawuchi, of Nigerian descent, ``I'll be
O.K.''

At least once, Meadowcreek pondered withdrawing from its region and declaring independence from the
state athletics association, which would allow it to set a softer schedule. Jason Dopson, the
athletic director, rejected the idea, although the school's four trophy cases contain no football
mementos.

``You're throwing in the towel at that point,'' he said.

Dopson prefers to measure the program by the number of players who advance to college, on
scholarship or not. Higher education is stressed throughout the school, which boasts that the class
of 2009 received more than \$2.5 million in aid.

``To me, that's the biggest reward,'' said Dopson, who added that Ward's main selling point was his
skill as a math teacher.

Ward, the Mustangs' seventh coach in 13 seasons, has endured far worse than at this school, where
game crowds are a couple of hundred and, by players' accounts, the last football pep rally was held
in 2006. At a previous job, he had only one assistant. ``And he sat in the bleachers and smoked
cigarettes during practice,'' Ward said.

``I'm one of those underdog guys,'' added Ward, whose r\'esum\'e also lists four degrees and
experience playing semipro football.

``You keep the kids enthusiastic,'' he said. ``That's when you know you've made it: You lose a game,
but you don't lose the team.''

He won't, if Wright speaks for most of his teammates.

``I never think,'' he said, ``that we're going to lose.''

\section{Following Trash and Recyclables on Their Journey}

\lettrine{W}{here} does all the trash go?

Karin Landsberg, 42, a self-described ``eco-geek'' in Seattle, was so curious that she invited
researchers from the Massachusetts Institute of Technology into her home last month to fish 12 items
out of her garbage and recycling bins -- a can of beans, a compact fluorescent light bulb -- and tag
them with small electronic tracking devices.

Her trash is now on its journey to the place where it goes to die or be reborn.

The Architectural League of New York went through a similar trash-tagging exercise as part of the
same project when it moved its offices from midtown Manhattan to SoHo two weeks ago. Among the
discarded items tagged were a coffee cup, a filing cabinet, a book shelf, a broken wine glass and an
empty plastic bottle that had held liquid soap.

``All they can tell me up to this point is that some of the stuff has gone through the Lincoln
Tunnel,'' said Gregory Wessner, director of digital programs and exhibitions for the league. ``It is
on the move. We're really excited to know what happens.''

Through the project, overseen by M.I.T.'s Senseable City Laboratory, 3,000 common pieces of garbage,
mostly from Seattle, are to be tracked through the waste disposal system over the next three months.
The researchers will display the routes in real time online and in exhibitions opening at the
Architectural League of New York on Thursday and the Seattle Public Library on Saturday.

One purpose of the project, said Carlo Ratti, director of the lab, is to give people a concrete
sense of their impact on the environment in a way that might lead them to change their habits.

``If you see where a plastic bottle ends up, a few miles down the road in a dump, you may want to
get tap water or some other container for the water,'' Mr.~Ratti said.

Collecting, transporting, storing and getting rid of garbage is a costly and often daunting task for
cities. Lynn Brown, a spokeswoman for Waste Management Inc., a company that runs both landfills and
recycling centers nationwide and is helping to underwrite the tracking project with \$300,000, said
garbage moved through a vast network of sites run by multiple contractors, which makes it
challenging to find the most efficient way to handle it.

It also means hundreds of possible journeys for trash.

``From a logistics standpoint, it's a very complicated situation,'' Ms.~Brown said. ``When you look
at how waste is handled in different cities, it's like snowflakes. It's all different.''

Other factors are also in play in the travel of recyclables like metal and plastic. Among them are
price fluctuations that may make it cheaper for a company to ditch items than to recycle them,
contamination that makes a can or paper useless, and human error in sorting or transporting
material.

Even when an item is headed where it is supposed to go, ``does it fall off the boat, or truck, or
whatever?'' said Ms.~Landsberg, a transportation planner for Washington State. ``Is the stuff
actually made into something useful in this country? Does it all end up shredded and shipped to
China, where who knows what happens to it?''

To answer some of those questions, the M.I.T. team is using battery-powered tags based on cellphone
technology.

The researchers say it will take several months to analyze the data generated by the cellular
signals. But they have already noticed that while some trash reaches its destination in a couple of
days, other items may take four or five weeks to wind their way to landfills or recycling and waste
processing plants.

In Seattle, where researchers recruited volunteers for the project through the Seattle Public
Library's Web site, the Seattle Public Utilities newsletter and other local publications, about 500
pieces have been tagged. One item, an aluminum can disposed of at a residence, traveled 2.5 miles to
a recycling facility in the city in just under two days.

In New York, where 50 items were tagged at the Architectural League's offices, a recyclable plastic
bottle picked up at Madison Avenue and 51st Street traveled 18.3 miles over four days to Kearny,
N.J., and is still en route, said Assaf Biderman, associate director of the M.I.T. lab.

The tracking has its limitations. Even though the tags have a battery life of two to six months and
can report back from overseas, they can easily be crushed in transit inside garbage trucks and at
processing facilities. Mr.~Biderman said a paper cup taken from a Seattle residence sent signals for
seven and a half days before it went silent and is assumed to have been destroyed.

But the researchers say most tags are likely to travel far enough to show which items go where and
how long it takes them to reach a destination, yielding information about inefficiencies in the
waste management system. In coming weeks the project is expected to gain an international component
when 50 items are tagged in London, Mr.~Biderman said.

Ms.~Brown of Waste Management said her company hoped that the experiment could eventually help
shorten or avoid overlaps in routes traveled by its 24,000 garbage trucks and to find more central
locations for transfer and disposal.

Ultimately, she said, ``we're looking for ways to recycle more and to do it all more efficiently.''

Brett Stav, a senior planning and development specialist for the Seattle Public Utilities, which
collects about 21,000 tons of trash and recyclables a day, said that aside from the help with
logistics, he saw ``tremendous educational value'' in the experiment.

``There is this hidden world of trash, and there are ramifications to the choices that people
make,'' Mr.~Stav said. ``People just take their trash and put it on the curb and they forget about
it and don't think about all the time and energy and money put into disposing of it.''

The point is well taken by Ms.~Landsberg of Seattle, who is so environmentally conscious that she
keeps a worm bin to compost her food waste.

``If I found out that it wasn't going where I think it does, if it is less recycled than I hoped,''
she said she ``might think about buying less of it or doing without.''

``Maybe it is more about the reduce than the re-use,'' she said.

\section{Obama Says Chicago Is 'Ready' to Win Bid to Hold 2016 Games}

\lettrine{W}{ith} 16 days left until the International Olympic Committee chooses a host city for the
2016 Olympics, President Barack Obama stood on the South Lawn of the White House on Wednesday and
made a pitch for Chicago's bid to win those Summer Games. He promised that if the I.O.C. chose
Chicago, the city would make the United States -- and the world -- proud.

``Chicago is ready,'' Obama said during an event featuring Olympians, Paralympians and local
schoolchildren. ``The American people are ready. We want these Games.''

``I promise you, we are fired up about this,'' he said of the possibility of the Games being awarded
to Chicago, where he lived for nearly 25 years before moving into the White House.

As it stands, though, Obama will not travel to Copenhagen for the I.O.C. vote on Oct.~2, in which
Chicago will go up against the other finalists: Madrid, Rio de Janeiro and Tokyo. While other bid
cities will have their heads of state or other top leaders at the vote, Chicago will not. Obama said
he must remain at home to deal with the pressing issue of health care reform.

``The good news is I'm sending a more compelling superstar to represent the city and country we
love, and that is our first lady, Michelle Obama,'' he said as the crowd cheered.

Michelle Obama, a lifelong Chicagoan, will lead the United States contingent at the meetings. On
Wednesday, she showed the crowd charisma that just may win over some I.O.C. members.

After taking to the podium, she encouraged the audience to cheer and show its Olympic spirit. She
then poked fun at her husband's attempt at a few of the Olympic sports that were on hand, causing
the crowd to roar with laughter.

``You should have seen the president in there fencing,'' she said of her husband, who said he had
always wanted to try the sport. ``It was pathetic.''

As she spoke, Mayor Richard M.~Daley of Chicago; Stephanie A.~Streeter, the United States Olympic
Committee acting chief executive; and Patrick G.~Ryan, the bid leader, smiled.

Michelle Obama said Chicago was the ``ideal home for the 2016 Games,'' not just for its landscape,
infrastructure or resources, but also for its people and their love of sports.

``You know, you have to admit, even White Sox fans are impressed by the fact that even though the
Cubs haven't won a World Series in centuries, Cubs games sell out,'' she said. ``Everybody's there.
It doesn't matter. Win or lose, we are going to watch the Cubs.''

Later, she added: ``I know that Barack and I would feel such tremendous pride to see the Olympic
torch burning brightly in the city that we love so much.''

In the past, heads of state have played an important role in the selection of an Olympic host city.
In 2005, Tony Blair, the British prime minister, traveled to the I.O.C. meeting and encouraged
voters to choose London for the 2012 Summer Games. He was successful. In 2007, Vladimir V.~Putin,
then the Russian president, addressed I.O.C. members, who eventually picked Sochi as host of the
2014 Winter Games.

With that in mind, Daley said Michelle Obama's presence would be crucial in helping Chicago win.

``It is very important and it will make a huge difference,'' he said. ``She is going to be an
excellent representative for us and the United States. And it helps that she is from Chicago. She
only lives a walk away from some of the venues.''

\section{Baucus Offers Health Plan but Lacks G.O.P. Support}

\lettrine{T}{he} chairman of the Senate Finance Committee on Wednesday unveiled his long-awaited
plan to remake the nation's health care system and insure millions of Americans, but he did not win
support from a single Republican. His proposal is less costly and extends the reach of government
less than other major health bills moving through Congress.

The chairman, Senator Max Baucus, Democrat of Montana, said the committee would start voting next
week on his proposal, which embodies President Obama's top domestic priority.

The bill closely resembles what Mr.~Obama said he wanted, except that it does not include a new
government insurance plan to compete with private insurers. Given reservations expressed Wednesday
by members of both parties, the Baucus bill is likely to be modified as it lurches through the
committee to the Senate floor and then, presumably, on to negotiations with the House, which is
working on a more costly, more expansive measure.

Mr.~Baucus described his plan as ``a good balanced bill that can pass the Senate.'' But the
prospects that Congress will deliver a comprehensive bill to Mr.~Obama by the end of the year, as he
has requested, are difficult to gauge. In seeking a middle ground, Mr.~Baucus disappointed many
liberal Democrats. Speaker Nancy Pelosi said the House bill, drafted by Democrats, was superior in
many ways and ``clearly does more to make coverage affordable for more Americans.''

With the release of Mr.~Baucus's proposal, a huge fight over the role of government in American life
-- comparable to the struggle over creation of Social Security in 1935 -- moves to a new phase, in
which hundreds of lobbyists and millions of average citizens will be engaged.

The Congressional Budget Office said Mr.~Baucus's bill would still leave 25 million people uninsured
in 2019. About one-third of them would be illegal immigrants. By contrast, the budget office said,
17 million people would be left uninsured under the House bill. At least 46 million people are now
uninsured.

For more than three months, Mr.~Baucus has been negotiating with three Republican senators, and some
of their ideas are incorporated in his bill. He predicted that some Republicans and most Democrats
would eventually support a version of his bill.

Senator Olympia J.~Snowe, Republican of Maine and one of the senators Mr.~Baucus has been
negotiating with, said she was not ready to endorse the plan, but would continue working with
Mr.~Baucus.

The need for action is urgent, Mr.~Baucus said, declaring: ``Our health care system is simply
unsustainable. It's breaking the bank for everyone, from families to businesses to governments.''

Mr.~Baucus's bill is the least expensive of five major health care bills moving through Congress.
The Congressional Budget Office said the expansion of coverage would cost \$774 billion over 10
years, compared with price tags of more than \$1 trillion for the other measures.

The budget office said the cost of Mr.~Baucus's bill would be fully offset by new taxes and fees,
along with savings squeezed from Medicare, so it could reduce the cumulative total of federal budget
deficits by \$49 billion over the next 10 years. At the end of that period, it said, new revenues
and savings would be growing faster than costs, so the bill could also help reduce deficits in the
decade after 2019.

Like the other bills, Mr.~Baucus's proposal would require most Americans to have health insurance,
with penalties for those who flout the requirement. But his plan differs from the others in
significant ways:

¶Instead of creating a new government health plan, Mr.~Baucus would set up nonprofit insurance
cooperatives in every state. The Congressional Budget Office said the co-ops ``seem unlikely to
establish a significant market presence in many areas of the country.'' This finding provides
ammunition to liberals who say the co-ops could not compete effectively with big insurance
companies.

¶The Baucus plan, like the other bills, offers subsidies to help low- and middle-income people buy
insurance. But eligibility is more limited, and the subsidies appear to be less generous than in the
other proposals, causing some Democrats to suggest that many people could still find insurance
unaffordable.

¶Unlike the other bills, the Baucus plan would impose a new excise tax on insurance companies that
sell high-end policies costing more than \$8,000 for individuals and \$21,000 for families.
Mr.~Baucus hopes the tax would put downward pressure on health costs. The American Federation of
State, County and Municipal Employees denounced the tax, saying it would hit ``health plans that
provide benefits for many middle-class families.''

¶The bill would not require employers to offer coverage. But employers with more than 50 workers
would have to reimburse the government for some or all of the cost of subsidies provided to
employees who buy insurance on their own.

The legislation would reshape the \$2.5 trillion-a-year health care industry, which accounts for
roughly one-sixth of the American economy. It would significantly expand Medicaid and substantially
reduce Medicare payments to hospitals and many other health care providers.

The bill would create state insurance marketplaces, or exchanges, where consumers could shop for
insurance and compare plans. Insurers could not deny coverage or charge higher premiums because of
health status or pre-existing medical conditions.

The White House issued a low-key statement saying the Baucus plan provided ``another boost of
momentum for the president's effort to reform the health system.''

The mixed feelings among Democrats were illustrated in two statements issued by the Senate majority
leader, Harry Reid of Nevada, who is up for re-election next year. In one statement, Mr.~Reid
praised the Baucus plan. But in the other, he expressed concern that the bill would force his state
to increase its Medicaid spending, and he said the measure ``needs improvement.''

Conservatives savaged the Baucus plan. ``Democrats should scrap their expensive plans and work with
Republicans on a fiscally responsible proposal,'' said the House Republican leader, Representative
John A.~Boehner of Ohio.

Mr.~Baucus would extend benefits to millions of the uninsured by expanding Medicaid, at a federal
cost of \$287 billion over 10 years, and by offering subsidies to individuals and families, at a
cost of \$463 billion over 10 years. He would offer tax credits to help small businesses buy
insurance, at a cost to the Treasury of \$24 billion over 10 years.

The budget office said the proposed tax on high-premium insurance plans would generate \$215 billion
in revenue over the next 10 years. People who go without insurance would pay penalties totaling \$20
billion, it said, while employers who do not offer coverage would pay penalties of \$27 billion. A
family without insurance could face a penalty up to \$3,800 a year.

\section{Yale Student Was Strangled, Authorities Say}

\lettrine{A}{s} a high school student in the shoreline town of Branford, Conn., Raymond Clark III
joined the Asian Awareness Club, which made spring rolls for a faculty lunch and organized a trip to
Chinatown for the Chinese New Year. He joined the Interact Club, which focused on community problems
like homelessness. And he played football and baseball, throwing long bombs as a quarterback and
knuckleballs as a pitcher.

His sister found him a job less than 10 miles from where they grew up, in the laboratory at Yale
where she worked. He, too, became a technician, a grunt in the rarefied world of medical research,
cleaning lab animals' cages and doing custodial chores while graduate students and their professors
did experiments that held the promise of life-changing breakthroughs.

On Wednesday, Mr.~Clark, 24, was at the center of the investigation into the killing of one of those
graduate students, Annie M.~Le, whose body was found on Sunday, the day she was to be married,
behind a wall in the building where they had worked.

The police said he was not the only person they were looking at in connection with her disappearance
and death, which the authorities revealed on Wednesday had been caused by strangulation. But he was
the only one from whom they had taken DNA samples, and the only one whose car -- a red Ford Mustang
-- they had hauled away to search.

The police said he was cooperative, and they released him early Wednesday morning after about five
hours. ``He provided what we needed,'' said Chief James Lewis of the New Haven police.

Mr.~Clark's future may depend on whether crime laboratory technicians find a match between the
samples from Mr.~Clark and hundreds of items from the Yale lab. No matter how long it took to
analyze the samples -- 30 minutes, or a matter of days -- Chief Lewis said, ``if we have one match
on a person we know was at that location,'' the police would make an arrest.

Chief Lewis said that investigators had taken more than 200 items from the crime scene with
potential DNA evidence that could be compared with the samples from Mr.~Clark. Earlier in the week,
Chief Lewis said the police had interviewed 150 people and watched 700 hours of surveillance video
from cameras in and around the building where Ms.~Le's body was found -- and where Mr.~Clark worked.

A law enforcement official said that Mr.~Clark had scratch marks on his chest and on parts of his
arms, where there were also bruises. This raised the suspicions of investigators, said the official,
who declined to be identified because the case was still open.

The police said they did not question Mr.~Clark on Tuesday night. They had questioned him earlier,
Chief Lewis said, but he had invoked his right to a lawyer. He was represented by David H.~Dworski
of Fairfield, Conn., and the New Haven Public Defender's Office was also providing counsel.

The public defender's office declined to comment. Mr.~Dworski said in a statement, ``We are
committed to proceeding appropriately with authorities with whom we are in regular communication.''

By Wednesday afternoon, about 12 hours after Mr.~Clark had been released, Connecticut's chief
medical examiner released the cause of Ms.~Le's death, a finding he had withheld earlier in the week
at the request of the state's attorney in New Haven. A spokeswoman for the medical examiner,
H.~Wayne Carver II, said Ms.~Le had been strangled -- she died from ``traumatic asphyxia'' caused by
``neck compression.''

Chief Lewis would not say whether she had been sexually assaulted.

It remained unclear whether there was any connection between Ms.~Le and Mr.~Clark beyond the fact
that they spent time in the same building, Mr.~Clark working as a technician, Ms.~Le watching
experiments with mice. A spokesman for the New Haven police, Officer Joe Avery, said there was no
indication that she had complained of being stalked or harassed by Mr.~Clark or anyone else in
recent weeks.

Mr.~Clark, who lives in Middletown with his girlfriend, also an animal research technician at Yale,
had been a focus of the police for several days. On Wednesday afternoon, after removing several
boxes of items from Mr.~Clark's apartment, investigators were still knocking on his neighbors' doors
seeking information about him.

Chief Lewis refused to discuss a police report from 2003, when Mr.~Clark was a senior at Branford
High School and the police were called about a dispute with his girlfriend at the time. According to
the police document, reported Wednesday in The New Haven Independent, a news Web site, she wanted to
break up but he did not, and tried to confront her. She complained to school authorities, who called
the police, but did not press charges, according to the Web site .

Mr.~Clark grew up in a rented gray house in a working-class neighborhood of aspirations when a
nearby factory was humming. Jim Garrett, 65, who lives two doors down, said the house the Clarks
lived in deteriorated as the years went by and the factory closed, and eventually Mr.~Clark's
parents moved to a condominium in Cromwell, Conn., north of Middletown, where Mr.~Clark's mother
works in the Wal-Mart across the street.

At Yale, Mr.~Clark's work habits were described by one researcher as ``very officious and very
demanding.'' But Mr.~Garrett remembered Mr.~Clark as ``industrious and busy.'' More than once,
Mr.~Clark rang the doorbell with Mr.~Garrett's runaway beagle at his heels. ``Your dog was loose,''
Mr.~Clark would tell him.

Mr.~Clark was pictured several times in the Branford High School yearbook in 2004, his senior year.
The first photograph was with the Asian Awareness Club; the second was with the Interact Club. He
was also shown with the Cheers for Charity Club, which gave baby showers for two pregnant women.

Teammates remembered Mr.~Clark as a talented and versatile athlete who was deeply competitive. Conor
Reardon, 23, who was on the team with Mr.~Clark, said that he and other teammates got Mr.~Clark to
succumb to a baseball rite of passage, chewing tobacco.

``But the next day, he came to school late because it had made him so sick that when he got out of
his car that night he could barely stand up,'' Mr.~Reardon said, adding that Mr.~Clark ``ended up
falling asleep in the bath tub in his baseball uniform after taking a shower fully clothed.''

\section{As Race Debate Grows, Obama Steers Clear of It}

\lettrine{P}{resident} Obama has long suggested that he would like to move beyond race. The question
now is whether the country will let him.

He woke up on Wednesday to a rapidly intensifying debate about how his race factors into the broader
discussion of civility in politics, a question prompted in part by former President Jimmy Carter's
assertion Tuesday that racism was behind a Republican lawmaker's outburst against Mr.~Obama last
week as the president addressed a joint session of Congress.

Even before that, several conservatives had accused their liberal counterparts of unfairly tainting
them as racists for engaging in legitimate criticism of the White House.

Mr.~Obama's response to all this, aides say, has been to tell his staff not to be distracted by the
charges and to focus on health care and the rest of his policy agenda.

``He could probably give a very powerful speech on race, just as he did in the course of the
campaign,'' said Valerie Jarrett, a senior adviser to Mr.~Obama. ``But right now his top domestic
priority is health care reform. It's difficult, challenging and complicated. And if he leads by
example, our country will be far better off.''

During the presidential campaign, when he disavowed the incendiary remarks of his pastor, the Rev.
Jeremiah A.~Wright Jr., he took the opportunity to explain his views on race in America and invite
reconciliation. And after he stumbled in July in accusing the police of ``acting stupidly'' by
arresting the Harvard professor Henry Louis Gates Jr., he used the occasion for what he called a
teachable moment.

But this time the White House has made clear that it does not want to engage on the topic, which
beyond threatening to distract attention from the health care push could also put further strain on
Mr.~Obama's broad but tenuous electoral coalition of liberals and moderates, Democrats and
independents.

Signaling that he had no intention of lending his voice to Mr.~Carter's accusation, the president
declined to answer a reporter's question on the subject in the Oval Office on Wednesday.

And his press secretary, Robert Gibbs, told reporters again and again at the daily White House
briefing that Mr.~Obama did not share Mr.~Carter's views on the motivations of Representative Joe
Wilson, the South Carolina Republican who shouted ``You lie!'' during Mr.~Obama's address.

Even as several leading Congressional Democrats distanced themselves from Mr.~Carter's comments,
some liberals pointed Wednesday to what they describe as an increasing number of racially tinged
attacks. At last weekend's conservative protest in Washington, there were Confederate flags,
references on placards to sending Mr.~Obama to Africa and pictures of him in whiteface, as the Joker
in the last ``Batman'' movie. There are also, of course, the continuing questions from some on the
right about his United States citizenship.

But a number of prominent conservatives say critics have been smeared by many of the president's
supporters.

On his radio program this week, Rush Limbaugh said, ``Today, it's all based in racism -- the
criticism of Obama's health care plan or whatever.'' On Fox News, former Speaker Newt Gingrich
added, ``I think it's very destructive for America to suggest that we can't criticize a president
without it being a racial act.''

It is difficult to gauge the extent to which the vitriol that has come Mr.~Obama's way is racially
motivated and the extent to which it is simply akin to that directed at his white predecessors.

Former officials who served under President George W.~Bush have been quick to recall this week that
protesters frequently compared him to Hitler and that the Senate majority leader, Harry Reid of
Nevada, called him a ``loser'' and a ``liar.''

In an interview Wednesday shortly after meeting privately with Mr.~Obama in the Oval Office, Colin
L.~Powell, secretary of state under Mr.~Bush, said: ``You can find pictures where Bush was called
all kinds of names, with all sorts of banners being held up and burned in effigy. I've seen it in
every presidency.''

Mr.~Powell said he believed that Mr.~Obama might be facing even more apparent hostility but that the
blame lay not necessarily with racial bias, but instead with the partisan culture of the Internet
and cable news and the way they amplify the more extreme voices.

``The issue there is not race, it's civility,'' Mr.~Powell said. ``This is not to say that we are
suddenly racially pure, but constantly talking about it and reducing everything to black versus
white is not helpful to the cause of restoring civility to our public dialogue.''

Yet other supporters of Mr.~Obama say they cannot help but see overt racism in some of the
conservative attacks.

``You cannot act like you don't have several hundred years of racial context here, where a painted
face has a racial context to it in this country,'' said Cornell Belcher, a Democratic pollster who
helped on Mr.~Obama's presidential campaign and has studied race extensively.

Mr.~Belcher and other Obama allies said that some race-based discomfort was inevitable, especially
among very conservative white voters who see Mr.~Obama's rise as reflecting a shift in the social
order that comes at their expense.

Representative Elijah Cummings, Democrat of Maryland and former chairman of the Congressional Black
Caucus, said: ``It's difficult, because people haven't seen this before: they've never seen a black
president wanting to talk to their children, a black man saying, 'We can do better.' ''

Mr.~Obama has many top aides who are white and have spent years dealing with race in the context of
politics. He also has a close group of African-American advisers and friends for whom a racial
conversation like the one bubbling up this week is not an abstract issue but a way of life. There
have been occasional tensions between the two groups in the past, but on Wednesday, at least, there
were no obvious signs of disagreement.

His goal, Mr.~Obama has told both camps, is to be seen as a president who happens to be black rather
than the nation's first black president.

\section{For Senate Democrats, 60 Is the Magic Number}

\lettrine{T}{he} unveiling of a compromise health care proposal has Senate Democrats pondering a
daunting mathematical challenge: how to keep all 59 Democrats united and attract at least one
Republican to pass an overhaul measure.

As many lawmakers on Wednesday got their first detailed look at a Finance Committee plan that was
months in the making, senators immediately began exploring whether the plan -- when combined with
elements of another, more liberal one -- could win enough senators to reach that magic procedural
number of 60.

``We have to meld a couple of things together and see where we are,'' said Senator Christopher
J.~Dodd, Democrat of Connecticut. ``I wouldn't say today with absolute certainty that you could get
to 60, but it would be just as foolish to say you can't get there either. This is the Senate.''

In trying to reach critical mass for legislative success, advocates of health care overhaul face an
extremely delicate balancing act. With the death of Senator Edward M.~Kennedy of Massachusetts,
Democrats control 59 seats, meaning they need at least one Republican to join them if they are to
proceed without employing a procedural shortcut that could cause havoc in the Senate.

And Senate Democrats have substantial differences of their own. More liberal members like Senator
John D.~Rockefeller IV of West Virginia and Senator Bernard Sanders, independent of Vermont, have
been strong advocates of a public insurance option; many of the more centrist Democrats have come
down just as strongly against one. Centrists are interested in holding costs down; progressives want
to bring more Americans under the insurance umbrella, a push that drives costs up.

As a result, changes intended to bring the centrists and conservatives in line could drive away
progressives, while any move to draw in the more liberal elements could end up alienating the
centrists. There is little margin for error.

``I continue to believe that you start at the center and then move to collect additional support on
both ends of the political spectrum,'' said Senator Ron Wyden, an Oregon Democrat and health policy
expert whose seat on the Finance Committee makes him a key player as that panel considers the new
plan from Senator Max Baucus, the Montana Democrat who leads the committee.

Juggling the individual party demands can be dizzying. Another Finance Committee member, Senator
Bill Nelson of Florida, is digging in against potential Medicare cuts. Senator Bob Casey of
Pennsylvania wants to make sure the plan does enough for poor children, while Senator Evan Bayh of
Indiana is worried about health costs for small businesses and middle-class families.

And Senator Ben Nelson of Nebraska, typically one of the hardest votes for Democratic leaders to
corral, is looming as a particularly tough sell. ``At the end of the day, I want to see everything
before I commit to anything,'' said Mr.~Nelson, who added that he would have trouble backing a bill
that did not have some Republican support.

Yet winning over just one Republican will take extraordinary effort, partly because of the intense
pressure from colleagues to deny Democrats a legislative victory. Senator Olympia J.~Snowe of Maine,
the Republican viewed as the most likely to back the proposal, did not endorse the plan on
Wednesday, though she is talking to Democrats.

Other Republicans considered potential allies under the right circumstances have their own
reservations. Senator Susan Collins, another Maine moderate who has joined with Democrats in the
past, has not seen enough emphasis on cutting health care costs over all.

Senator George V.~Voinovich of Ohio, who is retiring and thus may face fewer political constraints,
said he feared the expansion of coverage for the poor could bankrupt states. ``To be candid, I don't
know how you pay for it,'' he said.

Given the steep climb toward 60, Senate Democratic leaders have begun to make another argument to
lawmakers. They are pressing colleagues to vote with the party on procedural matters related to
health care legislation and against any filibuster -- a 60-vote issue -- even if they intend to
oppose the measure in the end when simple majority rules.

Senators are usually reluctant to clear the way for a bill they might vote against since they
relinquish their most powerful leverage, but the message is evidently reaching some.

``It is difficult to ask someone to facilitate the enactment of legislation with which they
disagree,'' Mr.~Bayh said. ``But to move the process forward, to improve things, to get to the point
where you can support it substantively, that of course I would be willing to do.''

Even with all the policy disputes compounded by political and procedural hurdles, Democrats said
they still believe they can get a bill through. But it will not be easy.

``Do I believe there is enough consensus around here to get a bill done in the end? Yes,'' Mr.~Casey
said. ``But I also believe it is going to be a difficult couple of weeks.''

\section{Recovery Picks Up in China as U.S.~Economy Remains Ailing}

\lettrine{J}{ust} eight months ago, thousands of Chinese workers rioted outside factories closed by
the global downturn.

Now many of those plants have reopened and are hiring again. Some executives are even struggling to
find enough temporary staff to fill Christmas orders.

The image of laid-off workers here returning to jobs stands in sharp contrast to the United States,
where even as the economy shows signs of improvement, the unemployment rate continues to march
toward double digits.

In China, even the hardest-hit factories -- those depending on exports to the United States and
Europe -- are starting to rehire workers. No one here is talking about a jobless recovery.

Even the real estate market is picking up. In this industrial town 90 miles northwest of Shanghai,
prospective investors lined up one recent Saturday to buy apartments in the still-unfinished Rose
Avenue complex. Many of them slept outside the sales office all night.

``The whole country's economy is back on track,'' said Shi Yingyi, a 34-year-old housewife who
joined the throng\cite{throng}. ``I feel more confident now.''

The confidence stems from China's three-pronged effort -- a combination of stimulus, liberal bank
lending and broad government support for exports.

The Chinese central bank said the country's economy surged\cite{surge} at an annualized rate of 14.9
percent in the second quarter. The United States economy shrank at an annual rate of 1 percent in
that period.

``So often China and the U.S.~are mixed together as being in the same situation, and that is totally
wrong,'' said Xu Xiaonian, an economist in Beijing with the China Europe International Business
School.

That does not mean the two nations are not connected, of course. China's rebound in growth may slow
if the American economy does not pick up. China needs the United States to buy its goods, and the
United States needs China to continue to buy its debt.

This mutual dependence makes it harder for either country to let the current dispute over Chinese
tires and American chicken and auto parts to grow into a trade war.

But with more economic planning than the United States, China has been able to disburse its stimulus
much faster, turning it into new rail lines and highways.

China's finance ministry announced in late June that half the \$173 billion in central government
spending had already been allocated to specific projects. The White House said in early July that a
quarter of the spending authority and tax cuts in the \$789 billion American stimulus had been
allocated or used.

But even more key to China's recovery, economists say, are two other government efforts that are
paying big dividends: looser lending and export supports.

The state-controlled banking system here -- which breezed through the global financial crisis with
minimal losses as American financial institutions reeled -- unleashed \$1.2 trillion in extra
lending to Chinese consumers and businesses in the first seven months of this year. That money is
financing everything from a boom in car sales, up 82 percent in August from a year earlier, to
frenzied\cite{frenzied} factory construction.

Beijing also has given huge tax breaks and other assistance to exporters. They include placing broad
restrictions on imports and intervening heavily in currency markets to hold down the value of the
renminbi, to keep Chinese exports competitive even in a weakened global economy.

Indeed, subsidies abound at all levels of government: the Wuxi municipal government just offered up
to \$146,000 to each local business that increases exports in the last three months of this year.

To be sure, not all the laid off workers throughout China have been hired back.

``Some plants reduced worker numbers by 20 to 30 percent, now they hire back 10 percent,'' said
Stanley Lau, deputy chairman of the Federation of Hong Kong Industries, which represents
export-oriented factories employing 10 million Chinese workers.

Even so, American trade data shows that imports from China only eroded\cite{erode} 14.2 percent in
the first seven months of this year while imports from the rest of the world plunged 32.6 percent.
China's trade surplus, already the world's largest, was \$108 billion for the seven-month period.

``We definitely see an upswing in sales orders in the second half of this year when compared to the
first half,'' said Gu Fung, the sales manager at the Wuxi Baolai Batteries Company.

China's well-capitalized banking system allowed for rapid investment.

Chinese banks came into the crisis with enormous excess reserves, the result of three years of tight
regulatory limits on lending to prevent the economy from overheating. When those limits were
removed, and authorities urged bank executives to lend, the total value of loans outstanding shot up
more in the first seven months of this year than in the previous 24 months.

By contrast, total loans and leases outstanding at financial institutions insured by the Federal
Deposit Insurance Corporation actually fell \$249 billion, or 3.2 percent, in the first half of this
year.

Though Washington has used taxpayer money to bail out American banks, it does not have Beijing's
power to force banks to lend that money to businesses and consumers.

As much as a third of the extra bank lending in China appears to have gone into real estate and
stock market speculation. But the bulk has gone into investments by companies and local governments,
with tangible\cite{tangible} results.

China's currency and trade policies, though highly effective, would be hard for the United States to
emulate.

For instance, government intervention in currency markets has prevented the renminbi from moving
appreciably against the dollar in more than 14 months, and has pushed the renminbi down by 18
percent against the euro since March.

Government agencies have been told not to buy imported goods with money from economic stimulus
programs unless no domestic alternative is available. Washington has imposed a less restrictive
rule, misleadingly known as ``Buy American,'' requiring that construction materials for the stimulus
program be bought from any of the 39 countries that have agreed to free trade in government
procurement -- which China has not.

Still, China's stimulus efforts could be sowing the seeds of future distress. With so much money
washing into the system so fast, regulators have voiced concerns about corruption in government
investment projects.

Cheap cash has a way of inflating bubbles -- just ask Wall Street -- that could damage China's
economy and its banks when they pop.

``You have to imagine the rigor\cite{rigor} and due diligence'' that mainland banks have been
showing in rushing out so many loans, said Benjamin Hung, the chief executive of the Hong Kong unit
of Standard Chartered Bank.

But such concerns are so 2008.

\section{Lab Technician Arrested in Murder of Yale Student}

\lettrine{B}{efore} there was blood, the high-tech lab at 10 Amistad Street at Yale University was a
model of efficiency. The mice and rats and rabbits stayed locked in cages. The technicians
responsible for their well-being circulated like emergency room nurses. Researchers hunched over the
cages for hours, intent on claiming a breakthrough.

The two groups interacted in professional if perfunctory\cite{perfunctory} ways, but on Thursday,
the authorities charged a technician, Raymond Clark III, with murdering one of the graduate
researchers, Annie M.~Le. Ms.~Le, 24, was strangled on Sept.~8, and her body was found on Sunday
hidden behind a wall, out of view from the immaculate\cite{immaculate} corridors of the lab.

Mr.~Clarke, 24, was arrested just after 8 a.m. Thursday in Room 214 of the Super 8 Motel in
Cromwell, Conn. He had been staying there with his father, at the end of several days in which the
authorities interviewed him, tailed him, took DNA samples and then kept him under surveillance. He
was charged with murder and driven back to New Haven, where he was arraigned but said little and did
not enter a plea. Bail was set at \$3 million.

The authorities said his DNA matched crime scene evidence, but did not elaborate.

Chief James Lewis of the New Haven police would not speak about a possible motive, but said, ``It is
important to note that this is not about urban crime, university crime, domestic crime, but an issue
of workplace violence, which is becoming a growing concern around the country.''

The chief sought to dispel any notion that Ms.~Le, who was about to be married, had been stalked by
Mr.~Clark or that Yale itself was unsafe. But the arrest opened a window into a
peculiar\cite{peculiar} work environment, populated by thousands of animals, driven researchers and
the technicians who perform the lab's menial but essential work.

Those technicians are given a special order: to serve as advocates for the animals and guardians of
regulation about how they should be treated.

``There is a certain stress that builds with the job,'' said David Russell, who worked as an animal
technician at Yale from 1997 to 2008. ``If there's something wrong, you are the one who is on the
hook.''

They come from a variety of backgrounds: former veterinary technicians; laid-off workers from
pharmaceutical companies; men and women fresh out of high school and college and looking for a
decent-paying job.

The jobs are competitive, and many get through the door with the help of a friend or relative.
Mr.~Clark's brother-in-law and sister also work as animal technicians, and she recommended Mr.~Clark
for a position in the washing center in 2004, the year he graduated from high school in nearby
Branford.

The university asks that technicians have familiarity with animals. Mr.~Clark confided in one
co-worker that he had listed on his r\'esum\'e that he had worked on a farm, even though he had not,
the co-worker said. The co-worker spoke on the condition of anonymity because Yale officials had
instructed employees not to speak with the news media.

Yale's Web site says it conducts criminal background checks of its employees, and since 2007, it has
required all educational and employment credentials be verified.

With its cutting-edge facilities, the Amistad building, which opened in 2007, is a place technicians
dream of working. It is home to about 4,000 mice alone, Mr.~Russell said, on a campus that also
keeps hamsters, gerbils, cats, dogs, pigs, sheep, fish and monkeys.

The washroom job is considered one of the toughest positions. It involves scraping dirty cages and
loading them onto a conveyor-belt washer, and lifting 40-pound bags of food and bedding.

On a daily basis, technicians must also make the rounds looking for green neon tags -- the mark
signifying that an animal needs to be euthanized. They take the animals to the basement, lock them
in a cage, and turn on the carbon dioxide machine.

``It is very easy to get attached to the animals,'' Mr.~Russell said. ``It wears on you.''

Mr.~Clark's co-worker said the technicians ``definitely do get a little desensitized.''

``But I don't know anyone who is bothered and upset on a daily basis,'' the co-worker said. The
university provides counseling to help employees cope with having to kill animals on a regular
basis.

Mr.~Clark was eventually promoted to full-time animal technician, a position that pays \$12 to \$25
per hour. That job required him to serve as caretaker for the animals, moving them into clean cages
and checking to make sure they were not dehydrated, sick or dead.

Animal technicians must also be watchdogs, making sure that in the bureaucratic world of animal
research, all documents have been filed and all ethical standards obeyed. They might remind a
student to put on a gown\cite{gown} before entering a room, or chide\cite{chide} a researcher for
failing to separate a litter of mice or clipping a mouse tail for a DNA sample, a practice the
university forbids.

They live in fear of being held responsible for somebody else's sloppiness; a single lapse like a
dehydrated animal or unsanitary work space could mean weeks of disciplinary hearings.

``A lot of them tend to view us as janitors\cite{janitor},'' the co-worker said. ``But we're more
than that. We are policemen. We are there to make sure everything is done humanely and ethically.''

Some thought Mr.~Clark went beyond the bounds of his position. A research team leader in the Amistad
building, where Ms.~Le's body was found, said that several of his researchers complained last year
that Mr.~Clark was rude to them, prompting the team leader to alert Mr.~Clark's supervisor.

``He would berate them for minor infractions,'' said the man, who requested anonymity. ``Everyone
enforces rules, but he enforced them in an officious manner.''

In a statement, Richard C.~Levin, the Yale president, said the supervisor ``reports that nothing in
the history of his employment at the university gave an indication that his involvement in such a
crime might be possible.''

Mr.~Russell and other technicians said Mr.~Clark had significant authority in his position, and said
it would not be strange for somebody in his position to directly call or send e-mail to students
about their lab work. He estimated there were about 200 employees in two dozen animal labs at Yale,
and that supervisors, particularly Mr.~Clark's, did not have time to get involved in the day-to-day
operations of the labs.

While most technicians do not forge close relationships with students, they do get to know each
other. They can spend anywhere from a few minutes to a few hours together daily, depending on how
many animals are in the lab room.

Students said that relations between researchers and technicians were cordial\cite{cordial}, for the
most part, although the technicians were clearly there for the more tedious aspects of any
scientific investigation.

``They clean and maintain lab supplies and prepare chemicals -- it's a job,'' said Frank Liu, a
postdoctoral associate in Yale's School of Medicine. ``We don't have that much interaction with
them.''

On Thursday, the authorities provided no details about what had prompted the killing. The arrest
warrant -- which Chief Lewis said ran to more than 1,000 pages -- was sealed for two weeks because
the investigation was continuing.

For Ms.~Le's family, the arrest brought some answers after nine days of uncertainty, going back to
when Ms.~Le was reported missing and her purse was found in her office, in another Yale building a
few blocks from the lab building. She was about to be married to Jonathan Widawsky, a graduate
student at Columbia she met when they were undergraduates at the University of Rochester.

Dennis Smith, the pastor\cite{pastor} of the New Haven Seventh-Day Adventist Church, who has been
speaking for the Le and Widawsky families, appeared on the ``Today'' show on Thursday and called the
arrest ``wonderful news.''

``They are very appreciative of the thoughts and prayers of everyone during this terrible time of
grief and loss,'' he said in an e-mail message later in the day.

Mr.~Clark came under suspicion last week, partly because of scratches on his chest and scratches and
bruises\cite{bruise} on his arms. Wearing khaki pants and a striped polo shirt that exposed tattoos
that circled both his arms, he appeared in court with two public defenders.

Mr.~Clark kept his head down during the five-minute arraignment in Superior Court in New Haven. When
Judge Jon C.~Blue asked Mr.~Clark to acknowledge that he had been read his constitutional rights,
Mr.~Clark nodded and said quietly, ``Yes, sir.''

Mr.~Clark was taken to a state prison described on its Web site as a ``high/maximum security level
multimission facility'' for men, including those who need ``pretrial protective custody.''

Many Yale students and employees seemed shaken that one of their own was accused of the crime. In
his statement, Mr.~Levin, the Yale president, echoed the police chief's concern about workplace
violence.

``We must reaffirm our deepest values as an institution -- our commitment to the search for truth,
undertaken in a spirit of openness, tolerance and civility,'' he said. ``The work of the university
requires us to engage with each other in the classroom, to collaborate in the laboratory and to
trust one another in workplaces across the campus.''

The killing ``could have happened in any city, in any university,'' he said. ``It says more about
the dark side of the human soul than it does about the extent of security measures.''

\section{Fans Hold On to 16-Inch Dreams}

\lettrine{F}{OREST} PARK, Ill. -- If you have never held a softball big enough to fill both hands
and then some, a ball that might mangle or at least jam a few fingers if heaved hard enough, a ball
that turns as mushy as an old shoe after enough whacks from a bat, you are probably not from around
here.

A small suggestion: Keep thoughts of wearing a glove to yourself.

For Chicagoans and those from its near suburbs, like this one, softball is a game played gloveless,
with a fat ball, 16 inches in circumference, not the ordinary 12-inch sphere found most everywhere
else. It has been this way here since sometime after 1887, when, as the legend goes, people in
Chicago fashioned a ball from a puffy boxing glove and a bat from a broomstick.

By now, the distinction is less a question of taste for some small segment of amateur athletes than
a matter of identity, connection, tradition in a city that revels in all it has created -- from
pizza to skyscrapers -- through impromptu breaking of rules.

``A hallmark of many of Chicago's greatest contributions to art and popular culture and sport
involves taking something and doing it differently,'' said Tim Samuelson, the city's cultural
historian, who, like many others here, carries childhood memories of 16-inch softball -- in his
case, a fly ball headed toward him in a schoolyard on the city's far North Side.

``You really want to catch it,'' Mr.~Samuelson said, ``but there it is, this behemoth of a ball, and
there's part of you that says, 'Uh oh.' ''

For more than a decade, people here had worked to create a Hall of Fame dedicated to the 16-inch
game, but even as people were being inducted into the hall, it long existed mainly as a concept, not
a place. At last, this summer, Inductee Park was built here, 11 miles west of Chicago's Loop.
Leaders of the Hall of Fame, with help from a local state legislator, have dreams of a museum
dedicated to the region's beloved sport nearby, in what village officials said was once a carwash.

In recent years, some had come to fear that 16-inch softball might disappear. The players -- who
number well into the tens of thousands -- began looking older, balder, rounder. Some schools and
clubs were teaching only baseball (or even ordinary softball), and computer games were booming.

``The game was dying out,'' said Rickey R.~Hendon, a state senator, and founder, manager and
starting pitcher for Senator Hendon's Fighting Fifth, a 16-inch team, who has pledged to find
\$500,000 in state financing for the museum, despite the state's budget crisis.

Al Maag, a co-founder of the Hall of Fame, said his hope, once the museum was built, was to keep the
sport alive (there are signs of its perking up, he says) and spread it far beyond Chicago.

``I'd like to see it played globally,'' said Mr.~Maag, 60, who grew up playing it and says his best
friends remain those who came through softball.

In a way, the sport's mystique is its reflection on the city: an unfussy sport for an unfussy place.
Though the game was invented years before, its popularity soared during the Depression. It was far
less expensive than ballgames that required gloves and other gear, and did not seem to rule out
those with ample bellies or graying hair. ``Anybody can play this game,'' said Tim Collins, as he
prepared to play on a recent afternoon.

All this took was a ball and bat, and, because of the heft of the ball, it could be played in
cramped, neighborhood parks. It is not easy to hit a 16-inch ball out of the park, any park.

Among better-known Chicagoans of yore with ties to this game are names rarely linked to athletic
achievement: Mike Royko, the sharp-tongued, Pulitzer Prize-winning newspaper columnist; Richard
J.~Daley, the former mayor and father of the current mayor (pictured in Inductee Park offering a
classic, underhanded, arching pitch while wearing a suit, tie and dress shoes); and Machine Gun Jack
McGurn, a member of Al Capone's gang who was the sponsor of a team. (Not surprisingly, there is said
to have once been betting on 16-inch around here.)

For a time, the sport bloomed and spread, gaining recognition over its better-known 12-inch cousin
in the 1960s and '70s. It gained popularity in pockets in Indiana, Wisconsin, Iowa, California,
Arizona and elsewhere, but Chicago remained its loyal home. (In one national tournament played since
1963, Chicago teams have won each year but one.) As the game's reach expanded, tense debates emerged
over whether to glove or not to glove.

Even in Chicago, gloves are now allowed for some, but many purists view them as anathema -- or, in
the words of one Chicago player, something ``no man with any pride would wear.''

``The whole magic behind this is having no glove,'' said Anthony T.~Calderone, the mayor of Forest
Park, which is home to the ``No Glove Nationals.''

This summer, in the fields of Grant Park in downtown Chicago, 81,000 people played 16-inch softball
in amateur leagues for lawyers, hotel workers, accountants, and on and on.

``I grew up watching my dad play here, and I hope to see my kids do the same,'' said Brian Duellman,
32, as he waited for his game to begin on a recent evening he described as ``perfect, the epitome of
16-inch.''

\section{House Votes to Ban Federal Funds for Acorn}

\lettrine{T}{he} House voted Thursday to deny federal money to the community-organizing group
Acorn\cite{acorn} after a video emerged in which employees of the group gave advice to two
conservative activists posing as a prostitute and a pimp who said they wanted buy a house in
Baltimore and start a brothel\cite{brothel}.

Republicans added the prohibition to a Democratic bill on college lending by a bipartisan vote of
345-75, showing that Acorn, the Association of Community Organizations for Reform Now, was
increasingly seen by Democrats as a political liability.

``Acorn has violated serious federal laws, and today the House voted to ensure that taxpayer dollars
would no longer be used to fund this corrupt organization,'' Representative Eric Cantor of Virginia,
the No.~2 House Republican, said.

Bolstered by conservative media outlets, Republicans have been on a crusade against the organization
for months, accusing it of improperly influencing elections and being protected by Democrats who
were the beneficiaries of the group's political activities. Matters came to a head in the past week
after the Acorn workers were videotaped.

The Census Bureau then announced it would not use Acorn in conducting the 2010 Census, and the
Senate voted on Monday to deny the group, which has received an estimated \$53 million in federal
aid since 1994, any housing or transportation money.

Despite that vote, Senator Michael Johanns, Republican of Nebraska, on Thursday offered the same ban
to a measure covering spending by the Interior Department.

The problem, according to Senator Dianne Feinstein, the California Democrat in charge of the bill,
was that there is nothing in the measure that sends money to Acorn.

But Ms.~Feinstein said she was willing to add his language to the measure without a roll call vote,
which she assumed would turn out similarly to the 83-7 vote against Acorn on Monday.

Mr.~Johanns, recognizing the political benefits in a vote, insisted on a roll call tally. ``This is
such an important issue,'' he said. ``This is an issue that people all across the country are
watching the Senate floor.''

He also said he was preparing a measure that would ban all federal funds for Acorn to avoid having
to attach the ban to a dozen different spending bills.

In the House, the 75 votes against the prohibition all came from Democrats. One opponent,
Representative Jerrold Nadler, Democrat of New York, said he considered the prohibition to be
unconstitutional.

Officials at Acorn have called the videotaped incidents isolated, begun their own investigation and
said workers at other offices reported the conservative imposters to local police.

The ban on funds was included in a measure that would still have to pass the Senate to become law.
The Senate prohibitions are also weeks if not months from becoming law.

House Republicans used one of the few procedural weapons available to the minority party to force
the vote. Earlier Thursday, Speaker Nancy Pelosi, Democrat of California, said the Appropriations
Committee should determine whether the group was entitled to federal aid.

``It's up to the Appropriations Committee to scrutinize Acorn or any other group that receives money
from the federal government that has such an allegation against it,'' she told reporters.

With their victories on Capitol Hill, Republicans quickly sought to turn the issue toward the White
House.

``Though today's vote indicates that the writing's on the wall for Acorn,'' Representative John A.
Boehner of Ohio, the Republican leader, said, ``President Obama must indicate whether he will join
the Congress in taking decisive action to break all government ties with this corrupt
organization.''

\section{Plan for a Kennedy Successor Advances}

\lettrine{A}{fter} hours of testy debate, the Massachusetts House of Representatives on Thursday
approved legislation allowing Gov. Deval Patrick to appoint an interim successor to Senator Edward
M.~Kennedy.

The House voted 95 to 58; the measure now goes to the State Senate, which could take up the proposal
on Friday. But Republicans have threatened to try to delay a vote, and, through procedural
maneuvers\cite{maneuver}, could do so for several days.

Mr.~Patrick, a Democrat, has said that if both chambers approved the measure, he would appoint a
temporary successor to Mr.~Kennedy within days. The appointee would serve until a special election
on Jan.~19, and could play a crucial role in the fate of the health care bill making its way through
Congress.

With the Kennedy seat empty, Democrats in the Senate are not assured the 60 votes necessary to pass
the legislation. They could use procedural shortcuts to get the bill passed, but that could cause
havoc\cite{havoc} in the Senate.

Therese Murray, a Democratic and president of the State Senate, has remained publicly noncommittal
on the proposal despite calls from the Obama administration and from Victoria Reggie Kennedy,
Mr.~Kennedy's widow. The State Senate, also heavily Democratic, is said to remain divided on the
issue.

Under current law, the Kennedy seat would remain empty until the special election. But shortly
before his death last month, Mr.~Kennedy asked the legislature to change the law and allow
Mr.~Patrick to name a temporary successor, arguing in a letter that Massachusetts needed full
representation in Congress.

Republicans and even a number of Democrats have attacked the proposal as overly partisan. Governors
here had the power to fill Senate vacancies until 2004, when the Democratic majority in the
legislature changed the law to require a special election. Democrats worried then that if Senator
John Kerry were elected president, Gov. Mitt Romney would appoint another Republican.

Representative Elizabeth Poirier, a Republican, said that if the legislature approved the change, it
would damage its reputation and become ``a pawn\cite{pawn} in the Washington game.''

But Representative Cory Atkins, a Democrat, said Massachusetts could not afford to let Mr.~Kennedy's
seat stay empty even for a few months.

``We are truly outnumbered on the floors of the Capitol in Washington, D.C.,'' she said. ``We need
every single vote.''

The appointee would not be bound by the legislation to stay out of the special election, but
Mr.~Patrick has said he would seek an ``explicit personal commitment'' from the appointee not to
seek Mr.~Kennedy's seat permanently.

Mr.~Patrick has been recovering from hip surgery for the last week and has not publicly discussed
possible candidates for a temporary appointment. But some Democrats have been lobbying for Michael
S.~Dukakis, the former governor and 1988 presidential nominee, who now teaches at Northeastern
University.

Other possible candidates include Scott Harshbarger, a former state attorney general, and Paul G.
Kirk Jr., chairman of the board of the John F.~Kennedy Library Foundation in Boston.

\section{Why Don't Doctors Wash Their Hands More?}

\lettrine{L}{ike} most medical and nursing students, my introduction to operating began not with
scalpels\cite{scalpel} and sutures\cite{suture} but with scrub\cite{scrub} sinks and
sterile\cite{sterile} gloves. Over the course of an entire day, my classmates and I learned the
meticulous art of scrubbing up -- generating a thick bacteria-killing lather from our fingertips to
above our elbows -- and maintaining sterile technique -- laying down and picking up sundry items
without contaminating a sterile field. Even during practice runs, the nurses and doctors modeled the
ritual of scrubbing up then donning sterile garb with such mindfulness and consistency that every
student believed it had to be important, even if it was, in fact, just washing our hands and getting
dressed.

That gravitas and germaphobic hypersensitivity sometimes led to situations bordering on slapstick.

One morning early on in medical school, I found myself assisting on a challenging
abdominal\cite{abdominal} operation. I came to the operating table only after having scrubbed my
hands and forearms to a bright pink and allowed the nurses to envelope my torso\cite{torso} and arms
in a layer of impenetrable Gore-Tex and latex sterility. Keenly aware of all the opportunities for
contaminating my newly sterile status, I trundled over to the sleeping patient and stood stiffly,
but quietly, next to the resident and across from the senior surgeon and scrub nurse.

The patient's abdomen was deep, and the surgeon and resident struggled to see what they were doing.
Despite their best efforts, loops of glistening bowels kept slithering into the operative field.
After craning his neck and twisting his torso yet again, the senior surgeon looked up and motioned
with his head toward the ceiling. ``Can someone move the lights?'' he asked. ``It's dark down in
here.''

I looked for someone to volunteer, but both doctors had their hands deep within the patient's belly,
and the scrub nurse was focused on preparing instruments for the surgeon's next move. My own left
hand was curled tightly around a retractor, but my right hand was free.

I reached for one of the lights.

``She's contaminated!'' shrieked the scrub nurse as I brushed my hand against the unsterile light
rim.

Horrified, I slapped my right hand against my uncovered, unscrubbed forehead.

``She did it again!'' boomed the surgeon, his eyes widening with panic.

I was sure that the surgeon and the scrub nurse would pounce on me next, their anxiety was so
palpable. As I staggered back, my oversized sterile gown wrapped and twisted around my legs. I
started to fall. In a desperate effort to break my descent, I reached for something to hold on to. I
grabbed the resident's gloved hand.

``Now she's contaminated me!'' the resident cried out, jumping away from the table and twisting his
hand free.

Within seconds the circulating nurse, the anesthesiologist and an operating room technician had
hustled me over to a corner of the room. Standing in a semi-circle around me, they clucked their
reprimands: You've got to be careful! Remember the sterile field! Just step away if you're
contaminated! Didn't you pay attention? And, perhaps most humiliating: You medical students never
learn!

If I hadn't known it before, I knew it then. Hand hygiene and sterile technique are so successfully
maintained in operating rooms not because of the reminders that hang over scrub sinks, but because
it is part of the culture and identity of those who work there. No self-respecting surgeon, nurse,
anesthesiologist or technician would ever dream of breaching those sterile protocols in the surgical
suites. Or of allowing any deviation from the aseptic norms to simply pass.

But such enthusiastic devotion to hand hygiene does not exist outside the operating room. And again
and again in discussions about quality and safety and the terrible infections that can ensue, one
issue continues to bedevil the patient-doctor relationship yet defies all reason: why don't doctors
wash their hands more?

Over the last 30 years, despite countless efforts at change, poor hand hygiene has continued to
contribute to the high rates of infections acquired in hospitals, clinics and other health care
settings. According to the World Health Organization, these infections affect as many as 1.7 million
patients in the United States each year, racking up an annual cost of \$6.5 billion and contributing
to more than 90,000 deaths annually.

If we can get every health care worker to be nearly maniacal about hand hygiene in surgery, the same
must be possible for the health care system as a whole. But it may require approaching the problem
as we do in the operating room.

Last week the Joint Commission, the most important hospital accrediting agency in the United States,
established a new center that would, in its first national project, attempt to change the way health
care organizations approach hand hygiene. Since last spring, the Joint Commission's Center for
Transforming Health Care has worked with eight hospital systems across the country and in
collaboration with private industry to identify the causes of hand hygiene failure and to institute
innovative solutions, like sensor tags that beep when a caregiver entering a patient's room has
failed to use sanitizing gel. Because every one of the eight hospital systems has its own specific
constellation of obstacles, each has been testing a different set of solutions. The Joint Commission
hopes that these varied solutions, once proven effective, will meet the needs of the thousands of
hospitals and health care organizations it accredits across the country.

``Health care organizations have been working on hand hygiene for a long time,'' said Dr.~Mark
R.~Chassin, president of the Joint Commission and the former Commissioner of the New York State
Department of Health. ``But no one has solved the problem to the degree that we all want. What has
been impressive so far about the Center's work is that hospitals, starting from the C.E.O.'s, have
come together in this new effort.''

Up until this point, the Joint Commission has focused primarily on identifying the highest priority
safety and quality problems. ``We put them in front of health care organizations and say, 'You have
to work on this,' '' Dr.~Chassin said. ``But what the Joint Commission has not done until now is
work with organizations to develop interventions that show them how to solve the problem, how they
can aspire to the highest possible levels of excellence.''

Although still in its early stages, the hand hygiene project has already made an impact in the eight
collaborating hospital systems. Average hand hygiene compliance rates have increased from less than
50 percent at baseline to 74 percent. ``The goal in the future,'' Dr.~Chassin said, ``is to be
consistently above 90 percent.''

The Joint Commission's approach could recapitulate the surgical experience. If successful, hand
hygiene could become as important to how all health care workers define their work and themselves as
sterile technique is to those who work in the operating room. And these changes could in turn
transform the patient-doctor relationship.

``There's an analogy that is not true in health care now but that I hope one day will be so,''
Dr.~Chassin reflected. ``When we get on an airplane, safety and good outcomes are someone else's
job. We don't go poring over accident and maintenance records. We don't go looking for judgments on
the pilots' safety records. All of that is guaranteed by someone else, so we are free to focus on
much more personally relevant issues like scheduling or layovers.''

Even more to the point, however, this effort may be the American health care system's best chance at
solving an intractable problem in the patient-doctor relationship. If the national agency that
defines excellence and standards of quality for more than 90 percent of hospitals cannot make good
hand hygiene a pivotal part of the culture and identity of health care, then who can?

\section{Unmasking the Ghosts}

\lettrine{T}{he} scientific integrity of medical research has been clouded in recent years by
articles that were drafted by drug company-sponsored ghostwriters and then passed off as the work of
independent academic authors.

Yet the leading medical journals have continued to rely largely on an honor system of disclosure to
detect such potential bias, asking authors to voluntarily report any industry ties or contributors
to their manuscripts.

But now, in light of recently released evidence that some drug makers have gone to great lengths to
turn scientific articles into marketing vehicles for their products, some influential medical
editors are cracking down on industry-financed ghostwriting. And they are getting help from some
members of Congress.

These editors are demanding that journals impose tougher disclosure policies for academic authors
and that the journals enforce their own rules by actively investigating the provenance of
manuscripts and by punishing authors who play down extensive contributions by ghostwriters.

In medical journal circles, the exorcism of industry-financed editorial assistance even has its own
name: ghostbusting.

In an editorial last week calling for a zero tolerance policy, the editors of the medical journal
PLoS Medicine, from the Public Library of Science, called for journals to identify and retract
ghostwritten articles and banish their authors.

``Any papers where this breach is substantiated should be immediately retracted,'' the editors
wrote. ``Authors found to have not declared such interest should be banned from any subsequent
publication in the journal and their misconduct reported to their institutions.''

In the past, researchers have raised allegations of ghostwriting in articles about quality-of-life
drugs like antidepressants, painkillers and diet pills. But the situation has become more serious
this year after a few editors said they had discovered ghostwriting in manuscripts about
life-and-death products like cancer and hematology drugs.

As Washington tries to revamp the health care system, concerns about ghostwriting are taking on new
urgency. One of the underlying assumptions of the health care overhaul effort is that money can be
saved and medical care improved by relying more heavily on research showing which drugs and
procedures are the most effective. But experts fear that the process could be corrupted if research
articles are skewed by the hidden influence of drug or medical device makers.

One senator on the trail of ghostwriting is Charles E.~Grassley, a Republican of Iowa and a member
of the Senate Finance Committee, which has taken a leading role in the health overhaul effort.

In July, Mr.~Grassley wrote letters asking eight leading medical journals about their ghostwriting
policies. He also asked whether, since 2004, the journals had taken action against any author who
had failed to report the involvement of a third party in the development of a manuscript.

None of the editors reported taking action against an author for ghostwriting. Their replies to the
senator, obtained by The New York Times, varied from assurances of editorial
diligence\cite{diligence} to the equivalent of ``don't ask, don't tell.'' One editor in chief, for
example, wrote that because his journal prohibited ghostwriting, the publication did not have a
specific policy on the practice.

Journals without explicit ghostwriting rules can expect to hear more from the senator.

``Objective research is really at the heart of public trust in medicine,'' Mr.~Grassley wrote in an
e-mail message to a reporter last Friday.

Allegations of ghostwriting first surfaced several years ago in the promotion of the diet drug
combination fen-phen, which was taken off the market because of safety concerns in 1997, and the
painkiller Vioxx, withdrawn in 2004. And last month, documents made public in litigation against the
pharmaceutical giant Wyeth showed that the company had paid a medical writing firm to draft
articles, published through 2005, favorable to its Premarin family of hormone drugs even as evidence
mounted that certain hormone drugs could increase the risk of breast cancer.

Some researchers say industry ghostwriting is widespread and continuing. Even with disclosure
policies already in effect at many publications, unnamed authors played a role in more than 40
articles published last year at six major medical journals, according to a study made public last
week. That study, conducted by an editorial team at The Journal of the American Medical Association,
or JAMA, defined ghostwriting broadly as any uncredited significant contribution to research or
writing, regardless of whether it was financed by industry.

Over the last few years, international associations of medical journal editors have developed
stricter disclosure criteria for authors of and contributors to scientific manuscripts. The
International Committee of Medical Journal Editors, for example, defines an author as a person who
makes a substantial contribution to developing a study or analyzing its results and in drafting a
manuscript, and who approves the final version of an article. Authors should identify other
contributors to an article and their financing sources, according to the group.

Drug companies say they are about to put these publication principles into effect for clinical
trials.

``The pharmaceutical industry is moving in lock step with the editors of medical journals,'' Jeffrey
K.~Francer, assistant general counsel of the Pharmaceutical Research and Manufacturers of America,
an industry trade group for drug makers, said in an interview last week. The new standards are to
take effect in October, he said.

But even though disclosure policies are already in place at many journals, the new JAMA study found
a ghostwriting rate of more than 7 percent at JAMA and PLoS Medicine, and nearly 11 percent at the
New England Journal of Medicine. Joseph S.~Wislar, who led the study, said in an interview last week
that The New England Journal of Medicine may have had a higher rate because the journal did not
require lead authors to list all other contributors.

Editors of The New England Journal of Medicine said that they were puzzled by and skeptical of the
JAMA data, but confirmed that the publication left such disclosures to the discretion of authors.

Experts who study disclosure said authorship policies might be inadequate in part because they asked
for incomplete information, but also because they typically had no teeth.

``Requiring someone to write a retraction or barring them from publishing in academic journals for
some period of time -- that would be an effective deterrent,'' said George Loewenstein, a professor
of economics and psychology at Carnegie Mellon University in Pittsburgh who has conducted research
on the effect of conflict-of-interest disclosures in medicine.

A few editors said they were already taking tougher stances after discovering their disclosure
policies had allowed authors to acknowledge writers financed by drug companies without explaining
that the paid writers played primary roles in creating the manuscripts.

The problem of incomplete disclosure is particularly worrisome for opinion pieces like review
articles, in which an author brings a personal perspective to a wide body of research, according to
an editorial in The Oncologist.

``These articles are likely to influence the direction of new investigation as well as the practice
of oncology,'' wrote Dr.~Bruce A.~Chabner, the clinical director of the cancer center at
Massachusetts General Hospital and the editor in chief of The Oncologist. ``It is critical that such
articles represent the unbiased views of the authors, and not those of a ghostwriter or a drug's
sponsor.''

The Oncologist plans to continue publishing clinical trials sponsored by drug companies, Dr.~Chabner
wrote. But the journal no longer accepts opinion pieces that involve writers with ties to companies
that have a commercial interest in an article's content -- nor will its editors correspond with
hired writers who are not named as the authors of manuscripts.

Mr.~Francer, of the Pharmaceutical Research and Manufacturers of America, said such measures could
be detrimental because they could ``chill research and chill support for research.''

But the trend may be too far along to deter.

In January, editors at Blood, the journal of the American Society of Hematology, discovered that an
unsolicited manuscript submitted by a prominent researcher involved significant contributions from a
pharmaceutical company employee named in the acknowledgments -- a major role in the manuscript that
should have qualified the employee to be listed as an author of the paper. Further detective work
quickly turned up two other ghostwritten manuscripts.

Editors decided to make their discoveries public in an editorial titled `` 'Ghostbusting' at Blood''
in which they wrote that the journal would henceforth reject opinion pieces that had industry ties.

In an interview last month, Dr.~Cynthia E.~Dunbar, the editor in chief of Blood, said that, in the
future, the journal would consider a ban of several years for authors caught lying about
ghostwriting, in addition to retracting their ghosted articles.

But, said Dr.~Dunbar, who is a hematologist at the National Institutes of Health in Bethesda, ``I
hope we don't have to do that.''

\section{Taming Your Digital Distractions}

\lettrine{I}{s} there any human invention more duplicitous than the personal computer? These
machines were manufactured and initially marketed as devices to help us at work. We were told they
would perform amazing feats of office derring-do -- adding up rows of numbers effortlessly, turning
our musings into beautiful magazine-quality documents, and letting us collaborate with one another
across continents.

Boy, that turned out well, didn't it? Sure, you could use your PC to analyze stats for the annual
sales report due in two days. But hey, look at this -- someone wants tobe your friend on Facebook!
And wait a second: A zany couple decided to start off their wedding by dancing down the aisle, and
lucky for everyone, they posted the video on YouTube. And did you hear what that ignorant
congressman just said about health care? Now you've got no choice but to spend the next five minutes
crafting an impassioned tweet to express your outrage.

And so it goes: You get to your PC every morning with hours of productive time ahead of you. Next
thing you know, it's 5 p.m. and you've frittered the day away on Digg, Hulu, Wikipedia and your
fantasy football league. And no wonder -- how can anyone expect to get anything done when you're
plying your trade on one of the most distracting machines ever invented? With so much available on
your PC -- your friends, blogs, games and even TV shows -- working in a modern office can often seem
as rattling as working on the floor of a Las Vegas casino.

During the last few weeks, I've been using a slate of programs to tame these digital distractions.
The apps break down into three broad categories. The most innocuous simply try to monitor my online
habits in an effort to shame me into working more productively. Others reduce visual bells and
whistles on my desktop as a way to keep me focused.

And then there are the apps that really mean business -- they let me actively block various parts of
the Internet so that when my mind strays, I'm prohibited from giving in to my shiftless ways. It's
the digital equivalent of dieting by locking up the refrigerator and throwing away the key.

The first category is epitomized by a program called RescueTime, which keeps track of everything
that happens on your computer, and then reports your habits in a series of charts and graphs. I
found the software's analysis tremendously illuminating. I learned, for instance, that during a
typical month I spend more than 70 hours surfing the Web, much of it on news and social-networking
sites.

By comparison, I spend only about half as much time in Microsoft Word, which, as a writer, is where
I do my work. Seeing these stats knocked me over; clearly, I wasn't using my time very wisely.
(RescueTime offers a free limited version; an upgraded plan with deeper usage stats costs \$8 a
month.)

So what to do? Over the years, several friends have suggested that I might stay more focused by
ditching Word in favor of a so-called minimalist writing tool. These programs, like Hog Bay
Software's WriteRoom for the Mac (\$24.95) or the free Dark Room for Windows, essentially take your
computer back two decades in time. Each presents you with a full-screen, monochrome window absent of
taskbars and menus; the experience is that of typing on an old-fashioned dedicated word processor,
with every other function of a modern PC hidden from view.

But the procrastinator's mind is not so quickly deceived. I found that I could easily switch from
working in Dark Room to wasting time in a Web browser -- and that was a problem.

Time for stronger medicine: I loaded up LeechBlock, a free add-on for Firefox whose main function is
to save you from yourself. LeechBlock works like a stern nanny: You tell it which Web sites to keep
you away from, and at the appointed hour, it stops you cold. Try to go to Facebook and you get back
a warning to go back to work.

The software is quite flexible. You can let it block out different sites at different times of the
day, or set a maximum daily or hourly limit for certain sites. For instance, I asked LeechBlock to
restrict my time on Twitter and Facebook to no more than five minutes an hour, and on news sites to
no more than 10 minutes an hour. This gave me a little bit of time to goof off, but not enough that
I'd lose sight of my larger purpose.

But LeechBlock suffers a crucial limitation: if you really want to get around it, all you have to do
is load up another browser. One Mac app that has found a way to solve this problem is called
Freedom, which blocks all of your computer's networking functions for a pre-determined number of
minutes. In other words, once you set it, you've got no Web, no instant messaging, no e-mail -- and
the only way to undo Freedom's block before the time runs out is to restart your machine.

I wish I could say that using these digital nannies has revolutionized the way I work. They didn't,
really. Though blocking time-sucking Web sites did keep me from goofing off on my computer, I found
that my brain quickly compensated by wasting time in other ways: As I'm writing this paragraph, for
example, I'm also eating a peach. But not just eating it without thinking -- I've been using a
paring knife to try to cut perfectly cubical pieces to pop into my mouth.

Perhaps this kind of unconscious fidgeting -- whether online or off -- is inevitable. The mind is a
restless place, and creative pursuits like writing seem unsustainable in long bursts; perhaps the
mind just needs frequent breaks.

But I did notice that net-blocking software was helpful in getting me to at least consider all the
ways that I was wasting my time. When LeechBlock threw up a roadblock in my path, it gave me pause;
when I went around it, I was at least conscious that it wasn't the right thing to do. Sometimes a
little shame is all you need.

\section{Palm Reports Bigger Loss Amid Declining Revenue}

\lettrine{S}{UNNYVALE}, Calif. (AP) -- The smartphone maker Palm posted a wider loss for its first
quarter on Thursday as revenue dropped, but adjusted earnings and sales topped Wall Street's
expectations.

For the three months ended Aug.~28, the company posted a loss after paying preferred dividends of
\$164.5 million, or \$1.17 a share, compared with a loss of \$41.9 million, or 39 cents a share, in
the period a year earlier.

Excluding one-time items, the company posted a much smaller loss of \$13.6 million, or 10 cents a
share, in the latest quarter.

Revenue tumbled to \$68 million from \$366.9 million. But adjusted sales, which exclude deferred
revenue and cost of sales from the Pre smartphone, totaled \$360.7 million.

Analysts on average were expecting a loss, excluding items, of 24 cents a share on adjusted sales of
\$297.7 million, according to a poll by Thomson Reuters.

``We're making significant progress with Palm's transformation, and our culture of innovation is
stronger than ever. We're launching more great Palm web OS products with more carriers, and turning
our sights toward growth,'' the chief executive, Jonathan Rubinstein, said in a statement.

Palm shipped 823,000 smartphones during the quarter, down 30 percent year over year. But compared
with the prior quarter, shipments more than doubled.

Palm said it planned to offer about 16 million of its shares and use the proceeds for working
capital and general corporate purposes. It also forecast second-quarter adjusted revenue of \$240
million to \$270 million and full-year revenue on the same basis of \$1.6 billion to \$1.8 billion.

\section{Canada Won't Review Sale of Nortel Unit}

\lettrine{T}{ony} Clement, the industry minister for Canada, said on Wednesday that his department
would not review the takeover of the wireless assets of bankrupt Nortel Networks by Ericsson of
Sweden.

Some Canadians, most prominently the co-chief executives of Research in Motion, the BlackBerry
maker, asked Mr.~Clement to examine the sale under foreign investment review laws. They argued that
because the company had benefited from Canadian government assistance, its assets should remain
under domestic control. R.I.M. also charged that it had been effectively and unfairly barred from
bidding.

But in a statement, Mr.~Clement praised the \$1.13 billion takeover of the unit which, unlike many
other parts of Nortel, is profitable.

``This deal is beneficial to Canada and will help Canada compete on the world stage,'' the minister
said.

A review by his department, he added, found that not enough of the wireless operation's assets are
in Canada to meet the threshold for review under the foreign investment regulations.

Speaking to reporters in Toronto, Mr.~Clement said that there was also no evidence that Nortel's
future wireless technology, the assets of most interest to Ericsson, were developed using government
money. Nortel itself has said that its lack of profits for the last several years made it unable to
exploit tax concessions.

In any case, Nortel is only licensing patents related to that technology under the terms of its sale
agreement with Ericsson.

Separately, Nortel said courts in Ontario and Delaware had approved the sale of another unit, which
makes networking equipment for corporations and governments, to Avaya. Verizon, a major Nortel
customer, had objected to the \$900 million deal in the Delaware bankruptcy court on national
security grounds. But that court concluded that Verizon's concerns were commercial, rather than
national, in nature.

\vocabulary
\end{document}
